% Working manuscript, 23 October 2026. Standalone pdfLaTeX source.
\documentclass[11pt,a4paper]{article}
\usepackage[T1]{fontenc}
\usepackage[utf8]{inputenc}
\usepackage{lmodern,microtype}
\usepackage[margin=25mm]{geometry}
\usepackage{amsmath,amssymb,amsthm,mathtools,mathrsfs}
\usepackage{booktabs,array,longtable,enumitem,listings}
\usepackage[dvipsnames]{xcolor}
\usepackage[colorlinks=true,linkcolor=MidnightBlue,citecolor=MidnightBlue,urlcolor=MidnightBlue]{hyperref}
\usepackage{bookmark,xurl}
\usepackage{fancyhdr}
\usepackage{needspace}
\pagestyle{fancy}\fancyhf{}
\fancyhead[L]{\small Untwisted cyclic Haagerup--Izumi centers}
\fancyhead[R]{}
\fancyfoot[C]{\thepage}
\setlength{\headheight}{14pt}
\setlength{\emergencystretch}{3em}
\setlength{\parskip}{2.5pt}
\setlength{\parindent}{1.2em}
\setlist{topsep=4pt,itemsep=3pt}
\lstset{language=Python,basicstyle=\ttfamily\scriptsize,
 breaklines=true,columns=fullflexible,keepspaces=true,
 numbers=left,numberstyle=\tiny\color{gray},numbersep=6pt,
 frame=single,rulecolor=\color{gray},showstringspaces=false}
\numberwithin{equation}{section}
\newtheorem{theorem}{Theorem}[section]
\newtheorem{proposition}[theorem]{Proposition}
\newtheorem{lemma}[theorem]{Lemma}
\newtheorem{corollary}[theorem]{Corollary}
\theoremstyle{definition}\newtheorem{definition}[theorem]{Definition}
\theoremstyle{remark}\newtheorem{remark}[theorem]{Remark}
\newcommand{\C}{\mathbb C}\newcommand{\Z}{\mathbb Z}\newcommand{\Q}{\mathbb Q}
\newcommand{\one}{\mathbf1}\newcommand{\B}{\mathcal B}\newcommand{\W}{W_N}
\newcommand{\FP}{\operatorname{FPdim}}\newcommand{\FPdim}{\operatorname{FPdim}}
\newcommand{\Ind}{\operatorname{Ind}}\newcommand{\Irr}{\operatorname{Irr}}
\newcommand{\Hom}{\operatorname{Hom}}\newcommand{\End}{\operatorname{End}}
\newcommand{\Rep}{\operatorname{Rep}}\newcommand{\Ext}{\operatorname{Ext}}
\newcommand{\rk}{\operatorname{rk}}\newcommand{\Tr}{\operatorname{Tr}}
\newcommand{\wt}{\operatorname{wt}}\newcommand{\loc}{\mathrm{loc}}
\newcommand{\ptr}{\operatorname{ptr}}\newcommand{\ch}{\operatorname{ch}}
\newcommand{\diag}{\operatorname{diag}}\renewcommand{\Vec}{\operatorname{Vec}}
\newcommand{\bo}{\boxtimes}\newcommand{\Cat}{\mathcal C}
\newcommand{\Par}{\mathcal P}\newcommand{\Neu}{\mathcal D}
\newcommand{\Mod}{\mathcal M}\newcommand{\HI}{\mathcal H}
\newcommand{\Frame}{\mathcal B}\newcommand{\mm}{\mathfrak m}
\newcommand{\eps}{\varepsilon}\newcommand{\NN}{\mathbb Z_{\geq0}}
\newcommand{\ds}{\displaystyle}
\newcommand{\parityrank}{p_N}
\setcounter{tocdepth}{1}
\hypersetup{pdftitle={Vertex operator algebra realizations of untwisted cyclic Haagerup--Izumi centers},pdfauthor={David E. Evans and Andrew Riesen},pdfsubject={Working manuscript, 23 October 2026},pdfkeywords={Haagerup--Izumi categories, vertex operator algebras}}
\title{\textbf{Vertex operator algebra realizations of\\untwisted cyclic Haagerup--Izumi centers}}
\author{David E. Evans\and Andrew Riesen}
\date{23 October 2026}
\begin{document}
\maketitle
\begin{abstract}
For every odd integer $N\geq3$, we construct an untwisted cyclic
Haagerup--Izumi fusion category $\mathcal H_N$ and a simple,
self-contragredient, regular, $C_2$-cofinite vertex operator algebra
$W_N$ of CFT type, with central charge $-5N$ and a ribbon equivalence:
\[
 \operatorname{Rep}(W_N)\simeq
 \mathcal P_N\boxtimes Z(\mathcal H_N).
\]
Here $\mathcal P_N$ is a pointed ribbon category of rank $2^N$, and
noninvertible simples of $\mathcal H_N$ have spherical dimension
$-(\sqrt{N^2+4}-N)/2$. An extension by shifted $E_7$ lattice modules
removes the pointed factor. The resulting simple, self-contragredient,
regular, $C_2$-cofinite VOA $E_N$ has central charge $-16N$ and
representation category $Z(\mathcal H_N)$, of rank $N^2+3$.
It is lower bounded and integer graded, with least conformal degree
$-(N-1)/2$.

The construction realizes Evans's cyclic Virasoro sum. Its regularity
follows from an explicit lattice frame and projectivity of the vacuum;
refined modular traces and branching then determine all simple modules.
Two condensations identify the ribbon center, whose full modular data
are the negative-branch datum
$\mathcal{MD}_{-}(\mathbb Z/N;\mathbb Z/(N^2+4),b\mapsto b^2/(N^2+4))$
of Gannon--Schopieray--Yadav. This comparison holds uniformly, including
at composite discriminant. An equivalence with an independently
prescribed Haagerup--Izumi category is not asserted.
\end{abstract}


\section{Introduction and main theorem}
\label{sec:intro}

The representation category of a rational vertex operator algebra,
under the standard finiteness and duality hypotheses, records its
fusion, braiding and conformal twists. The converse realization problem
asks for a vertex operator algebra with prescribed categorical data.
Haagerup--Izumi categories provide a family of test cases: their simple
objects consist of a group of invertibles and a single orbit of
noninvertibles, while their Drinfeld centers are modular tensor categories.

We construct vertex operator algebra realizations for every odd cyclic
order $N\geq3$, with the negative spherical dimension specified below.
The first algebra, $W_N$, is a CFT-type extension of a tensor product
of $N$ irrational Virasoro algebras. Its representation category
contains a nondegenerate pointed factor $\mathcal P_N$ of rank $2^N$.
We construct an untwisted cyclic HI category $\mathcal H_N$ and prove:
\[
 \operatorname{Rep}(W_N)\simeq
 \mathcal P_N\boxtimes Z(\mathcal H_N),
 \qquad c(W_N)=-5N.
\]
An explicit extension using shifted $E_7$ lattice modules cancels
$\mathcal P_N$ and gives a second algebra $E_N$, with central charge
$-16N$ and representation category $Z(\mathcal H_N)$. Both algebras
are simple, self-contragredient, regular and $C_2$-cofinite. Their
grading conventions and complete invariants are stated in
Theorem~\ref{thm:main}.

Two features of the proof are essential. First, the Virasoro factors
are not rational, so regularity of their extension requires a separate
argument. We prove $C_2$-cofiniteness and split extensions of the
vacuum before using modular traces. Second, a single conformal-weight residue does not distinguish
all exceptional labels when $N^2+4$ is composite. Keeping the partial conformal weights resolves these
labels and makes the branching calculation sufficient to exhaust
the simple modules. The center equivalence then follows from two
specified condensation algebras.

\paragraph{Origin of the construction.}
Evans \cite[Section 6.2]{evanssfold} proposed the cyclic Virasoro sum
$V_\varphi$ with parameters $\psi_i+\psi_{i-1}^{-1}=k_i$.
For $(k_0,k_1,\ldots,k_{N-1})=(N+2,2,\ldots,2)$, that sum is exactly
the underlying frame module of $W_N$ used here. The cited section
presents a heuristic cyclic closure of the duality-wall gluing of
Creutzig--Gaiotto \cite{cgduality}. The work below supplies the uniform
VOA construction, regularity proof, module exhaustion and ribbon
center identification for this cyclic family. The earlier physical
proposals of Gang--Kim, Gang--Kim--Lee and Jeong--Sohn
\cite{gangkim,gangkimlee,jeongsohn} provide the S-fold and character
context; the hyperbolic-word description is developed in
\cite{evanswords}. Flohr's screened one-variable algebras and theta-difference
characters \cite{flohr1993,flohr1994} are an earlier precedent for
these methods. Their character calculations do not supply the
$N$-variable regularity argument proved below.

\paragraph{Relation to categorical existence results.}
The categorical background includes Izumi's endomorphism and tube
algebras \cite{izumi}, the Evans--Gannon grafting of modular data
\cite[Section 3.3]{evansgannon2011}, and their near-group and
nonunitary constructions \cite{evansgannon2014,evansgannon2017}.
Existence for every odd cyclic order is already established in
\cite{huangcyclic} and \cite[Theorem A(1)]{gsy}; the latter includes
the pseudo-unitary Galois conjugate. Its Theorem 5.1 identifies the
untwisted center modular data. Thus the result here is a VOA and
condensation realization, rather than a new categorical existence
claim.

We also identify the full modular data of our center with the
negative cyclic class $(N^2+4)^-_1$ of Gannon--Schopieray--Yadav,
including a bijection of simple labels. This is distinct from an
equivalence with a prescribed category. At order fifteen, their
examples have identical center modular data but are pairwise not
categorically Morita equivalent \cite[Corollary 5.9]{gsy}.
Our center equivalence retains the condensation algebras and their
multiplication; the modular-data comparison is a subsequent theorem.
The announced $c=8$ Haagerup VOA \cite{riesenc8} concerns a separate
positive-central-charge realization.

\subsection{Conventions and statement}

Throughout, $N\geq3$ is odd. All categories are over $\mathbb C$. A fusion category is finite,
semisimple, rigid, and has simple unit. A modular tensor category carries
its specified ribbon structure; dimensions below are spherical dimensions
unless marked $\operatorname{FPdim}$. For a fusion category $\mathcal H$,
$Z(\mathcal H)$ denotes its Drinfeld center with the ribbon structure
induced by the stated spherical structure.

A VOA is lower bounded and integer graded, with finite-dimensional
homogeneous spaces. It is of CFT type when its negative degrees vanish
and its degree-zero space is $\mathbb C\mathbf1$. A VOA is \emph{regular}
when every weak module is a direct sum of irreducible ordinary modules.
Regularity implies rationality, meaning complete reducibility of
admissible modules, for both $W_N$ and $E_N$. The regularity of $E_N$
is proved directly, without assuming CFT type. We use $C_2$-cofinite
in the sense
\[
 C_2(V)=\operatorname{span}_{\mathbb C}\{a_{-2}b:a,b\in V\},
 \qquad \dim V/C_2(V)<\infty.
\]
This condition depends on the vertex operators, not on the choice
of conformal vector. Self-contragredience means an isomorphism
$V\cong V^\vee$ with the graded contragredient module for the stated
conformal vector. The notation $\operatorname{Rep}(V)$ denotes the
category of grading-restricted generalized modules; for the regular
VOAs considered here it consists of finite direct sums of ordinary
simple modules.
For a fusion category, its rank is the number of isomorphism classes of
simple objects. We write $\boxtimes$ for Deligne products of categories
and, with a VOA subscript when needed, for fusion products of modules.

\begin{definition}
An untwisted cyclic Haagerup--Izumi (HI) fusion category of odd order $N$
has simple objects $g$ and $g\rho$, $g\in G=\mathbb Z/N$, with:
\[
 gh=g+h,\qquad \rho g=(-g)\rho,\qquad
 \rho^2=\mathbf1\oplus\bigoplus_{g\in G}g\rho,
\]
and its pointed subcategory is tensor equivalent to
$\operatorname{Vec}_{G}$ with trivial associator. This condition does
not prescribe the remaining associator entries.
\end{definition}

Set:
\[
 \mu=N^2+4,\qquad
 t=\frac{\sqrt\mu-N}{2},\qquad
 \delta=t^{-1}=\frac{N+\sqrt\mu}{2},\qquad \parityrank=2^N.
\]
Thus $0<t<1$, $t(t+N)=1$, and $\delta^2=1+N\delta$.
The two roots of the HI dimension equation $d^2=1+Nd$ are $-t$ and
$\delta$. The term \emph{negative spherical branch} refers to the root
$-t$; it does not assert that both roots arise from different pivotal
structures on one fixed fusion category.
The symbol $\mathcal P_N'$ in the theorem denotes the centralizer of
$\mathcal P_N$ in $\operatorname{Rep}(W_N)$: its objects have trivial
double braiding with every object of $\mathcal P_N$.

\begin{theorem}\label{thm:main}\label{r:thm:main}
For every odd integer $N\geq3$, the constructions in this paper give:
\begin{enumerate}
\item A simple, self-contragredient, regular, $C_2$-cofinite VOA $W_N$
of CFT type, with $c(W_N)=-5N$ and $(W_N)_1=0$.
Its modular tensor category has rank $2^N(N^2+3)$ and contains a
nondegenerate pointed ribbon subcategory $\mathcal P_N$ of rank $2^N$.
\item An untwisted cyclic HI fusion category $\mathcal H_N$,
obtained by the specified two-stage condensation, with
$\dim(g)=1$ and $\dim(g\rho)=-t$, and a ribbon equivalence:
\[
 \mathcal P_N'\simeq Z(\mathcal H_N).
\]
\item A simple, self-contragredient, regular, $C_2$-cofinite VOA $E_N$,
given by the explicit finite extension in
Theorem~\ref{thm:parity-cancellation-voa}, with:
\[
 c(E_N)=-16N,\qquad
 \operatorname{Rep}(E_N)\simeq Z(\mathcal H_N),\qquad
 |\operatorname{Irr}(E_N)|=N^2+3.
\]
Its least conformal degree is exactly $-(N-1)/2$.
\item The normalized ribbon modular data of
$\mathcal P_N'$, $Z(\mathcal H_N)$, and $\Rep(E_N)$ are
\[
 \mathcal{MD}_{-}\bigl(\mathbb Z/N;\mathbb Z/\mu,
                            b\longmapsto b^2/\mu\bigr),
 \qquad \mu=N^2+4,
\]
in the conventions of Gannon--Schopieray--Yadav
\cite[Appendix F.1]{gsy}. The explicit bijection of labels and all
entries of $S$ and $T^{\mathrm{rib}}$ are given in
Theorem~\ref{thm:gsy-modular}.
\end{enumerate}
Both VOAs are rational. The VOA $E_N$ is lower bounded and integer
graded, with finite-dimensional homogeneous spaces; its regularity
is in the weak-module sense specified above. Its negative least
degree excludes CFT type, but not $C_2$-cofiniteness or
self-contragredience.
\end{theorem}

The category $\mathcal H_N$ in part (2) is the one constructed by
condensation; no uniqueness of an HI category with these fusion rules
is assumed. Part (4) identifies its modular data, with the tensor unit
fixed and $S_{\one,\one}=1/\dim\mathcal H_N$. The negative spherical
dimensions make these realizations nonunitary.
For ordinary VOA characters, $T^{\mathrm{an}}=e^{-2\pi i c/24}T^{\mathrm{rib}}$;
this scalar is accounted for in Section~\ref{sec:gsy-modular}.

\subsection{Frame and lattice notation}

We collect the conventions used in the construction. The proof outline
in Section~\ref{sec:proof-structure} can be read independently of
these formulas.

Indices $i,j$ are cyclic modulo $N$ unless a signed antiperiodic
extension is stated. Put:
\[
 \psi_i=\frac{t+i+1}{t+i},\qquad
 c_i=13-6(\psi_i+\psi_i^{-1}),\qquad
 h_i(R,S)=\frac{(\psi_iR-S)^2-(\psi_i-1)^2}{4\psi_i}.
\]
Each $\psi_i$ is positive and irrational. The key identities are:
\[
 \psi_i+\psi_{i-1}^{-1}=k_i,\qquad
 k_0=N+2,\quad k_i=2\ (i>0).
\]
The frame is $\mathcal B=\bigotimes_iL(c_i,0)$, and
$L_i(R,S)$ abbreviates the simple Virasoro module of central charge
$c_i$ and highest weight $h_i(R,S)$. For positive integral $R,S$,
it is the irreducible quotient at the degenerate Kac weight.

The positive cycle form $K=K_+$ is defined without an ambiguous
small-cycle matrix convention by:
\[
 x^{\mathsf t}K_+x=\sum_i(x_i-x_{i+1})^2+Nx_0^2.
\]
The antiperiodic form $K_-$ replaces the closing difference
$(x_{N-1}-x_0)^2$ by $(x_{N-1}+x_0)^2$.
Their determinants are $N^2$ and $\mu$.
Use the even positive lattice $L=(\mathbb Z^N,2K_+)$ with basis $\beta_i$.

Write $B=\sum_i\beta_i$, $\varepsilon=e^B$, and $D=L(-1)$.
In character formulas, $q=e^{2\pi i\tau}$ and
$q_i=e^{2\pi i\tau_i}$. Quantum-group parameters denoted by $q$
are introduced locally; the parity rank is always $\parityrank$.

For $m\in\mathbb Z_{\geq0}^N$, write:
\[
 M_m=\bigotimes_iL_i(2m_i+1,2m_{i+1}+1),\qquad R_j=M_{e_j}.
\]
The state and module called $C$ in the vacuum calculation correspond
to $m=\mathbf1$, whereas $\mathcal C=\operatorname{Rep}(W_N)$ is
the full representation category. A proper interval containing site
$0$ is called \emph{heavy}; an \emph{ordinary} proper interval avoids
$0$. We also write $W=W_N$ when $N$ is fixed. Bold $\mathbf1$ denotes
the vacuum or the all-ones coordinate vector according to its ambient
space; the distinction is explicit whenever both occur together.

The main categorical notation is collected here for reference.
\begin{center}\small
\begin{tabular}{@{}ll@{}}
\toprule
Symbol & Meaning\\
\midrule
$\mathcal B$ & The tensor product of the $N$ Virasoro frame algebras\\
$A$ & The deformed lattice VOA, also its algebra object over $W_N$\\
$\mathcal C$ & $\operatorname{Rep}(W_N)$ after regularity is established\\
$P=(\mathbb Z/2)^N$, $\parityrank$ & The parity-label group and its order $2^N$\\
$\mathcal P=\mathcal P_N$ & The pointed ribbon subcategory on the $J_\epsilon$\\
$\mathcal D=\mathcal P'$ & The centralizer of $\mathcal P$ in $\mathcal C$\\
$\mathcal M=\mathcal D_A$ & Right $A$-modules in $\mathcal D$\\
$\mathcal D_A^{\mathrm{loc}}$ & The full subcategory of local $A$-modules\\
$\mathcal H_N$ & The fusion category obtained by the second condensation\\
\bottomrule
\end{tabular}
\end{center}
\subsection{Structure of the proof}
\label{sec:proof-structure}

Regularity precedes the module classification and modular traces.

\emph{Construction (Part~I).}
Cyclic Virasoro gluing gives $W_N$ (Theorem~\ref{c:cyclic-gluing}).
Theorem~\ref{c:cyclic-c2} proves $C_2$-cofiniteness, and the explicit
commuting lattice frame gives a conformal embedding
(Theorem~\ref{c:screening-frame} and Corollary~\ref{c:cyclic-embedding}).
Theorem~\ref{c:currents} constructs the parity simple currents.

\emph{Regularity (Part~II).}
Proposition~\ref{neutral:prop:two-lowest-weights} reduces vacuum
extensions to lowest weights zero and one.
Theorems~\ref{r:thm:h0} and~\ref{r:thm:h1} exclude these cases,
making the vacuum projective. Lemma~\ref{lem:projective-unit} then
gives rigidity and semisimplicity, and
Theorem~\ref{r:thm:regularity} passes to all weak modules.
The parity factorization follows in Section~\ref{sec:parity-factor}.

\emph{Modules and traces (Part~III).}
We classify the periodic lattice sectors and prove the needed
uniqueness results for rational frame sectors. Formal refined
traces retain the joint partial weights and detect every exceptional
label, including those not separated by a single conformal-weight
residue when $N^2+4$ is composite.

\emph{Center identification and realization (Part~IV).}
Exact branching and the dimension identity in
Section~\ref{r:sec:budget} exhaust the simple modules. Two
condensations use the lattice algebra $A$ and the regular algebra
$D_H$ of a cyclic Lagrangian subgroup:
\[
 \mathcal D=\mathcal P_N'
 \quad\longmapsto\quad \mathcal M=\mathcal D_A
 \quad\longmapsto\quad \mathcal H_N=\mathcal M_{D_H}.
\]
The second local category is trivial.
Corollary~\ref{cor:our-center} identifies the ribbon center;
Section~\ref{sec:gsy-modular} computes its full modular data.
Theorem~\ref{thm:parity-cancellation-voa} cancels the parity factor by an explicit
$E_7$ extension and proves its finiteness, duality, weak-module
regularity and grading assertions.

The appendices supply the interval and Zhu calculations and the
symbolic frame certificate. Section~\ref{c:app:frame} explains the
reduction to that finite certificate. No theorem from a companion
twisted-case manuscript is assumed.

\paragraph{AI disclosure.}
Generative AI was used extensively in author-directed proof
exploration, symbolic calculations, code development, and drafting,
and produced substantial portions of the initial arguments and
exposition. The authors formulated the research objective, guided
the mathematical strategy, and revised the arguments, notation,
and presentation.

The authors take responsibility for the mathematical claims and
final text. The accompanying calculations and code make the finite
checks inspectable; their role in the proof is explained in the
manuscript. Independent human verification of the complete argument
remains pending. Corrections and feedback on the mathematics,
exposition, attribution, and historical account are welcome.

\paragraph{Acknowledgments.}
The work of the first author was supported in part by a grant of access to OpenAI models through the ChatGPT for Academic Researchers programme.

% BEGIN COMPACT CONTENTS
\clearpage
\section*{Contents}
\begingroup
\small
\setlength{\parskip}{0pt}
\makeatletter
\renewcommand{\l@section}{\@dottedtocline{1}{0pt}{2.5em}}
\renewcommand{\l@part}[2]{%
  \addpenalty{-\@highpenalty}%
  \addvspace{0.8em}%
  {\bfseries\raggedright #1\par}\nobreak}
\@starttoc{toc}
\makeatother
\endgroup
\clearpage
% END COMPACT CONTENTS
\part{Construction and lattice realization}

\section{Cyclic gluing of generic Virasoro algebras}
\label{sec:cyclic-gluing}
Fix an odd integer $N\ge3$ and retain the parameters and Virasoro
modules introduced in Section~\ref{sec:intro}. Thus
$\psi_i=(t+i+1)/(t+i)$ and $t(t+N)=1$. We construct $W_N$ as an
extension of the frame $\mathcal B$ by transporting a canonical
diagonal algebra.

\begin{lemma}[Signs of the quantum parameter on integer spins]\label{c:quantum-sign}
Let $q\in\mathbb C^\times$ be not a root of unity. The type-I
integer-spin ribbon categories for $q$ and $-q$ are ribbon equivalent.  The corresponding
category for $q^{-1}$ is braid-reversed equivalent to that for $q$.
\end{lemma}
\begin{proof}
Use $\Delta E=E\otimes K+1\otimes E$ and
$\Delta F=F\otimes1+K^{-1}\otimes F$.
The map from $U_{-q}(\mathfrak{sl}_2)$ to $U_q(\mathfrak{sl}_2)$ given by
$K'\mapsto K$, $E'\mapsto-E$, and $F'\mapsto F$ is a Hopf isomorphism.
On even highest weights it preserves the type-I condition.  In the
standard universal $R$-matrix its $n$th nilpotent coefficient is
\[
 q^{n(n-1)/2}\frac{(q-q^{-1})^n}{[n]_q!}E^n\otimes F^n.
\]
Since $[n]_{-q}=(-1)^{n-1}[n]_q$, changing $q$ to $-q$ changes the
scalar coefficient by $(-1)^n$, which is cancelled by $E'\mapsto-E$.
The diagonal factor $q^{H\otimes H/2}$ is unchanged on even weights:
the product of two even weights divided by two is even.  The ribbon
twist is unchanged as well, since its exponent on highest weight
$2m$ is $(2m)(2m+2)/2=2m(m+1)$.
For inversion, the Hopf isomorphism
$K'\mapsto K$, $E'\mapsto KF$, $F'\mapsto EK^{-1}$ from
$U_{q^{-1}}$ to $U_q$ carries the standard $R$-matrix to
$R_{21}^{-1}$ and the ribbon twist to its inverse; this follows
term by term from the same $R$-matrix formula.  These identities
give the stated ribbon equivalences.
\end{proof}

\begin{theorem}[Cyclic gluing]\label{c:cyclic-gluing}
There is a simple VOA $W_N$ of CFT type and central charge $-5N$,
with frame decomposition:
\[
 W_N=\bigoplus_{r_0,\ldots,r_{N-1}\in2\Z_{\ge0}+1}
       \bigotimes_{i=0}^{N-1}L_i(r_i,r_{i+1}).
\]
It is self-contragredient, $(W_N)_1=0$, and $\dim(W_N)_2=N(N+1)/2$.
\end{theorem}
\begin{proof}
At an irrational Virasoro parameter, the category of finite-length
$C_1$-cofinite modules is a semisimple rigid vertex tensor category
\cite{virasoro}. Its odd first-row and first-column subcategories are generic integer-spin quantum $\mathfrak{sl}_2$ ribbon categories \cite[Theorem 1.1]{c:Koshida}. They centralize one another: the unique product is $L_i(r,s)$ and
\[
 h_i(r,s)-h_i(r,1)-h_i(1,s)=-\frac{(r-1)(s-1)}2,
\]
so balancing gives trivial monodromy for odd $r,s$. Consequently, the tensor-product functor is braided on these odd-label
subcategories. It maps distinct simple pairs bijectively to the odd-label
simples, by the fusion rules in \cite[Theorem 5.2.5]{virasoro}.

The cyclic relations are
\[
 \psi_0+\psi_{N-1}^{-1}=N+2,\qquad
 \psi_i+\psi_{i-1}^{-1}=2\quad(1\le i<N).
\]
For the second equality, substitute the displayed fractions. For the first, use $1/t=t+N$ and $1/(t+N)=t$. The odd first-row category at site $i$ is therefore braid-reversed equivalent to the odd first-column category at site $i-1$: their squared quantum parameters are mutual inverses. Lemma~\ref{c:quantum-sign} removes the possible sign of a quantum parameter on integer spins. The twists agree with this reversal because, for odd $r$, the sum of the two paired conformal weights is an integer.

Let $\mathcal R$ be the product of these $N$ odd first-row categories. Its canonical diagonal algebra in $\mathcal R\bo\mathcal R^{\mathrm{rev}}$ exists in the direct-sum completion and is simple, by \cite[Theorem 3.5 and Remark 3.7]{gluing}. Transport it through the preceding braid-reversed equivalences and reorder factors. This produces exactly the displayed multiplicity-free sum $W_N$, with nonzero multiplication in every allowed fusion channel. The odd row and column subcategories at each site centralize each other, so this transport is a braided equivalence onto the relevant full subcategory of the tensor-product Virasoro category.

To check grading restriction, write $r_i=2m_i+1$ with $m_i\ge0$. The lowest weight of its summand is
\begin{equation}\label{c:chain-weight}
 h(\mathbf m)=\sum_{i=0}^{N-1}(m_i-m_{i+1})^2+Nm_0(m_0+1).
\end{equation}
It is a nonnegative integer, vanishing only at $\mathbf m=0$. At bounded weight, $m_0$ is bounded and so are every successive difference and every $m_i$; hence only finitely many summands contribute. Each has finite-dimensional graded pieces. The twist is the identity and the unit is injective. Theorem 5.7 of \cite{gluing} therefore turns this commutative algebra into a conformal vertex algebra, and the grading calculation makes it a VOA. Simplicity follows either from simplicity of the canonical algebra, or by multiplying any nonzero simple summand in an ideal by its dual to reach the vacuum.

Summing the cyclic relations gives $\sum_i(\psi_i+\psi_i^{-1})=3N$, hence $\sum_i c_i=-5N$. There is no weight-one state: a nonconstant cyclic integer sequence has squared-difference sum at least two, while a nonzero constant sequence has $Nm_0(m_0+1)\ge2N$. For weight two, $m_0=0$ and the sequence consists of ones on a nonempty interval among the $N-1$ other sites and zeros elsewhere. There are $N(N-1)/2$ such intervals. Together with the $N$ Virasoro stress tensors this gives $N(N+1)/2$. Finally a simple CFT-type VOA with zero weight-one space has a nondegenerate invariant bilinear form, so is self-contragredient.
\end{proof}


\section{Uniform \texorpdfstring{$C_2$}{C2}-cofiniteness}\label{c:app:c2}
We prove finiteness in two steps. First, the $C_2$ quotient is finite
over the algebra generated by the stress tensors. Second, the
level-three null relations transport the values of these generators
around the cycle. The resulting monodromy forces every such value
to vanish, so the stress-tensor algebra is itself finite-dimensional.

\begin{theorem}\label{c:cyclic-c2}
For every odd $N\ge3$, the VOA $W_N$ of
Theorem~\ref{c:cyclic-gluing} is $C_2$-cofinite. It admits a finite set
of strong generators of conformal weights at most $10N$.
\end{theorem}


\subsection*{Finite generation over the stress tensors}
Write $M_{\mathbf m}$ for the summand in \eqref{c:chain-weight}, and $v_{\mathbf m}$ for its highest vector. Put
\[
 R=W_N/C_2(W_N),\quad x_i=[T_i],\quad
 R_{\mathrm{st}}=\C[x_0,\ldots,x_{N-1}]\subset R,\quad D_i=(T_i)_{(0)}=L^{(i)}_{-1}.
\]
Here $R_{\mathrm{st}}$ denotes the image of the polynomial algebra. Each $D_i$ preserves $C_2$, since zero modes are derivations. In $R$, $D_i x_j=0$, the modes $L^{(i)}_{-k}$ for $k\ge3$ act by zero, and $L^{(i)}_{-2}$ acts by multiplication by $x_i$.

\begin{lemma}\label{c:stress-finite}
The algebra $R$ is finite as a $R_{\mathrm{st}}$-module.
\end{lemma}
\begin{proof}
For a nonempty binary vector $u$ with $m-u\ge0$, the top-channel exponent is
\begin{align*}
 \Delta&=h(m)-h(u)-h(m-u)\\
 &=2\sum_i(u_i-u_{i+1})\big((m_i-u_i)-(m_{i+1}-u_{i+1})\big)
       +2Nu_0(m_0-u_0).
\end{align*}
Choose $u$ to be a superlevel set of $m$. Each boundary contribution in the sum is $|m_i-m_{i+1}|-1\ge0$. It is positive at a boundary gap of at least two. The last term is positive when $u_0=1$ and $m_0\ge2$. Consequently, whenever
\[
 m\notin\mathcal L_N:=\{m:m_0\le1,\ |m_i-m_{i+1}|\le1\text{ for every }i\},
\]
one can choose such a superlevel set with $\Delta>0$. Both $u$ and $m-u$ are then nonzero.

The coefficient $(v_u)_{(-\Delta-1)}v_{m-u}$ belongs to $C_2(W_N)$. Its projection to $M_m$ is a nonzero multiple of $v_m$: the canonical multiplication has a nonzero top fusion channel, and its highest-input, highest-output coefficient is nonzero. To see the latter assertion directly, pair its coefficients with all descendants of the highest output vector. Virasoro Ward identities determine these pairings recursively from the leading highest-weight coefficient. If that coefficient vanished, nondegeneracy of the irreducible target's contravariant form would make the entire intertwiner vanish. Every other summand in this coefficient has labels coordinatewise at most $m$ and strictly smaller total label. Thus $[v_m]$ reduces to lower-label summands. Their images are stable under all $D_i$ and $x_i$, so the entire image of $M_m$ reduces as well. Induction on $\sum_i m_i$ leaves only $m\in\mathcal L_N$.

This is a finite set, since $m_i\le1+\lfloor N/2\rfloor$. In each surviving module put $r_i=2m_i+1$. The singular vector at site $i$ has level $d_i=r_i r_{i+1}$ and can be normalized to have leading term $(L^{(i)}_{-1})^{d_i}$; see \cite[Proposition 2.2.2]{virasoro}. Modulo $C_2$, it gives a monic relation for $D_i$ with coefficients in $\C[x_i]$. All remaining terms have smaller $D_i$-degree. Commutators with $L^{(i)}_{-2}$ vanish modulo $C_2$. Hence the images
\[
 \prod_i D_i^{a_i}v_m,\qquad m\in\mathcal L_N,\quad 0\le a_i<d_i,
\]
generate $R$ over $R_{\mathrm{st}}$.
\end{proof}

For completeness, a smaller bound for strong generators follows directly from $C_1$. If $\max m_i\ge2$, let $u$ be the indicator of one connected maximal plateau (all sites when $m$ is constant). The same formula gives $\Delta\ge0$, enough to put the coefficient in $C_1$. The zero modes $D_i$ preserve $C_1$, while $L^{(i)}_{-k}$, $k\ge2$, on a positive-weight state belongs to $C_1$. Induction leaves binary $u$ and the vacuum summand. The singular-vector bounds give
\[
 \dim W_N/C_1(W_N)\le N+\sum_{u\in\{0,1\}^N}9^{|u|}=N+10^N.
\]
For a binary summand the highest required weight is
\[
 h(u)+\sum_i(r_ir_{i+1}-1)
 =6|u|+2\sum_i u_i u_{i+1}+2Nu_0\le10N.
\]

\begin{lemma}\label{c:partial-translations}
For every algebra homomorphism $\ell:R\to\C$ and every $v\in W_N$,
\[
 \ell(D_i v)=0\qquad\text{for all }i.
\]
\end{lemma}
\begin{proof}
In characteristic zero, a derivation of a finite commutative $R_{\mathrm{st}}$-algebra that annihilates $R_{\mathrm{st}}$ annihilates its reduced algebra. Indeed derivations preserve the nilradical and each minimal prime. On each integral component, pass to fraction fields; the resulting finite algebraic extension of the fraction field of the image of $R_{\mathrm{st}}$ is separable, so a derivation over that base field is zero. Apply this to $D_i$ using Lemma~\ref{c:stress-finite}.
\end{proof}
Thus separate Virasoro Ward identities can be evaluated at any point
$\ell$ of the $C_2$ algebra, even though the VOA translation operator
is the sum of the $D_i$.

\subsection*{A Laurent-polynomial lemma}
For a positive irrational number $a$, write $h_a=2a-1$ and define
\[
 C_{a,X}(z)=
 \begin{pmatrix}
 0&1&0\\0&0&1\\
 4ah_a^2z^{-3}&4aX+2ah_a z^{-2}&-4az^{-1}
 \end{pmatrix}.
\]
This is the companion connection of
\begin{equation}\label{c:scalar-ode}
 f'''+4az^{-1}f''-(4aX+2ah_a z^{-2})f'-4ah_a^2z^{-3}f=0.
\end{equation}

\begin{lemma}\label{c:laurent-obstruction}
Let $a>b>0$ be irrational. If
\[
 F'=(C_{a,X}\otimes I+I\otimes C_{b,Y})F
\]
has a nonzero vector solution with Laurent-polynomial entries, then $aX=bY$.
\end{lemma}
\begin{proof}
First suppose $X,Y\ne0$ and put $A=aX$, $B'=bY$. If $A\ne B'$, the constant coefficient matrix has a simple zero eigenvalue: its two constituent spectra are $\{0,\pm2\sqrt A\}$ and $\{0,\pm2\sqrt{B'}\}$. The right zero eigenvector is $e_0\otimes e_0$. Comparing the highest Laurent power and projecting the next equation to the left zero eigenspace shows that this power must be zero. The residue matrix annihilates the right zero eigenvector. The formal solution with prescribed constant coefficient is therefore unique: the recursion is invertible off the zero eigenspace, and its scalar equation at order $z^{-n-1}$ has coefficient $-n$ for $n\ge1$.

The normalized scalar solution with zero exponential part is
\[
 f_{a,X}(z)=\sum_{n\ge0}u_nz^{-2n},\qquad
 u_n=\frac{(a)_n(\tfrac12-a)_n(1-2a)_n}{n!A^n}.
\]
Substitution into \eqref{c:scalar-ode} proves the recurrence; explicitly the coefficient factorization is
\[
 -2j(2j+1)(2j+2)+8aj(2j+1)+4ah_aj-4ah_a^2
 =-4(a+j)(2a-2j-1)(2a-j-1).
\]
Consequently the first component of the unique tensor-product solution is $f_{a,X}f_{b,Y}$. None of its Pochhammer parameters is a nonpositive integer, and the gamma-function asymptotic gives
\begin{equation}\label{c:factorial-asymptotic}
 u_n=\frac{(n!)^2 n^{-3/2-2a}}{\Gamma(a)\Gamma(\tfrac12-a)\Gamma(1-2a)A^n}
       (1+O(n^{-1})).
\end{equation}
Let $v_n$ be the analogous coefficients for $b,Y$. Their convolution cannot terminate. The following estimates apply also when $X$ and $Y$ are complex. If $|A|<|B'|$, then
\[
 \sum_{j=0}^n u_{n-j}v_j=u_n(1+o(1));
\]
if $|A|>|B'|$, the analogous statement holds with $v_n$. For the first assertion,
split the sum at $n/2$. After division by $|u_n|$, the terms with $1\le j\le n/2$ are bounded by a constant times $\binom nj^{-2}(|A|/|B'|)^j$, because both power exponents in \eqref{c:factorial-asymptotic} are negative. Their sum is $O(n^{-2})$. The other half has the exponentially small factor $(|A|/|B'|)^{n/2}$ times a fixed power of $n$. Initial finitely many coefficients are absorbed in the constants.

If $|A|=|B'|$, split at $n/2$ and compare each half with its adjacent
endpoint. For $1\le j\le n/2$, the first half is bounded by
$C|v_n|\binom nj^{-2}$ after absorbing the fixed negative powers
of $j$ and the bounded ratio $(n-j)/n$; the other half is bounded
by $C|u_n|\binom nj^{-2}$ after reversing $j$.
Since $\sum_{j=1}^{\lfloor n/2\rfloor}\binom nj^{-2}=O(n^{-2})$,
we obtain
\[
 \sum_{j=0}^n u_jv_{n-j}=u_n+v_n+
 O(n^{-2})(|u_n|+|v_n|).
\]
Since $a>b$, $|u_n/v_n|=O(n^{-2(a-b)})\to0$, so again the convolution is nonzero for all sufficiently large $n$. Thus $f_{a,X}f_{b,Y}$ is not a Laurent polynomial, a contradiction.

If $X=0$ and $Y\ne0$, rescale the first-factor coordinates to $(f,zf',z^2f'')$. Its connection becomes $E_a/z$, where
\[
 E_a=\begin{pmatrix}0&1&0\\0&1&1\\4ah_a^2&2ah_a&2-4a\end{pmatrix},\qquad
 \operatorname{Spec}(E_a)=\{2a,1-2a,2-4a\}.
\]
The second constant matrix $C_0$ has a simple zero eigenspace.
If the leading coefficient is $z^m v$, the highest-power equation
places $v$ in $\C^3\otimes\ker C_0$. Projecting the next equation
to this kernel gives $m v=(E_a\otimes1)v$: the projected residue
of the second factor is zero. Thus the integer $m$ would be an
eigenvalue of $E_a$, impossible because all three are irrational. The case $Y=0$, $X\ne0$ is identical with $b$. This proves the lemma.
\end{proof}

\subsection*{Null vectors transport the stress values around the cycle}
Fix an algebra homomorphism $\ell:R\to\C$ and an index $i$. Let $v=v_{e_i}$, with $(3,1)$ label at site $i$, $(1,3)$ at site $i-1$, and vacuum labels elsewhere. Set
\[
 a=\psi_i>1,\quad b=\psi_{i-1}^{-1}\in(0,1),\quad
 X=\ell(T_i),\quad Y=\ell(T_{i-1}),
\]
and abbreviate the two translations by $D_A,D_B$. Define
\[
 F_{pq}(z)=\ell\big(Y(D_A^pD_B^qv,z)v\big),\qquad 0\le p,q\le2.
\]
Every entry is a Laurent polynomial with nonpositive powers: all positive-power coefficients belong to $C_2$. Moreover $F_{00}\ne0$, since its most singular coefficient is the nonzero vacuum pairing of $v$.

The level-three null relation is
\[
 D_A^3v-4aL^A_{-2}D_Av+2ah_aL^A_{-3}v=0,
\]
and likewise for $b$. For any $D_B$-descendant $u$ of $v$, normal ordering and Lemma~\ref{c:partial-translations} give
\begin{align*}
 \ell(Y(L^A_{-2}u,z)v)
 &=X\ell(Y(u,z)v)-z^{-1}\ell(Y(D_Au,z)v)+h_az^{-2}\ell(Y(u,z)v),\\
 \ell(Y(L^A_{-3}u,z)v)
 &=z^{-2}\ell(Y(D_Au,z)v)-2h_az^{-3}\ell(Y(u,z)v).
\end{align*}
Indeed, on the left only $L^A_{-2}$ survives; on the right only $L^A_{-1}$ and $L^A_0$ act nontrivially on the highest vector, and the $L^A_{-1}$ action is moved to the inserted vector using $\ell D_A=0$. The null relation therefore gives
\[
 F_{3q}=4aXF_{1q}-4az^{-1}F_{2q}+2ah_az^{-2}F_{1q}+4ah_a^2z^{-3}F_{0q},
\]
with the analogous relation for $F_{p3}$. Total translation gives $F'_{pq}=F_{p+1,q}+F_{p,q+1}$. Thus the nine entries satisfy precisely the tensor-sum system of Lemma~\ref{c:laurent-obstruction}. It follows that
\begin{equation}\label{c:stress-transport}
 \psi_i\ell(T_i)=\psi_{i-1}^{-1}\ell(T_{i-1})\qquad\text{for every }i.
\end{equation}

Finally,
\[
 \prod_i\psi_i=\frac{t+N}{t}=t^{-2}.
\]
Iterating \eqref{c:stress-transport} gives $\ell(T_0)=t^4\ell(T_0)$. Since $0<t<1$, every $\ell(T_i)$ is zero. By Lemma~\ref{c:stress-finite}, $R$ is a finitely generated commutative algebra; the Nullstellensatz makes each $x_i$ nilpotent. Thus $R_{\mathrm{st}}$ is finite-dimensional, and the finite $R_{\mathrm{st}}$-module $R$ is finite-dimensional. This proves $C_2$-cofiniteness and completes the proof of
Theorem~\ref{c:cyclic-c2}.


\section{The ambient lattice and its conformal deformation}\label{c:ambient}
Choose orthogonal real currents $A_i$ with
$(A_i,A_j)=2\psi_i\delta_{ij}$ and put
\[
 \beta_i=A_i-\psi_{i-1}^{-1}A_{i-1},\qquad
 L=\bigoplus_{i=0}^{N-1}\mathbb Z\beta_i,\qquad
 B=\sum_i\beta_i.
\]
The cyclic parameter relations give the Gram matrix $2K$, where
\begin{equation}\label{c:lattice-quadratic}
 m^TKm=\sum_i(m_i-m_{i+1})^2+Nm_0^2.
\end{equation}
It is positive definite.  All lattice pairings are even, so we choose
the trivial cocycle for exponential fields.  Also
\[
 B=\sum_i\frac{\psi_i-1}{\psi_i}A_i,
 \qquad (B,B)=2N.
\]
Let $D$ be the lattice translation operator and set
\[
 \varepsilon=e^B,\qquad
 \omega_0=\omega_L+\tfrac12 B(-2)\one
 =\sum_i\left(\frac{A_i(-1)^2}{4\psi_i}
       +\frac{\psi_i-1}{2\psi_i}A_i(-2)\right)\one.
\]

\begin{proposition}[The ambient conformal vector]\label{c:ambient-deformation}
For every $\lambda\in\mathbb C$, the vector
$\omega_\lambda=\omega_0+\lambda D^2\varepsilon$ is a conformal vector
on the simple lattice vertex algebra $V_L$, with central charge $-5N$.
Its translation and grading operators agree with those of $\omega_0$.
The grading is nonnegative, has finite-dimensional homogeneous spaces,
and has weight-zero space $\mathbb C\one\oplus\mathbb C\varepsilon$.
The resulting vertex operator algebra $A_N=(V_L,\omega_\lambda)$ is regular.
\end{proposition}
\begin{proof}
The shift by $B/2$ changes the central charge to
$N-12(B/2,B/2)=-5N$.  The ground weight of charge
$k=\sum_i k_i\beta_i$ is
\begin{equation}\label{c:ambient-weight}
 h_0(k)=k^TKk-Nk_0
 =\sum_i(k_i-k_{i+1})^2+Nk_0(k_0-1).
\end{equation}
For integral $k$ this is a nonnegative even integer and vanishes only
for $k=0$ and $k=(1,\ldots,1)$.  Positive definiteness gives finiteness
at bounded weight.  Thus $\varepsilon$ is a primary of weight zero.
The self-bracket of $\varepsilon$ is zero, since $(B,B)=2N>0$.
Writing $[\omega_0{}_z\varepsilon]=D\varepsilon$, sesquilinearity
and skew-symmetry give
\[
 [\omega_0{}_z D^2\varepsilon]+[D^2\varepsilon{}_z\omega_0]
 =(D+z)^2D\varepsilon-z^2D\varepsilon
 =(D+2z)D^2\varepsilon.
\]
The self-bracket of $D^2\varepsilon$ also vanishes.  Hence
$\omega_\lambda$ satisfies the Virasoro relation with the same central
charge.  The mode formula
$(D^2\varepsilon)_{(n)}=n(n-1)\varepsilon_{(n-2)}$ shows that the
translation and grading operators are unchanged.

The positive-definite lattice VOA is regular
\cite[Theorem 3.16]{c:DLMRegular}.  Weak modules depend only on the
underlying vertex algebra, whose translation operator has not changed.
Every weak module is therefore a direct sum of the standard lattice-coset
modules.  On each such module the new grading is the positive quadratic
lattice energy minus the linear term $(B,k)/2$; it is bounded below,
diagonalizable, and has finite-dimensional eigenspaces.  The
$D^2\varepsilon$ term does not change this grading.  Thus these simple
modules are ordinary for $\omega_\lambda$ and $A_N$ is regular.
\end{proof}


\section{The uniform screening frame and its certificate}\label{c:app:frame}
We construct commuting Virasoro vectors in the ambient lattice algebra.
Each vector has a correction linear in the coupling $u$; their sum
first changes at order $u^N$.

Put $r_i=t+i$, $p_i=(r_i+1)^{-1}$, and
\[
 P_- =\prod_{i=0}^{N-1}r_i,\qquad P_+=\prod_{i=0}^{N-1}(r_i+1),\qquad
 s=(-1)^N,\qquad C_N=\frac{s}{2}(P_-+P_+).
\]
Thus $P_+/P_-=(t+N)/t$. Pairing $r_jr_{N-j}=1+j(N-j)$ gives the useful closed formula
\[
 C_N=-\frac{\sqrt{N^2+4}}2
           \prod_{j=1}^{(N-1)/2}\bigl(1+j(N-j)\bigr)\ne0.
\]
Use the orthogonal Heisenberg currents $A_i=\sqrt{2\psi_i}\,H_i$, with $(A_i,A_j)=2\psi_i\delta_{ij}$. For a charge $m=\sum m_i\beta_i$, set
\[
 p_i(m)=m_i-\frac{m_{i+1}}{\psi_i},\qquad
 h_i(m)=\psi_i p_i(m)^2-(\psi_i-1)p_i(m).
\]
In these coordinates
\[
 T_i=\frac{A_i(-1)^2\one}{4\psi_i}
       +\frac{\psi_i-1}{2\psi_i}A_i(-2)\one,
 \qquad h(m)=\sum_i h_i(m)=\sum_i(m_i-m_{i+1})^2+Nm_0(m_0-1).
\]
All lattice pairings are even, so the exponential fields below can use the trivial lattice cocycle.

\begin{theorem}[Uniform Virasoro frame]\label{c:screening-frame}
For every odd $N\ge3$, there are explicitly defined finite sums $\widetilde T_i(u)\in V_L[u]$ satisfying
\[
 [\widetilde T_i(u){}_{z}\widetilde T_j(u)]
 =\delta_{ij}\big((D+2z)\widetilde T_i(u)+c_i z^3\one/12\big),
 \qquad
 \sum_i\widetilde T_i(u)=\omega_0+C_Nu^N D^2\varepsilon.
\]
They have $\widetilde T_i(0)=T_i$, and their degree-one corrections are
\[
 u\big(e^{\beta_i}-e^{\beta_{i+1}}\big).
\]
Consequently, every nonzero $\lambda\in\C$ admits a conformal embedding
\[
 \bigotimes_i L(c_i,0)\ \hookrightarrow\ (V_L,\omega_0+\lambda D^2\varepsilon).
\]
This gives the Virasoro frame embedding. The extension of this map to
$W_N$ is established in Corollary~\ref{c:cyclic-embedding}.
\end{theorem}

\subsection*{Finite formulas for the frame}
A proper cyclic interval is an oriented consecutive set $I=[a,\ldots,b]$ of length $\ell<N$. Also write $I$ for its binary charge vector. If $k\in I$, define
\[
 F_{I,k}=(-1)^{\ell-1}
   \prod_{\substack{j\in I\\j<_I k}}r_j
   \prod_{\substack{j\in I\\k\le_I j<_I b}}(r_j+1),
 \qquad c_{i,I}=F_{I,i}-F_{I,i+1},
\]
where $F_{I,k}=0$ for $k\notin I$. Empty products are one and site indices in $c_{i,I}$ are cyclic. In particular $c_{i,\{a\}}=\delta_{ia}-\delta_{i,a-1}$ and $\sum_i c_{i,I}=0$.

For the charge-$B$ term write $a_j=A_j(-1)$, $b_j=A_j(-2)$ and define the quadratic oscillator polynomial $S_i$ by its coefficients:
\[
\begin{array}{c|c}
\text{monomial}&\text{coefficient in }S_i\\\hline
b_i&C_Np_i\\
b_j\ (j\ne i)&0\\
a_i^2&\dfrac{s}{2}\big(P_+p_i-P_-r_ip_i^2\big)\\[3pt]
a_j^2\ (j<i)&-\dfrac{sP_+r_jp_j}{2r_i(r_i+1)}\\[3pt]
a_j^2\ (j>i)&-\dfrac{sP_-r_jp_j}{2r_i(r_i+1)}\\[3pt]
a_ia_j\ (j<i)&sP_+p_ip_j\\
a_ia_j\ (j>i)&sP_-p_ip_j\\
a_ja_k\ (j\ne k,\ i\notin\{j,k\})&0
\end{array}
\]
The order $<$ in this table is the ordinary order $0<1<\cdots<N-1$. Telescoping $1/[r_i(r_i+1)]=1/r_i-1/(r_i+1)$ gives
\begin{equation}\label{c:screening-total}
 \sum_i S_i\varepsilon=C_ND^2\varepsilon.
\end{equation}

There are further scalar terms $d_{i,I}e^{B+I}$ for intervals $I\subseteq\{1,\ldots,N-1\}$. Define them in increasing order of $|I|$. In the already constructed partial frame at $u=1$, let $q_{i,B+I}$ be the coefficient of $z e^{B+I}$ in
\[
 [\widetilde T_i{}_{z}\widetilde T_i]-(D+2z)\widetilde T_i-c_i z^3\one/12.
\]
Set
\begin{equation}\label{c:screening-tail}
 d_{i,I}=-\frac{q_{i,B+I}}{2(h_i(B+I)-1)}.
\end{equation}
This is a finite, explicit recursion in lattice OPE coefficients. Its denominators never vanish, as checked below. The complete formula is
\begin{equation}\label{c:screening-full}
 \widetilde T_i(u)=T_i+
 \sum_{\text{proper cyclic }I}u^{|I|}c_{i,I}e^I
 +u^NS_i\varepsilon+
 \sum_{I\subseteq\{1,\ldots,N-1\}\text{ interval}}
       u^{N+|I|}d_{i,I}e^{B+I}.
\end{equation}

\subsection*{Proof and the complete charge bound}
If a nonnegative charge has weight at most three, its weight is zero or two: the displayed formula for $h(m)$ is an even nonnegative integer. If $m_0\ge2$, then $h(m)\ge2N\ge6$. If $m_0=0$, squared differences totaling two force a binary interval avoiding zero. If $m_0=1$, the possibilities are $B$, a proper binary cyclic interval containing zero, or $B+I$ with $I$ avoiding zero. Weight zero gives precisely $0,B$.

For the singular OPE of two weight-two fields, the coefficient of $z^k$ has total weight $3-k\le3$. The preceding list therefore exhausts every positive charge at which a frame equation can fail. Hence the finite ansatz \eqref{c:screening-full} contains every
possible correction.

For proper intervals, the only nonlinear contributions come from decompositions $I=J\sqcup K$ into adjacent intervals. For $v,w$ in the oriented interval, direct product telescoping gives, when $v<_I w$,
\[
 \sum_{J|K}F_{J,v}F_{K,w}=F_{I,w}-F_{I,v}.
\]
The sum is over cuts separating $v$ and $w$. It is the elementary identity for $\prod(r_j+1)-\prod r_j$, expanded by the place at which the factor is changed. The version with an extra momentum factor is
\[
 \sum_{J|K}F_{J,v}F_{K,w}p_l(J)=
 \begin{cases}
 p_l(I)(F_{I,w}-F_{I,v}),&l<_I v,\\
 -F_{I,v}/(r_l+1),&v\le_I l<_I w,\\
 0,&l\ge_I w.
 \end{cases}
\]
Here sites are ordered along the interval with its left boundary; outside its support and boundary the expression is zero. For the reversed ordered pair, use $p_l(K)=p_l(I)-p_l(J)$.

Taking discrete differences in $v,w$ produces exactly the two required coefficient identities
\begin{align*}
 \sum_{J|K}(c_{i,J}c_{j,K}+c_{i,K}c_{j,J})
 &=-h_i(I)c_{j,I}-h_j(I)c_{i,I}+2\delta_{ij}c_{i,I},\\
 \sum_{J|K}\big(c_{i,J}c_{j,K}p_l(J)+c_{i,K}c_{j,J}p_l(K)\big)
 &=-c_{j,I}p_i(I)\delta_{li}+c_{i,I}p_j(I)\delta_{lj}\\
 &\quad-c_{i,I}h_j(I)p_l(I)+\delta_{ij}c_{i,I}p_l(I).
\end{align*}
They are respectively the coefficients of $z e^I$ and $A_l(-1)e^I$ in the frame equations. The singleton equations hold directly, and induction proves every proper-interval equation.

At charge $B$, the possible residuals are cubic oscillator polynomials
and their $z$-multiples. Verification that the displayed $S_i$ cancel
these residuals is the computer-assisted step in the proof. The complete
program in Appendix~\ref{app:certificate} checks the resulting rational
identities with $N,t$ and the site positions left indeterminate. The
finite reduction, described below, has 172 coefficient types.
Equation~\eqref{c:screening-total} follows separately by telescoping.

It remains to prove that the recursion \eqref{c:screening-tail} solves every equation at the remaining charges. For $\mu=B+[a,b]$, with $1\le a\le b<N$, its nonzero individual weights are
\[
 h_{a-1}(\mu)=\frac{r_{a-1}-1}{r_{a-1}+1},\qquad
 h_i(\mu)=\frac{2}{r_i(r_i+1)}\ (a\le i<b),\qquad
 h_b(\mu)=1+\frac2{r_b}.
\]
All other weights are zero. None equals one, since $t$ is irrational and positive. The momenta are also nonzero: outside the interval they equal
$1/(r_i+1)$, in its interior they equal $2/(r_i+1)$, and at its
left and right boundaries they equal $(1-r_i)/(r_i+1)$ and
$(r_i+2)/(r_i+1)$, respectively. Irrationality excludes $r_i=1$.

\begin{lemma}[Automatic mixed equations]\label{c:automatic-mixed}
At a first unresolved charge $\mu$ of minimum weight two, suppose $h_i(\mu)\ne1$ and $p_i(\mu)\ne0$ for all $i$. If the self-Virasoro equations have been repaired at that charge, all mixed equations vanish there for $N\ge3$.
\end{lemma}
\begin{proof}
For $i\ne j$, the mixed residual has the form
\[
 [\widetilde T_i{}_{z}\widetilde T_j]_\mu
   =\left(\sum_k v_kA_k(-1)+zw\right)e^\mu.
\]
Apply $[L_0^i,L_{-1}^i]=L_{-1}^i$ to $\widetilde T_j$ and take charge $\mu$. All contributions involving earlier mixed residuals vanish. The original Fock action then gives
\[
 (h_i(\mu)+\delta_{ki}-1)v_k=w p_i(\mu)\delta_{ki}.
\]
Thus $v$ is supported at $i$. Skew-symmetry says that the reversed constant coefficient is $wp(\mu)-v$; the same argument shows that it is supported at $j$. A third site $k$ has $p_k(\mu)\ne0$, forcing $w=0$, then $v=0$.
\end{proof}

Adding $d e^\mu$ changes a self-equation's $z e^\mu$ coefficient by $2(h_i(\mu)-1)d$. Therefore \eqref{c:screening-tail} kills that coefficient. Skew-symmetry kills the constant coefficient too: any self-bracket at this charge has constant term equal to one half the derivative of its $z$ coefficient. Higher powers of $z$ are impossible by weight. The lemma now kills every mixed equation, completing induction over the tail charges.

Finally the sum of the repaired frame is Virasoro. Its proper-interval corrections sum to zero, and its charge-$B$ correction is \eqref{c:screening-total}. Compare it with the Virasoro vector $\omega_0+C_ND^2\varepsilon$ from Proposition~\ref{c:ambient-deformation}. A first nonzero tail difference would be a scalar $d e^\mu$ of total weight two. Its self-equation has linear coefficient $2d z e^\mu$, so $d=0$. Hence all tail coefficients sum to zero. Charge homogeneity reinstates the powers of $u$ in \eqref{c:screening-full} and proves the theorem. Each parameter $\psi_i$ is irrational, so the corresponding Virasoro vacuum algebra is simple; the nonzero homomorphism from their simple tensor product is injective.

\subsection*{Reduction to finitely many symbolic identities}
Every complementary pair of proper intervals consists uniquely of $I=[a,b]\subseteq\{1,\ldots,N-1\}$ and its cyclic complement $J$. Put
\[
 m_k=\mathbf1_{a\le k\le b},\quad H_k=\mathbf1_{b<k},\quad
 R=\frac{t+N}{t},\quad G=sP_-.
\]
To see precisely which factors are extracted, put
\[
 \alpha_I=(-1)^{|I|-1}\prod_{j\in I}r_j,\qquad
 \alpha_J=(-1)^{|J|-1}\prod_{j\in J}r_j.
\]
The identities $r_j+1=r_{j+1}$ give
$F_{I,k}=\alpha_I m_k/r_k$ and
$F_{J,k}=\alpha_J[1-m_k+(R-1)H_k]/r_k$.
For the second formula the factor $R$ appears exactly when the
product defining $F_{J,k}$ crosses the closing edge. In particular
$\alpha_I\alpha_J=(-1)^{N-2}P_-=G$.
Thus $c_{i,I}=\alpha_Iu_i$, $c_{i,J}=\alpha_Jv_i$, and the
normalized coefficients and momenta are
\begin{align*}
 u_i&=\frac{m_i}{r_i}-\frac{m_{i+1}}{r_i+1},&
 v_i&=\frac{g_i}{r_i}-\frac{g_{i+1}}{r_i+1},&
 g_i&=1-m_i+(R-1)H_i,\\
 p_i(I)&=m_i-\frac{r_i}{r_i+1}m_{i+1},&
 p_i(J)&=p_i-p_i(I).
\end{align*}
In these formulas use $m_0=m_N=0$, $H_0=0$, $H_N=1$. For a Boolean monomial $\prod_{k\in M}m_k\prod_{l\in H}H_l$, summation over all $1\le a\le b<N$ gives exactly
\[
 \begin{cases}
 (\min M)\max(0,U-\max M),&M\ne\varnothing,\\
 U(U-1)/2,&M=\varnothing,
 \end{cases}
 \qquad U=\min(H\cup\{N\}).
\]
Indeed, for $M\ne\varnothing$ the allowed choices are
$1\le a\le\min M$ and $\max M\le b<U$; for $M=\varnothing$
there are $b$ choices of $a$ for each $1\le b<U$.
Boolean products are reduced using $m_k^2=m_k$, $H_k^2=H_k$.
Successor sites are represented as $j+1$ before counting. When
$j+1$ coincides with another site, the same minimum and maximum
values give the same count; an empty endpoint range gives zero.
At the closing edge $g_N=R$, so
$g_N/(t+N)=g_0/t=1/t$. Thus the successor convention gives the
cyclic coefficient at site zero exactly. The factors at $m_0,m_N$
and $H_0,H_N$ are also enforced by the endpoint count.
This replaces every rank-dependent interval sum by a polynomial
on each equality and order pattern. The charge-$B$ identities
hold with $N,t$ independent; the relation $t(t+N)=1$ is still
used in the preceding lattice and charge-bound arguments.

The complementary exponential bracket is
\[
 [e^I{}_{z}e^J]=
 \left(\frac{P_1^3}{6}+\frac{P_1P_2}{2}+\frac{P_3}{3}
       +\frac z2(P_1^2+P_2)+\frac{z^2}2P_1+\frac{z^3}{6}\right)\varepsilon,
 \quad P_k=\sum_l p_l(I)A_l(-k).
\]
Thus a coefficient involves at most the two frame sites and three oscillator sites. The script enumerates every equality and order type of at most five distinct sites, keeping their positions $j_0<\cdots<j_{m-1}$ symbolic. The seven shapes are $z^3$, $z^2A(-1)$, $zA(-2)$, $zA(-1)^2$, $A(-3)$, $A(-2)A(-1)$ and $A(-1)^3$. There are $7,40,69,46,10$ types with one through five distinct sites,
respectively, for a total of 172. For each type, the program computes
the source and the linear correction, clears denominators, and expands
the numerator of their sum. The claimed identity is precisely that
this numerator is zero. The complete Python/SymPy implementation is
printed in Appendix~\ref{app:certificate}; it makes no numerical
substitution for $N$, $t$, or the site positions. This certifies the
finite algebraic calculation, subject to the reduction proved above;
it is not a proof-assistant verification of the theorem.

For clarity, the linear correction is computed with the undeformed
Fock modes
\[
 L_i(n)=\frac1{4\psi_i}\sum_{k\in\Z}:A_i(k)A_i(n-k):
 -\frac{(n+1)(\psi_i-1)}{2\psi_i}A_i(n),
\]
where $A_i(0)e^B=2\psi_ip_i e^B$ and
$A_i(k)=2\psi_i k\,\partial/\partial A_i(-k)$ for $k>0$.
The first linear bracket is
$\sum_{k=0}^3 z^k L_i(k-1)(S_je^B)/k!$; the reversed one is
obtained by skew-symmetry, replacing $z$ by $-z-D$ and negating.
These formulas explain the mode action in the certificate.
As an independent finite cross-check, the supplied program
\nolinkurl{direct_charge_B_n3.py} constructs these normally ordered
modes and all six proper cyclic intervals directly over
$\Q(\sqrt{13})$. For $N=3$ it verifies all nine ordered frame pairs,
comprising 315 coefficient slots, and detects a perturbation of the
$A_0(-2)e^B$ coefficient. This check does not use the Boolean
interval encoding and does not replace the uniform argument.

\subsection*{A nondegenerate invariant form on the ambient algebra}
\begin{proposition}\label{c:screening-form}
For $\lambda\ne0$, $(V_L,\omega_\lambda)$ is regular and has a unique normalized nondegenerate invariant bilinear form. Its weight-zero space is two-dimensional, so it is not of CFT type.
\end{proposition}
\begin{proof}
The total grading is unchanged by the $D^2\varepsilon$ correction. It has
\[
 (V_L)_0=\C\one\oplus\C\varepsilon,\qquad
 (V_L)_1=\bigoplus_i\big(\C A_i(-1)\one\oplus\C A_i(-1)\varepsilon\big).
\]
The degree-one-to-zero action is
\[
 L_1^{(\lambda)}A_i(-1)\one
   =-2(\psi_i-1)(\one+2\lambda\varepsilon),\qquad
 L_1^{(\lambda)}A_i(-1)\varepsilon=0.
\]
Therefore $(V_L)_0/L_1(V_L)_1$ is one-dimensional, and the vacuum has nonzero image exactly when $\lambda\ne0$. The invariant-form theorem \cite{c:LiForm} gives a unique form with $\langle\one,\one\rangle=1$. Its radical is an ideal of the underlying simple lattice vertex algebra, hence is zero. On weight zero its matrix is
\[
 \begin{pmatrix}1&-1/(2\lambda)\\-1/(2\lambda)&0\end{pmatrix}.
\]
Regularity follows from Proposition~\ref{c:ambient-deformation}.
\end{proof}


\section{Fock modules and signed Kac labels}\label{c:fock-orientation}
We fix the free-field conventions used below, since the distinction
between a Verma module and its contragredient is essential.
Let $[a_m,a_n]=m\delta_{m+n,0}$ and
\[
 Q=\frac{\psi-1}{\sqrt{2\psi}},\qquad
 T=\tfrac12 a(-1)^2\one+Q a(-2)\one,
 \qquad c=1-12Q^2,
\]
where $\psi>0$ is irrational.  The Fock space $F_\xi$ has ground
momentum $a_0=\xi$, and its highest weight is
$h(\xi)=\xi(\xi-2Q)/2$.
Put
\[
 \alpha_+=\sqrt{2\psi},\qquad \alpha_-=-\sqrt{2/\psi},\qquad
 \xi_{r,s}=\frac{1-r}{2}\alpha_++\frac{1-s}{2}\alpha_-.
\]
The classical free-field determinant formula \cite[Proposition 2.7]{c:TK}
for the ground-vector map $M(c,h(\xi))\longrightarrow F_\xi$ is,
at level $n$ and up to a nonzero basis constant,
\begin{equation}\label{c:fock-determinant}
 \det\Phi_{\xi,n}
 =C_n\prod_{\substack{r,s\ge1\\rs\le n}}
          (\xi-\xi_{r,s})^{p(n-rs)},\qquad C_n\ne0.
\end{equation}
Here $p(n)$ is the partition number. To check conventions against
\cite{c:TK}, set $w=\xi-Q$, $\lambda=-Q$, and
$s_\pm=\alpha_\pm$. Their raising operator $L_n$ is our $L_{-n}$;
their determinant factor
$w+rs_+/2+ss_-/2$ is exactly $\xi-\xi_{r,s}$.
At level one this gives
$L_{-1}|\xi\rangle=\xi a_{-1}|\xi\rangle$.
Fock contragredience is $F_\xi^\vee\simeq F_{2Q-\xi}$; the pairing
is obtained by transposing nonzero oscillator modes as
$a_n\mapsto-a_{-n}$ and the zero mode as $a_0\mapsto2Q-a_0$.

\begin{lemma}[Signed Kac coordinates]\label{c:signed-fock}
For nonzero integers $r,s$ and irrational $\psi$ the following hold.
If $r,s>0$, then $F_{\xi_{r,s}}$ is the contragredient Verma module
at $h_{r,s}$, with unique simple socle $L(r,s)$.
If $r,s<0$, it is the Verma module at $h_{-r,-s}$; its simple socle
is the irreducible partner Verma module and is not $C_1$-cofinite.
If $rs<0$, it is an irreducible Verma module and is not $C_1$-cofinite.
\end{lemma}
\begin{proof}
Irrationality implies that $\xi_{r,s}=\xi_{u,v}$ only when
$(r,s)=(u,v)$.  Thus \eqref{c:fock-determinant} identifies
$F_{\xi_{r,s}}$ with its Verma module whenever at least one of $r,s$
is negative.  Reflection sends $(r,s)$ to $(-r,-s)$, so
contragredience proves the assertion when both are positive.
The generic Verma-module classification gives one primitive singular
vector at level $rs$ for a positive pair and an irreducible partner
Verma module below it.  A mixed-sign pair is not a positive Kac
weight, even after reflection, so its Verma module is irreducible.
The $C_1$-cofinite simple modules are exactly the positive Kac
modules \cite[Theorems 2.2.4 and 3.1.4]{virasoro}, which proves the
remaining assertions.
\end{proof}

At site $i$ of our lattice, the momentum of charge $k$ is
$\xi=\alpha_+k_i+\alpha_-k_{i+1}$, hence its signed Kac coordinates
are $(1-2k_i,1-2k_{i+1})$.  In particular the charges $-m$ and
$\mathbf1+m$ have, respectively, contragredient-Verma and Verma
orientation for the positive pair $(2m_i+1,2m_{i+1}+1)$.


\section{Simple generator modules for the deformed frame}\label{c:sec:uniform-socle}

Retain the lattice $L$, the vector $B$, and the exponential
$\varepsilon=e^B$ from Section~\ref{c:ambient}. Fix a nonzero coupling
in the frame of Theorem~\ref{c:screening-frame}, so that:
\[
 \sum_i\widetilde T_i=\omega_0+\lambda D^2\varepsilon,
 \qquad \lambda\ne0.
\]
Each correction to $T_i$ has a strictly positive charge in the cone
$\sum_i\mathbb Z_{\geq0}\beta_i$.  The ambient normalized invariant form
satisfies $\langle1,\varepsilon\rangle=-1/(2\lambda)$.
We write $L_i(n)$ for the modes of $\widetilde T_i$.

For $m\in\mathbb Z_{\geq0}^N$, set
\[
 r_i=2m_i+1,\quad
 \eta_i=h_i(r_i,r_{i+1}),\quad
 \ell_i=r_ir_{i+1},\quad
 H=\sum_i \eta_i=\sum_i(m_i-m_{i+1})^2+Nm_0(m_0+1).
\]
The two distinguished charges are $\mu_-=-\sum_i m_i\beta_i$ and
$\mu_+=B+\sum_i m_i\beta_i$.  Their undeformed Fock modules are,
respectively, the tensor product of the contragredient Verma modules
$M_i^\vee$ and the tensor product of the Verma modules $M_i$, by Lemma~\ref{c:signed-fock}.
Here $M_i=M(c_i,\eta_i)$ has one primitive singular vector, at level
$\ell_i$, and its simple quotient is $L_i(r_i,r_{i+1})$.

\begin{proposition}[Simple generator socles]\label{c:generator-socles}
For each $m=e_j$ with $0\leq j<N$, the ambient lattice contains exactly one
simple frame submodule of type
\[
 M_m=\bigotimes_i L_i(2m_i+1,2m_{i+1}+1).
\]
Its highest vector can be normalized to have lowest charge $\mu_-$,
with coefficient one at $e^{\mu_-}$.
\end{proposition}

\subsection*{The finite spectral calculation}
For a lattice charge $k=(k_i)$, its undeformed partial highest weights are
\[
 h_i(k)=\psi_i k_i(k_i-1)+\psi_i^{-1}k_{i+1}(k_{i+1}-1)
             -2k_ik_{i+1}+k_i+k_{i+1}.
\]
Put $a_i=1+i(N-i)$, including $a_0=a_N=1$.  Since
\[
 \psi_i=1+\frac{N-i+t}{a_i},\qquad
 \psi_i^{-1}=1-\frac{N-i-1+t}{a_{i+1}},
\]
the condition that every $h_i(k)-\eta_i$ is an integer implies
\begin{equation}\label{c:socle-common-C}
 \frac{k_i(k_i-1)-m_i(m_i+1)}{a_i}=\kappa\quad\text{for every }i.
\end{equation}
This follows by comparing coefficients of the irrational number $t$.
The total Fock ground weight is
\[
 h(k)=\sum_i(k_i-k_{i+1})^2+Nk_0(k_0-1).
\]

For $m=e_j$, $j>0$, all grades needed below are at most $H+3=5$.
Equation \eqref{c:socle-common-C} gives $\kappa=k_0(k_0-1)\geq0$;
if $\kappa>0$, then $h(k)\geq2N>5$.  Hence $\kappa=0$.

For the cyclic elementary generator $m=e_0$, the required grades are
at most $H+3=2N+5<6N$.  Here $\kappa=k_0(k_0-1)-2$.
If $\kappa=-2$, every site $i>0$ would satisfy
$k_i(k_i-1)=-2a_i<0$, which is impossible.  If $\kappa>0$, then
$k_0(k_0-1)\geq6$ contradicts the grade bound.  Thus $\kappa=0$ for the closing generator as well.

Thus throughout these finite ranges $k_i\in\{-m_i,m_i+1\}$.
Write $k_i=(1+\epsilon_i r_i)/2$, $\epsilon_i\in\{\pm1\}$.
Direct substitution gives
\begin{equation}\label{c:socle-domain-wall}
 h_i(k)=\eta_i+
 \begin{cases}0,&\epsilon_i=\epsilon_{i+1},\\
 \ell_i,&\epsilon_i\ne\epsilon_{i+1}.
 \end{cases}
\end{equation}
A nonconstant cyclic sign sequence has at least two domain walls.
Consequently, at the joint weight $\eta=(\eta_i)$, and at each joint weight
$\eta+\ell_i e_i$, the only Fock charge quotients are $\mu_-$ and
$\mu_+$.  Below $\eta$ in any one coordinate there are no quotients.
Oscillators add nonnegative integers separately in each coordinate, so
they do not change these conclusions.

\subsection*{Spectral projections, without a semisimplicity assumption}
All arguments take place in finite-dimensional total-weight spaces.
The commuting operators $L_i(0)$ are triangular for the positive-charge
filtration; their associated graded operators are the undeformed Fock
$L_i(0)$ operators.  Their joint generalized spectral projections are
polynomials in these commuting operators and preserve the filtration.
One can equivalently choose a linear combination separating all the
finitely many joint eigenvalue tuples and use its ordinary spectral
projections.  This construction does not require the individual $L_i(0)$ to be
semisimple.

The joint generalized eigenspace at $\eta$ has dimension two.  Its
intersection with the cone-filtered subspace supported at charges
$\mu_++\sum_i\mathbb Z_{\geq0}\beta_i$ has dimension one.
To justify the dimension count, restrict the polynomial spectral
projector to this invariant filtered subspace. Its associated graded
image is the sum of the Fock eigenspaces with tuple $\eta$ in that
cone. Of the two possible charges, only $\mu_+$ belongs to the
cone, and only its ground vector has that tuple. Thus the image
has dimension one. The commuting nilpotent parts vanish on a
one-dimensional space. It is therefore spanned by a genuine joint eigenvector $e$ with lowest charge
$\mu_+$ and coefficient one at $e^{\mu_+}$.  Projecting $e^{\mu_-}$ gives a vector $f$ with
lowest charge $\mu_-$ and coefficient one there.  Thus
\[
 L_i(0)e=\eta_i e,\qquad L_i(0)f=\eta_i f+b_i e.
\]
The absence of spectral targets below $\eta$ implies $L_i(n)e=L_i(n)f=0$
for all $n>0$.  The submodule $E$ generated by $e$ is the tensor product
of the Verma modules $M_i$: its map from that tensor product is
injective, since its lowest-charge projection is exactly the
isomorphism to the positive Fock module at $\mu_+$.

Let $S_i$ be the primitive singular operator in $M_i$, with coefficients
fixed at $\eta_i$.  The leading $\mu_-$ part of $S_i f$ vanishes, since
that Fock module is contragredient Verma.  The spectral calculation at
$\eta+\ell_i e_i$ therefore places $S_i f$ in $E$.  In more detail, inside
the cone-filtered subspace above $\mu_-$, the kernel of projection to
the $\mu_-$ Fock quotient has generalized-block dimension
$p(\ell_i)$, supplied entirely by the $\mu_+$ Fock quotient; the
corresponding ordinary weight space of $E$ has that same dimension.
They are equal.  This remains valid before the coefficients $b_i$ have
been shown to vanish.

\subsection*{Eliminating the possible Jordan coupling}
Each frame vector is quasi-primary for the total conformal vector, so
the ambient invariant form satisfies
$\langle L_i(n)x,y\rangle=\langle x,L_i(-n)y\rangle$.
The module $E$ is isotropic.  Indeed $e$ is a total primary of weight
$H$, and every charge in it is at least $\mu_+\geq B$.  Therefore the
weight-zero product $e_{(2H-1)}e$ cannot have charge $0$ or $B$, the only
weight-zero lattice charges.  This proves $\langle e,e\rangle=0$,
and contravariance proves $\langle E,E\rangle=0$.

On the other hand,
\begin{equation}\label{c:socle-cross-pair}
 \langle e,f\rangle=(-1)^H\langle1,e_{(2H-1)}f\rangle
                  =-\frac1{2\lambda}\ne0.
\end{equation}
Here $H$ is even, the lowest-charge product is
$e^{\mu_+}_{(2H-1)}e^{\mu_-}=\varepsilon$, and every correction has
strictly larger charge and hence cannot contribute to weight zero.

Let $P_i(z)$ be the Shapovalov norm polynomial of the fixed vector
$S_i$: on a generalized highest vector, $S_i^\dagger S_i$ acts as
$P_i(L_i(0))$.  The Kac determinant has a simple zero at $\eta_i$ in
level $\ell_i$, because $\psi_i$ is irrational and the only positive
Kac pair with this highest weight is $(r_i,r_{i+1})$.
Thus the level-$\ell_i$ Gram matrix has a one-dimensional kernel and
\[
 P_i(\eta_i)=0,\qquad P_i'(\eta_i)\ne0.
\]
The last assertion follows directly by differentiating the determinant:
for a symmetric matrix with a one-dimensional kernel, its adjugate is
a nonzero scalar times the outer product of a kernel vector with
itself.  A simple determinant zero is therefore equivalent to a
nonzero derivative of that vector's norm.

Since $S_i f\in E$, isotropy and contravariance now give
\[
 0=\langle S_i f,S_i f\rangle
   =\langle f,P_i(L_i(0))f\rangle
   =b_iP_i'(\eta_i)\langle f,e\rangle.
\]
Equation \eqref{c:socle-cross-pair} proves $b_i=0$ for every site.

\subsection*{The common singular-vector kernel}
Now $f$ is a genuine joint highest vector.  Thus $S_i f$ is a joint
singular vector in $E$ at weight $\eta+\ell_i e_i$, and
\[
 S_i f=c_i S_i e
\]
for some scalar $c_i$.  Operators belonging to different Virasoro
factors commute, and $E$ is a tensor product of Verma modules.  Hence
\[
 (c_i-c_j)S_iS_j e=0,\qquad S_iS_j e\ne0,
\]
so all $c_i$ have a common value $c$.  The vector $v=f-ce$ is killed
by every primitive singular operator.  It therefore generates the
simple tensor product $M_m$.  Its leading negative exponential has
coefficient one.  No other such simple submodule can occur in this
two-dimensional bottom space, since a kernel vector in $E$ would
contradict its Verma structure.  This proves the proposition.


\section{Closure and classification of the frame socle}

We first show that the sum of simple $C_1$-cofinite frame submodules
is closed under vertex products. The category of lower-bounded
$C_1$-cofinite modules for $\mathcal B=\bigotimes_iL(c_i,0)$ is the
Deligne product of the generic Virasoro $C_1$ categories, and is
therefore semisimple. This follows by applying McRae's tensor-product
theorem inductively \cite[Theorems 3.7 and 4.13]{c:McRaeDeligne}.
The required closure under contragredients holds in each generic
Virasoro factor by \cite{virasoro}.

For two simple $C_1$-cofinite frame submodules $M,N$ of the ambient
algebra, let $P$ be the span of the coefficients of $Y(M,z)N$.
The intertwining-operator Jacobi identity makes $P$ a frame submodule.
It inherits a lower-bounded, finite-dimensional total-weight grading
from the lattice.  The intertwining operator with target $P$ is
surjective by this definition, so Miyamoto's image theorem, in the form of \cite[Corollary 2.14]{cmy}, makes
$P$ $C_1$-cofinite.  The preceding tensor-product theorem then makes
$P$ a finite direct sum of simple $C_1$ frame modules.  Thus the sum
of all such simple submodules is closed under vertex products.  This
argument does not assume that the full ambient lattice module is
locally $C_1$-cofinite.

\begin{lemma}[Aligned socle and multiplicity]\label{c:aligned-C1-socle}
Every simple $C_1$-cofinite frame submodule of the ambient lattice is
of type $M_m$, for a unique $m\in\mathbb Z_{\geq0}^N$, and its
highest vector has a nonzero component at $e^{-\sum_i m_i\beta_i}$.
Each such type occurs with multiplicity at most one.
\end{lemma}
\begin{proof}
Let $v$ be the highest vector of a simple $C_1$ frame submodule $M$.
Filter by the scalar charge height $d(k)=\sum_i k_i$.
The vector $v$ has finite charge support; let $d_0$ be its minimum
height.  Every deformed frame mode preserves or strictly increases
this height, since every correction charge lies in the positive
cone.  The same lower bound therefore holds on the module generated
by $v$.  For any charge $k$ at height $d_0$ with nonzero component in
$v$, projection to that charge is a nonzero module homomorphism
\[
 M\longrightarrow F(k)=\bigotimes_iF_i(k).
\]
Indeed a correction could contribute at $k$ only from a charge of
height strictly less than $d_0$, which is absent.  Simplicity of $M$
makes this map injective.

Here the two signed Kac coordinates of $F_i(k)$ are
\[
 R_i=1-2k_i,\qquad S_i=1-2k_{i+1}.
\]
They are nonzero odd integers.  Lemma~\ref{c:signed-fock} gives the following three possibilities.  If
$R_i,S_i>0$, then $F_i(k)$ is a contragredient Verma module, whose
unique simple socle is $L_i(R_i,S_i)$ and is $C_1$-cofinite.  If
$R_i,S_i<0$, it is the Verma module at the reflected positive Kac
label, whose simple socle is the irreducible partner Verma module;
that partner is not $C_1$-cofinite.  If the signs are opposite, the
Fock module itself is an irreducible Verma module and is not
$C_1$-cofinite.  Irrationality of $\psi_i$ excludes any additional
positive Kac label in the last two cases.

By the tensor-product description above, $M$ is an external tensor
product of simple $C_1$ Virasoro modules.  Its nonzero embedding in
$F(k)$ forces such a simple socle in every factor.  This can also be
seen directly on its highest vector: the highest-weight and
primitive-null equations act separately on the finite-dimensional
site weight spaces, so their joint solution space is the tensor
product of the corresponding site solution spaces.  Consequently
$R_i,S_i>0$ at every site, equivalently $k_i\leq0$ for every $i$.
Writing $m=-k$, the embedded socle is exactly $M_m$.  Its highest
vector in $F(k)$ is the Fock ground exponential, so the coefficient
at $e^{-m}$ is nonzero.

The label uniquely determines $m$: equality of two positive Kac
weights at irrational $\psi_i$ forces equality of their positive
Kac pairs.  Evaluation of the highest-vector coefficient at this
one ground exponential is therefore an injective linear map
\[
 \operatorname{Hom}_{\mathcal B}(M_m,V_L)\longrightarrow\mathbb C.
\]
If a nonzero linear combination had coefficient zero, applying the
preceding argument to its nonzero simple image would require that
same coefficient to be nonzero.  This contradiction proves the
multiplicity assertion.
\end{proof}


\section{Nonzero self-couplings of all elementary generators}

The elementary generators have an explicit nonzero self-channel. The calculation is uniform for all edges $j$, including $j=0$.

Set
\[
 a=\psi_j,\quad b=\psi_{j-1}^{-1},\quad k=a+b,
 \qquad A=2a,\quad B'=2b,\quad H=2k-2.
\]
For an ordinary edge $k=2$; for the cyclic edge $k=N+2$.  The symbol
$B'$ in this section is a scalar and is not the lattice vector $B$.
Use the two orthogonal currents $X=A_j$ and $Z=-bA_{j-1}$; their
norms are $A$ and $B'$ and $\beta_j=X+Z$.  The corresponding
exponential $E=e^{\beta_j}$ has individual frame weights $(1,1)$.
The charge-$\beta_j$ frame corrections are exactly $+uE$ and $-uE$.

Define the Schur states by
\[
 \sum_{p\geq0}S_p(X)z^p
 =\exp\left(\sum_{n\geq1}\frac{X(-n)}n z^n\right),
\]
and similarly for $Z$.  The normalized simple generator $v_j$ has
lowest term $e^{-\beta_j}$.  At charge zero its state is exactly
\begin{equation}\label{c:edge-neutral-state}
 uP,\qquad
 P=\sum_{p=0}^H\frac{S_p(X)S_{H-p}(Z)}{A-1-p}.
\end{equation}
Indeed, $E_{(1)}e^{-\beta_j}$ is the Schur state of total degree
$H=\beta_j^2-2$, and the original partial $L_j(0)$ acts by $p$ on
its summand of $X$-degree $p$.  The eigenvalue equation at charge
zero therefore gives \eqref{c:edge-neutral-state}.  Its denominators
are nonzero because $a$ is irrational.  No other charge in the
global frame can enter this equation: both $v_j$ and the frame
corrections are supported in the nonnegative cone above their
respective lowest charges.

Let $\sigma_j$ denote the coefficient of the normalized highest
vector $v_j$ in $v_{j(H-1)}v_j$.  The leading charge $-\beta_j$
can come only from the ordered charge pairs $(-\beta_j,0)$ and
$(0,-\beta_j)$.  Contracting $X(-n)$ against $e^{-\beta_j}$ replaces
it by $(-1)^n A$, and similarly replaces $Z(-n)$ by $(-1)^nB'$.
Thus the Schur states evaluate to
$(-1)^p(A)_p/p!$ and $(-1)^{H-p}(B')_{H-p}/(H-p)!$.
Since $H$ is even, their total sign is positive.  Skew-symmetry gives
the same coefficient for the reversed charge pair; derivative terms
cannot change the leading highest vector at its minimum weight.
Consequently
\begin{equation}\label{c:edge-sigma-sum}
 \frac{\sigma_j}{u}
 =2\sum_{p=0}^H
 \frac{(A)_p(B')_{H-p}}{p!(H-p)!(A-1-p)},
 \qquad A+B'=H+2.
\end{equation}
Vacuum descendants have only nonnegative charge and cannot contribute
to this leading negative exponential.  The next possible diagonal
channel, with edge spin two, has weight $6k-4>H$, so the coefficient
in \eqref{c:edge-sigma-sum} is indeed the self-channel scalar.

\begin{lemma}[An elementary finite-sum identity]
For every integer $H\geq0$, as an identity of rational functions of
$A$,
\begin{equation}\label{c:edge-sigma-product}
 2\sum_{p=0}^H
 \frac{(A)_p(H+2-A)_{H-p}}{p!(H-p)!(A-1-p)}
 =
 2^{H+1}\frac{\displaystyle\prod_{j=1}^H(2A-2j-1)}
                  {\displaystyle\prod_{j=1}^{H+1}(A-j)}.
\end{equation}
\end{lemma}
\begin{proof}
The only possible finite poles on the left are simple poles at
$A=p+1$, $0\leq p\leq H$.  The residue there is
\[
 2\binom{2p}{p}\binom{2H-2p}{H-p}.
\]
The right side has the same residue: evaluating its numerator at
$A=p+1$ gives the product of the odd integers up to $2p-1$ and,
with sign $(-1)^{H-p}$, the odd integers up to $2(H-p)-1$;
the derivative of its denominator is
$(-1)^{H-p}p!(H-p)!$.  The double-factorial formula gives precisely
the displayed central-binomial residue.

It remains to check that the left side is proper.  For $0\leq p\leq H$,
the finite Newton expansion is
\[
 \frac1{A-1-p}
 =\sum_{r=0}^p
 \frac{p(p-1)\cdots(p-r+1)}{(A-1)(A-2)\cdots(A-r-1)}.
\]
Vandermonde convolution gives, for $A+B'=H+2$,
\[
 \sum_{p=0}^H p(p-1)\cdots(p-r+1)
       \frac{(A)_p(B')_{H-p}}{p!(H-p)!}
 =\frac{(A)_r(A+B'+r)_{H-r}}{(H-r)!}.
\]
After substitution, every term has numerator of degree $r$ in $A$
and denominator of degree $r+1$, since $A+B'$ is constant.  The
left side is therefore $O(A^{-1})$.  The right side is also proper.
Their difference is a polynomial vanishing at infinity, hence zero.
\end{proof}

All zeros and poles of \eqref{c:edge-sigma-product} occur at rational
$A$.  Therefore $\sigma_j\ne0$ for every edge in the cyclic
frame and every nonzero common coupling $u$.


\subsection{Exact ordinary-root coefficients}
\begin{proposition}\label{c:ordinary-root-coefficients}
Let $1\le j<N$, put $h=2\psi_j-1$, and normalize the elementary
highest vector $v_j$ to have coefficient one at $e^{-\beta_j}$.
For the common screening coupling $u\ne0$, its vacuum and self-channel
coefficients are
\begin{align}
 \nu_j&=u^2\frac{3(3h-5)(3h-1)}{h(h-2)(h-1)^2},\label{c:ordinary-root-nu}\\
 \sigma_j&=u\frac{8(2h-1)(2h-3)}{h(h-1)(h-2)}.\label{c:ordinary-root-sigma}
\end{align}
Both coefficients are nonzero.
\end{proposition}
\begin{proof}

We compute the invariant norm of the ordinary generator $v_j$.
Set $a=\psi_j$, $b=\psi_{j-1}^{-1}=2-a$,
$X=A_j$, $Y=A_{j-1}$, and $E_j=e^{\beta_j}$.
Use the common screening coupling $u$, so the charge-$\beta_j$ parts
of the two frame vectors are $+uE_j$ and $-uE_j$.
Only the charges $-\beta_j,0,\beta_j$ in the normalized highest vector
can contribute to the charge-zero vacuum coefficient in its
self-product.  Joint primarity fixes those terms as
\[
 v_j=e^{-\beta_j}+uP_j+u^2c_jE_j+
       \text{terms at other charges},
\]
where
\[
 P_j=\alpha X(-1)^2+\beta X(-1)Y(-1)+\gamma Y(-1)^2
                         +\delta X(-2)+\epsilon Y(-2),
\]
\[
 \alpha=\delta=-\frac1{2(3-2a)},\qquad
 \beta=\frac{b}{2(1-a)},\qquad
 \gamma=\frac{b^2}{2(2a-1)},\qquad
 \epsilon=-\frac{b}{2(2a-1)},
\]
and
\[
 c_j=\frac{10a^2-20a+9}{(a-1)^2(2a-3)(2a-1)}.
\]
For example, the $L_j(0)$ equation at charge zero gives these five
coefficients directly, since $L_j(0)$ has eigenvalues $2,1,0,2,0$
on the five displayed oscillator monomials, and its source is
$\tfrac12((X-bY)(-1)^2+(X-bY)(-2))$.
The $L_j(1)$ and $L_j(2)$ equations then hold identically.

The contractions are $(X,X)=2a$, $(Y,Y)=2/b$, $(X,Y)=0$.
The coefficient $E_{j(1)}P_j$ is
\[
 4a^2\alpha-4a\beta+4\gamma-2a\delta+2\epsilon,
\]
giving the stated $c_j$.  Finally the charge-zero coefficient in
$v_{j(3)}v_j$ is
\begin{align*}
 \nu_j&=u^2\left(2c_j+8a^2\alpha^2+
       \frac{4a}{b}\beta^2+\frac8{b^2}\gamma^2
       -12a\delta^2-\frac{12}{b}\epsilon^2\right)\\
 &=u^2\frac{3(3a-4)(3a-2)}{(a-1)^2(2a-3)(2a-1)}\ne0.
\end{align*}
Every possible zero or pole in the last expression would require a
rational value of $a$, whereas $a=\psi_j$ is irrational.
The $C_1$ image theorem puts this self-product in the frame vacuum
submodule, whose weight-zero space is $\mathbb C1$; there is therefore
no $\varepsilon$ term.  Thus $\nu_j$ is also the invariant norm.
Substituting $a=(h+1)/2$ gives \eqref{c:ordinary-root-nu}.
For an ordinary edge, $H=2$ and $A=2a=h+1$ in
\eqref{c:edge-sigma-product}. Hence
\[
 \frac{\sigma_j}{u}
 =8\frac{(2A-3)(2A-5)}{(A-1)(A-2)(A-3)}
 =\frac{8(2h-1)(2h-3)}{h(h-1)(h-2)},
\]
which is \eqref{c:ordinary-root-sigma}. Irrationality again excludes
all zeros and poles.
\end{proof}


\section{Recognizing the cyclic algebra from its generator socles}

We distinguish three questions: closure of the simple socle, survival of its
generators in a simple quotient, and uniqueness of the resulting diagonal
multiplication.  The following lemmas address the last two questions without
assuming that every fusion-allowed operator product is already nonzero.

\subsection{Object-fixing automorphisms of the generic integer-spin category}

Write $\mathcal R_q$ for the semisimple category of integer-spin type-I
$U_q(\mathfrak {sl}_2)$ modules, with its usual ribbon structure, where
$q$ is not a root of unity.  Its simple objects are $X_m$, $m\geq0$,
with $X_0=\mathbf1$ and $X_1$ the three-dimensional representation.

\begin{lemma}\label{c:object-fixing-trivial}
Every $\mathbb C$-linear braided tensor autoequivalence of
$\mathcal R=\boxtimes_{i=0}^{N-1}\mathcal R_{q_i}$ that fixes the isomorphism
class of every simple object is monoidally naturally isomorphic to the
identity.  The parameters $q_i$ need not be equal.
\end{lemma}
\begin{proof}
First consider one factor.  The generic $SO(3)_q$ category is the additive
idempotent completion of the pivotal category generated by a symmetrically
self-dual object $X=X_1$, its cup and cap, and a rotationally invariant
trivalent vertex $v:\mathbf1\longrightarrow X^{\otimes3}$.
All morphisms are linear combinations of trivalent diagrams.  This is the
standard even Temperley--Lieb realization; see \cite[Proposition 4.3]{c:MPS}.  In their normalization the bigon coefficient is one, so the closed
theta evaluation is $\dim X=q^2+1+q^{-2}\ne0$.

Let $F$ be the autoequivalence.  Since every simple object is fixed and its
twist is scalar, $F$ preserves the twist.  Being braided, it consequently
preserves the associated pivotal structure.  Choose an isomorphism
$F(X)\simeq X$.  After multiplying this isomorphism by a nonzero scalar,
the transported cup equals the original cup; this is possible over
$\mathbb C$ since the cup scales quadratically.  The transported cap then
also agrees, by the snake identities.  The transported trivalent vertex is
$\alpha v$ for some $\alpha\ne0$.  Preservation of the nonzero theta
evaluation gives $\alpha^2=1$.  Replacing the chosen isomorphism
$F(X)\simeq X$ by its negative leaves cup and cap unchanged and changes the
sign of the trivalent vertex.  Thus it can also be arranged that $F(v)=v$.

Tensor the chosen isomorphism on words in $X$.  Because all morphisms between
these words are generated by cup, cap, and $v$, the resulting maps are natural;
by construction they are monoidal.  They extend to finite direct sums and
idempotent summands, proving the claim for one factor.  For a Deligne product,
perform the same normalization in each factor.  All remaining interchanges
between generators from distinct factors are braidings, which $F$ preserves.
The same argument therefore proves naturality and monoidality for the product.
\end{proof}

The lemma controls the tensor structure, including its coherence.
In particular, an object-fixing braided autoequivalence cannot produce
a different diagonal algebra through a residual tensor $2$-cocycle.

\subsection{Recovering an algebra from the mirror functor}

\begin{lemma}\label{c:mirror-recovers-algebra}
Let $\mathcal U,\mathcal V$ be semisimple ribbon categories with simple units,
and let
\[
 A=\bigoplus_{i\in I}U_i\boxtimes V_i
\]
be a simple commutative algebra in their direct-sum completion, with trivial
twist.  Suppose the $U_i$ and the $V_i$ are complete sets of distinct simple
objects, the unit summand occurs once, and $U_0=\mathbf1_{\mathcal U}$, $V_0=\mathbf1_{\mathcal V}$ for the unit index $0$.  Then $A$, including its multiplication, is the canonical diagonal
algebra associated to the braid-reversed equivalence constructed from $A$
in \cite[Proposition 4.4 and Theorem 4.5]{gluing}.
\end{lemma}
\begin{proof}
We reconstruct the multiplication from the mirror equivalence.  Let $F$ be free
$A$-module induction and let $F_{\mathcal U}$ and $F_{\mathcal V}$ denote its
restrictions to the two factors.  The proof of the cited theorem constructs
a tensor category $\mathcal M\subset\operatorname{Rep}A$ in which both
restrictions are tensor equivalences onto $\mathcal M$.  Its unit is $A$.
The central lifts of the two functors identify their braidings oppositely.
If $\tau=F_{\mathcal V}^{-1}F_{\mathcal U}$ is the resulting mirror functor,
then, on transporting $\mathcal M$ to $\mathcal U$, the action of
$\mathcal U\boxtimes\mathcal V$ on its unit is the tensor functor
\[
 (U,V)\longmapsto U\otimes\tau^{-1}(V).
\]
The coherence here is inherited from the central lifts, so it is the same
coherence used for the canonical diagonal algebra.

For finite objects $Z$ of $\mathcal U\boxtimes\mathcal V$, Frobenius
reciprocity gives the natural isomorphism
\[
 \operatorname{Hom}(Z,A)
   \simeq\operatorname{Hom}_{\operatorname{Rep}A}(F(Z),A).
\]
Thus $A$ represents the internal endomorphisms of the unit for that action.
Moreover this identification preserves multiplication.  Indeed, if
$f:Z\to A$ and $g:T\to A$, their free $A$-linear extensions
$\widehat f:F(Z)\to A$ and $\widehat g:F(T)\to A$ have tensor product
\[
 F(Z\otimes T)\simeq F(Z)\otimes_A F(T)
 \xrightarrow{\widehat f\otimes_A\widehat g}
 A\otimes_A A\simeq A.
\]
Under Frobenius reciprocity this is exactly $\mu_A\circ(f\otimes g)$.
The unit is preserved as well.  The representing algebra is therefore the
canonical internal-End algebra of the action, as in
Creutzig--Kanade--McRae, Theorems 3.2--3.5.  Uniqueness of representing
objects proves the algebra assertion.  For infinitely many simple labels, test the representing property on
finite objects $Z$. The argument then identifies every component of
the multiplication in the direct-sum completion.
\end{proof}

\subsection{Recognition from a generating family with nonzero vacuum channels}

Fix an edgewise braid-reversed equivalence between $\mathcal R$ and a ribbon
category $\mathcal D$.  Write $D_m$ for its image of $X_m^*$, so the canonical
algebra is
\[
 A_{\mathrm{can}}=\bigoplus_{m\in\mathbb Z_{\geq0}^N}X_m\boxtimes D_m.
\]
All these labels are self-dual.  We write $M_m=X_m\boxtimes D_m$.

\begin{theorem}[Diagonal recognition]\label{c:diagonal-recognition}
Let $S$ be a connected commutative unital algebra of trivial twist in
$\operatorname{Ind}(\mathcal R\boxtimes\mathcal D)$.  Suppose its underlying
object is a subobject of $A_{\mathrm{can}}$: only the indicated diagonal
simples occur, each with multiplicity at most one.  Suppose a collection
$\{M_g:g\in G\}$ occurs in $S$, its $\mathcal R$-labels tensor-generate
$\mathcal R$, and each vacuum channel
\[
 M_g\otimes M_g\ \xrightarrow{\mu_S}\ S
        \xrightarrow{\pi_0}\mathbf1
\]
is nonzero.  Then $S$ is simple, has every diagonal summand, and is isomorphic
as an algebra to $A_{\mathrm{can}}$.
Moreover, the unital subalgebra generated by the summands $M_g$, $g\in G$,
is all of $S$.

For example, either of the following generating collections suffices:
\[
 G=\{e_1,\ldots,e_{N-1},\mathbf1\},\qquad
 G=\{e_0,\ldots,e_{N-1}\},
\]
where $\mathbf1=(1,\ldots,1)$.
\end{theorem}
\begin{proof}
The ind-completion of a semisimple category is semisimple as an abelian
category.  A proper ideal of $S$ has no unit summand: an ideal containing
the unit is all of $S$.  Because the unit has multiplicity one, the sum of
all proper ideals still has no unit summand.  This sum $J$ is therefore a
proper ideal and is the unique maximal ideal.  Put $Q=S/J$.

No $M_g$ is contained in $J$.  If it were, multiplying it by itself and
using its nonzero vacuum channel would put the unit in $J$.  Hence all the
chosen generators survive in the simple algebra $Q$.  The diagonal decomposition gives the commuting-pair hypothesis:
$\operatorname{Hom}_{\mathcal R}(\mathbf1,X_m)$ and
$\operatorname{Hom}_{\mathcal D}(\mathbf1,D_m)$ vanish for $m\ne0$.  Apply
\cite[Proposition 4.4]{gluing} to $Q$: its simple-label support
is closed under tensor products and taking summands.  Since it contains
the generating family, it is all of $\mathcal R$.  Thus $Q$ has every
diagonal summand.  But $S$ has no off-diagonal summands and has multiplicity
at most one on each diagonal summand.  The quotient $S\twoheadrightarrow Q$
can consequently have no kernel, and $S=Q$ is simple with full support.

Theorem 4.5 of the same paper supplies a mirror equivalence for $S$.
Relative to the fixed edgewise equivalence, it is a braided autoequivalence
of $\mathcal R$ fixing every simple label.  Lemma~\ref{c:object-fixing-trivial}
makes it monoidally naturally isomorphic to the identity.  By
Lemma~\ref{c:mirror-recovers-algebra}, this identifies the entire algebra
structure of $S$ with the canonical diagonal one.

For the stated generating set, $X_{\mathbf1}\otimes
\bigotimes_{j=1}^{N-1}X_{e_j}$ has $X_{e_0}$ as a summand: choose the
vacuum channel in each of the factors $j>0$.  Then the spin-one objects
$X_{e_i}$ generate every integer-spin factor, hence their product category.
This argument takes place in the support of the already simple quotient;
it does not presume that those operator-product channels were nonzero in $S$.
Finally the unital subalgebra generated by the $M_g$ satisfies the same
hypotheses: it is again a multiplicity-free diagonal subobject, and it has
the same nonzero generator vacuum channels.  Applying the proof to it shows
that it too has every diagonal summand, so its inclusion in $S$ is an
isomorphism.
\end{proof}

\paragraph{Necessity of the two-point hypothesis.}
Multiplicity-free diagonal support and associativity alone are not enough.
For example, a Cartan-product degeneration can retain the same diagonal
object while killing lower channels, including a nontrivial generator's
vacuum channel.  The proof above uses precisely the vacuum channels to rule
out such degenerations before applying the mirror theorem.

\subsection{A nonzero self-channel forces the vacuum channel}

\begin{lemma}[The spin-one associativity test]\label{c:spin-one-vacuum}
In the setup of Theorem~\ref{c:diagonal-recognition}, let $M=M_{e_i}$ occur in
$S$.  If the projection of $\mu_S:M\otimes M\to S$ to $M$ is nonzero,
then its projection to the unit is nonzero.  No simplicity assumption on
$S$ and no prior occurrence of $M_{2e_i}$ are needed.
\end{lemma}
\begin{proof}
Only one generic $SO(3)_q$ factor is involved.  Write its spin-one object as
$X$ and its three tensor-square summands as $X_0=\mathbf1,X_1=X,X_2$.
Use the nonzero canonical diagonal multiplication to fix bases for every
allowed diagonal multiplication channel.  Express the multiplication of
$S$ in these bases as scalars $a_{ij}^{\ell}$, assigning scalar zero if the
output summand is absent.  Unit normalization gives $a_{01}^1=a_{10}^1=1$.
Because the reference multiplication is itself commutative, commutativity
of $S$ gives $a_{1j}^1=a_{j1}^1$; this statement includes any braiding
scalars and does not assume that a naked fusion basis has trivial braiding.

Put $T=\operatorname{Hom}(X^{\otimes3},X)$, which has dimension three.
Dual bases in the second factor identify the corresponding nine-dimensional
space of diagonal-category morphisms with $\operatorname{End}(T)$.
The left-associated multiplication of $M^{\otimes3}$ into $M$ is the
diagonal endomorphism
\[
 D=\operatorname{diag}\bigl(a_{11}^0,\ (a_{11}^1)^2,
                                      a_{11}^2a_{21}^1\bigr)
\]
in the left fusion-tree basis.  The right-associated multiplication is the
same diagonal endomorphism in the right fusion-tree basis.  If $F$ is the
spin-one Racah change-of-basis matrix, the second factor changes
by its inverse transpose: the evaluation pairing of the two
fusion spaces must remain the identity. Explicitly, if
$v'_a=\sum_bF_{ba}v_b$, then
$(v'^a)=\sum_b(F^{-1})_{ab}v^b$.
Consequently a diagonal channel tensor
$\sum_a d_av'_a\otimes v'^a$ has matrix $FDF^{-1}$ in the
unprimed bases. Associativity therefore says
\[
 D=FDF^{-1},\qquad FD=DF.
\]
The vacuum-to-spin-one entry $F_{01}$ is nonzero.  For completeness, in the
usual normalized pivotal fusion-tree bases the vacuum row is, up to
nonzero signs,
\[
 F_{0j}=\frac{\sqrt{\dim X_j}}{\dim X},\qquad j=0,1,2.
\]
This follows by closing the vacuum cup-cap pair against the two trivalent
vertices; the remaining closed diagram is the nonzero theta evaluation.
All quantum dimensions in this row are nonzero for generic $q$.
In particular $F_{01}=\pm1/\sqrt{\dim X}\ne0$.
The $(0,1)$ entry of $FD=DF$ now yields
$a_{11}^0=(a_{11}^1)^2$.  The implication in the statement is independent
of all the nonzero choices of channel bases.
\end{proof}


\subsection{Application to the cyclic frame}

Let $\mathcal B=\bigotimes_i L(c_i,0)$ denote the Virasoro frame subalgebra
in the deformed lattice construction, and let
\[
 \mathscr S(V_L)=\sum\{M\subset V_L:
              M\text{ is a simple }C_1\text{-cofinite }\mathcal B\text{-module}\}.
\]
Miyamoto's image theorem (in the form of Creutzig--McRae--Yang,
\emph{Direct limit completions}, Corollary 2.14), together with generic
Virasoro semisimplicity, makes this a vertex subalgebra.  By Lemma~\ref{c:aligned-C1-socle},
its $\mathcal B$-module decomposition is a subobject of
\[
 \bigoplus_{m\geq0}\bigotimes_i
                L_i(2m_i+1,2m_{i+1}+1).
\]
On odd labels the row/column interchange has trivial double braiding, so
the edgewise rearrangement identifies this with the diagonal categorical
setup of Theorem~\ref{c:diagonal-recognition}.

\begin{corollary}[Uniform cyclic embedding]\label{c:cyclic-embedding}
For every odd $N\geq3$ and every nonzero coupling, the cyclic VOA
$W_N$ embeds conformally in the deformed lattice VOA. More precisely:
\[
 \mathscr S(V_L)\simeq W_N
\]
as an extension of $\mathcal B$, with canonical diagonal multiplication.
\end{corollary}
\begin{proof}
The uniform calculation \eqref{c:edge-sigma-product} gives a nonzero
self-channel for every $M_{e_i}$. Lemma~\ref{c:spin-one-vacuum} therefore
gives a nonzero vacuum channel for every elementary generator. These
objects tensor-generate the product of integer-spin categories. Theorem~\ref{c:diagonal-recognition}
applies and identifies the full socle, with its multiplication, with $W_N$.
In particular the frame and the elementary highest vectors $v_{e_i}$
vertex-generate $W_N$, and all fusion-allowed diagonal channels are nonzero.
\end{proof}


\subsection{The orthogonal splitting and a second weight-zero simple module}
\begin{proposition}\label{c:ambient-splitting}
For $\lambda\ne0$ the conformal embedding just constructed gives
a $W_N$-module decomposition
\[
 A_N=W_N\oplus X,\qquad X=W_N^\perp.
\]
The line $X_0$ is spanned by $\zeta=\one+2\lambda\varepsilon$.
There is a simple $W_N$-module $T\not\simeq W_N$ of lowest weight
zero, obtained as a quotient of the cyclic submodule generated by
$\zeta$.  If $C$ denotes the highest vector of the all-ones frame
summand, normalized to have leading charge $e^{-B}$, then
$o(C)\zeta=2\lambda\zeta$.
\end{proposition}
\begin{proof}
The normalized ambient invariant form restricts to a nonzero invariant
form on the simple CFT-type subalgebra $W_N$.  Its restriction is
therefore nondegenerate.  Both algebras have finite-dimensional
homogeneous spaces, so grade-by-grade orthogonal projection gives
the asserted module splitting.  The weight-zero Gram matrix in
Proposition~\ref{c:screening-form} identifies $X_0$ as stated.

The all-ones summand is present by the diagonal recognition theorem.
Lemma~\ref{c:aligned-C1-socle} supplies its normalized lowest charge
$-B$.  Its total highest weight is $2N$.
Consequently the charge-zero coefficient of $o(C)\varepsilon$ is
one, while weight considerations give
$o(C)\varepsilon=\one+b\varepsilon$ for some $b$.
Since $C$ is primary of even weight and $o(C)\one=0$, invariance gives
\[
 0=\langle o(C)\varepsilon,\one\rangle
   =1-\frac{b}{2\lambda}.
\]
Thus $b=2\lambda$ and $o(C)\zeta=2\lambda\zeta$.
The cyclic submodule generated by $\zeta$ has finite length: it is
grading restricted over the CFT-type $C_2$-cofinite algebra $W_N$,
so the finite-category theorem applies \cite{huang}.
Any simple quotient has nonzero image of its cyclic generator and
therefore has lowest weight zero.  Its $C$ zero mode acts there by
$2\lambda\ne0$, whereas it acts by zero on the vacuum top.  The
simple quotient is consequently different from $W_N$.
\end{proof}


\section{Parity simple currents}
\label{sec:parity-currents}
We extend the odd-label gluing to all parity sectors. Continue to write
$\mathcal B=\bigotimes_iL(c_i,0)$ and let $K=K_+$ be the cycle form
from Section~\ref{sec:intro}. Its diagonal entries are
$k_0=N+2$ and $k_i=2$ for $i\ne0$, and its cycle off-diagonal
entries are $-1$. For $N=3$, each distinct off-diagonal entry is
$-1$. The identities needed below are:
\[
 x^TKx=\sum_i(x_i-x_{i+1})^2+Nx_0^2,
 \qquad K\mathbf1=Ne_0,\qquad \det K=N^2.
\]
The canonical odd-label gluing is:
\[
 W=\bigoplus_{r_i\geq1,\ r_i\text{ odd}}
       \bigotimes_iL_i(r_i,r_{i+1}).
\]
Here $W=W_N$ is the extension constructed in
Theorem~\ref{c:cyclic-gluing}.

The ambient category used throughout is
$\mathscr V_{\mathcal B}$, the category of $C_1$-cofinite grading-restricted
generalized $\mathcal B$-modules. The generic Virasoro results
\cite[Theorems 3.1.4,4.2.6,5.1.2,5.5.3]{virasoro}, followed by
\cite[Theorem 1.2, equivalently Theorem 4.13]{c:McRaeDeligne}, show that
$\mathscr V_{\mathcal B}$ is exactly the semisimple rigid Deligne product of the
corresponding generic Virasoro categories; it is closed under contragredients.
Thus \cite[Theorems 7.1,7.5,7.7]{cmy} applies to its direct-limit completion,
to algebra extensions there, and to their local modules. These theorems do
not require $W$ to be rational or $C_2$-cofinite.

\subsection{Construction and unrestricted fusion}
For $\epsilon\in P=(\mathbb Z/2\mathbb Z)^N$, define
\[
 J_\epsilon=\bigoplus_{\substack{r_i\geq1\\
                         r_i\equiv1+\epsilon_i\pmod2}}
                  \bigotimes_iL_i(r_i,r_{i+1}).
\]
Represent the entries of $\epsilon$ by $0,1$ when evaluating integral
quadratic expressions.

\begin{theorem}\label{c:currents}
Each $J_\epsilon$ is a simple local grading-restricted $W$-module, and
\[
 J_\epsilon\bo_W J_\eta\simeq J_{\epsilon+\eta}.
\]
These fusion products have the intertwining-operator universal property
against every grading-restricted generalized $W$-module, not only against
modules whose $\mathcal B$-restriction belongs to $\Ind(\mathscr V_{\mathcal B})$.

The full finite pointed ribbon subcategory $\mathcal P$ on these modules has
self-braiding quadratic form, conformal twist, and categorical dimension
\begin{align*}
 Q(\epsilon)&=\exp\left(\frac{\pi i}{2}\epsilon^TK\epsilon\right),\\
 \theta_\epsilon&=\exp\left(\frac{\pi i}{2}
                    (\epsilon^TK\epsilon+2N\epsilon_0)\right),\\
 d_\epsilon&=(-1)^{\epsilon_0}.
\end{align*}
Its monodromy is $(-1)^{\epsilon^TK\eta}$, hence is nondegenerate.
\end{theorem}

\begin{proof}
We first work in $\Ind(\mathscr V_{\mathcal B})$. The generic first-row category is
braided equivalent to type-I $U_q(\mathfrak{sl}_2)$ modules with
$q=e^{\pi i\psi_i}$; see \cite[Theorem 1.1, Theorems 5.10--5.11]{c:Koshida}.
The first-column category uses $e^{\pi i\psi_i^{-1}}$. At edge $i$, the
parameters of the row at $i$ and column at $i-1$ sum to $k_i$.
Moriwaki's integer parameter shift theorem for type A
\cite[Theorem 2.16]{c:Moriwaki} and quantum-coordinate-ring functor
\cite[Lemma 3.4, Proposition 3.5]{c:Moriwaki} therefore transport the
$\mathbb Z/2$-graded pieces of $\mathcal O_q(SL_2)$ to a lax braided
functor whose source quadratic form has value $i^{k_i}$ on the odd grade.

Reordering the edge factors and multiplying the row and column factors at
each site gives a strong monoidal equivalence onto $\mathscr V_{\mathcal B}$;
the monoidal structure is the usual braiding interchange, and
$L_i(r,1)\bo L_i(1,s)\simeq L_i(r,s)$ by
\cite[Theorem 5.2.5]{virasoro}. This functor is not braided for the naive
product braiding. Indeed, balancing gives
\[
 h_i(r,s)-h_i(r,1)-h_i(1,s)
       =-\frac{(r-1)(s-1)}2,
\]
so the row/column double braiding is $(-1)^{(r-1)(s-1)}$.
One must multiply the source braiding by the bicharacter
\[
 \beta(\epsilon,\eta)=(-1)^{\sum_i\epsilon_{i+1}\eta_i}.
\]
The exponent is bilinear over $\mathbb F_2$, so
$\beta(\epsilon+\eta,\zeta)=\beta(\epsilon,\zeta)\beta(\eta,\zeta)$
and likewise in the second argument. Multiplying a braiding by
$\beta$ therefore preserves both hexagons with the same associator.
Its polarization is
$\beta(\epsilon,\eta)\beta(\eta,\epsilon)
=(-1)^{\sum_i(\epsilon_{i+1}\eta_i+\eta_{i+1}\epsilon_i)}$,
exactly the row--column monodromy introduced by the interchange.
This is the correction in the braided compatibility diagram for
the multiplication functor. Transposing the convention corresponds
to reversing that interchange. The resulting source self-braiding is
\[
 \prod_i i^{k_i\epsilon_i}\,
        (-1)^{\sum_i\epsilon_i\epsilon_{i+1}}
       =\exp\left(\frac{\pi i}{2}\epsilon^TK\epsilon\right).
\]
The zero-grade image is the canonical odd-label algebra $W$, and the
$\epsilon$-grade image is $J_\epsilon$. Normalized abelian cocycles have
trivial components whenever one entry is zero. The lax braided functor
therefore supplies ordinary local $W$-module structures on every
$J_\epsilon$ and balanced multiplication maps between them.

These balanced maps are isomorphisms, not merely nonzero maps. Before the
transports, $\mathcal O_q(SL_2)$ is strongly $\mathbb Z/2$-graded. In its
standard presentation $a,b,c,d$ are odd and $ad-qbc=1$, so
$1\in\mathcal O_1\mathcal O_1$. Strong grading implies
\[
 \mathcal O_a\otimes_{\mathcal O_0}\mathcal O_b
          \xrightarrow{\sim}\mathcal O_{a+b}.
\]
For example a finite decomposition $1=\sum_jx_jy_j$ in odd elements
supplies the inverse to multiplication. The map is equivariant, hence so
is its inverse. Taking $N$ factors and transporting through the grading
twists and strong monoidal equivalences preserves these relative tensor
isomorphisms. Thus $J_\epsilon\bo_WJ_\eta\simeq J_{\epsilon+\eta}$ in
the local direct-limit category. Since $W$ is simple, invertibility also
implies simplicity of all these modules.

The total top weight of a $\mathcal B$-summand is
\[
 \sum_i h_i(r_i,r_{i+1})=\frac{r^TKr-N}{4}.
\]
For $r=\mathbf1+v$, with $v_i\geq0$, this equals
\[
 \frac14\sum_i(v_i-v_{i+1})^2+\frac N4v_0(v_0+2).
\]
Positive definiteness gives a lower bound and finite-dimensional weight
spaces for each parity sum. Reduction modulo integers gives the displayed
conformal twist. Dividing the twist by the self-braiding quadratic form
gives the displayed dimension character. The polarization of $Q$ is
$(-1)^{\epsilon^TK\eta}$; $\det K$ is odd, so it is nondegenerate.

It remains to verify that the calculated fusion products are universal
outside the chosen ambient category. Let $\mathcal Y$ be any $W$-intertwining
operator of type $\binom X{J_\epsilon\ J_\eta}$, where $X$ is a
grading-restricted generalized $W$-module. Restrict to finite
$\mathscr V_{\mathcal B}$-subobjects $U\subset J_\epsilon$ and $V\subset J_\eta$.
By \cite[Corollary 2.14]{cmy}, the image of this $\mathcal B$-intertwining operator
is $C_1$-cofinite and grading restricted. That corollary permits an
arbitrary generalized target; it is the direct-limit version of
Miyamoto's image theorem \cite{c:Miyamoto}. Consequently the full image
$\operatorname{Im}\mathcal Y$ is a union of $\mathscr V_{\mathcal B}$-submodules and
belongs to $\Ind(\mathscr V_{\mathcal B})$. The intertwining Jacobi identity makes it
a $W$-submodule. Factor $\mathcal Y$ through that image and use the
already-established local-category tensor universal property. Composing
with the inclusion into $X$ proves the claimed unrestricted universal
property.
\end{proof}

\begin{remark}[Comparison of braiding conventions]
The braided identification above uses the first-row equivalence of
\cite{c:Koshida}. With the Kac weights fixed here, balancing gives
row--column monodromy $(-1)^{(r-1)(s-1)}$. In particular, it equals
$-1$ when both highest-weight parities are odd. Thus the trivial
row--column braiding asserted in \cite[Proposition 6.3]{c:Koshida}
does not agree with these weights. The monoidal multiplication
functor of its Proposition 6.2 is sufficient for the construction;
the displayed bicharacter supplies the required braiding correction.
\end{remark}

\Needspace{12\baselineskip}
\part{Vacuum extensions and regularity}

\section{Reduction of vacuum extensions to lowest weights zero and one}
\label{neutral:sec:neutral-reduction}

We prove regularity by splitting every extension of the vacuum
by a simple module.  The first step reduces this problem to simple
modules of lowest weight zero or one.

Let $W=W_N$ be the vertex operator algebra of
Theorem~\ref{c:cyclic-gluing}.  We use its conformal frame
$\mathcal B=\bigotimes_{i=0}^{N-1}L(c_i,0)$, its frame decomposition,
and generation by the frame and the elementary summands $R_j=M_{e_j}$;
see Corollary~\ref{c:cyclic-embedding}.
Thus
\[
 h_i(r,s)=\frac{(r\psi_i-s)^2-(\psi_i-1)^2}{4\psi_i},\qquad
 \psi_i=\frac{t+i+1}{t+i},\qquad t(t+N)=1,
\]
and, with cyclic indices,
\[
 \psi_j+\psi_{j-1}^{-1}=k_j,
 \qquad k_0=N+2,\quad k_j=2\quad(1\leq j<N).
\]
Each $\psi_i$ is irrational.  Throughout this section, modules are
grading-restricted generalized $W$-modules.  Their restrictions to
$\mathcal B$ need not lie in the $C_1$-cofinite Virasoro category;
the block argument below establishes the finiteness properties that
we need.

\subsection{Exact frame-block projections}

For an irrational parameter $\psi$, the Feigin--Fuchs structure theorem
says that a Virasoro Verma module is simple unless its highest weight
is $h(r,s)$ for positive integers $r,s$.  In that case its unique
nonzero proper submodule is the simple Verma module of highest weight
$h(r,-s)=h(r,s)+rs$; see \cite[Proposition 2.2.1 and
Theorem 2.2.4]{virasoro}.  Hence every Verma module has length at most two.
The corresponding blocks are singletons and the disjoint pairs
\[
                     \{h(r,s),h(r,-s)\},\qquad r,s>0.
\]
The pairs are disjoint because equality of two signed Kac weights at
irrational $\psi$ forces equality of their signed labels up to
simultaneous reversal.

\begin{lemma}\label{neutral:lem:frame-blocks}
The restriction of a grading-restricted generalized $W$-module $E$
to $\mathcal B$ is the direct sum of its joint Virasoro block
components.  The projections onto these components are canonical
$\mathcal B$-module maps and are exact on short exact sequences.
Every fixed joint block component has finite length.
\end{lemma}

\begin{proof}
The subspace generated by finitely many vectors under the joint
positive Borel algebra is finite-dimensional.  Indeed, positive modes lower
total conformal weight, partial zero modes preserve total weight,
the total grading is bounded below, and each total generalized weight
space is finite-dimensional.  The image of this Borel algebra is
solvable: the positive modes strictly lower the finite set of relevant
total weights and the zero modes commute.  Lie's theorem therefore
gives a finite flag with one-dimensional quotients, on which positive
modes vanish.  PBW induction is exact.  Consequently the frame module
generated by the chosen vectors is a quotient of a module with a
finite filtration by tensor products of ordinary Virasoro Verma
modules.  It has finite length.

We justify block separation also for generalized partial zero modes.
In a finite generalized highest-weight module for one Virasoro factor,
the nilpotent part of $L(0)$ commutes with all modes.  In an extension
of two nonisomorphic ordinary simple modules, that nilpotent part
therefore factors through a homomorphism from the quotient simple to
the submodule simple, and is zero.  Such an extension is ordinary.

An ordinary extension between two simples in different displayed
blocks splits.  If their highest weights are incongruent modulo
$\mathbb Z$, weight cosets split it.  Otherwise, dualize if necessary
so that the quotient has the smaller highest weight.  Its highest
vector lifts to a highest vector of the extension, since the
submodule has no weights below its larger highest weight.  The
resulting Verma-module map has zero intersection with the submodule:
the Verma structure theorem lists no constituent of that highest
weight in a different block.  Its image is thus a section.  Restricted
contragredience is exact, so this also proves the other orientation.
Finite-length induction now gives the canonical block decomposition
for finite generalized highest-weight modules.

The restriction of all of $E$ to one factor can have infinite partial
weight multiplicities.  For each vector, however, its one-factor
positive-Borel orbit is finite by the total-weight argument above.
The induced module has finite partial weight spaces and finite Verma
length.  Thus this restriction is a union of finite generalized
highest-weight submodules.  Their canonical block decompositions
agree under inclusions.  They define canonical projections on the
union.  The remaining factors commute with the chosen one and
preserve its block components.  Iterating gives the joint projections.
Their functoriality and the direct-sum decompositions prove exactness.

A fixed joint block contains at most $2^N$ simple frame types.  Let
$D$ be their finite set of total lowest weights.  Every finite
submodule of this component has length at most
\[
                         \sum_{d\in D}\dim E_{[d]}.
\]
This follows by additivity of generalized weight-space dimensions
along a composition series: each constituent contributes at least one
dimension at its own total lowest weight.  The lengths of the finitely generated submodules are therefore
bounded integers.  Choose one of maximal length.  Adjoining any
further vector gives a finite-length submodule of the same length,
so the inclusion is an equality.  Thus the chosen submodule is the
entire joint block component.
\end{proof}

\subsection{The aligned odd collection}

Call a joint frame block \emph{aligned and odd} if it has a simple
constituent
\[
 \bigotimes_iL_i(r_i,s_i),\qquad
 r_i,s_i\in2\mathbb Z+1,\qquad r_i^2=s_{i-1}^2.
\]
Linkage preserves this condition because it changes only signs of
the Kac labels.  Let $\mathscr P$ be this collection of blocks, and
let $E_{\mathscr P}$ be the sum of their components in $E$.

\begin{lemma}\label{neutral:lem:aligned-odd-stable}
If the total generalized conformal weights of $E$ are integers,
then $E_{\mathscr P}$ is a $W$-submodule and has no negative total
weights.  The restriction of $W$ to the frame belongs to $\mathscr P$.
\end{lemma}

\begin{proof}
We first record the necessary local null constraint, including the
case of signed Kac labels.  For a nonzero intertwining operator with
first input $L(3,1)$, second input $L(r,s)$, and a simple highest-weight
output of weight $x$, one has
\[
 (x-y)\bigl[(x-y-2\psi+1)(x-y-1)-4\psi y\bigr]=0,
                    \qquad y=h(r,s).
\]
Its three roots are $h(r-2,s),h(r,s),h(r+2,s)$.  The reciprocal formula
for first input $L(1,3)$ changes $s$ by $-2,0,2$.
The polynomial is the image of the elementary level-three singular
vector in the Virasoro Zhu bimodule
\cite[Proposition 5.2.3]{virasoro}.  Necessity for arbitrary simple
highest-weight second and output modules follows from the corrected
injectivity theorem \cite[Lemma 5.5 and Propositions 5.6--5.8]{huangyang}:
both the second module and the output contragredient are generated by
their lowest spaces, and the first Zhu bimodule is cyclic on its
highest class.  Thus a nonzero intertwiner cannot have zero highest
coefficient.  No sufficiency assertion or fusion decomposition for
signed-label modules is being used.  For the joint frame these same
necessary relations hold separately at each factor.

Project the action of an elementary summand to a fixed source and
target frame block.  Both blocks have finite length by
Lemma~\ref{neutral:lem:frame-blocks}.  If this projected intertwiner is
nonzero, restrict it to the first nonzero source-filtration term;
it factors through that term's simple quotient.  The span of all
coefficients is a frame submodule, by the intertwiner commutator
formula.  Choose the smallest target-filtration term containing this
coefficient span and project to its simple quotient.  Minimality makes
the resulting intertwiner between simple subquotients nonzero.
The preceding necessary constraints therefore determine
the possible target blocks even when the original blocks are not
semisimple.

The summand $R_j$ changes only $r_j$ and $s_{j-1}$, each by $0,\pm2$.
Let $\Delta_r$ and $\Delta_s$ be the respective changes of their
squares.  The possibly irrational part of the change of total
highest weight is
\[
 \frac14\bigl(\psi_j\Delta_r+\psi_{j-1}^{-1}\Delta_s\bigr)
 =\frac14\bigl(\psi_j(\Delta_r-\Delta_s)+k_j\Delta_s\bigr).
\]
All other terms are rational.  Both highest weights occur in the
integer-graded module $E$, so irrationality of $\psi_j$ forces
$\Delta_r=\Delta_s$.  Alignment and odd parity are preserved.
The frame and every $R_j$ therefore preserve $E_{\mathscr P}$;
the vertex-generation theorem proves $W$-stability.

For nonnegativity, reverse both labels at each site if necessary to
write $r_i=\rho_i>0$ and $s_i=\epsilon_i\rho_{i+1}$, where every
$\rho_i$ is odd and $\epsilon_i\in\{1,-1\}$.  Direct substitution gives
\begin{equation}\label{neutral:eq:aligned-positive-weight}
 4\sum_i h_i(r_i,s_i)
 =\sum_i(\rho_i-\rho_{i+1})^2+N(\rho_0^2-1)
       +4\sum_{i:\epsilon_i=-1}\rho_i\rho_{i+1}\geq0.
\end{equation}
Frame descendants only increase total weight.  Finally, the
constructed summands $M_m$ of $W$ have $r_i=s_{i-1}=2m_i+1>0$,
so belong to this collection.
\end{proof}

\begin{lemma}[Vacuum-like vectors]\label{neutral:lem:vacuum-like}
If a vector $v$ of a weak $W$-module is killed by the total translation
operator $D=L(-1)$, then $a_nv=0$ for every $a\in W$ and every
$n\geq0$.  The map $a\mapsto a_{-1}v$ is a $W$-module homomorphism
from $W$ to the module containing $v$.
\end{lemma}

\begin{proof}
Lower truncation gives $a_nv=0$ for all sufficiently large integers
$n$.  Descending induction using $[D,a_n]=-n a_{n-1}$ proves the
vanishing for $n\geq0$.  The module Jacobi identity then gives
$b_m(a_{-1}v)=(b_ma)_{-1}v$, which is the asserted homomorphism.
\end{proof}

\begin{proposition}\label{neutral:prop:two-lowest-weights}
Let $S$ be a simple grading-restricted $W$-module.  Every nonsplit
grading-restricted generalized extension
\[
                0\longrightarrow S\longrightarrow E
                         \longrightarrow W\longrightarrow0
\]
has $E=E_{\mathscr P}$, has $S\not\cong W$, and has lowest weight
$h(S)\in\{0,1\}$.  On its lowest space, the possible joint generalized
partial spectra are $\{0\}$ if $h(S)=0$, and
$\{e_0,\ldots,e_{N-1}\}$ if $h(S)=1$.  In the weight-zero case the
joint generalized spectrum on all of $E_0$ is also zero and $L(0)=0$
there exactly.
\end{proposition}

\begin{proof}
An irreducible lower-bounded generalized module is ordinary: the
kernel of the nilpotent part of $L(0)$ is a nonzero submodule, hence
the entire module.  Its weights lie in one coset modulo $\mathbb Z$.
If the coset differs from the vacuum coset, the two total weight
cosets split the extension.

In the remaining case $E$ is integer graded.  The map $E\to W$
restricts to a surjection $E_{\mathscr P}\to W$, by exactness of the
frame-block projections.  Its kernel is $S\cap E_{\mathscr P}$,
which is either zero or all of $S$ by simplicity and
Lemma~\ref{neutral:lem:aligned-odd-stable}.  If the kernel is zero, this restriction is an isomorphism onto $W$
and gives a section.  If the kernel is $S$, then $E_{\mathscr P}$
contains $S$ and surjects onto $E/S$, so $E=E_{\mathscr P}$.
Lemma~\ref{neutral:lem:aligned-odd-stable} then excludes negative
weights.

If $h(S)>1$, the unique vacuum lift $v\in E_0$ has $Dv=0$ because
$E_1=0$.  Lemma~\ref{neutral:lem:vacuum-like} splits the extension.  The same
argument handles $S\cong W$: here $E_1=0$ because $W_1=0$, and any
lift of the vacuum is killed by $D$.  Hence only the two stated
lowest weights remain, with $S\not\cong W$.

To determine the partial spectra, apply
\eqref{neutral:eq:aligned-positive-weight}.  If the
positive odd sequence $(\rho_i)$ is nonconstant, its cyclic differences
include at least two nonzero even integers, so its total lowest weight
is at least two.  If it is constant and greater than one, that weight
is at least $2N$.  At weights zero and one, therefore, all $\rho_i=1$.
No negative sign gives the vacuum tuple $(0,\ldots,0)$; exactly one
negative sign at site $j$ gives the tuple $e_j$.  The simple vacuum
frame module has no total weight-one descendants.  Composition
series of the finite frame blocks consequently give exactly the
stated possible generalized spectra, with no assertion of partial
semisimplicity.

Finally, the nilpotent part of total $L(0)$ on $E$ vanishes.  It is a
module endomorphism, kills $S$, is zero on the quotient, and so factors
through a homomorphism $W\to S$.  That homomorphism is zero because
$S\not\cong W$.  In particular $L(0)=0$ on $E_0$ in the weight-zero
case.
\end{proof}

\subsection{Exact identities for a weight-one vacuum defect}

Suppose now $h(S)=1$, choose the unique vacuum lift $v\in E_0$, and set
\[
 H_i=L_i(0),\qquad w_i=L_i(-1)v,\qquad w=\sum_iw_i.
\]
Every positive-weight zero mode and every degree-lowering mode kills
$v$.  The $w_i$ are jointly highest and satisfy $H_iw_j=\delta_{ij}w_j$.
A nonsplit extension has $w\ne0$ by Lemma~\ref{neutral:lem:vacuum-like}.

For a positive-weight joint frame eigenvector $a\in W$, of total
weight $p$ and partial weights $h_i$, put $r_a=a_{p-2}v$.  The
following identities hold without a semisimplicity assumption on
the partial zero modes acting on $S_1$:
\begin{align}
 o(a)w_i&=(h_i-H_i)r_a,&
 o(a)w&=(p-1)r_a,\label{neutral:eq:neutral-first-image}\\
 (L_i(-n)a)_{p+n-2}v&=(-1)^n o(a)w_i
                         &&(n\geq1).\label{neutral:eq:neutral-desc-first-image}
\end{align}
Indeed, the iterate formula for
$(\omega_i{}_{1-n}a)_{p+n-2}v$ leaves only the term containing
$\omega_i{}_0v=w_i$: all $a$-modes of degree at most zero on $v$
vanish and $\omega_i{}_kv=0$ for $k\geq1$.  This gives
\eqref{neutral:eq:neutral-desc-first-image}.  The commutator with $H_i$ at
mode $p-2$ gives the first identity in
\eqref{neutral:eq:neutral-first-image}; summing it uses $\sum_iH_i=1$ on $S_1$.

\begin{lemma}\label{neutral:lem:first-image-cubic}
If the local factor of a joint primary $a$ is $(3,1)$ or $(1,3)$,
with partial weight $h$, then its first image satisfies
\[
                         H_i(H_i-1)(H_i-h)r_a=0.
\]
In particular, its components at partial eigenvalues zero and one
are ordinary at that site.
\end{lemma}

\begin{proof}
The local singular vector is
\[
 (L_i(-1)^3-2(h+1)L_i(-2)L_i(-1)+h(h+1)L_i(-3))a=0.
\]
Write $X=H_i$.  Equations
\eqref{neutral:eq:neutral-first-image}--\eqref{neutral:eq:neutral-desc-first-image}
give, respectively, the following polynomials on $r_a$ for its three
terms:
\[
 (X-h)(X-h-1)(X-h-2),\quad
 -(X-h)(X-h-1),\quad X-h.
\]
Their stated linear combination is $(X-h)X(X-1)$.  Here $h$ is
irrational, so these roots are distinct.  The assertion follows on
each finite generalized eigenspace.  This is a direct first-image
mode computation, rather than an ordinary Zhu calculation applied
to a degree-raising mode.
\end{proof}

For homogeneous states $a,b$ of weights $p,q\geq2$, the usual mode
commutator and \eqref{neutral:eq:neutral-first-image} give
\begin{equation}\label{neutral:eq:neutral-defect-product}
 o(a)o(b)w=(q-1)\sum_{\ell=0}^{p-1}
 \binom{p-1}{\ell}\frac{o(a_\ell b)w}{p+q-\ell-2}.
\end{equation}
To see this, replace $o(b)w$ by $(q-1)b_{q-2}v$, commute $o(a)$ past
$b_{q-2}$, and use $o(a)v=0$.  The summand is the first image of
$a_\ell b$, whose weight is $p+q-\ell-1$.
For two weight-two states, skew symmetry and $W_1=0$ reduce this to
\[
                       \{o(a),o(b)\}w=o(a_1b)w.
\]
Together with the ordinary-root identities proved below, this
anticommutator relation shows that
$U=\operatorname{span}\{w_i\}$ is invariant under every ordinary
interval zero mode: one adjoins adjacent disjoint intervals
inductively.  We will use this invariant subspace without requiring
a description of all of $S_1$.

We also use the exact partial-conformal Zhu identity
\begin{equation}\label{neutral:eq:neutral-partial-Zhu}
 [L_i(-n)a]=(-1)^n\bigl(-[\omega_i]*[a]+n[a]*[\omega_i]+h_i[a]\bigr),
                          \qquad n\geq1.
\end{equation}
For $n=1$ this is the Zhu commutator formula; $n=2$ follows by expanding
$[\omega_i]*[a]$.  The higher circle relations give the recurrence
$[L_i(-n-2)a]+2[L_i(-n-1)a]+[L_i(-n)a]=0$, proving all remaining
cases.  When it is iterated, $h_i$ is the weight of the descendant
currently being reduced, not the weight of its original primary.

\subsection{An algebraic criterion for second-order nilpotence}

The following form of Nakayama's lemma controls the partial
nilpotent operators on the first cyclic image.

\begin{lemma}\label{neutral:lem:cut-Nakayama}
Let $V$ be finite-dimensional, let commuting operators $H_i$ have
joint generalized spectrum $e_j$ on a component, and suppose
$\sum_i H_i=1$.  Let $r$ be a vector in that component.  Suppose
polynomials $K_0,\ldots,K_{N-2}$ vanish at $e_j$, their gradients are
independent modulo $\sum_i(H_i-\delta_{ij})=0$, and every $K_a(H)r$
is an ordinary simultaneous eigenvector.  Then
\[
 (H_i-\delta_{ij})(H_k-\delta_{kj})r=0\qquad\text{for all }i,k.
\]
\end{lemma}

\begin{proof}
On the cyclic space generated by $r$, let $\mathcal A$ be the finite
commutative algebra generated by $n_i=H_i-\delta_{ij}$.  It is local
with nilpotent maximal ideal $\mathfrak m=(n_i)$.  Invertibility of
the gradients says that the $K_a$ span $\mathfrak m/\mathfrak m^2$;
Nakayama's lemma makes them generators of $\mathfrak m$.  Ordinariness
of their values gives $\mathfrak mK_a(H)r=0$, so
$\mathfrak m^2r=0$.  The maximal ideal is nilpotent because $\mathcal A$ is a finite
local algebra.  The argument therefore requires no additional bound
from a $(3,3)$ singular vector.
\end{proof}


\subsection{The ordinary-root defect calculation}

The ordinary roots determine the initial defect space without any
classification of its ambient module.  We give the calculation because
it is used in every interval equation below.

\begin{lemma}\label{neutral:lem:ordinary-root-defect}
For an ordinary root $R=R_j$, $1\leq j<N$, write
\[
 a=1+\frac2{t+j},\qquad b=2-a,\qquad
 r=R_0v,\quad O=R_1,\quad w_a=w_j,\quad w_b=w_{j-1}.
\]
In the root normalization with nonzero scale $u$, one has the formulas
below.  For the root with leading charge coefficient one, $u$ equals the
screening parameter.  Rescaling $R$ rescales $u$ by the same factor,
so the formulas apply to independently normalized ordinary roots.
\begin{align*}
 r&=u\left(\frac{a+1}{(a-2)(a-1)}w_a
                +\frac{a-3}{a(a-1)}w_b\right),\\
 Ow_a&=u\left(\frac{a+1}{a-2}w_a+
                         \frac{a-3}{a-1}w_b\right),\\
 Ow_b&=u\left(-\frac{a+1}{a-1}w_a-
                         \frac{a-3}{a}w_b\right),
 \qquad Ow_i=0\quad(i\notin\{j-1,j\}).
\end{align*}
If the vacuum extension is nonsplit, all $w_i$ are nonzero and
$U=\operatorname{span}\{w_i\}$ has dimension $N$.
\end{lemma}

\begin{proof}
Off the two active sites $L_i(-1)R=0$, so
\eqref{neutral:eq:neutral-desc-first-image} gives $H_ir=0$ there.
At the active sites Lemma~\ref{neutral:lem:first-image-cubic} and
Proposition~\ref{neutral:prop:two-lowest-weights} make $r$ an ordinary
sum $r_a+r_b$ with spectra $(1,0)$ and $(0,1)$.  The first-image
commutator gives
\[
 Ow_a=(a-1)r_a+ar_b,\qquad
 Ow_b=br_a+(b-1)r_b.
\]

The explicit ordinary-root OPE, in the stated normalization, has
vacuum and self-channel coefficients
\begin{equation}\label{neutral:eq:ordinary-root-constants}
 \nu=u^2\frac{3(3a-5)(3a-1)}{a(a-2)(a-1)^2},\qquad
 \sigma=u\frac{8(2a-1)(2a-3)}{a(a-1)(a-2)}.
\end{equation}
These are the ordinary-edge specializations of the root product
calculation in Proposition~\ref{c:ordinary-root-coefficients}.
The next aligned channel has lowest weight eight, so only these
two channels occur through weight four.

For a local degenerate weight $h$, put
\[
 c_h=\frac{-3h^2+7h-2}{h+1},\qquad
 A_h=\frac{h(h+1)}{(h+2)(2h-1)},\qquad
 B_h=\frac{(h-2)(h+1)}{4(h+2)(2h-1)},
\]
\[
 \alpha_h=\frac{3h(c_h-2h+4)}{c_h(5c_h+22)},\qquad
 \beta_h=\frac{2h(5h+1)}{c_h(5c_h+22)}.
\]
The level-two self-channel is
$A_hL(-2)R+B_hL(-1)^2R$, and the level-four vacuum channel is
$\alpha_hL(-4)\mathbf1+\beta_hL(-2)^2\mathbf1$.
Applying $L(1),L(2)$ to the former and $L(2),L(4)$ to the latter
gives the displayed coefficients.

The identity $R_2R=0$, together with skew symmetry, gives
\[
 Or=\nu\left(\frac a{c_a}w_a+\frac b{c_b}w_b\right)
                        +\frac\sigma2r.
\]
The iterate formula also gives $(R_{-1}R)_2v=Or$.  Its vacuum
channel contributes
$\nu((\alpha_a+\beta_a)w_a+(\alpha_b+\beta_b)w_b)$.
Its self-channel contributes $\sigma F(H_j,H_{j-1})r$, where
\begin{align*}
 F(x,y)={}&A_a(a-x)+B_a(x-a-1)(x-a)\\
         &+A_b(b-y)+B_b(y-b-1)(y-b)
                         +\tfrac14(x-a)(y-b).
\end{align*}
Here the mixed vacuum descendant contributes zero, and the mixed
self descendant has coefficient $1/4$.  These statements follow
directly from \eqref{neutral:eq:neutral-first-image}--
\eqref{neutral:eq:neutral-desc-first-image}.
The two scalar denominators in comparing the formulas are
\[
 F(1,0)-\tfrac12
 =-\frac{a(a-2)(3a-5)}{4(a-4)(2a-3)(2a-1)},\qquad
 F(0,1)-\tfrac12
 =-\frac{a(a-2)(3a-1)}{4(a+2)(2a-3)(2a-1)}.
\]
They are nonzero because $a$ is irrational.  Substituting
\eqref{neutral:eq:ordinary-root-constants} yields the asserted
$r_a,r_b$, and hence the two matrices.

Every off-diagonal coefficient is nonzero.  If one $w_i$ vanished,
the equations on each adjacent ordinary edge would make its neighbor
vanish.  Connectedness of the ordinary path would give $w=0$,
splitting the extension.  Thus all $w_i$ are nonzero.  Their distinct
joint spectra make them linearly independent.
\end{proof}

The anticommutator identity now gives invariance of $U$ under every
ordinary interval.  If adjacent disjoint intervals $J,K$ are shorter
than their union, then
$(P_J)_1P_K=\kappa P_{J\cup K}$ with $\kappa\ne0$, so
$o(P_{J\cup K})w\in U$.  Equation
\eqref{neutral:eq:neutral-first-image} then puts
$o(P_{J\cup K})w_i$ in $U$ for every $i$.



\section{Proper intervals and the cyclic first image}
\label{r:sec:cyclic}

Let \(H_*=2N+2\). Every connected proper interval containing root
zero has primary \(P_J\) of weight \(H_*\); ordinary intervals have
weight two. Write
\[
                   r_J=(P_J)_{H_*-2}v.
\]
For the full cyclic primary \(C\), set
\[
      H=2N,\qquad h_i=\wt_i(C),\qquad
      r_C=C_{H-2}v,\qquad H_i=L_i(0).
\]
The universal frame commutators give
\begin{equation}\label{r:eq:Cwi}
        o(C)w_i=(h_i-H_i)r_C,\qquad
        o(C)w=(H-1)r_C.
\end{equation}

\subsection{Semisimplicity of proper-interval first images}

Appendix~\ref{neutral:app:proper-intervals} proves that every heavy
proper-interval first image \(r_J\) is a sum of simultaneous
eigenvectors of the \(H_i\), with eigenvalue tuples among the \(e_i\).
Thus the partial nilpotent operators vanish on these images.
The proof uses growth by adjacent weight-two intervals.
For a prefix \(J=[0,k]\), cutting at \(a<k\) gives the operator
\[
        D_a=\sum_{i=a+1}^kH_i-(t+a)H_a-1
\]
on \(r_{[0,k]}\), with right-hand side a linear combination of
\((h_i-H_i)r_{[0,a]}\). On the old eigencomponents and the fixed
endpoint \(e_{N-1}\), \(D_a\) has nonzero scalar part; induction forces
ordinariness. On the new endpoint \(e_k\), the right-hand sides
vanish for all cuts. The endpoint \((3,1)\) null relation makes
\(H_k\) semisimple, and the cut equations then successively kill
the internal nilpotents.

The reversed growth proves the suffix statement. For a wrapping
proper interval with growth on both sides of zero, a cut on the
side opposite a given eigencomponent has nonzero scalar part
and reduces to a shorter interval. These arguments prove the
claim for every proper heavy interval. The full operator identities and both directions of growth are given
in Appendix \ref{neutral:app:proper-intervals}.

\subsection{The closing-root product through level three}

For the ordinary complementary interval \(P=P_{[1,N-1]}\) and
the closing root \(R_0\), canonical fusion gives
\[
 (R_0)_3P=\kappa C,\qquad
 (R_0)_jP=0\ (j>3),\qquad \kappa\ne0.
\]
The other singular coefficients are level-one, level-two, and
level-three descendants of \(C\).
For a rank-one parameter \(p\), the \((3,1)\) followed by
\((1,3)\) descendant coefficients are
\begin{align*}
 v_1&=\frac p{p-1}L_{-1}C,\\
 v_2&=-\frac{2p}{(p-2)(2p-1)}L_{-2}C
       +\frac{p^2(2p-3)}{2(p-2)(p-1)(2p-1)}L_{-1}^2C,\\
 v_3&=-\frac{p+2}{3(p-2)}L_{-3}C
       -\frac{2p}{3(p-2)(p-1)}L_{-2}L_{-1}C
       +\frac{p^2}{6(p-2)(p-1)}L_{-1}^3C.
\end{align*}
Reversing the inputs gives
\begin{align*}
 \widetilde v_1&=-\frac1{p-1}L_{-1}C,\\
 \widetilde v_2&=-\frac{2p}{(p-2)(2p-1)}L_{-2}C
       +\frac{3p-2}{2(p-2)(p-1)(2p-1)}L_{-1}^2C,\\
 \widetilde v_3&=\frac{2p+1}{3(2p-1)}L_{-3}C
       -\frac{2p}{3(p-1)(2p-1)}L_{-2}L_{-1}C
       +\frac1{6(p-1)(2p-1)}L_{-1}^3C.
\end{align*}
The Ward equations for \(L_1,L_2,L_3\) determine these coefficients
uniquely, since the output has no singular vector below level nine.

\subsection{The all-cut polynomial system}

For each prefix cut \(a=0,\ldots,N-2\), let
\(B_a=P_{[0,a]}\), \(A_a=P_{[a+1,N-1]}\), and normalize
\[
 (B_a)_3A_a=\kappa_a C,\qquad
 v_a=\kappa_a^{-1}o(B_a)o(A_a)w,\qquad \kappa_a\ne0.
\]
Put
\[
 r=t+a,\qquad
 Z=\sum_{i<a}H_i,\quad X=H_a,\quad Y=H_{N-1}.
\]
Tensoring the preceding local coefficients with the vacuum
descendants at the prefix sites and applying \eqref{neutral:eq:neutral-defect-product}
gives
\begin{equation}\label{r:eq:K}
                        K_a(H)r_C=v_a.
\end{equation}
Each \(v_a\) is a sum of simultaneous eigenvectors.  Indeed,
\(o(A_a)w\in U\), and the analogue of \eqref{r:eq:Cwi} for \(B_a\)
expresses its image under \(o(B_a)\) in terms of the semisimple
proper-interval first image.
Appendix \ref{neutral:app:cut-polynomials} gives a finite explicit
polynomial prescription for every \(K_a\).

The polynomial \(K_a\) vanishes at each allowed joint eigenvalue.
Its gradient with respect to \((Z,X,Y)\) has four forms:
\begin{align*}
 G_{\mathrm{suffix}}&=
 \left(-\frac16,-\frac{r+1}6,-\frac{t-1}{6t}\right),\\
 G_{\mathrm{prefix}}&=
 \left(\frac13,
 \frac{(r+1)(r^2+4r-2)}{3(r-1)(r+2)},
 \frac{(t-1)(2t^2+4t-1)}{3t(t+1)(2t-1)}\right),\\
 G_{\mathrm{cut}}&=
 \left(\frac{3r+2}6,-\frac{r+1}{3(r-1)},
 \frac{(t-1)(6rt^2+9rt-3r+4t^2+8t-2)}
      {6t(t+1)(2t-1)}\right),\\
 G_{\mathrm{last}}&=
 \left(\frac{2t-3}{6t},
 \frac{(r+1)(2r^2t-3r^2+8rt-9r-4t+6)}
      {6t(r-1)(r+2)},
 \frac{t-1}{3(t+1)}\right).
\end{align*}
For the component \(e_j\), use prefix if \(j<a\), cut if \(j=a\),
suffix if \(a<j<N-1\), and last if \(j=N-1\).
All denominators are nonzero for \(r=t+a\).

\begin{theorem}[Uniform second-order nilpotent bound]\label{r:thm:nil}
In the hypothetical weight-one extension above, decompose the
first cyclic image as \(r_C=\sum_jr_{C,j}\), with joint generalized
eigenvalue \(e_j\). Put \(N_i=H_i-\delta_{ij}\) on that component.
Then, for every odd \(N\ge3\),
\[
                         N_iN_\ell r_{C,j}=0
                         \qquad\text{for all }i,\ell,j.
\]
The statement does not assert \(N_ir_{C,j}=0\).
\end{theorem}

\begin{proof}
The operators \(N_i\) commute and are nilpotent on the finite
joint generalized eigenspace, and \(\sum_iN_i=0\) because total
\(L(0)\) is the scalar one on \(S_1\).
Appendix \ref{r:app:jac} proves that the Jacobian formed from all
the displayed gradients is invertible after imposing this sum.

Lemma \ref{neutral:lem:cut-Nakayama}, applied to the ordinary
right-hand sides of \eqref{r:eq:K}, now gives the assertion.
\end{proof}

The theorem kills products of two partial nilpotents, but may leave
first-order nilpotent images.  To eliminate these, we next prove that
the proper heavy interval images on the right-hand sides of
\eqref{r:eq:K} vanish.

\subsection{All first partial nilpotents vanish}

\begin{theorem}[Uniform first-image elimination]\label{r:thm:first}
For every odd \(N\ge3\), in the hypothetical weight-one extension,
\[
 r_B=0\quad\text{for every proper heavy cyclic interval }B,
 \qquad
 (H_i-\delta_{ij})r_{C,j}=0\quad\text{for all }i,j.
\]
Consequently every proper heavy interval zero mode annihilates
\(U\), and \(r_C\) is a sum of simultaneous partial eigenvectors.
The final contradiction will use the consequence \(o(R_0)U=0\).
\end{theorem}

\begin{proof}
Appendix \ref{r:app:first} proves that the two possible components of
the closing-root first image \(r_{R_0}\) vanish. For \(N\ge5\),
triangle differences of complementary-interval equations give
the result uniformly; for \(N=3\), two explicit matrices have
nonzero determinants in \(\mathbb Q[t]/(t^2+3t-1)\).

An ordinary interval zero mode preserves \(U\). If a heavy interval
\(B\) annihilates \(U\), its commutator with an adjoining ordinary
interval therefore annihilates \(U\). The unique order-two OPE
identifies this commutator with
\(\kappa[Q,o(B')]\), where \(B'\) is their proper union and \(Q\)
is an explicit linear combination of the \(H_i\). The first image
of \(B'\) is ordinary, and \(Q\) has at least two distinct
coefficients on its active support. Evaluating on suitable source
vectors \(w_i\) forces every component of \(r_{B'}\) to vanish.
Starting at \(R_0\) proves this for every proper heavy interval.

All right sides of \eqref{r:eq:K} consequently vanish. The same
invertible Jacobian now gives \(\mm r_{C,j}=0\), proving the claim.
The detailed scalar coefficients and propagation identity are
given in Appendix \ref{r:app:first}.
\end{proof}

\section{A second ground character from the lattice embedding}
\label{neutral:sec:second-ground-character}

Here a \emph{ground character} means a one-dimensional representation
of the Zhu algebra on the lowest conformal-weight space. We construct
one that differs from the vacuum character; it will determine the
second root of the cyclic Zhu relation in the next section.

The construction gives a conformal embedding $W\subset A$ into the
deformed lattice vertex operator algebra, with
\[
 A_0=\mathbb C\mathbf1\oplus\mathbb C\varepsilon,
 \qquad \varepsilon=e^B,\qquad
 \langle\mathbf1,\mathbf1\rangle=1,
 \quad\langle\varepsilon,\mathbf1\rangle=-\frac1{2\lambda},
 \quad\lambda\ne0.
\]
The frame zero modes kill both ground vectors.  For $\mathbf1$ this is
the vacuum property.  For $\varepsilon$, the undeformed partial
weights are zero, while every correction to a frame vector has
strictly positive charge.  Such a correction has no degree-zero
target above charge $B$ and hence also kills $\varepsilon$.
Normalize the full cyclic primary $C$ by its leading charge term
$e^{-B}$.

Each elementary highest vector was constructed in the cone above
its charge $-\beta_j$.  For disjoint binary root supports, diagonal
fusion has only their union as output.  Taking the highest
coefficient of successive disjoint products therefore places every
binary highest vector in the cone above
$-\sum_{j\in J}\beta_j$.  In root coordinates, a primary with binary
support $J$ thus has charges componentwise at least $-\mathbf1_J$.

\begin{lemma}\label{neutral:lem:second-ground-character}
There is a simple grading-restricted $W$-module $T$ of lowest weight
zero whose ground character has
\[
 H_i=0,\qquad o(P_J)=0\quad(J\text{ a proper interval}),
                     \qquad o(C)=2\lambda.
\]
Its existence does not require semisimplicity of the restriction of
$A$ to $W$.
\end{lemma}

\begin{proof}
For a proper interval $J$, choose a coordinate absent from its
support.  Applied to $\varepsilon$, its primary zero mode can only
have charge $B$ in degree zero: charge zero is excluded by that
coordinate.  Thus $o(P_J)\varepsilon$ is a scalar multiple of
$\varepsilon$.  This scalar is zero.  Indeed, $P_J$ is a primary of
even total weight, so its zero mode is self-adjoint for the invariant
form, and
\[
 \langle o(P_J)\varepsilon,\mathbf1\rangle
       =\langle\varepsilon,o(P_J)\mathbf1\rangle=0,
                     \qquad\langle\varepsilon,\mathbf1\rangle\ne0.
\]
For $C$, extraction of its leading charge gives
$o(C)\varepsilon=\mathbf1+d\varepsilon$.  The same pairing gives
$0=1-d/(2\lambda)$, and hence
\[
             o(C)\varepsilon=\mathbf1+2\lambda\varepsilon.
\]
Thus the nonzero one-dimensional line $(A/W)_0$ has the stated Zhu
character.  Let $M$ be the $W$-submodule of $A/W$ generated by that
line.  Any proper submodule misses its degree-zero line, because that
line generates $M$.  The sum of all proper submodules still misses
degree zero and is therefore proper.  The quotient by this sum is
simple, grading restricted, and retains the same ground character;
call it $T$.  Since its $C$-eigenvalue differs from the vacuum value
zero, $T\not\cong W$.
\end{proof}



\section{Lowest weight zero and the Zhu algebra}
\label{r:sec:h0}

We now treat lowest weight zero.  A spanning result for the Zhu
algebra controls the number of factors involving the distinguished
root.  Appendix~\ref{neutral:app:filtered-spanning} gives the full
maximum-plateau argument; we record its filtered consequence below.

\subsection{A filtered spanning lemma}

The partial frame classes, all proper cyclic interval highest
classes, and the full-cycle class \([C]\) generate \(A(W)\).
Disconnected binary supports factor by regular normal products;
higher labels reduce by the maximum-plateau induction in
Appendix \ref{neutral:app:filtered-spanning}.
Let \(A_0\subset A(W)\) be the ordinary open-chain image, including
the partial frame classes, and put
\[
                 F_pW=\bigoplus_{m_0\le p}M_m.
\]
Aligned fusion gives \((F_pW)_n(F_qW)\subset F_{p+q}W\).
If the first input has label zero at the distinguished root, fusion
preserves the second input's label at that root.

\begin{lemma}[Filtered Zhu span]\label{r:lem:filtered}
One has
\begin{equation}\label{r:eq:filteredspan}
 [F_1W]\subseteq A_0+
 \sum_{P\in\{C,\ \text{proper heavy intervals}\}}A_0[P]A_0.
\end{equation}
\end{lemma}

\begin{proof}
If \(m_0\le1\) and \(\max_i m_i\ge2\), peel a connected
maximal plateau with indicator \(u\). It avoids root zero.
The top-channel weight formula gives
\(\Delta=h(m)-h(u)-h(m-u)\ge0\), and
\((v_u)_{-\Delta-1}v_{m-u}\) has a nonzero \(v_m\) component.
Every other aligned output has strictly smaller total label and
the same root-zero label. Induction on total label therefore
retains at most one root-zero-bearing binary factor in a
normal-product expression. Partial Virasoro descendants retain
this count: each \(L_i(-1)\) is a derivation, and higher negative
Virasoro modes add only a frame factor. A terminal binary support
has at most one connected component containing zero, which is
heavy or the full cycle. Its other components lie in the open chain.

Pass to Zhu classes by a second induction on total conformal
weight. Replacing a normal product by its Zhu product leaves only
lower-weight singular terms, still in \(F_1W\), so the count is
preserved. Higher negative modes are handled by total derivatives.
Partial descendants reduce using
\[
 [L_i(-n)a]=(-1)^n\bigl(-H_i*[a]+n[a]*H_i+h_i[a]\bigr),
 \qquad n\ge1,
\]
which adds only left and right frame classes. This proves
\eqref{r:eq:filteredspan} without commuting its noncommutative factors.
\end{proof}

\subsection{The two-character quotient}

\begin{theorem}[Closure of the weight-zero branch]\label{r:thm:h0}
Let \(M\) be a finite-dimensional Zhu representation with
generalized partial joint spectrum zero and exact total
conformal operator \(\sum_iH_i=0\). Then \(H_i=0\) for every
\(i\), every proper interval primary acts by zero, and
\[
             C(C-2\lambda)=0,\qquad \lambda\ne0.
\]
In particular,
\[
 A(W)/\langle H_0,\ldots,H_{N-1}\rangle
 \cong\C[X]/(X(X-2\lambda))\cong\C\oplus\C.
\]
The relevant aligned lowest-weight-zero simples are exactly
\(W,T\), and every vacuum extension in this branch splits.
\end{theorem}

\begin{proof}
Appendix \ref{r:app:h0-open-chain} first proves that ordinary
roots commute with all the partial nilpotents. Comparing a
root square with its derivative commutator gives
\[
 (t+j+1)H_j-(t+j-1)H_{j-1}\in\mm^2,
\]
where \(\mm\) is the nilpotent ideal generated by the \(H_i\).
Together with \(\sum_iH_i=0\), the positive recurrence coefficients
give \(\mm=\mm^2\), hence \(\mm=0\).
A difference of endpoint commutators then kills every ordinary
interval; the possible height-two output channels cancel explicitly.

Appendix \ref{r:app:h0-completion} eliminates maximal proper heavy
intervals using two independent coefficients in an ordinary-root
product with \(C\), calibrated on the ground character constructed in
Lemma \ref{neutral:lem:second-ground-character}.
A nonzero level-two coefficient descends to every heavy interval.
Lemma \ref{r:lem:filtered} now implies \(C^2=\alpha+\beta C\):
the square has weight at most \(4N\), below the \(6N\) threshold
for root-zero label two. The actual vacuum and \(T_0\) characters
give \(\alpha=0\), \(\beta=2\lambda\), and both distinct
characters occur. To make this an identity of the quotient algebra,
apply the preceding operator relations to its left regular module;
it is finite-dimensional because \(A(W)\) is finite-dimensional.
This proves the exact quotient. Zhu's lowest-space correspondence
\cite[Theorem 2.2.2]{zhu} identifies its two simple characters with
the already constructed simple modules \(W\) and \(T\).

Proposition~\ref{neutral:prop:two-lowest-weights} verifies these
hypotheses for the aligned lowest-weight-zero tops and the vacuum
extensions under consideration.
For an extension of \(W\) by a distinct simple of weight zero,
the nilpotent part of total \(L(0)\) factors through
\(\Hom(W,S)=0\), so the total operator on its top is also zero.
Lemma~\ref{r:lem:lift-top} below lifts the splitting of the lowest
space to a splitting of the module extension.
\end{proof}

\subsection{Why a split top splits a vacuum extension}

\begin{lemma}[Lifting a splitting of the lowest space]
\label{r:lem:lift-top}
For a simple \(S\not\simeq W\) of lowest weight zero, restriction to
degree zero induces an injection
\[
 \Ext_W^1(W,S)\hookrightarrow
 \Ext_{A(W)}^1(\C_{\rm vac},S_0).
\]
The left-hand side is computed in grading-restricted generalized
modules, and \(\C_{\rm vac}\) is the vacuum Zhu character.
\end{lemma}

\begin{proof}
Suppose the induced extension of lowest spaces splits.  Choose a lift
\(v\in E_0\) with the vacuum Zhu character. Every
positive-weight zero mode kills \(v\). The commutator
\[
 a_nL(-1)v=L(-1)a_nv+n\,a_{n-1}v
\]
shows that all degree-lowering modes kill \(L(-1)v\).
PBW degree ordering puts its generated submodule in weights
at least one. It lies in the simple \(S\), whose weight-zero
space is nonzero, so \(L(-1)v=0\). The lift is vacuum-like
and splits the original extension.  Thus the map on extension
classes has zero kernel.
\end{proof}

The semisimple quotient in Theorem~\ref{r:thm:h0} supplies the splitting
of the lowest space.  Vacuum self-extensions were already handled in
Proposition~\ref{neutral:prop:two-lowest-weights}.

The exact-total-operator hypothesis matters: the theorem is
not an assertion that every generalized zero-eigenvalue
block of an arbitrary module is semisimple. The argument
needs precisely the simple tops and vacuum extensions
specified above.

\section{The closing-root square excludes weight-one extensions}
\label{r:sec:h1}

\begin{theorem}[Closure of the weight-one branch]\label{r:thm:h1}
There is no nonsplit extension of \(W\) by a simple module
of lowest weight one.
\end{theorem}

\begin{proof}
Assume such an extension exists. Its nonzero defect space
belongs to the aligned odd collection by
Proposition \ref{neutral:prop:two-lowest-weights}. That space
\(U=\operatorname{span}\{w_j\}\) has
\(H_iw_j=\delta_{ij}w_j\), with every \(w_j\ne0\).
Theorem \ref{r:thm:first} gives \(o(R_0)U=0\).
The partial Zhu descendant formula then shows that every
partial Virasoro descendant of \(R_0\) also annihilates \(U\).
This uses only \(H_iU\subset U\); it does not require a
classification of the full target top.

Put
\[
 K=\wt(R_0)=2N+2,\qquad a=1+2/t,\qquad b=1-2t,
 \qquad a+b=K.
\]
The only aligned channels of \(R_0\times R_0\) are the
vacuum, \(R_0\), and \(2e_0\). The latter has lowest weight
\(6N+8>2K=4N+4\), so it cannot occur in the Zhu square.
Every contribution of the self-channel annihilates \(U\).

At a local degenerate weight \(h\), evaluation of the normalized
vacuum channel on Virasoro highest weight one has exact series
\[
 V_h(z)={}_3F_2\!\left(
 \begin{matrix}-h-1,-h,2\\1-h/2,(1-3h)/2\end{matrix};
 -\frac{z^2}{4(1+z)}\right).
\]
Appendix \ref{r:app:heavy-vacuum-obstruction} derives this formula
from the level-three null relation and proves its uniqueness
as the normalized vacuum solution. On \(w_0\), the active
partial weights are \((1,0)\), so the other endpoint vacuum
series evaluates to \(1\). The resulting Zhu coefficient is
\begin{equation}\label{r:eq:finalnonzero}
 \begin{split}
 P_K(a)&=[z^{2K}](1+z)^K V_a(z)\\
 &=\left(-\frac14\right)^K
 \frac{(-a-1)_K(-a)_K(2)_K}
 {(1-a/2)_K((1-3a)/2)_K\,K!}\ne0.
 \end{split}
\end{equation}
To see the extraction directly, write
\(V_a(z)=\sum_nd_nz^{2n}(1+z)^{-n}\).
For \(n<K\), the corresponding term after multiplication
by \((1+z)^K\) is a polynomial of degree \(K+n<2K\).
For \(n>K\), its lowest degree is \(2n>2K\).
Only \(d_K\) remains. Irrationality of \(a\) excludes all
zeros and poles of the finite Pochhammer products.

Let \(\nu_0\ne0\) be the vacuum-channel coefficient. Its
nonvanishing follows from the nondegenerate invariant form on the
multiplicity-one self-contragredient frame summand \(M_{e_0}\).
Applying \(R_0*R_0\) to \(w_0\) gives
\[
              0=\nu_0P_K(a)w_0\ne0,
\]
a contradiction. This proves the theorem.
\end{proof}

Only the closing-root first-image vanishing is used in this proof;
no description of repeated cyclic action on the full target top is required.

\section{Vacuum projectivity and regularity}
\label{r:sec:projectivity}

The passage from vacuum projectivity to semisimplicity is categorical.
We isolate it from the extension calculations, using the rigidity
criterion of Etingof--Penneys \cite[Theorem 1.1]{ep}.

\begin{lemma}[A projective-unit criterion]\label{lem:projective-unit}
Let $\mathcal C$ be a finite $\mathbb C$-linear abelian braided
monoidal category with simple unit $\one$ and right-exact tensor
product. Suppose it has a contravariant duality $X\mapsto X'$ with
natural isomorphisms
\[
 \operatorname{Hom}(Y\boxtimes X,\one)
       \cong\operatorname{Hom}(Y,X').
\]
If $\one$ is projective, then $\mathcal C$ is rigid and semisimple.
\end{lemma}
\begin{proof}
For every nonzero $X$, the identity of $X'$ corresponds to a nonzero
map $e_X:X'\boxtimes X\to\one$. Since the unit is simple, $e_X$ is
an epimorphism; projectivity supplies a section. Thus $\one$ is a
direct summand of $X'\boxtimes X$. Every indecomposable $X$ is
therefore non-negligible in the sense of Etingof--Penneys. This
summand argument does not assume the snake identities.

We verify moderate growth without using tensor exactness. Write
$\ell$ for composition length, let $S_i$ run through the finitely
many simple objects, and set $c_X=\max_i\ell(S_i\boxtimes X)$.
Right exactness and induction on a composition series give:
\[
 \ell(Y\boxtimes X)\le c_X\ell(Y),\qquad
 \ell(X^{\boxtimes n})\le\ell(X)c_X^{n-1}.
\]
Every simple endomorphism algebra is $\mathbb C$. Induction on
length in the two variables, using left exactness of Hom, gives
$\dim\operatorname{Hom}(A,B)\le\ell(A)\ell(B)$. Consequently:
\[
 \dim\operatorname{End}(X^{\boxtimes n})
     \le\ell(X)^2c_X^{2n-2}<n!
\]
for sufficiently large $n$. The category is idempotent complete,
has finite-dimensional Hom spaces and finite direct sums, and
$\operatorname{End}(\one)=\mathbb C$. Etingof--Penneys therefore
makes each indecomposable object rigid. Finite direct sums of rigid
objects are rigid, so $\mathcal C$ is rigid.

The tensor functors now have both adjoints and are exact. For every
$X$, rigidity gives:
\[
 \operatorname{Hom}(X,-)
   \cong\operatorname{Hom}(\one,X^*\boxtimes-).
\]
The right-hand side is exact because the unit is projective. Every
object is therefore projective, and the category is semisimple.
\end{proof}

\Needspace{6\baselineskip}
\begin{theorem}[Regularity]\label{r:thm:regularity}
The vertex operator algebra $W_N$ is regular.
\end{theorem}
\begin{proof}
Let $\mathcal C$ be the category of grading-restricted generalized
$W_N$-modules. The CFT-type, $C_2$-cofiniteness and self-contragredience
properties established in Part~I allow us to apply Huang's
vertex tensor-category theorems \cite{huang} and the duality theorem
\cite[Theorem 2.12]{alsw}. Thus $\mathcal C$ is a finite abelian
braided Grothendieck--Verdier category, with right-exact tensor
product and dualizing object $W_N'\cong W_N$.

Proposition~\ref{neutral:prop:two-lowest-weights} reduces a nonsplit
extension of the vacuum by a simple module to lowest weight zero or
one. Theorems~\ref{r:thm:h0} and~\ref{r:thm:h1} exclude these two
cases. Hence:
\begin{equation}\label{r:eq:allvacext}
 \operatorname{Ext}^1_{\mathcal C}(W_N,S)=0
 \qquad\text{for every simple }S.
\end{equation}
Induction on length and the long exact sequence for Ext extend
this vanishing to every object. The vacuum is projective, and
Lemma~\ref{lem:projective-unit} makes $\mathcal C$ rigid and
semisimple.

It remains to pass from grading-restricted modules to all weak
modules. By \cite[Proposition 2.11]{mcrae}, every finitely generated
weak module for an $\mathbb N$-graded $C_2$-cofinite VOA is
grading-restricted generalized. Every cyclic weak $W_N$-module is
therefore a finite sum of ordinary simple modules. An arbitrary
weak module is the sum of its cyclic submodules, hence a sum of
simple submodules. Such a module is semisimple. Thus every weak
$W_N$-module is a direct sum of irreducible ordinary modules, as
required. In particular $W_N$ is rational; compare \cite{abd}.
\end{proof}


\section{The parity factor after regularity}
\label{sec:parity-factor}

By Theorem~\ref{r:thm:regularity}, the category
$\Cat=\operatorname{Rep}(W_N)$ is semisimple and rigid.
McRae's factorizability theorem~\cite[Theorem 5.3]{mcrae}, applied to
the simple, self-contragredient, $\mathbb N$-graded, $C_2$-cofinite
VOA $W_N$, makes it a modular tensor category. The parity modules constructed earlier generate
a pointed braided subcategory $\Par$. Its self-braiding form is:
\[
 Q(\epsilon)=i^{\epsilon^{\mathsf t}K_+\epsilon},
 \qquad \epsilon\in(\mathbb Z/2)^N.
\]
The monodromy pairing is $(-1)^{\epsilon^{\mathsf t}K_+\eta}$.
Since $\det K_+=N^2$ is odd, this pairing is nondegenerate.
The centralizer factorization for a nondegenerate fusion subcategory
\cite{dgno} consequently gives a ribbon equivalence:
\[
 \Cat\simeq\Par\boxtimes\Neu,\qquad \Neu=\Par'.
\]
Every simple has a unique expression $J_\epsilon\boxtimes Z$ with
$Z\in\operatorname{Irr}\Neu$. In particular the parity action on
simples is free, every orbit has $2^N$ members, and it has a unique
representative in $\Neu$. This use of finite semisimple category
theory occurs after regularity, and is not an input to its proof.

\Needspace{12\baselineskip}
\part{Simple modules and modular traces}

\section{Frame decomposition of the nonzero periodic lattice sectors}
\label{g:sec:periodic-frame}
Let $L=(\mathbb Z^N,2K_+)$ and $B=\sum_i\beta_i$.  Choose a
class $\xi\in L^*/L$ with nonzero odd-order component $\gamma$.
In the coordinate basis $\beta_i$, write:
\[
 \gamma_i=\frac{a+ib}{N}\pmod{\mathbb Z},\qquad
 \xi=\gamma+\frac\eta2,\qquad \eta\in\{0,1\}^N.
\]
Here $a,b\in\mathbb Z/N\mathbb Z$ are not both zero.  We write
$M_\xi=V_{L+\xi}\downarrow W_N$ for restriction to the conformally
embedded algebra.  A charge vector $k=(k_i)$ denotes the lattice
vector $\sum_i k_i\beta_i$; thus $e^k$ is its lattice exponential.

The frame decomposition follows from the block projections of
Lemma~\ref{neutral:lem:frame-blocks} and the triangularity of the
screening deformation: every correction strictly raises the total
coordinate height $\sum_i k_i$.  We first show that the undeformed charge Focks
occupy distinct singleton blocks.

\begin{lemma}[Nondegeneracy and separation of charge Focks]
\label{g:lem:periodic-focks}
For every charge $k\in\xi+\mathbb Z^N$, the undeformed tensor Fock space
$F(k)$ is an irreducible tensor product of nondegenerate Virasoro Verma
modules. Its joint frame linkage block is a singleton, and no other charge
in the same coset has the same joint block.
\end{lemma}
\begin{proof}
A nonzero affine progression modulo $N$ cannot have two consecutive zero
entries: their difference would give $b=0$, and then $a=0$.
The Kac labels at site $i$ are
$R_i=1-2k_i$ and $S_i=1-2k_{i+1}$. At least one of these is nonintegral,
because their integrality would make $\gamma_i=\gamma_{i+1}=0$.
If its weight were a degenerate Kac weight $h_i(r,s)$ with positive
integers $r,s$, the difference-of-squares identity and irrationality of
$\psi_i$ would force $(R_i,S_i)=\pm(r,s)$, a contradiction. The Verma
module at that weight is therefore irreducible. The charge Fock space has
exactly its character, so the highest-vector homomorphism from the
irreducible Verma is an isomorphism. The irrational Virasoro
linkage theorem makes its block a singleton.

If $F(k)$ and $F(k')$ have the same joint block, then at every site
$(R_i',S_i')=\sigma_i(R_i,S_i)$ with $\sigma_i=\pm1$. A negative sign
would force $2\xi_i,2\xi_{i+1}\in\mathbb Z$, hence
$\gamma_i=\gamma_{i+1}=0$, because $k,k'$ lie in the same lattice coset.
This is impossible, so every sign is positive and $k'=k$.
\end{proof}

\begin{proposition}[Deformed frame decomposition]
\label{g:prop:periodic-frame}
There are canonical simple $\mathcal B$-submodules $E_k\subset M_\xi$
indexed by $k\in\xi+\mathbb Z^N$, and
\[
M_\xi\downarrow\mathcal B=\bigoplus_{k\in\xi+\mathbb Z^N}E_k,
\qquad E_k\cong F(k).
\]
Projection to the lowest charge identifies $E_k$ with $F(k)$. In
particular its highest vector $f_k$ can be normalized to have lowest
charge $k$ and leading term $e^k$.
\end{proposition}
\begin{proof}
Filter $M_\xi$ by scalar charge height
$d(k)=\sum_i k_i$, using the descending subspaces
$F_{\ge d}$ of vectors supported in heights at least $d$. These are
$\mathcal B$-submodules because every frame correction strictly raises
height. The associated quotient at height $d$ is the direct sum of the
undeformed charge Focks at that height.

Fix the singleton block of $F(k)$. Any nonzero vector of the corresponding
deformed block has a minimum charge height. Projecting its generated
finite-length frame module to that minimum height gives a nonzero map
to the undeformed Focks at that height. Since the only such Fock in this
block is $F(k)$, that minimum must be $d(k)$, and the image lies in $F(k)$.
Consequently projection to $F(k)$ is injective on the entire deformed
block: a vector in its kernel would have strictly larger minimum height.
It is surjective by exact block projection applied to
$F_{\ge d(k)}\twoheadrightarrow F_{\ge d(k)}/F_{>d(k)}$.
Thus the deformed block is $F(k)$. The same minimum-height argument
excludes all other blocks.
\end{proof}

\begin{proposition}[Downward charge transitions]
\label{g:prop:downward-links}
For every elementary generator $R_j$ and every $k$, the component of its
vertex action from $E_k$ to $E_{k-\beta_j}$ is nonzero.
\end{proposition}
\begin{proof}
The normalized highest vector $v_j\in R_j$ has lowest charge $-\beta_j$
and leading term $e^{-\beta_j}$, including for the closing generator
$j=0$. All its other charges have strictly larger scalar height.
The normalized $f_k$ has lowest charge $k$. Therefore the minimum-charge
part of $Y(v_j,z)f_k$ is exactly
$Y(e^{-\beta_j},z)e^k$, which is nonzero. Projecting to the canonical
frame block of $F(k-\beta_j)$ retains this leading term. Hence the
asserted block component is nonzero.
\end{proof}


A nonzero submodule could a priori be supported on a proper subset
closed under downward charge transitions.  To prove simplicity, we
also need the upward transitions.  They are obtained from the
monodromy of the level-three BPZ equation in the next section.

\section{BPZ monodromy and simplicity of the periodic sectors}
\label{g:sec:bpz-periodic-simplicity}

We prove that every elementary generator acts nontrivially in both
charge directions.  The argument combines the frame decomposition in
Proposition~\ref{g:prop:periodic-frame} with Mimachi's irreducibility
criterion for the Dotsenko--Fateev equation
\cite[Theorem 11]{mimachi}.  At this stage we use only the frame
restrictions of the lattice modules; the corresponding statement for
arbitrary modules is proved in Section~\ref{g:sec:generic-frame-semisimple}.

\subsection{Frame labels and vacuum coefficients}
Write $\mathcal B=\bigotimes_iL(c_i,0)$ and
$M_\xi=V_{L+\xi}\downarrow W_N$, where the odd-order component of
$\xi\in L^*/L$ is nonzero. The exact frame decomposition is
\[
 M_\xi=\bigoplus_{k\in\xi+\mathbb Z^N}E_k.
\]
Every $E_k$ is an irreducible tensor product of generic Virasoro Verma
modules, occurs once, and has highest vector $f_k$ with rational Kac labels
\[
 R_i=1-2k_i,\qquad S_i=1-2k_{i+1}.
\]
All these labels have odd denominator. At every site at least one of
$R_i,S_i$ is nonintegral. The summands $E_k$ are pairwise nonisomorphic.
The elementary generator $v_j\in R_j$ has $(3,1)$ at site $j$,
$(1,3)$ at site $j-1$, and the vacuum elsewhere.

Put
\[
 a=\psi_j,\qquad b=\psi_{j-1}^{-1},\qquad
 \kappa_j=a+b=\begin{cases}2,&j\ne0,\\N+2,&j=0.\end{cases}
\]
Both $a,b$ are positive irrational numbers. At site $j-1$ use the
equivalent parametrization $b=\psi_{j-1}^{-1}$, interchanging the two
Kac labels. Thus both active factors of $v_j$ are of type $(3,1)$,
at parameters $a$ and $b$ respectively. Their source momenta are
\[
 p=aR_j-S_j,\qquad q=bS_{j-1}-R_{j-1}.
\]
In particular
\[
 p+q=\kappa_j-2-2(K_+k)_j\in\mathbb Z.
\]

The generator has a nonzero vacuum two-point coefficient:
\[
 (v_j)_{2\operatorname{wt}(v_j)-1}v_j=\nu_j\mathbf1,
 \qquad \nu_j\ne0.
\]
This holds also for $j=0$. Indeed the cyclic construction is a
multiplicity-free sum of pairwise distinct self-dual simple frame
modules. Its nondegenerate invariant form restricts nondegenerately to
each of them; distinct summands are orthogonal. Its restriction to the
one-dimensional highest-weight line of $R_j$ is therefore nonzero.
Equivalently, one may use directly the nonzero vacuum channel of the
canonical diagonal algebra in the cyclic-gluing theorem.

\subsection{Exact identification of the third-order local system}
For an irrational parameter $a$, write
\[
 d=2a-1,\qquad H=\frac{p^2-(a-1)^2}{4a}.
\]
The level-three singular relation for the $(3,1)$ primary $X$ is
\[
 \bigl(L_{-1}^3-4aL_{-2}L_{-1}+2a(2a-1)L_{-3}\bigr)X=0.
\]
Consider the Virasoro four-point blocks
$F(z)=\langle H|X(1)X(z)|H\rangle$. Ward identities give
\[
 F'''-4aU F''+(-4aV+2adT)F'+2adW F=0,
 \tag{BPZ}\label{g:eq:bpz}
\]
where
\begin{align*}
 U&=\frac{1-2z}{z(z-1)},\\
 V&=\frac{H}{z^2}+\frac{d}{(z-1)^2}
      -\frac{2d+1}{z(z-1)},\\
 T&=\frac{3z^2-3z+1}{z^2(z-1)^2},\\
 W&=-\frac{2H}{z^3}-\frac{2d}{(z-1)^3}
       -\frac{2d}{z^2}+\frac{2d}{(z-1)^2}.
\end{align*}
For example, translation and dilation give
$\partial_0F=(z-1)F'+2dF$ and
$\partial_1F=-zF'-2dF$. Applying these identities to the
$L_{-1}X$ insertion gives the displayed $V$ coefficient. Thus the
operator includes its accessory parameter, rather than being inferred
only from its local exponents.

To compare this equation with \cite[equation (1.1)]{mimachi},
denote the Dotsenko--Fateev parameters by $A,B,C,G$ and set:
\[
 A=a-1-p,\qquad B=C=-2a,\qquad G=2a.
\]
His equation (1.1) is
\begin{align*}
 &z^2(z-1)^2I'''+(K_1z+K_2(z-1))z(z-1)I''\\
 &\qquad +(L_1z^2+L_2(z-1)^2+L_3z(z-1))I'
 +(M_1z+M_2(z-1))I=0,
\end{align*}
with
\begin{align*}
 K_1&=-G-3B-3C,\qquad K_2=-G-3A-3C,\\
 L_1&=(B+C)(2B+2C+G+1),\\
 L_2&=(A+C)(2A+2C+G+1),\\
 L_3&=(B+C)(2A+2C+G+1)+(A+C)(2B+2C+G+1)\\
    &\quad +(C-1)(A+B+C)+(3C+G)(A+B+C+G+1),\\
 M_1&=-C(2B+2C+G+1)(2A+2B+2C+G+2),\\
 M_2&=-C(2A+2C+G+1)(2A+2B+2C+G+2).
\end{align*}
The substitution
\[
 F(z)=z^{-A}(1-z)^{2a}I(z)
\]
turns this equation exactly into \eqref{g:eq:bpz}. One checks this by setting
$w=-A/z+2a/(z-1)$ and replacing
\begin{align*}
 I'&=z^A(1-z)^{-2a}(F'-wF),\\
 I''&=z^A(1-z)^{-2a}(F''-2wF'+(w^2-w')F),\\
 I'''&=z^A(1-z)^{-2a}
  (F'''-3wF''+(3w^2-3w')F'+(-w^3+3ww'-w'')F).
\end{align*}
Substitution recovers each of the three monic coefficients in
\eqref{g:eq:bpz}.  The two local systems therefore differ by a
rank-one twist, which preserves irreducibility.

\begin{lemma}[A sufficient generic irreducibility criterion]
\label{g:lem:bpz-irreducible}
The rank-three local system \eqref{g:eq:bpz} is irreducible if
\[
 p\notin(a+\mathbb Z)\cup(-a+\mathbb Z)
       \cup(2a+\mathbb Z)\cup(-2a+\mathbb Z)
       \cup(\tfrac12+\mathbb Z).
\]
If $p=aR-S$, with $R,S$ rational of odd denominator and not both
integral, this condition holds. Under this additional rational-label
hypothesis, its three exponents at zero
\[
 \lambda_+=1-a+p,\qquad
 \lambda_0=1-2a,\qquad
 \lambda_-=1-a-p
\]
are pairwise distinct modulo $\mathbb Z$.
\end{lemma}
\begin{proof}
Mimachi's Theorem 11 requires the following eleven quantities to be
nonintegral:
\begin{gather*}
 A,\ A+G/2,\ B,\ B+G/2,\ C,\ C+G/2,\\
 A+B+C+G/2,\ A+B+C+G,\ (G+1)/2,\\
 A+C+(G+1)/2,\ B+C+(G+1)/2.
\end{gather*}
After substitution the nonautomatic conditions are exactly
$a-1-p,2a-1-p,-2a-1-p,-a-1-p,-p-1/2\notin\mathbb Z$.
All other conditions follow from irrationality of $a$.

If $p=ma+n$ with $m,n\in\mathbb Z$, irrationality gives
$R=m$ and $S=-n$, contrary to the hypothesis. If $2p\in\mathbb Z$,
irrationality gives $R=0$ and $2S\in\mathbb Z$; the odd-denominator
condition then makes $S$ integral, again a contradiction. These
observations prove that the rational-label hypothesis implies the
criterion and also gives nonresonance of the three exponents, whose
differences are $a+p,a-p,2p$. Nonresonance is not asserted from the
irreducibility criterion alone: that criterion need not exclude
$2p\in\mathbb Z$.
\end{proof}

\subsection{Factorization and a finite-jet invariance lemma}
Fix $k,j$ and let $f_k^*$ be the grade-zero functional on the frame
summand $E_k$, normalized by $f_k^*(f_k)=1$ and extended by zero on
the other summands. This is a homogeneous vector in the restricted
dual of $M_\xi$, and is a highest vector for the contragredient frame
action. Consider
\[
 \mathcal F(z)=f_k^*Y(v_j,1)Y(v_j,z)f_k.
\]
It is a rational function of $z$ by the ordinary VOA module Jacobi
identity. In particular, it is single-valued about both zero and one.

The exact frame decomposition justifies the following factorization
without a rationality assumption on the frame or on $W_N$. Project
the inner action onto an intermediate $E_l$. The singular relation
for each active $(3,1)$ factor restricts its intermediate highest
weight to the three values with momentum $p,p\pm2a$ at site $j$,
and $q,q\pm2b$ at site $j-1$. This is also obtained directly from
the indicial polynomial of \eqref{g:eq:bpz}. Every other site is unchanged.
All these intermediate Verma modules are generic: changing a rational
Kac label by an even integer does not make a nonintegral pair integral.
Their nondegenerate Shapovalov forms and the Virasoro Ward identities
uniquely determine each projected four-point block from its leading
three-point constants. The PBW bases and the inverse Shapovalov forms
are tensor products, so these projected blocks factor into the product
of the two rank-three Virasoro blocks. This is a direct formal-series
factorization; regular-singular ODE convergence makes it analytic
near zero.

Let $\boldsymbol f(z)=(f_+(z),f_0(z),f_-(z))^T$ and
$\boldsymbol g(z)=(g_+(z),g_0(z),g_-(z))^T$ be the normalized
Frobenius bases of the two equations. There is therefore a constant
matrix $Q$ such that
\[
 \mathcal F(z)=\boldsymbol f(z)^TQ\boldsymbol g(z).
\]
More generally the two-variable expression
\[
 \mathcal H(z_1,z_2)=\boldsymbol f(z_1)^TQ\boldsymbol g(z_2)
\]
has the property
\begin{equation}
 \left.\partial_{z_1}^r\partial_{z_2}^s
          \mathcal H(z_1,z_2)\right|_{z_1=z_2=z}
 =f_k^*Y(v_j,1)
    Y\bigl(L_j(-1)^rL_{j-1}(-1)^sv_j,z\bigr)f_k.
 \label{g:eq:nine-descendants}
\end{equation}
Here the componentwise $L(-1)$ derivative identity is applied to
the factorized frame intertwiners. Its constants are the same $Q$:
the descendants are obtained by differentiating the same moving
insertion. Each right-hand side is again a rational VOA module
matrix coefficient.

\begin{lemma}[Nine rational derivatives force tensor invariance]
\label{g:lem:nine-jets}
The matrix $Q$ defines a monodromy-invariant tensor in the tensor
product of the two rank-three local systems.
\end{lemma}
\begin{proof}
After continuation along a loop about zero or one, write the new
coefficient matrix, in the original local bases, as $Q'$. Rationality
of \eqref{g:eq:nine-descendants} for $0\le r,s\le2$ implies
\[
 \boldsymbol f^{(r)}(z)^T(Q'-Q)\boldsymbol g^{(s)}(z)=0
 \qquad(0\le r,s\le2).
\]
At any common ordinary point of the two differential equations, their
two Wronskian matrices are invertible. Stacking these nine equations
therefore gives $Q'=Q$. The two loops generate the fundamental group.

The nine descendant coefficients are essential here.  The scalar
identity for $\mathcal F$ alone need not determine $Q$, because
multiplication of the two solution spaces can have a kernel.
\end{proof}

\subsection{All elementary charge channels are nonzero}
\begin{proposition}
\label{g:prop:all-periodic-channels}
For every $k\in\xi+\mathbb Z^N$ and every elementary index $j$,
both components
\[
 R_j\times E_k\longrightarrow E_{k-\beta_j},\qquad
 R_j\times E_k\longrightarrow E_{k+\beta_j}
\]
of the $W_N$ vertex action are nonzero. This includes the closing
index $j=0$.
\end{proposition}
\begin{proof}
The two source Kac pairs satisfy Lemma~\ref{g:lem:bpz-irreducible},
so both rank-three monodromy representations are irreducible.
The invariant tensor $Q$ is nonzero. Indeed the vacuum channel of
$Y(v_j,1)Y(v_j,z)$ gives
\[
 \mathcal F(z)=\nu_j(1-z)^{-2\operatorname{wt}(v_j)}
                 +\text{terms with smaller pole order},
 \qquad \nu_j\ne0.
\]
The remaining $W_N$ states have positive conformal weight, and
$f_k^*(f_k)=1$. Thus $\mathcal F$ is not zero.

An invariant tensor between two irreducible three-dimensional
representations is a homomorphism from the dual of one to the other.
Since it is nonzero, it is an isomorphism. Hence $Q$ is invertible.

The intermediate frame summands in the fixed lattice coset are limited to:
\[
 E_{k-\beta_j},\qquad E_k,\qquad E_{k+\beta_j}.
\]
Indeed, equality of
rational Kac weights at an irrational parameter identifies the two
Kac pairs up to simultaneous sign.
A negative sign between two pairs whose charges are in the same
coset would force both adjacent odd-discriminant coordinates to be
zero, which is excluded. Thus at all inactive sites the labels agree
literally, and at the two active sites their changes must be the same
$0,2,$ or $-2$. In the two Frobenius bases, $Q$ is consequently
diagonal, with entries belonging to these three intermediate summands.

Equivalently, this diagonality follows from local monodromy. The
exponents for the second factor are $1-b+q,1-2b,1-b-q$.
Since $p+q\in\mathbb Z$, the matched pairs have integral sums.
The unmatched corner pairs would require $2p\in\mathbb Z$;
a pair involving just one middle exponent would require
$p\in\pm a+\mathbb Z$. Both are impossible by the arithmetic in
Lemma~\ref{g:lem:bpz-irreducible}.

Invertibility makes each diagonal entry nonzero. The entry for
$E_{k\pm\beta_j}$ is the product of the two corresponding
three-point coefficients, for the passage from $E_k$ to that
intermediate summand and back. Its nonvanishing implies that each
of those two actions is nonzero, proving the claim.
\end{proof}

\begin{theorem}[Simplicity of every nonzero odd lattice restriction]
\label{g:thm:periodic-simple}
For every odd $N\ge3$ and every $\xi\in L^*/L$ whose odd-order
component is nonzero, the restriction
\[
 V_{L+\xi}\downarrow W_N
\]
is simple. Its character is therefore the genuine simple-module
character
\[
 \operatorname{ch}M_\xi(\tau)
  =\frac1{\eta(\tau)^N}
    \sum_{k\in\xi+\mathbb Z^N}
     q^{(k-\mathbf1/2)^TK_+(k-\mathbf1/2)}.
\]
Here $k$ is the coordinate column in the basis $(\beta_i)$, so the
lattice vector $B/2$ has coordinate column $\mathbf1/2$.
\end{theorem}
\begin{proof}
A nonzero submodule is a direct sum of some of the pairwise
nonisomorphic simple frame summands $E_k$. Proposition
\ref{g:prop:all-periodic-channels} makes its charge support closed
under both $k\mapsto k+\beta_j$ and $k\mapsto k-\beta_j$, for
every $j$. The lattice graph is connected, so its support is the
entire coset. The submodule is all of $M_\xi$. The character is
the ambient lattice character computed above.
\end{proof}

\section{Frame semisimplicity in generic rational sectors}
\label{g:sec:generic-frame-semisimple}

We next pass from the lattice modules to simple $W_N$-modules in
nonzero rational sectors.  The main issue is that semisimplicity of
the total conformal operator does not imply semisimplicity of its
commuting partial operators.  The level-three singular vectors give
polynomial relations for those partial operators separately; their
simple roots eliminate the possible nilpotent parts.

\begin{lemma}[Structure of a generic generalized frame block]
\label{g:lem:generic-generalized-Verma}
Let $\Lambda$ be a joint frame block whose unique simple object is the
irreducible Verma tensor $F_\lambda$ of joint highest weight $\lambda$.
Every finite-length module $Y$ in this block has the form
\[
 Y\cong U(\mathfrak n_-)\otimes T,
 \qquad T=Y_{[\lambda]},
\]
where all positive modes annihilate $T$ and
$L_i(0)|_T=\lambda_i+N_i$ for commuting nilpotent $N_i$.
Its ordinary frame socle is
\begin{equation}
 \operatorname{soc}_{\mathcal B}(Y)
  =U(\mathfrak n_-)\otimes\bigcap_i\ker N_i.          \label{g:eq:generic-socle}
\end{equation}
Both $Y$ and its restricted contragredient are generated by their
joint highest spaces.  Moreover the Zhu annihilation space is exactly
$\Omega(Y)=T$: there are no additional vectors annihilated by all
positive Virasoro modes at a strictly higher joint degree.
\end{lemma}
\begin{proof}
All simple constituents have weights in
$\lambda+\mathbb Z_{\geq0}^N$, so the same is true of $Y$, and positive
modes kill $T$.  Induction from this highest space gives a generalized
Verma tensor $\mathcal V(T)$ and a map $\mathcal V(T)\to Y$.
The commuting nilpotents on $T$ admit a flag with one-dimensional trivial
nilpotent quotients.  Induction is exact by PBW.  Hence $\mathcal V(T)$
has a finite flag of irreducible ordinary Verma tensors $F_\lambda$.
The cokernel of the displayed map has zero joint highest space; since
every nonzero object in the block has a constituent $F_\lambda$, exact
highest-space extraction makes that cokernel zero.  The kernel has
zero joint highest space by construction and vanishes for the same
reason.  This proves the asserted description.

A vector in $\bigcap_i\ker N_i$ generates an ordinary irreducible
Verma tensor, and every simple submodule must start at joint weight
$\lambda$.  This proves \eqref{g:eq:generic-socle}.  Restricted contragredience preserves
finite length and the single simple type, so the same argument applies
to $Y'$.

For the last assertion, the Shapovalov matrix at every positive joint
degree is invertible at the ordinary generic highest weight $\lambda$.
Replacing its entries by their values at the commuting operators
$\lambda_i+N_i$ changes this invertible matrix by a nilpotent matrix
over the finite local algebra generated by the $N_i$.  It remains
invertible.  Therefore a positive-degree vector killed by all positive
Virasoro modes must be zero.  The highest space is killed by every
degree-lowering mode of the frame VOA, proving $\Omega(Y)=T$.
\end{proof}

\begin{lemma}[The elementary generalized top constraint]
\label{g:lem:generic-top-constraint}
Let $F$ be an ordinary irreducible Verma tensor with rational Kac pairs
of odd denominator, with no Kac pair consisting of two integers.
Let $Y$ be a finite-length generic joint block, ordinary for the
\emph{total} conformal zero mode.  Every frame intertwining operator
obtained by restricting the $W_N$ vertex action,
\[
       \mathcal I:\quad R_j\times F\longrightarrow Y,
\]
has image contained in $\operatorname{soc}_{\mathcal B}(Y)$.
\end{lemma}
\begin{proof}
Write $T=Y_{[\lambda]}$ and $L_i(0)|_T=\lambda_i+N_i$ as in
Lemma~\ref{g:lem:generic-generalized-Verma}.  Since total $L(0)$ is
semisimple, $\sum_iN_i=0$, but individual $N_i$ need not vanish a
priori.  Let $f$ be the highest vector of $F$ with partial weights
$h_i$.  For a homogeneous $u\in R_j$, project
$\mathcal I(u,z)f$ to $T$ and take its unique total-weight-compatible
coefficient.  Denote it by $t(u)\in T$.

The standard necessary leading-coefficient Zhu map, obtained directly
by the intertwining Jacobi identity, sends the left and right actions of
the $i$th conformal vector to $\lambda_i+N_i$ and $h_i$, respectively.
This assertion is valid for a nonsemisimple partial $L_i(0)$: the
residue proof uses only the zero-mode operators and the annihilation
of the source and target highest spaces by positive modes.  In
particular, with $x_i,y_i$ denoting the two Zhu actions, the identities
for a state of partial weight $d_i$ are
\begin{equation}\label{g:eq:partial-zhu}
\begin{aligned}
 [L_i(-1)u]&=(x_i-y_i-d_i)[u],\\
 [L_i(-2)u]&=(2y_i-x_i+d_i)[u],\\
 [L_i(-3)u]&=(x_i-3y_i-d_i)[u].
\end{aligned}
\end{equation}
For example, the first follows already from
\[
 [L_i(0),\mathcal I(u,z)]
   =z\mathcal I(L_i(-1)u,z)+\mathcal I(L_i(0)u,z).
\]
The remaining formulas follow from the left/right Zhu products and
the relation $[L_i(-3)u]+2[L_i(-2)u]+[L_i(-1)u]=0$.
Successive applications of these relations give the PBW reduction
of the Zhu bimodule of $R_j$: it is generated by its highest class
under the commuting $x_i,y_i$.  The argument uses this necessary
leading-coefficient map, rather than a converse fusion-rule statement.

The reduction is performed separately at each Virasoro factor.
Indeed:
\[
 A(\mathcal B)=\mathbb C[x_0,\ldots,x_{N-1}],\qquad x_i=[e_i],
\]
and modes from distinct factors commute.  Each factor of the root
bimodule is generated by its highest class under its own left and
right zero modes, so the tensor bimodule $A(R_j)$ is generated by
the full root highest class.  All partial variables remain
independent.  Lemma~\ref{g:lem:generic-generalized-Verma} identifies
the target annihilation space $\Omega(Y)$ with $T$.

At an active factor use the parameter $a$, so that its root highest
weight is $d=2a-1$ and its singular vector is
\[
 (L_{-1}^3-4aL_{-2}L_{-1}+2a(2a-1)L_{-3})u=0.
\]
Set $s=x-y-d$.  Its Zhu image, from \eqref{g:eq:partial-zhu}, is the polynomial
\begin{equation}
 P_a(x,y)=s(s-1)(s-2)-4a(y-s+1)s
                      +2a(2a-1)(s-2y).              \label{g:eq:zhu-cubic}
\end{equation}
For $h_a(p)=(p^2-(a-1)^2)/(4a)$, direct factorization gives
\begin{equation}
 P_a(x,h_a(p))
   =\prod_{m=-1}^1\bigl(x-h_a(p+2ma)\bigr).         \label{g:eq:zhu-factorization}
\end{equation}
If $p=aR-S$ with the stated rational labels, the differences of
these three roots are $a+p,a-p,2p$, all nonzero; indeed each is
nonintegral by irrationality and the odd-denominator condition.

Apply \eqref{g:eq:zhu-factorization} to $t(u_j)$, where $u_j$ is the highest vector of $R_j$.
If $\lambda_i$ is not one of the three roots, the polynomial in
$\lambda_i+N_i$ is invertible and $t(u_j)=0$.  If $\lambda_i$ is
a root, it is simple, and
\[
 P_a(\lambda_i+N_i,h_i)=N_i U_i(N_i),\qquad U_i(0)\ne0.
\]
The nilpotent polynomial $U_i(N_i)$ is invertible, so
$N_i t(u_j)=0$.  At the other active factor apply the same
calculation with reciprocal parameter, converting $(1,3)$ to
$(3,1)$.  At every inactive vacuum factor the Zhu relation is
$x_i=y_i$; it likewise makes the top coefficient zero unless
$\lambda_i=h_i$, and otherwise forces $N_it(u_j)=0$.
Thus $t(u_j)\in T_0:=\bigcap_i\ker N_i$.

Cyclicity of the root Zhu bimodule now puts \emph{every} $t(u)$ in
$T_0$, because each partial zero mode preserves $T_0$.  Pass to the
quotient $\overline Y=Y/(U(\mathfrak n_-)T_0)$.  All highest
coefficients of the resulting intertwiner vanish.  This implies that
the intertwiner itself vanishes: $F$ is generated by its highest
vector, and $\overline Y'$ is generated by its highest space by
Lemma~\ref{g:lem:generic-generalized-Verma}.  Move their negative
Virasoro modes across the matrix coefficient by the intertwining
Jacobi identity.  Every matrix coefficient is thereby reduced to a
highest-target/highest-source coefficient with some descendant in
the first input, all of which vanish.  Equivalently, the generic
Shapovalov matrices remain invertible after the nilpotent highest
weight substitution, giving the same induction at each degree.

Thus the original image lies in $U(\mathfrak n_-)T_0$, the ordinary
frame socle in \eqref{g:eq:generic-socle}.
\end{proof}

\begin{theorem}[Semisimple frame restriction in a nonzero rational sector]
\label{g:thm:generic-sector-frame-semisimple}
Let $X$ be a simple grading-restricted generalized $W_N$-module
containing an ordinary simple frame module
\[
 F=\bigotimes_iL(c_i,h_i(R_i,S_i)),
\]
where every $R_i,S_i$ is rational with odd denominator and no pair
$(R_i,S_i)$ consists of two integers.  Then
\[
       X\big|_{\mathcal B}
\]
is a direct sum of ordinary irreducible Verma tensors, with finite
multiplicity at each joint highest weight.  In particular the
conclusion holds for every simple module in any of the
nonzero periodic or antiperiodic odd coarse sectors.
\end{theorem}
\begin{proof}
The nilpotent part of total $L(0)$ is a natural $W_N$-endomorphism.
On a simple grading-restricted generalized module it vanishes, so
$X$ is ordinary for total $L(0)$.

Let $X^{\mathrm{ord}}_{\mathrm{rat}}$ be the sum of the ordinary
simple frame submodules whose Kac pairs satisfy the hypotheses.
It contains $F$, and hence is nonzero.  We show that every elementary
generator preserves this sum.

Take a simple frame summand $F_0$ in this sum and project the action
of $R_j$ on $F_0$ to a joint frame block $Y$.  By
Lemma~\ref{neutral:lem:frame-blocks}, the block $Y$ has finite
length.  If the projected intertwiner is nonzero, the frame submodule
spanned by its coefficients has a simple quotient $Q$, and the
induced intertwiner to $Q$ is nonzero.  Both $F_0$ and $Q'$ are
generated by their highest spaces.  The necessary leading-coefficient
Zhu map is therefore injective, and the relations
\eqref{g:eq:partial-zhu}--\eqref{g:eq:zhu-factorization} apply
to its highest coefficient.  At each active factor the target
momentum is $p$, $p+2a$, or $p-2a$; at each inactive factor the
highest weight is unchanged.  Thus the Kac pairs of $Q$ are
obtained from those of $F_0$ by even integral shifts and simultaneous
sign changes.  Each still has odd denominators and at least one
nonintegral entry.

Consequently $Q$ lies in a generic singleton block.  Since $Q$
is a constituent of $Y$, the target block $Y$ is that same generic
block.  Lemma~\ref{g:lem:generic-top-constraint} now puts the entire
projected image in the ordinary frame socle of $Y$, all of whose
simple summands satisfy the required rationality condition.
Each coefficient has only finitely many block components, so the
action of $R_j$ preserves $X^{\mathrm{ord}}_{\mathrm{rat}}$.

The frame and the elementary generators vertex-generate $W_N$.
Thus $X^{\mathrm{ord}}_{\mathrm{rat}}$ is a nonzero
$W_N$-submodule of $X$, and simplicity gives
$X^{\mathrm{ord}}_{\mathrm{rat}}=X$.  Finite-dimensional total
conformal-weight spaces bound the multiplicity of each fixed
joint highest weight.
\end{proof}

\begin{remark}[The leading-coefficient map]
The Zhu-map input above is the necessary leading-coefficient map of
Huang--Yang \cite[Lemma 5.5 and Propositions 5.6--5.8]{huangyang},
in their corrected formulation at level $0$.  Its injectivity applies
when the second input and the contragredient of the target are
generated by their highest spaces, as verified above.  This use does
not require an unrestricted Frenkel--Zhu fusion-rule isomorphism.
The direct residue calculation also keeps the partial conformal
operators separate.  Frame semisimplicity is the conclusion of this
section; multiplicity one and uniqueness require the transport
argument that follows.
\end{remark}
\section{Multiplicity one and uniqueness in rational sectors}
\label{g:sec:generic-sector-classification}

We now determine the multiplicities and the module structure on each
regular rational grid.  Besides frame semisimplicity and the BPZ
calculation, we use the following property of the canonical diagonal
algebra: multiplication is nonzero in every allowed diagonal frame
fusion channel. In particular, with $R_j=M_{e_j}$ in its usual diagonal
labeling,
\begin{equation}
 R_j\cdot R_j\longrightarrow\mathbf1,\quad R_j,\quad M_{2e_j}
 \quad\hbox{are all nonzero},                         \label{g:eq:root-self-products}
\end{equation}
and, for $i\ne j$,
\begin{equation}
 R_i\boxtimes_{\mathcal B}R_j=M_{e_i+e_j},\qquad
 R_i\cdot R_j\longrightarrow M_{e_i+e_j}\ne0.         \label{g:eq:mixed-root-product}
\end{equation}
The fusion in \eqref{g:eq:mixed-root-product} has just one simple summand. At a shared Virasoro
site it is $(3,1)\boxtimes(1,3)=(3,3)$; at every other site one
factor is vacuum. The canonical diagonal tensor-category description
also proves this directly. Finally, the frame and the $R_j$
\emph{vertex-generate} $W_N$, by the construction above.  This
vertex-generation statement suffices for the submodule arguments
below.

\subsection{Regular rational label grids}
Choose signs $\sigma_i\in\{1,-1\}$, and define
\[
 (K_\sigma r)_j=\kappa_jr_j
                  -\sigma_{j-1}r_{j-1}-\sigma_jr_{j+1},
 \qquad \kappa_0=N+2,\quad\kappa_j=2\ (j\ne0).
\]
After a diagonal sign change, this is $K_+$ when
$\prod_i\sigma_i=1$, and $K_-$ when $\prod_i\sigma_i=-1$.
Its determinant is respectively $N^2$ or $N^2+4$, so it is invertible.
For $r\in\mathbb Q^N$ put
\[
 E(r)=\bigotimes_i L\bigl(c_i,h_i(r_i,\sigma_i r_{i+1})\bigr).
\]
Call $r$ regular here if its entries have odd denominators and
no two cyclically adjacent entries are both integral. Every factor
of $E(r)$ is then a generic irreducible Verma module in a singleton
frame block. The term ``regular'' in this section refers only to
this explicit arithmetic condition, not to regularity of a VOA.

Fix such an $r^{(0)}$ with $K_\sigma r^{(0)}\in\mathbb Z^N$, and set
\[
 \mathcal O(r^{(0)})=r^{(0)}+2\mathbb Z^N.
\]
Every point of this grid has the same regularity property. The frame
simples $E(r)$, for $r\in\mathcal O(r^{(0)})$, are pairwise distinct.
Indeed equality of weights forces simultaneous sign equality of the
two rational Kac labels at each irrational Virasoro parameter. A
negative sign at a site, for two points in the same grid, would force
both adjacent coordinates to be integral. This is excluded. The
argument permits isolated integral or zero coordinates.

The nonzero odd periodic and antiperiodic discriminant classes are
regular in this sense.  Periodic residue vectors are the nonzero
affine progressions modulo $N$.  The antiperiodic representatives,
verified below, have coordinates $s(N-2i)/(N^2+4)$ modulo
$\mathbb Z$, where $s\ne0$ in $\mathbb Z/(N^2+4)\mathbb Z$.
Two adjacent periodic entries can vanish only when the progression
is zero.  In the antiperiodic case, adjacent numerators have
difference $2$, except across the closing edge, where their sum is
$2$.  Since $N^2+4$ is odd, two vanishing adjacent entries would
again force $s=0$.

\begin{lemma}[The root action stays on one grid]
\label{g:lem:generic-grid-closure}
Suppose a simple $W_N$-module $X$ contains the frame simple
$E(r^{(0)})$ with the above hypotheses. Then its frame restriction
is semisimple, and every frame summand belongs to
$\{E(r):r\in\mathcal O(r^{(0)})\}$.
\end{lemma}
\begin{proof}
The generic frame-semisimplicity theorem applies to the nonintegral
rational coarse tuple. Alternatively apply its top-polynomial argument
successively from the given simple frame submodule: shifting a Kac
label by $0,2,-2$ preserves its nonintegral residue, so every
possible target block remains generic. The frame socle is stable
under each root, hence under their vertex-generated VOA.

At a point $r$ set
\[
 a=\psi_j,\quad b=\psi_{j-1}^{-1}=\kappa_j-a,
 \quad p=ar_j-\sigma_jr_{j+1},
 \quad q=b\sigma_{j-1}r_j-r_{j-1}.
\]
Then $p+\sigma_{j-1}q=(K_\sigma r)_j\in\mathbb Z$.
All nine elementary combinations of the two three-channel null
polynomials are available as necessary frame fusion possibilities.
An actual module mode changes the total weight by an integer. The
mixed middle/corner channels would require
$p\in\pm a+\mathbb Z$ or $q\in\pm b+\mathbb Z$; the
wrongly paired corner channels would require $2p\in\mathbb Z$.
All are excluded by regularity and irrationality, as in the BPZ
irreducibility lemma. Thus the only possible targets are
\[
                 E(r-2e_j),\quad E(r),\quad E(r+2e_j).
\]
The grid is closed under all root actions. Its sum of frame summands
in $X$ is a nonzero $W_N$-submodule, so simplicity makes it all of $X$.
\end{proof}

Write the resulting frame decomposition as
\[
 X\big|_{\mathcal B}
       =\bigoplus_{r\in\mathcal O(r^{(0)})}E(r)\otimes T_r,
 \qquad\dim T_r<\infty,
\]
initially allowing $T_r=0$. Fix normalized highest vectors of all
$E(r)$ and normalized frame intertwiners for every allowed channel.
Each such intertwiner space has dimension at most one, by the
highest-weight Ward recursion and the nondegenerate generic
Shapovalov forms. Denote
the three multiplicity maps of the root action by
\[
 A_j^\epsilon(r):T_r\longrightarrow T_{r+2\epsilon e_j},
                 \qquad\epsilon\in\{-1,0,1\}.
\]

\subsection{The two-root invariant tensor with multiplicity spaces}
\begin{lemma}[Return maps and middle maps are scalar]
\label{g:lem:scalar-root-maps}
For every $r$ with $T_r\ne0$, all adjacent spaces $T_{r\pm2e_j}$
are nonzero and the maps $A_j^\pm(r)$ are isomorphisms. Their
oppositely oriented compositions are nonzero scalar multiples of
the identity. Moreover $A_j^0(r)$ is a scalar multiple of the
identity. All these scalars depend only on the normalized $W_N$
multiplication and the frame labels, and not on $X$ or on the
dimension of $T_r$.
\end{lemma}
\begin{proof}
Use the operator-valued four-point coefficient
\[
 \Phi_j(z)=\operatorname{pr}_{E(r)^{\rm top}\otimes T_r}
                Y(v_j,1)Y(v_j,z)\big|_{E(r)^{\rm top}\otimes T_r}.
\]
By choosing a vector and a functional in $T_r$ one obtains ordinary
rational matrix coefficients. The factorization and nine-derivative
Wronskian proof apply to each matrix entry. Thus its coefficient
tensor belongs to
\[
 (V_a\otimes V_b)^{\pi_1}\otimes\operatorname{End}(T_r),
\]
where $V_a,V_b$ are the two irreducible rank-three BPZ local systems.
The existence of the nonzero vacuum coefficient shows these systems
are dual; Schur's lemma gives
$\dim(V_a\otimes V_b)^{\pi_1}=1$.

Normalize a generator of this one-dimensional invariant space by
its vacuum-channel coefficient. The vacuum term of $\Phi_j$ is
the fixed nonzero number $\nu_j$ times $\operatorname{id}_{T_r}$.
It follows that the entire invariant tensor equals this normalized
scalar tensor multiplied by $\nu_j\operatorname{id}_{T_r}$.

In the channel at $z=0$ the invariant tensor has the three allowed
entries, one for each $E(r+2\epsilon e_j)$. It is invertible as a
$3\times3$ tensor, so all three entries are nonzero. The two extreme
entries are, with the fixed three-point normalizations,
\[
 A_j^{-\epsilon}(r+2\epsilon e_j)A_j^\epsilon(r)
          =\lambda_{j,r,\epsilon}\operatorname{id}_{T_r},
 \qquad \lambda_{j,r,\epsilon}\ne0.
\]
Applying the same assertion at the neighboring point gives the
reverse composition as a nonzero scalar identity. Hence these maps
are isomorphisms and every neighboring multiplicity is present.

For the middle map use the channel at $z=1$, where the two roots
fuse first. By \eqref{g:eq:root-self-products} the three possible frame intermediates are
$\mathbf1,R_j,M_{2e_j}$. The individual-factor local exponents
are distinct modulo integers, making the three paired tensor
channels independent: in one factor they are
$2-4a,1-2a,2a$, whose differences are nonintegral because $a$ is
irrational. The other factor has parameter $b=\kappa_j-a$.
The tensor in this basis is again diagonal and invertible. Its
$R_j$ entry is the nonzero multiplication constant in
$R_j\cdot R_j\to R_j$ times $A_j^0(r)$. Since every tensor
entry is a fixed scalar times $\operatorname{id}_{T_r}$, the
middle map is scalar too. The scalar is nonzero.
\end{proof}

\subsection{Mixed generators and scalar square transport}
\begin{lemma}[Every elementary square has scalar transport]
\label{g:lem:scalar-squares}
For $i\ne j$ there is a nonzero number $\gamma_{ij}(r)$, depending
only on the normalized frame and algebra data, such that
\begin{equation}
 A_i^+(r+2e_j)A_j^+(r)
     =\gamma_{ij}(r)A_j^+(r+2e_i)A_i^+(r).           \label{g:eq:scalar-square}
\end{equation}
\end{lemma}
\begin{proof}
Project a two-root matrix coefficient from the highest space of
$E(r)$ to that of $E(r+2e_i+2e_j)$. In the product ordering where
$R_j$ acts first, there is just one possible intermediate frame
simple, namely $E(r+2e_j)$. The other ordering similarly has only
$E(r+2e_i)$. The three-channel grid closure and the injectivity of the label
map give these unique intermediate summands.

By \eqref{g:eq:mixed-root-product}, the product of the two root states belongs to the single
simple frame module $P=M_{e_i+e_j}$, and its multiplication
intertwiner is nonzero. The frame intertwiner space
\[
 \mathcal I_{\mathcal B}\binom{E(r+2e_i+2e_j)}{P\quad E(r)}
\]
has dimension at most one: all modules are highest-weight modules,
the target Verma factors are generic, and Ward recursion determines
all three-point coefficients from the highest coefficient. Let
$D_{ij}(r):T_r\to T_{r+2e_i+2e_j}$ be the multiplicity map for
this channel of the $P$ action, if present.

Associativity expresses each ordered product through the same
$P$ action and the same one-dimensional frame intertwiner space.
Thus each ordered composition in \eqref{g:eq:scalar-square} is a scalar multiple of
$D_{ij}(r)$. Those scalars are the ordinary frame change-of-channel
constants together with the fixed multiplication in \eqref{g:eq:mixed-root-product}; they do
not depend on the module. Both ordered compositions are invertible
by the preceding lemma. Hence $D_{ij}(r)$ and both scalar constants
are nonzero. Dividing gives \eqref{g:eq:scalar-square}, with $\gamma_{ij}(r)\ne0$.

In particular the relevant frame intertwiner exists: otherwise
associativity would make an invertible ordered product vanish.
\end{proof}

\begin{theorem}[One copy of every grid point]
\label{g:thm:generic-grid-multiplicity-one}
If a simple $W_N$-module contains $E(r^{(0)})$ as above, then
\[
 X\big|_{\mathcal B}
       =\bigoplus_{r\in r^{(0)}+2\mathbb Z^N}E(r)
\]
with multiplicity one at every point.
\end{theorem}
\begin{proof}
The scalar-return lemma makes every grid point present, all
multiplicity spaces have a common finite dimension, and transport
along every edge is invertible. Every closed lattice walk is a
product of elementary squares and backtracking pairs. By \eqref{g:eq:scalar-square} and
the return identities, its transport is a scalar automorphism.
These scalar holonomies need not equal one.

Choose a nonzero subspace $U\subseteq T_{r^{(0)}}$, and transport
it to $U_r\subseteq T_r$ along any grid path. The result is
path-independent as a \emph{subspace}, because scalar
automorphisms preserve every subspace. The sum
\[
                       \bigoplus_r E(r)\otimes U_r
\]
is stable under the frame and under each root: the two extreme
maps are the edge transports and their scalar inverses, while
the middle map is scalar. Vertex generation makes it a
$W_N$-submodule. Simplicity implies that $U$ cannot be a proper
nonzero subspace. Hence $\dim T_{r^{(0)}}=1$ and all
multiplicities equal one.
\end{proof}

\begin{theorem}[Uniqueness on a fixed regular grid]
\label{g:thm:generic-grid-uniqueness}
Any two simple $W_N$-modules containing the same regular frame
simple $E(r^{(0)})$, with $K_\sigma r^{(0)}\in\mathbb Z^N$,
are isomorphic.
\end{theorem}
\begin{proof}
They have the same multiplicity-one frame decomposition by the
preceding theorem. Fix the same normalized frame intertwiners for
both. The scalar return identities, middle maps, and square
constants are determined by the same frame local systems and
the same fixed $W_N$ multiplication, so they agree in both modules.

Choose a nonzero linear isomorphism between the multiplicity lines
at one base point and extend it by the positive and negative edge
maps. It is independent of the path: the ratio of the two
transports around any square or backtracking pair is one,
because the scalar holonomies agree. It intertwines the extreme
root maps by construction and the middle maps because their
scalars agree. It therefore intertwines every root and the frame,
hence their vertex-generated VOA $W_N$.
\end{proof}

\subsection{Theta characters and the global sign}
The simple module in a regular grid, whenever it exists, has character
\begin{equation}
 \boxed{\qquad
 \operatorname{ch}X_{[r^{(0)},\sigma]}(\tau)
    =\frac1{\eta(\tau)^N}
      \sum_{m\in\mathbb Z^N}
       q^{\frac14(r^{(0)}+2m)^TK_\sigma(r^{(0)}+2m)}.
 \qquad}                                                     \label{g:eq:grid-character}
\end{equation}
Indeed the oscillator-normalized exponent of $E(r)$ is
\[
 \sum_i\frac{(\psi_i r_i-\sigma_i r_{i+1})^2}{4\psi_i}
                         =\frac14r^TK_\sigma r.
\]
Consequently, once the modular-support argument establishes the
existence of a nonzero antiperiodic grid module, its entire character
is given by \eqref{g:eq:grid-character}.

The global sign gives the same set of frame simples:
$E(r)=E(-r)$ and negation bijects the two grids. The uniqueness
theorem consequently identifies their simple modules. For the
periodic lattice restriction $r=\mathbf1-2k$, this global sign is
\[
                       k\longmapsto\mathbf1-k.
\]
On lattice vectors this is the reflection in $B/2$.
It sends the lattice coset $\xi$ to $-\xi$. Thus, for nonzero odd
component,
\[
                 M_\xi\cong M_{-\xi}
                 \qquad\hbox{and hence}\qquad M_\xi'\cong M_\xi.
\]
Different regular grids can agree as sets of frame simples only
under the sign/gauge identifications described next. A coarse label
alone does not identify different parity lifts modulo $2\mathbb Z^N$.

\subsection{Alignment and canonical sign gauges}
The next proposition derives rational alignment from local
monodromy.  It applies to generic blocks with the stated
nonresonance conditions; integral and resonant blocks are treated
separately later.

\begin{proposition}[Alignment from the local monodromy spectra]
\label{g:prop:monodromy-alignment}
Suppose a generic singleton-block frame simple in a semisimple frame
restriction has momenta $p_i$, so that the preceding highest-weight
factorization applies, and at each elementary edge both rank-three BPZ local systems
are irreducible with nonresonant local exponents. Suppose in addition
that their middle eigenvalues cannot match a corner eigenvalue;
it suffices that
\[
 p_j\notin\pm\psi_j+\mathbb Z,\qquad
 -p_{j-1}/\psi_{j-1}
                \notin\pm\psi_{j-1}^{-1}+\mathbb Z.
\]
Then signs $\sigma_i$ and integers $n_i$ exist such that
\[
 p_j-\sigma_{j-1}p_{j-1}/\psi_{j-1}=n_j.
\]
There is a rational vector $r$ with
\[
 p_i=\psi_i r_i-\sigma_i r_{i+1},\qquad K_\sigma r=n.
\]
In particular its denominators are odd, and its sign holonomy is
periodic or antiperiodic. If this $r$ has no adjacent integral pair,
the preceding classification applies.
\end{proposition}
\begin{proof}
The same rational-derivative argument gives a nonzero invariant
tensor between the two irreducible local systems at each edge.
Their local eigenvalues at zero therefore match inversely.
The middle eigenvalues match each other because
$\psi_j+\psi_{j-1}^{-1}=\kappa_j\in\mathbb Z$.
The stated exclusions rule out matching middle with corner.
Matching the two remaining eigenvalues gives
$p_j\pm q_j\in\mathbb Z$, with
$q_j=-p_{j-1}/\psi_{j-1}$, proving the first assertion.

Define linear maps
\[
 (Dr)_i=\psi_i r_i-\sigma_i r_{i+1},\qquad
 (Lp)_j=p_j-\sigma_{j-1}p_{j-1}/\psi_{j-1}.
\]
The cyclic parameter identity gives $LD=K_\sigma$.
Since $K_\sigma$ is invertible, both square matrices $L,D$ are
invertible. Setting $r=K_\sigma^{-1}n$ gives $Dr=p$.
The matrix $K_\sigma$ is integral with odd determinant, so
$r\in\mathbb Q^N$ has odd denominators.
\end{proof}

\begin{corollary}[Exhaustion of the regular rational sectors]
\label{g:cor:regular-rational-exhaustion}
Let $X$ be a simple $W_N$-module containing a frame simple with
rational Kac pairs of odd denominator, at least one entry of each
pair being nonintegral. Then $X$ is one of the regular periodic
or antiperiodic grid modules just classified.
\end{corollary}
\begin{proof}
The generic frame-semisimplicity theorem applies. For both the row
and reciprocal-column parametrizations, the rational-label
irreducibility and nonresonance criterion holds at every site.
Proposition~\ref{g:prop:monodromy-alignment} therefore gives a
canonical periodic or antiperiodic rational vector
$r\in K_\pm^{-1}\mathbb Z^N$. Irrationality identifies its Kac
pairs with the original pairs up to simultaneous sign, so no
adjacent coordinates are both integral. Its odd-discriminant
component is consequently nonzero. The periodic Smith form gives
one of the nonzero affine-progressions modulo $N$; the antiperiodic
Smith form gives $s(N-2i)/(N^2+4)$ modulo $\mathbb Z$ with
$s\ne0$. The grid classification and uniqueness theorems apply.
\end{proof}

For the regular rational labels used in this section,
$2p_i$ and $2p_i/\psi_i$ are nonintegral. Hence precisely one
of the two possible edge-sign matching conditions can hold.
Changing the independently chosen momentum signs
$p_i\mapsto\varepsilon_ip_i$ consequently changes the unique
edge signs by
\[
                  \sigma_i\mapsto
                    \varepsilon_i\varepsilon_{i+1}\sigma_i.
\]
Their product is invariant, including at points where a coordinate
$r_i$ is zero or integral. Once $\sigma$ is fixed to the standard
periodic or antiperiodic representative, the only remaining gauge
is global sign. This establishes the sign identification using
local monodromy and integrality; it does not assume sign propagation
through nonzero coordinates.


\section{Refined modular traces after regularity}
\label{sec:ordinary-jets}

Regularity gives finitely many simple $W_N$-modules and ordinary
trace functions on each of them.  We refine these traces by recording
the commuting partial conformal weights.  To avoid a separate
convergence problem in several variables, we work with formal
Taylor germs along the diagonal.  The Ward identities determine
all their coefficients from convergent one-variable traces.

Let $V$ be a regular, $C_2$-cofinite CFT-type VOA with a commuting conformal frame
$\omega=\sum_{i=0}^{N-1}e_i$, and write $L_i(n)$ for its modes and $c_i$
for its central charges.  For a simple $V$-module $M$, put
\[
 F_M(v;\boldsymbol\tau)=
 \operatorname{Tr}_M\left(o(v)
          \prod_i e^{2\pi i\tau_i(L_i(0)-c_i/24)}\right).
\]
Initially this expression is interpreted coefficientwise in the total
$q$-grading and as a formal Taylor germ at
$\boldsymbol\tau=(\tau,\ldots,\tau)$.  The Ward recursion below
proves that each Taylor coefficient is holomorphic for $\tau$ in the
upper half-plane.  All the partial zero modes commute and preserve
each finite-dimensional total-weight space.  For a vector that is not
ordinary homogeneous, $o(v)$ is extended linearly from its homogeneous
components. The analytic theorem in \cite[Theorem 1]{km}
assumes convergence on the full product of upper half-planes.
We use its Ward recursion at the formal diagonal germ, and prove
that version below; no such global convergence hypothesis is
silently applied to the unknown modules.

For ordinary homogeneous $v$, use the square-bracket coordinate
\[
 Y[v,z]=Y(v,e^{2\pi iz}-1)e^{2\pi iz\operatorname{wt}(v)},
 \qquad Y[v,z]=\sum_{m\in\mathbb Z}v[m]z^{-m-1},
\]
extended linearly, and set
\[
 \widetilde e_i=(2\pi i)^2(e_i-c_i\one/24),\qquad
 L_i[m]=\widetilde e_i[m+1].
\]
The Eisenstein normalization is
\[
 G_{2k}(\tau)=2\zeta(2k)+
 \frac{2(2\pi i)^{2k}}{(2k-1)!}
 \sum_{n\geq1}\sigma_{2k-1}(n)e^{2\pi in\tau}.
\]
For a bracket-frame homogeneous vacuum descendant $v$ of partial
weights $w_i$, the trace Ward identities are
\begin{align}
 F_M(L_i[-1]v)&=0,\notag\\
 F_M(L_i[-2]v)&=2\pi i\partial_iF_M(v)
       +w_iG_2(\tau_i)F_M(v)
       +\sum_{k\ge2}G_{2k}(\tau_i)F_M(L_i[2k-2]v).
 \label{eq:ordinary-ward}
\end{align}
Only finitely many terms occur for each $v$.  The residue proof uses
mode commutators and cyclicity of trace on finite total-weight spaces;
it is valid as an identity of formal germs.  For a Taylor coefficient of a fixed order, solve this identity for
$\partial_iF_M(v)$ and iterate.  Only finitely many frame descendants
are introduced at that order, and the diagonal values are convergent
ordinary one-point functions.  Derivatives of the Eisenstein series
are holomorphic as well.  This proves the claimed convergence of
individual Taylor coefficients, without convergence of the full
multivariable series.  The identities are
\cite[Lemmas 4--5 and equation (2.5)]{km}.

\begin{lemma}[Formal covariance with the ordinary modular matrix]
\label{lem:ordinary-jets}
Let $v$ be a bracket-frame homogeneous vector in the vacuum
subalgebra generated by the commuting $e_i$, with partial weights
$w_i$. Let $S_{M,X}^{\mathrm{an}}$ be the coefficients in the full
insertion-dependent one-point $S$-transformation of $V$.
Then, as formal germs along the diagonal,
\[
 F_M(v;-1/\tau_0,\ldots,-1/\tau_{N-1})
    =\prod_i\tau_i^{w_i}
       \sum_{X\in\Irr(V)} S_{M,X}^{\mathrm{an}}
                              F_X(v;\boldsymbol\tau).
\]
\end{lemma}
\begin{proof}
We prove the analogous identity for each
$\gamma=\left(\begin{smallmatrix}a&b\\c&d\end{smallmatrix}\right)
\in\operatorname{SL}_2(\mathbb Z)$.
Put $j_i=c\tau_i+d$ and subtract the proposed right-hand side from
$\prod_i j_i^{-w_i}F_M(v;\gamma\boldsymbol\tau)$, calling the result
$H_v$.  Full one-point modularity makes $H_v$ vanish on the diagonal
for every frame vacuum descendant $v$.  In this normalization,
\[
 G_2(\gamma\tau)=(c\tau+d)^2G_2(\tau)-2\pi ic(c\tau+d).
\]
The anomaly cancels the derivative of $j_i^{-w_i}$ in the covariant
operator $2\pi i\partial_i+w_iG_2$.  Higher Eisenstein series have
their usual modular weights.  Thus \eqref{eq:ordinary-ward} implies
\[
 2\pi i\partial_iH_v
 =H_{L_i[-2]v}-w_iG_2(\tau_i)H_v
       -\sum_{k\ge2}G_{2k}(\tau_i)H_{L_i[2k-2]v}.
\]
Induct simultaneously for every frame descendant $v$ on the total
order of a multi-index of partial derivatives.  Order zero is the
ordinary modular identity on the diagonal.  Apply a derivative of
order $m$ to the displayed recursion.  Every term on its right is an
Eisenstein coefficient times a derivative of some $H_u$ of order at
most $m$, and hence vanishes on the diagonal by the induction
hypothesis.  Thus every derivative of order $m+1$ vanishes there.
All transverse Taylor coefficients therefore vanish, proving the
formal identity.  This is the formal Ward-recursion argument
underlying \cite{km}; it requires convergence of the individual
one-variable coefficients only.
\end{proof}

\begin{lemma}[Ordinary trace support in a frame block]
\label{lem:ordinary-support}
For a grading-restricted $W_N$-module, multiplication of its refined
ordinary character by $\eta_\Pi=\prod_i\eta(\tau_i)$ leaves only
the joint endpoints of the irrational Virasoro linkage blocks in
its frame restriction.  In a singleton block with irreducible
Verma tensor of joint highest weight $\lambda$, its contribution is
\[
 m_\lambda\prod_iq_i^{\lambda_i-(c_i-1)/24},
 \qquad m_\lambda\in\mathbb Z_{\ge0}.
\]
For a finite linear combination of refined characters of simple
$W_N$-modules, a nonzero coefficient at such a \emph{singleton-block}
endpoint implies that at least one simple module with a nonzero
coefficient in that linear combination contains the corresponding
frame simple.  This assertion is not made for a subtracted endpoint
of a two-element linkage block.
\end{lemma}
\begin{proof}
The frame-block lemma gives a finite-length restriction in each
fixed joint block, with finitely many simple types.  Ordinary trace
is additive on a composition series, including when partial zero
modes are not semisimple: in a simultaneous upper triangular basis
the exponential has the prescribed diagonal eigenvalues, and the
strictly upper triangular terms have zero trace.  Its block trace
is therefore the sum of the ordinary frame-simple characters with
their composition multiplicities.

At an irrational Virasoro parameter, an ordinary simple is either
an irreducible Verma module or the quotient of a Verma module by
its single irreducible singular submodule.  Its character is
accordingly one Verma character or the difference of two.  Multiplication
by eta cancels the oscillator factor, leaving just the one or two
linkage endpoints.  Tensoring gives the stated support description.
A singleton block with positive multiplicity is a nonzero
finite-length module and has a simple socle of that unique type.
Only the same linkage block can contribute that endpoint.  Hence,
even in a linear combination with complex coefficients, a nonzero
singleton endpoint requires a nonzero block in one of its summands,
proving the stated submodule conclusion.

The possibly infinite sum of blocks causes no ambiguity.  At each
fixed total $q$-degree only finitely many joint weight spaces occur.
To see that the formal germ distinguishes them, fix a total
exponent and let $\lambda^{(1)},\ldots,\lambda^{(r)}$ be the distinct
joint exponent tuples occurring there. They have the same coordinate
sum. Choose a transverse direction $u$ with $\sum_i u_i=0$ such that
$u\cdot\lambda^{(1)},\ldots,u\cdot\lambda^{(r)}$ are distinct;
this avoids finitely many proper hyperplanes. The directional
derivatives of orders $0,\ldots,r-1$ of
\[
 \sum_{a=1}^r c_a\exp(2\pi i z\,u\cdot\lambda^{(a)})
\]
at $z=0$ form a Vandermonde system in the $c_a$, so its vanishing
forces every $c_a=0$. The same argument applies after multiplication
by eta, which contributes only finitely many terms at each bounded
degree. Coefficients can therefore be compared degree by degree in
the formal germs of Lemma~\ref{lem:ordinary-jets}.
\end{proof}

The theta transforms used below have real, positive-definite
quadratic exponents.  Comparing all transverse derivatives
separates joint weight tuples even when their total conformal
weights coincide.  This distinction is necessary at composite
$N^2+4$: a nonzero discriminant label can have square zero
modulo $N^2+4$, so its conformal residue alone need not separate
it from the zero label.


\subsection{The trace normalization and categorical dimensions}
\label{sec:analytic-normalization}

We next relate the analytic modular matrix to categorical dimensions,
fixing compatible conventions for the Hopf link.  The argument applies
to a simple, self-contragredient,
regular, $C_2$-cofinite, nonnegatively graded VOA $V$; it does not require
$V_0=\mathbb C\mathbf1$. Its finite module category is modular by
\cite[Corollary 5.19]{mcrae}.

Let $S^{\rm an}$ be the matrix acting on the full insertion-dependent
ordinary traces. Its square is the contragredient permutation $C$.
Here this is an identity of full trace functions, not an inference from
scalar characters. Indeed the defining contragredient vertex operator
uses the coordinate inversion $x\mapsto x^{-1}$. At the insertion
$x=1$, the coordinate $f(z)=e^{2\pi iz}-1$ satisfies
\[
 f(-z)=-\frac{f(z)}{1+f(z)}.
\]
Thus inversion in the exponential coordinate becomes $z\mapsto-z$.
For a bracket-homogeneous insertion of weight $w$, the contragredient
trace is consequently $(-1)^w$ times the original trace. This is
exactly the slash action of $-I\in\operatorname{SL}_2(\mathbb Z)$.
Linear independence with all insertions identifies its matrix with
$C$, proving $(S^{\rm an})^2=C$.

Write $d_X=\dim X$ and $\Delta=\sum_Xd_X^2$.
Choose the orientation of the categorical Hopf-link matrix
$\widetilde S$ compatible with the analytic generator. For clarity,
McRae's geometrically modified convention uses
$(-1/\tau)^{L(0)}$, whereas the ordinary bracket convention uses
$\tau^{-w}$. These represent the two inverse lifts of
$\tau\mapsto-1/\tau$; their matrices differ by $C$.
Reversing the Hopf-link orientation likewise dualizes one index.
Both choices of $\widetilde S$ are symmetric and satisfy
\[
 \widetilde S_{\mathbf1,X}=d_X,\qquad
 \widetilde S^2=\Delta C.
\]
These are the usual ribbon-category identities
\cite[Section 4]{ko}, with no positivity condition on the $d_X$.

In the matching convention, apply the open-Hopf-link formula
\cite[Theorem 5.2, equation (5.5)]{mcrae} directly to $V$.
All simple $V$-modules are rigid, and regularity makes the modular
transform of the vacuum trace a sum of ordinary simple-module traces.
On a simple $Y$, the open Hopf link is the scalar
$\widetilde S_{X,Y}/d_Y$. Linear independence of the full
insertion-dependent traces therefore gives
\[
 S^{\rm an}_{X,Y}
  =S^{\rm an}_{\mathbf1,Y}\frac{\widetilde S_{X,Y}}{d_Y}.
\]
This argument applies separately to $V=W_N$ and to the deformed
lattice VOA $V=A$: both are simple, self-contragredient, regular,
$C_2$-cofinite and nonnegatively graded, as proved above. The theorem
requires no one-dimensional degree-zero space. In particular this
application to $A$ does not use the orbifold-specific extension
setting of \cite[Remark 5.21]{mcrae}. Set
\[
 D=\operatorname{diag}\left(
       S^{\rm an}_{\mathbf1,Y}/d_Y\right)_Y.
\]
Spherical dimensions of simple objects are nonzero, so $D$ is
well-defined.  The preceding identity reads
$S^{\rm an}=\widetilde SD$.
Using the two square identities gives
\[
 D\widetilde SD=\Delta^{-1}\widetilde S.
\]
Its $(\mathbf1,\mathbf1)$ entry gives
$D_{\mathbf1}^2=\Delta^{-1}$; its $(\mathbf1,Y)$ entries then give
$D_Y=D_{\mathbf1}$ for every $Y$.
Consequently
\begin{equation}\label{eq:analytic-vacuum-row}
 S^{\rm an}_{\mathbf1,X}
    =d_X S^{\rm an}_{\mathbf1,\mathbf1},\qquad
 S^{\rm an}_{\mathbf1,\mathbf1}\ne0.
\end{equation}
In particular the analytic matrix is symmetric, in either compatible
orientation. If $J$ centralizes $X$, its closed Hopf link is
$d_Jd_X$ in both orientations, and hence
\begin{equation}\label{eq:analytic-centralizer-row}
 \frac{S^{\rm an}_{J,X}}{S^{\rm an}_{\mathbf1,X}}=d_J.
\end{equation}
These two identities are independent of the compatible orientation
chosen.  They provide the lattice-index calculation and the
categorical dimension ratios extracted from the refined parity
rows, without requiring a unique module of minimal conformal
weight.

\section{The multivariable theta identity and its Fourier transform}
\label{sec:multivariable-theta}
\label{g:sec:theta}
The difference of two theta quotients is the multivariable analogue
of Flohr's one-variable character formula
\cite[equations (4.2)--(4.5) and Table 4.2]{flohr1993}.
The periodic and antiperiodic metric sectors are the two sectors
in the Evans--Gannon grafting \cite[Section 3.3]{evansgannon2011},
proved for untwisted HI centers in \cite[Theorem 5.1]{gsy}.
Here the separate frame variables are retained to resolve labels
that the single total conformal weight does not separate.
Let $\sigma=(\sigma_i)$ have all signs positive, or just its closing
sign negative. Write these two choices as $+$ and $-$, and define
\[
 (B_\sigma r)_i=\sqrt{\psi_i}\,r_i
                  -\frac{\sigma_i r_{i+1}}{\sqrt{\psi_i}}.
\]
Direct multiplication, using the cyclic parameter relations,
gives
\[
 B_\sigma^TB_\sigma=K_\sigma,\qquad
 \det B_+=t^{-1}-t=N,\qquad
 \det B_-=t^{-1}+t=\sqrt\mu.
\]
For $v\in\frac12\Z^N$, put
\[
 \Theta_{\sigma,v}(\boldsymbol\tau)
 =\sum_{x\in\Z^N+v}
       \exp\left(2\pi i\sum_i\tau_i(B_\sigma x)_i^2\right),
 \qquad
 \eta_\Pi(\boldsymbol\tau)=\prod_i\eta(\tau_i).
\]

\begin{lemma}[All-parity refined identity]\label{g:lem:theta}
For $v_\epsilon=(\one+\epsilon)/2$ modulo $\Z^N$,
\[
 \ch J_\epsilon(\boldsymbol\tau)
   =\frac{\Theta_{+,v_\epsilon}(\boldsymbol\tau)
                  -\Theta_{-,v_\epsilon}(\boldsymbol\tau)}
                  {2\eta_\Pi(\boldsymbol\tau)}.
\]
\end{lemma}
\begin{proof}
Expand the product of generic Virasoro character differences in the
known frame decomposition of $J_\epsilon$. Its positive labels have
parity $1+\epsilon_i$. Each term is indexed by edge signs. Independent
changes of signs of the labels gauge all edge signs except their
product; they preserve \emph{each} partial exponent, so separate
$\tau_i$ are permitted. Overall label reversal gives the factor two.

Completing positive labels to full signed theta sums introduces zero
coordinates for even labels. They cancel. If $j$ is the least zero
coordinate, reverse the labels at $0,\ldots,j-1$. This toggles the
closing sign and the sign on an edge incident to that zero, preserving
each partial square. If $j=0$ the closing cross term already vanishes.
This is an involutive, parity-preserving cancellation of the extra
terms.
\end{proof}

Define
\[
 \mathcal F_a(\boldsymbol\tau)=
 2^{-N}\sum_{\epsilon\in(\Z/2)^N}
       (-1)^{\epsilon\cdot a}
       \ch J_\epsilon(-1/\tau_0,\ldots,-1/\tau_{N-1}).
\]
Poisson summation with the matrix
$B_\sigma^T\diag(\tau_i)B_\sigma$, followed by finite Fourier
orthogonality, gives
\begin{equation}\label{g:eq:poisson}
\begin{split}
 \eta_\Pi\mathcal F_a
 ={}&(-1)^{|a|}2^{-N/2-1}
 \left[
 \frac1N\sum_{\substack{z\in\Z^N\\z\equiv a\ (2)}}
   e^{\,2\pi i\sum_i\tau_i(B_+^{-T}z)_i^2/4}\right.\\[-2pt]
 &\hspace{34mm}\left.
 -\frac1{\sqrt\mu}
 \sum_{\substack{z\in\Z^N\\z\equiv a\ (2)}}
   e^{\,2\pi i\sum_i\tau_i(B_-^{-T}z)_i^2/4}
 \right].
\end{split}
\end{equation}
The determinant factors cancel the eta automorphy factors.
The shift phase is $(-1)^{z\cdot(\one+\epsilon)}$, so the Fourier
sum keeps exactly $z\equiv a\pmod2$.
All coefficients within either bracketed branch have the same sign.
The theta expressions converge on $\mathbb H^N$; for comparison to
ordinary traces, only their formal diagonal germs are needed.

\section{Exceptional modules for all nonzero discriminant labels}
\label{g:sec:exceptional}
Set
\[
 v_i=\frac{N-2i}{\mu},\qquad
 v_N=-v_0=-\frac N\mu.
\]
Direct multiplication gives
\[
 K_-v=e_0.
\]
The interior equations are second differences, the final equation is
$2(N-2N+2)-(N-2N+4)+N=0$, and the zeroth equation is
$(N+2)N-(N-2)+(2-N)=\mu$.
The class of $v$ has order $\mu$, since
$v_{(N-1)/2}=1/\mu$. It generates the odd discriminant group.

For $s\not\equiv0\pmod\mu$ define the joint highest weights
\begin{equation}\label{g:eq:exceptionaltuple}
 \lambda_i(s)=h_i(sv_i,sv_{i+1}),\qquad 0\leq i<N,
\end{equation}
where the definition of $v_N$ includes the closing minus sign.
Let
\[
 \mathcal Q(s)=
 \left(\left[(sv_i\bmod\Z,sv_{i+1}\bmod\Z)\right]_{\pm}\right)_i.
\]
The sign quotient at each site means simultaneous sign reversal of
both entries.

\begin{lemma}[Genericity and separation]\label{g:lem:arithmetic}
For $s\not\equiv0\pmod\mu$, every factor in
$\bigotimes_i M(c_i,\lambda_i(s))$ is a nondegenerate irreducible
Verma module. Moreover
\[
 \mathcal Q(s)=\mathcal Q(s')
       \quad\Longleftrightarrow\quad s'\equiv\pm s\pmod\mu.
\]
Every such tuple is disjoint from all periodic lattice coarse tuples.
\end{lemma}
\begin{proof}
No adjacent entries $sv_i,sv_{i+1}$ can both be integral: their
difference is $2s/\mu$, also at the closing edge with the signed
extension. Since $\mu$ is odd, both integral would imply $\mu\mid s$.
For rational $R,S$, an equality
$h_i(R,S)=h_i(m,n)$ with positive integral $m,n$ forces
$(R,S)=\pm(m,n)$, by taking square roots and using irrationality of
$\psi_i$. Thus every factor is nondegenerate.

At the middle edge $j=(N-1)/2$, the pair is exactly
\[
       (sv_j,sv_{j+1})=(s/\mu,-s/\mu).
\]
Equality of this coarse pair already forces $s'\equiv\pm s$.
The converse holds at every site.
Periodic odd lattice coordinates have denominators dividing $N$.
Their coarse Kac pairs therefore cannot contain $s/\mu\notin\Z$,
because $\gcd(N,\mu)=1$. Parity shifts change Kac labels by integers
and preserve these coarse tuples.
\end{proof}

\begin{theorem}[All exceptional labels]\label{g:thm:alllabels}
For every integer $s$ with $\mu\nmid s$ there is a simple
grading-restricted $W_N$-module $X_s$ whose frame restriction contains
\[
       \bigotimes_{i=0}^{N-1} M(c_i,\lambda_i(s)).
\]
At least $(\mu-1)/2$ distinct parity orbits are thereby forced,
indexed by nonzero $s$ modulo global sign, and disjoint from the
periodic orbits.
\end{theorem}
\begin{proof}
In \eqref{g:eq:poisson} take $z=se_0$ and $a=z\bmod2$.
Since
$B_\sigma^{-T}=B_\sigma K_\sigma^{-1}$,
its antiperiodic joint exponent is
\[
       \left((B_-sv)_i^2/4\right)_i
       =\left(\lambda_i(s)-(c_i-1)/24\right)_i.
\]
It cannot cancel inside the antiperiodic sum, where all coefficients
have one sign.

Suppose a periodic term had the same joint exponent. Its rational
labels are the entries of $K_+^{-1}z'$ with denominator dividing
$N^2$. Equality of each squared Kac expression, by irrationality
of $\psi_i$, forces simultaneous-sign equality of the corresponding
rational pairs. The middle antiperiodic pair then forces
$s/\mu\in\frac1{N^2}\Z$, contrary to $\mu\nmid s$ and
$\gcd(N^2,\mu)=1$. Thus the tuple survives the entire transform.
This argument also works when $s^2\equiv0\pmod\mu$.

Apply full ordinary one-point modularity and
Lemma~\ref{lem:ordinary-jets} to the parity-module traces.
The result is a finite linear combination of ordinary simple-module
traces with identical formal frame jets. Lemma~\ref{lem:ordinary-support}
and finite exponential independence at each total $q$ exponent force
the surviving tuple to be a linkage endpoint of one of these modules.
By Lemma~\ref{g:lem:arithmetic} its joint block is a singleton. Its
finite-length block component has a nonzero simple socle, necessarily
the displayed irreducible Verma tensor. Thus that simple $W_N$-module
contains the stated frame module.

The necessary level-three BPZ constraints for the elementary
generators change rational Kac labels by $0,\pm2$; the parity-current
constraints change them by integers. The ordered coarse tuple modulo simultaneous sign is preserved by
these operations and by frame linkage. It therefore labels an
entire simple $W$-module and its parity orbit.
Lemma~\ref{g:lem:arithmetic} supplies all the claimed distinctions.
\end{proof}


The same proof gives every exceptional parity lift directly. For
$u\in\mathbb Z^N$, take $z=se_0+K_-u$ in \eqref{g:eq:poisson}.
Its antiperiodic label vector is $r=sv+u$. The no-adjacent-integral
argument and the middle-pair noncancellation proof are unchanged.
Thus a simple module contains the corresponding generic frame tensor
for every $u$. Reducing $u$ modulo $2\mathbb Z^N$ supplies all parity
grids. The multiplicity-one and uniqueness theorems give their full frame
decompositions and theta characters. Fusion with $J_\epsilon$ shifts
the rational grid by $\epsilon$ modulo $2\mathbb Z^N$, whereas
the elementary roots shift it by even integers. For a nonzero
odd-order discriminant label, a parity shift cannot identify the
grid with its global negative: that would make twice the odd-order
label zero. Thus the parity action is free, and the $2^N$ lifts
of a fixed nonzero label form one parity orbit. The local monodromy sign
argument identifies two such grids precisely by global sign after
fixing the antiperiodic gauge. There are consequently
$2^N(\mu-1)/2$ exceptional simples, or $(\mu-1)/2=(N^2+3)/2$
parity orbits, in the detected family. Exhaustion of all $W_N$-simples
will follow from the nonnegative dimension budget below.

For example, the middle pair separates labels $1$ and $16$ at $N=9$
even though their squares agree modulo $85$. At $N=11$ the label
$25$ is nonzero modulo $125$ and has square zero, but its middle
pair $(1/5,-1/5)$ remains nonintegral. The argument covers both
phenomena uniformly and does not count square residues alone.

\Needspace{12\baselineskip}
\part{Condensation and the Haagerup--Izumi center}

\section{The lattice extension and its categorical index}
\label{r:sec:extension}

By Section~\ref{r:sec:projectivity}, the category
$\Cat=\Rep(W_N)$ is modular. We now regard the conformal lattice
extension $A$ as an algebra object in $\Cat$. The first task is to
show that its categorical dimension is nonzero; separability will
then follow from the extension theorems.

\subsection{The deformation preserves the braided category}

Let \(B=\sum_i\beta_i\) and \(\eps=e^B\). The two conformal vectors
have the form
\[
 \omega_A=\omega_L+Dv,\qquad
 v=\frac12B(-1)\one+\lambda D\eps .
\]
Here \(v_0\omega_L=0\): the Heisenberg zero mode annihilates
\(\omega_L\), and \((D\eps)_0=0\). Both conformal gradings are
nonnegative by the construction. McRae's Proposition 2.17
\cite{mcrae} therefore identifies the entire braided tensor
structures for the two conformal vectors, including associators.
The twist can change, as explained in his Remark 2.18.

The possible failure of CFT type for \(A\) causes no modularity
problem. McRae's Corollary 5.19 applies to a simple,
self-contragredient, \(C_2\)-cofinite, nonnegatively graded VOA
with semisimple Zhu algebra. Regularity supplies the last property.
Thus \(\Rep(A)\) is modular and its braided category is the
pointed lattice category.

Its simple objects are indexed by
\[
 D_L=L^*/L,\qquad |D_L|=\parityrank N^2,\qquad
 D_L[2]=(\tfrac12\Z^N)/\Z^N,\qquad |D_L[2]|=\parityrank.
\]
Every entry of its categorical $S$-matrix is nonzero, because all
simple objects are invertible. Duality acts by $\xi\mapsto-\xi$:
the conformal shift changes the contragredient charge by $B\in L$,
which does not change the discriminant class.

\subsection{Disjointness of the integral branches}

The lattice charge labels at coset \(\xi\) are
\[
                    r_i=1-2(\xi_i+n_i),\qquad n_i\in\Z.
\]
For \(\xi=u/2\in D_L[2]\), all labels are integral, with parity
\(r=\one-u\pmod2\). Irrational Virasoro linkage changes their signs
but preserves their parities. The root label shifts are \(0,\pm2\),
so they preserve the same invariant.

The exact block projections and triangular charge
argument consequently separate the simple supports of the
restrictions of all distinct two-torsion lattice modules. They
also separate these restrictions from the nonzero odd branches.
On the latter branches, the periodic-grid classification gives
\[
 M_\xi\big|_W\text{ simple},\qquad
 M_\xi\big|_W\simeq M_\eta\big|_W
       \Longleftrightarrow \eta=\xi\text{ or }-\xi.
\]
The two-primary coordinate is retained in this equivalence.

In particular, \(W\) occurs only in the restriction of the zero
lattice module \(A\). Its multiplicity there is one. For example,
a homomorphism \(W\to A\) sends \(\one\) to a vector annihilated by
translation; in this underlying positive lattice vertex algebra,
\(\ker D=\C\one\).

\subsection{A nonzero index before invoking separability}

Let \(b_{\xi,Z}\) be the multiplicity of the \(W\)-simple \(Z\) in
the restriction of \(M_\xi\). The inclusion is conformal, so each
restriction is grading restricted. The $C_2$-cofinite module category
has finite length, and regularity makes these restrictions semisimple;
thus $b$ is a finite matrix. The preceding paragraph gives
\[
                         b_{\xi,\one}=\delta_{\xi,0}.
\]
For every insertion \(w\in W\), restriction of ordinary one-point
traces gives
\[
 F^A_\xi(w;\tau)=\sum_Zb_{\xi,Z}F^W_Z(w;\tau).
\]
The inclusion is conformal, so its coordinate changes agree.
Modular transformation and linear independence of the full
insertion-dependent \(W\)-trace functions imply
\begin{equation}\label{r:eq:analytic}
                  S_L^{\mathrm{an}}b=bS_W^{\mathrm{an}}.
\end{equation}
Scalar characters alone would not justify this inference.

Apply the normalization argument of
Section~\ref{sec:analytic-normalization} separately to $W$ and $A$.
Symmetry and \eqref{eq:analytic-vacuum-row} give, for either VOA,
\[
 (S^{\mathrm{an}})_{M,V}
       =\dim_{\Rep(V)}(M)(S^{\mathrm{an}})_{V,V}.
\]
Both vacuum entries are nonzero by the same argument.
Taking the zero row and vacuum column of \eqref{r:eq:analytic} yields
\begin{equation}\label{r:eq:index}
       \dim_{\Cat}A
         =\frac{(S_L^{\mathrm{an}})_{0,0}}
                 {(S_W^{\mathrm{an}})_{0,0}}\ne0.
\end{equation}
This is a categorical index identity. It uses no
high-temperature limit and no minimal-weight conjecture.

\subsection{Connected \'etale algebra and local modules}

The conformal extension gives a commutative algebra object \(A\)
in \(\Cat\), connected by the vacuum multiplicity above and with
trivial twist by integer grading. It is simple as its own algebra
module because the vertex algebra \(A\) is simple. Equation
\eqref{r:eq:index} and Kirillov--Ostrik's Lemma 1.20 \cite{ko}
make it a rigid algebra. Their Theorem 3.3 supplies an algebra
bimodule splitting of multiplication, so \(A\) is separable.

The induced category \(\Cat_A\) is spherical, and induction
preserves spherical dimensions by \cite[Theorem 1.18 and
Remark 1.19]{ko}. The local module category is the
VOA category \(\Rep(A)\), with its braiding and tensor structure,
by \cite[Theorem 3.65]{ckm}; the twists agree on the underlying
modules. As a further check, the modern extension theorem
\cite[Theorem 4.10]{cmsy} independently gives fusion-category
semisimplicity in this nonnegatively graded setting.

\section{Branching determines the rank and parity locality}
\label{sec:branching-rank}

\subsection{Determining the branching Gram matrix}

For \(\xi\in D_L[2]\), set
\[
 d_\xi=\dim\End_{\Cat}(M_\xi|_W).
\]
Disjointness and the opposite-coset identification give
\begin{equation}\label{r:eq:gram}
 bb^t=I+C_L+D_0,
\end{equation}
where \(C_L\) is inversion on \(D_L\), and \(D_0\) is diagonal,
supported on \(D_L[2]\), with entries \(d_\xi-2\).
This formula includes all unknown integral \(W\)-simples in the
columns of \(b\).

Use the analytic matrices in \eqref{r:eq:analytic}; both are symmetric
by Section~\ref{sec:analytic-normalization}. Transposing that identity
gives $b^tS_L^{\rm an}=S_W^{\rm an}b^t$, and hence
\[
 S_L^{\rm an}bb^t=bS_W^{\rm an}b^t=bb^tS_L^{\rm an}.
\]
Write $S_L=S_L^{\rm an}$ in the following calculation.
Since \(I+C_L\) commutes with \(S_L\), so does \(D_0\).
Thus
\[
 (D_{0,\eta}-D_{0,\xi})(S_L)_{\xi,\eta}=0
                     \qquad(\xi,\eta\in D_L).
\]
Every pointed \(S\)-entry is nonzero. Hence \(D_0\) is scalar.
It vanishes on the nonzero odd labels, so it vanishes everywhere.
Therefore
\[
                         d_\xi=2\quad(\xi\in D_L[2]).
\]
At \(\xi=0\), connectedness and the simple top-zero constituent \(T\)
constructed from \(A/W\) now force
\begin{equation}\label{r:eq:A1T}
                              A=\one\oplus T.
\end{equation}
No further constituent of the lattice restriction remains.

\subsection{The full module-category rank}

We use right modules with the tensor product convention of
\cite[Section 3.3]{dmno}: a right module acquires its left action
using the inverse braiding. This fixes the orientation in the center
equivalences below.

\begin{lemma}\label{lem:branching-rank}
Let $\mathcal C$ be a nondegenerate braided fusion category over
$\mathbb C$, let $A$ be connected \'etale, and put
$\mathcal B=\mathcal C_A$ and $\mathcal E=\mathcal C_A^0$.
If $U:\mathcal E\to\mathcal C$ is restriction and
$b_{e,X}=\dim\operatorname{Hom}_{\mathcal C}(X,U(e))$, then
\[
 \operatorname{rank}\mathcal B
 =\sum_{e\in\operatorname{Irr}\mathcal E}
       \dim\operatorname{End}_{\mathcal C}U(e)
 =\operatorname{Tr}(bb^{\mathsf t}).
\]
No unitarity or positivity of categorical dimensions is required.
\end{lemma}
\begin{proof}
Write $I_{\mathcal B}$ for the right adjoint of
$\operatorname{For}:Z(\mathcal B)\to\mathcal B$.
The center equivalence
\cite[Corollary 3.30 and Remark 3.31(i)]{dmno}
identifies $Z(\mathcal B)$ with
$\mathcal C\boxtimes\mathcal E^{\mathrm{rev}}$.
Under its forgetful functor, $X\boxtimes e$ goes to
$F_A(X)\otimes_A e$.
Adjunction and rigidity therefore give
\[
 \operatorname{Hom}_{\mathcal B}(F_A(X)\otimes_A e,A)
 \cong\operatorname{Hom}_{\mathcal B}(F_A(X),e^*)
 \cong\operatorname{Hom}_{\mathcal C}(X,U(e^*)).
\]
Consequently
\[
 I_{\mathcal B}(A)
 \cong\bigoplus_{e\in\operatorname{Irr}\mathcal E}U(e^*)\boxtimes e,
 \qquad
 \dim\operatorname{End}_{Z(\mathcal B)}I_{\mathcal B}(A)
 =\sum_e\dim\operatorname{End}_{\mathcal C}U(e).
\]
On the other hand, the ordinary center-adjunction formula
\cite[Proposition 9.2.2]{egno} gives
\[
 \operatorname{For}I_{\mathcal B}(A)
 \cong\bigoplus_{M\in\operatorname{Irr}\mathcal B}M\otimes_A M^*.
\]
Adjunction now identifies the preceding endomorphism dimension with
\[
 \dim\operatorname{Hom}_{\mathcal B}
 (\operatorname{For}I_{\mathcal B}(A),A)
 =|\operatorname{Irr}\mathcal B|.
\]
The last equality in the statement follows by decomposing each $U(e)$.
\end{proof}


Applying Lemma~\ref{lem:branching-rank} to the algebra $A$ gives:
\begin{equation}\label{r:eq:rankfull}
  \operatorname{rk}(\Cat_A)=\operatorname{Tr}(bb^t)=\parityrank(N^2+1).
\end{equation}
The local category has $\parityrank N^2$ invertible simples, so
exactly $\parityrank$ simples of $\Cat_A$ are nonlocal. Here
``nonlocal'' means that the right $A$-action fails the trivial-monodromy
condition; this count is taken in the full fusion category of
right $A$-modules.

\subsection{An orbit argument forces the parity factor to be local}

If a nonlocal simple \(Q\) in \(\Cat_A\) were invertible, tensoring
it by all \(\parityrank N^2\) local invertibles would give \(\parityrank N^2\) distinct
nonlocal simples. Distinctness follows by cancellation; locality
of any one of them would imply locality of \(Q\). This contradicts
the count \(\parityrank\), since \(N\ge3\).

Every induced parity current \(F(J_\epsilon)\) is invertible
because induction is a tensor functor. Hence it is local.
Locality of the free module \(A\otimes J_\epsilon\) implies that
\(J_\epsilon\) centralizes \(A\): inserting the unit into the inner
\(A\)-factor of the locality equation gives precisely the trivial
monodromy equation. Thus
\begin{equation}\label{r:eq:central}
                              A\in\Neu=\Par'.
\end{equation}

Induction is fully faithful on \(\Par\). Indeed,
\[
 \Hom_{\Cat_A}(F(J_\epsilon),F(J_\eta))
   =\Hom_{\Cat}(J_\epsilon,A\otimes J_\eta)
   =\delta_{\epsilon,\eta},
\]
since \(A=\one\oplus T\) lies in \(\Par'\) and
\(\Par\cap\Par'=\mathrm{Vec}\). Its local image is the full
two-primary pointed subgroup of \(\Rep(A)\).

Algebra-module formation now factors:
\[
 \Cat_A\simeq\Par\boxtimes\Neu_A,\qquad
 \Cat_A^{\loc}\simeq\Par\boxtimes\Neu_A^{\loc}.
\]
Equation \eqref{r:eq:rankfull} gives
\begin{equation}\label{r:eq:rankneutral}
                    \rk(\Neu_A)=N^2+1.
\end{equation}

\subsection{The odd local ribbon category}

The integral map
\[
 m\longmapsto\bigl((m_i-m_{i+1})_i,m_0\bigr)
\]
identifies the lattice with
\[
                  L\simeq \sqrt2 A_{N-1}\oplus\langle2N\rangle.
\]
The difference vector has sum zero and, together with \(m_0\),
reconstructs \(m\), so this is an integral isometry, not merely a
rational diagonalization.

The two odd discriminant factors have quadratic values
\(-b^2/N\) and \(a^2/N\) modulo \(\Z\). Thus the odd metric group is
\[
 E=(\Z/N)^2,\qquad
 q_E(a,b)=\exp\!\left(\frac{2\pi i(a^2-b^2)}N\right).
\]
Because \(N\) is odd, \(x=a+b,\ y=a-b\) identifies this with
\(\exp(2\pi ixy/N)\).
The conformal shift pairs integrally with odd representatives:
\[
             (B/2,\xi)=\one^TK_+\xi=N\xi_0\in\Z.
\]
This follows directly from the odd discriminant representatives
\(\xi_i=a+ib\), \(Na=Nb=0\) modulo \(\Z\).
The \(D^2\eps\) term does not alter \(L(0)\).
Consequently \(\Neu_A^{\loc}\) is the standard hyperbolic
ribbon category, including its cyclic Lagrangian subgroup and its
complementary subgroup with trivial associator.

\section{Modular coefficients, multiplicities, and exhaustion}
\label{r:sec:budget}

\subsection{The remaining near-group multiplicity}

By \eqref{r:eq:rankneutral}, the category \(\Mod=\Neu_A\) has the
\(N^2\) local invertibles \(E\) and one additional simple \(R\).
The object \(R\) is fixed on both sides by \(E\), is self-dual,
and obeys
\[
                        R^2=\sum_{e\in E}e+mR
\]
for an integer \(m\ge0\). Rigidity makes each invertible
multiplicity exactly one.

Set \(r=\FPdim R\) and \(a=\FPdim A\). The standard dimension
formulas for algebra-module categories give
\[
 r^2=N^2+mr,\qquad a=1+r^2/N^2.
\]
For the local unit, restriction gives $U(\one)=A=\one\oplus T$;
for a nonzero local invertible $e$, restriction gives the simple
$U(e)=Y_{[e]}$. These constituents are disjoint by the branching
calculation. Consequently Frobenius reciprocity gives
\[
 [F(T):e]=[T:U(e)]=\delta_{e,\one}.
\]
There is only one nonlocal simple $R$, so
\[
 F(T)=\one+kR,\qquad k\in\NN.
\]
Since induction preserves FP dimension,
\[
 a-1=1+kr,\qquad a=2+kr,\qquad
 r^2=N^2+N^2kr,\qquad m=N^2k.
\]
For the periodic simple indexed by \([e]=\{e,-e\}\), reciprocity
similarly gives exactly the two local summands $e$ and $-e$, each
once. They are distinct because $N$ is odd and $e\ne0$.
The module-dimension formula gives
$\operatorname{FPdim}_{\Neu}U(e)=a\operatorname{FPdim}_{\Mod}e=a$.
Thus $\operatorname{FPdim}Y_{[e]}=a=2+kr$, and the remaining
multiplicity is $k$:
\begin{equation}\label{r:eq:induction}
 F(Y_{[e]})=e\oplus(-e)\oplus kR.
\end{equation}
Every other simple \(Z\) of \(\Neu\) has
\begin{equation}\label{r:eq:otherind}
                            F(Z)=n_ZR,\qquad n_Z\in\Z_{>0}.
\end{equation}
There is no local summand, and a tensor functor between fusion
categories cannot annihilate a nonzero object.

\subsection{A sum-of-squares dimension identity}

The global dimension identity is
\[
 N^2a^2
 =1+(a-1)^2+\frac{N^2-1}{2}a^2+r^2\sum_Zn_Z^2.
\]
Substituting \(a=2+kr\) and \(r^2=N^2(1+kr)\) gives
\begin{equation}\label{r:eq:budget}
                \boxed{\sum_Zn_Z^2
                    =2+\frac{N^2-1}{2}k^2.}
\end{equation}
The sum runs over all neutral simples other than $\mathbf1$, $T$,
and the periodic family. In particular, it includes both the
exceptional modules constructed in Section~\ref{g:sec:exceptional}
and any additional simples. Thus the identity is available before
exhaustion has been proved.

\subsection{Extracting the modular coefficient ratio}

For a regular module \(Z\), let \(a_Z=z\bmod2\) be its parity Fourier
label in the refined dual-lattice formula. The refined
Poisson identity has coefficient
\[
 2^{-N/2-1}N^{-1}(-1)^{(\one+\epsilon)\cdot a_Z}
 \quad\hbox{or}\quad
 -2^{-N/2-1}\mu^{-1/2}(-1)^{(\one+\epsilon)\cdot a_Z}.
\]
More explicitly, for a regular label choose $z_0$ in the
appropriate discriminant class and let
$\Gamma=z_0+2K_\pm\Z^N$ and $-\Gamma=-z_0+2K_\pm\Z^N$.
These grids are distinct: equality would make the nonzero odd
part of the class two-torsion. Their parity labels agree because
$-z_0\equiv z_0\pmod2$, and $z\mapsto-z$ preserves every squared
partial exponent. Consequently their contributions combine as
\[
 \frac{c_\pm}{2}
 \left(\sum_{z\in\Gamma}q^{\mathbf h(z)}
       +\sum_{z\in-\Gamma}q^{\mathbf h(z)}\right)
 =c_\pm\sum_{z\in\Gamma}q^{\mathbf h(z)},
\]
where $q^{\mathbf h(z)}$ denotes the multivariable exponential in
\eqref{g:eq:poisson}. Here
$c_+=2^{-N/2}N^{-1}(-1)^{(\one+\epsilon)\cdot a_Z}$ and
$c_-=-2^{-N/2}\mu^{-1/2}(-1)^{(\one+\epsilon)\cdot a_Z}$.
This accounts for the factor two before a module coefficient is
read off.
The exact nonintegral frame support uniquely isolates its regular
simple and excludes every unknown integral module. Hence
\begin{align}
 S^W_{J_\epsilon,Y}
    &=2^{-N/2}N^{-1}(-1)^{(\one+\epsilon)\cdot a_Y},\label{r:eq:SY}\\
 S^W_{J_\epsilon,X}
    &=-2^{-N/2}\mu^{-1/2}(-1)^{(\one+\epsilon)\cdot a_X}.
       \label{r:eq:SX}
\end{align}
Rationality ensures that the full one-point modular action here
is the ordinary modular action; the formal frame-germ covariance
uses the same coefficients. This argument does not depend on
linear independence of scalar characters.

For a representative \(Z\in\Neu\), trivial monodromy with \(\Par\)
implies
\[
             \frac{S^W_{J_\epsilon,Z}}{S^W_{\one,Z}}
                       =\dim J_\epsilon.
\]
Equations \eqref{r:eq:SY}--\eqref{r:eq:SX} therefore force the same
\(a_Z\) for every neutral representative in both regular families.
Removing the common parity vacuum entry gives the exact ratio
\begin{equation}\label{r:eq:Sratio}
                \boxed{\frac{S^\Neu_{\one,X_s}}
                              {S^\Neu_{\one,Y_{[e]}}}
                       =-\frac N{\sqrt{N^2+4}}.}
\end{equation}

\subsection{Determining the integer multiplicities}

Let \(x=\dim_{\Mod}R\). The local invertibles have spherical
dimension one: in a spherical category their dimensions are
\(\pm1\), and their group has odd order. Dimension-preserving
induction gives
\[
 x^2=N^2+N^2kx,\qquad
 \dim Y_{[e]}=2+kx,\qquad
 \dim Z=n_Zx.
\]
The number \(x\) may be negative. The quadratic relation implies
\[
        \left(\frac{x}{2+kx}\right)^2
                         =\frac{N^2}{N^2k^2+4}.
\]
For each of the \(M=(N^2+3)/2\) known exceptional simples,
equation \eqref{r:eq:Sratio} yields
\begin{equation}\label{r:eq:nconstraint}
                    n_s^2(N^2+4)=N^2k^2+4.
\end{equation}
All \(n_s\) are positive integers, hence have the same value \(n\).
Subtract their contribution from \eqref{r:eq:budget}:
\begin{equation}\label{r:eq:unused}
 \boxed{\sum_{\text{additional }Z}n_Z^2
 =2+\frac{N^2-1}{2}k^2-\frac{N^2+3}{2}n^2
 =\frac{2(1-k^2)}{N^2+4}.}
\end{equation}
Every term on the left is the square of a positive integer.
Nonnegativity therefore gives $k\leq1$. If $k=0$,
\eqref{r:eq:nconstraint} gives \(n^2=4/(N^2+4)<1\), impossible.
Hence $k=n=1$, and the sum over additional simples vanishes.
No additional simple can occur, since it would contribute a positive
integer square. This proves both the asserted multiplicities and
exhaustion.

The negative sign of \eqref{r:eq:Sratio} selects
\[
 x=-Nt,\qquad
 \dim T=t^2,\qquad
 \dim Y=1+t^2,\qquad
 \dim X=-Nt.
\]
The positive FP root is \(r=N\delta\), giving
\[
 \FPdim T=\delta^2,\qquad
 \FPdim Y=1+\delta^2,\qquad
 \FPdim X=N\delta.
\]

\begin{center}
\begin{tabular}{llll}
\toprule
Neutral simple & Number & Induction to \(\Neu_A\) & Spherical dimension\\
\midrule
\(\one\) & \(1\) & \(\one\) & \(1\)\\
\(T\) & \(1\) & \(\one\oplus R\) & \(t^2\)\\
\(Y_{[e]}\) & \((N^2-1)/2\) & \(e\oplus(-e)\oplus R\) & \(1+t^2\)\\
\(X_s\) & \((N^2+3)/2\) & \(R\) & \(-Nt\)\\
\bottomrule
\end{tabular}
\end{center}


\section{Reconstruction of the untwisted ribbon center}
\label{sec:actual-center}

\subsection{Cyclic de-equivariantization gives the HI rules}
\label{sec:cyclic-reconstruction}

The near-group setting originates in Izumi's work \cite{izumi}
and the constructions and doubles of Evans--Gannon
\cite{evansgannon2014,evansgannon2017}.
The equivariantization construction in \cite[Section 6]{gsy}
relates cyclic HI categories to near-group categories. For the
converse used here, the near-group fusion rules must be accompanied
by a specified central Tannakian algebra. The next theorem records
these hypotheses and reconstructs the fusion rules after
de-equivariantization.

\begin{theorem}\label{thm:cyclic-converse}
Let $N\geq3$ be odd. Let $\mathcal B$ be a complex fusion category
whose simple objects are the group
$E=(\mathbb Z/N)^2$ of invertibles and one object $R$, with
\[
 R^2=\bigoplus_{e\in E}e\oplus N^2R.
\]
Suppose a central Tannakian functor
$\operatorname{Rep}(\Gamma)\to Z(\mathcal B)$, with
$\Gamma\cong\mathbb Z/N$, has an underlying fully faithful functor
identifying its simple objects with a subgroup $H\subset E$ of order
$N$. Let $D_H$ be the image of its regular algebra and put
$\mathcal H=\mathcal B_{D_H}$, using its specified central structure.
Then $\mathcal H$ has precisely the simple objects
$g$ and $g\rho$, $g\in G=E/H\cong\mathbb Z/N$, and
\[
 gh=g+h,\qquad \rho g=(-g)\rho,\qquad
 \rho^2=\mathbf1\oplus\bigoplus_{g\in G}g\rho.
\]
In particular, $\mathcal H$ is a fusion category with the
cyclic Haagerup--Izumi fusion rules.
\end{theorem}
\begin{proof}
De-equivariantization by the stated central Tannakian algebra is a fusion
category, with a tensor action of $\Gamma$ whose equivariantization is
$\mathcal B$; the underlying full faithfulness is precisely what ensures
a simple unit \cite[Theorem 4.18, Proposition 4.19, Remark 4.21]{dgno}.
Let $F:\mathcal B\to\mathcal H$ be free induction.
The subgroup $H$ is cyclic of order $N$, because its simple objects
are the characters of the cyclic group $\Gamma$. A generator of
$H\subset(\mathbb Z/N)^2$ is a primitive vector modulo $N$ and
extends to a basis; hence $E/H\cong\mathbb Z/N$.
For invertibles $e,f$,
\[
 \operatorname{Hom}_{\mathcal H}(F(e),F(f))
 \cong\bigoplus_{h\in H}\operatorname{Hom}_{\mathcal B}(e,fh).
\]
Thus the $F(e)$ give $N$ distinct simple invertibles, indexed by $G$.
Moreover,
\[
 \operatorname{End}_{\mathcal H}F(R)
 \cong\bigoplus_{h\in H}\operatorname{Hom}_{\mathcal B}(R,Rh).
\]
Every summand is one-dimensional. Choosing an isomorphism $R\to Rh$
in each degree, composition multiplies such basis vectors by nonzero
scalars. The endomorphism algebra is therefore a twisted group algebra
of $H$. A cyclic group has trivial $H^2(H,\mathbb C^\times)$: after
choosing a lift of its generator, rescale that lift so that its $N$th
power is $1$. Hence this algebra is $\mathbb C[H]\cong\mathbb C^N$.
It follows that $F(R)=Q_1\oplus\cdots\oplus Q_N$, with all $Q_i$
simple, distinct and of multiplicity one. Adjunction excludes the
$F(e)$ among the $Q_i$. Dominance of $F$ shows that this list exhausts
the simples of $\mathcal H$.

Put $\delta=(N+\sqrt{N^2+4})/2$. Then $\operatorname{FPdim}R=N\delta$. Tensor functors preserve Frobenius--Perron dimensions,
and de-equivariantization divides the global dimension by $N$.
Writing $d_i=\operatorname{FPdim}Q_i$ gives
\[
 \sum_i d_i=N\delta,\qquad
 \sum_i d_i^2
 =\frac{N^2+N^2\delta^2}{N}-N=N\delta^2.
\]
Equality in Cauchy--Schwarz forces every $d_i=\delta>1$.
For any $Q_i$, write the total invertible and noninvertible
multiplicities of $Q_iQ_i^*$ as $a,b$. Since
$\delta^2=1+N\delta$ is irrational,
$a+b\delta=1+N\delta$ gives $a=1$, $b=N$.
Frobenius reciprocity identifies $a$ with the order of the left
$G$-stabilizer. Thus the left action is free, and hence transitive;
the same argument with $Q_i^*Q_i$ handles the right action.
Choose $\rho=Q_1$ and write
$\rho g=\varphi(g)\rho$, where $\varphi\in\operatorname{Aut}(G)$.

The deck action fixes the isomorphism class of each $F(e)$: twisting
the regular-algebra action is undone by rescaling its $H$-graded
free factor. It is transitive on the $Q_i$. Indeed each deck orbit
of noninvertible simples supplies a noninvertible equivariant simple,
whereas $\mathcal B$ has just one. Since there are $N$ deck elements
and $N$ such simples, the action is free and transitive.
As it fixes every invertible, it commutes with left translation by
$G$. Thus the deck action on the noninvertible labels consists of
all translations $g\rho\mapsto(g+a)\rho$.

Write
$\rho^2=g_0\oplus\bigoplus_g m_g\,g\rho$.
The same dimension calculation shows there is exactly one
invertible term and $\sum_gm_g=N$.
Applying the deck translation $a$, and comparing unique invertible
terms with
$(a\rho)^2=(a+\varphi(a))\rho^2$, gives
$a+\varphi(a)=0$ for every $a$. Hence $\varphi=-1$.
Now $(a\rho)^2=\rho^2$, so the same comparison of noninvertible
terms gives $m_{g+a}=m_g$; consequently all $m_g=1$.
Finally associativity of $\rho^3$ gives
$g_0\rho=(-g_0)\rho$ after cancelling the common sums.
The free $G$-action yields $2g_0=0$. Oddness of $N$ gives $g_0=0$,
as required.
\end{proof}

\subsection{The complementary subgroup fixes the pointed associator}

\begin{corollary}\label{cor:untwisted}
In Theorem \ref{thm:cyclic-converse}, suppose the specified central
Tannakian category is one Lagrangian factor of a braided embedding
\[
 \mathcal E^{\mathrm{rev}}\longrightarrow Z(\mathcal B),\qquad
 \mathcal E\simeq\mathcal C(G\times\widehat G,q),\quad
 q(g,\chi)=\chi(g),
\]
whose underlying functor identifies all $N^2$ invertibles of
$\mathcal B$. Then the pointed subcategory of $\mathcal H$ is
tensor equivalent to $\operatorname{Vec}_G$ with trivial associator.
\end{corollary}
\begin{proof}
The hyperbolic category has a skeletal realization with associator $1$
and braiding
$c_{(g,\chi),(h,\psi)}=\psi(g)$; replacing this braiding by its
inverse-transpose realizes the reverse category.
Both factors $H=\{0\}\times\widehat G$ and
$K=G\times\{1\}$ are therefore tensor subcategories with trivial
associator. Use the regular algebra on $H$.
The underlying central embedding and free induction restrict to a
tensor functor $\operatorname{Vec}_K\to\mathcal H$.
The adjunction calculation in the theorem shows that its $N$ simple
images are distinct and exhaust the invertibles of $\mathcal H$;
it is fully faithful and hence an equivalence onto
$\mathcal H_{\mathrm{pt}}$. This proves the associator assertion.
It does not require $K$ to remain a central subcategory after induction.
\end{proof}

\subsection{Preservation of the ribbon structure}

\begin{theorem}\label{thm:ribbon-center}
Let $\mathcal D$ be a modular tensor category, and let $A$ be a
connected \'etale algebra with $\theta_A=1$ and
$\dim_{\mathcal D}A\ne0$. Suppose
$\mathcal B=\mathcal D_A$ has the near-group rules in
Theorem \ref{thm:cyclic-converse} and its local category is the
standard hyperbolic ribbon category of order $N^2$, with all its
objects invertible in $\mathcal B$.
Then there is a connected \'etale algebra $A'$ over $A$ such that
$\mathcal H=\mathcal D_{A'}$ has the untwisted HI rules and
\[
                  \mathcal D\simeq Z(\mathcal H)
\]
as ribbon categories. If $\dim_{\mathcal B}R=-Nt$, then
$\dim_{\mathcal H}\rho=-t$.
\end{theorem}
\begin{proof}
The central embedding
$(\mathcal D_A^0)^{\mathrm{rev}}\to Z(\mathcal D_A)$ is part of
the center equivalence \cite[Corollary 3.30]{dmno}.
It is fully faithful after forgetting to $\mathcal D_A$ on the local
subcategory. Choose either cyclic Lagrangian factor $H$ and its
regular algebra. By \cite[Section 3.6 and Proposition 3.16]{dmno},
this algebra in $\mathcal D_A^0$ gives a connected \'etale algebra
$A'$ over $A$ in $\mathcal D$, and algebra-module formation in
stages identifies $\mathcal D_{A'}$ with the corresponding central
de-equivariantization. Our inverse-braiding right-module convention
agrees with the reverse local factor in the displayed center embedding.
Theorem \ref{thm:cyclic-converse} and Corollary \ref{cor:untwisted}
apply to it.

The local modules after the second condensation have group
$H^\perp/H=0$, so $\mathcal D_{A'}^0\simeq\operatorname{Vec}$.
Equivalently $\operatorname{FPdim}(A')^2=\operatorname{FPdim}\mathcal D$.
The central lift of free induction is therefore a braided equivalence
$\mathcal D\to Z(\mathcal D_{A'})$
\cite[Corollary 4.1]{dmno}.

For completeness, this equivalence preserves the required ribbon
structures, not only the braiding.
Local restriction preserves twists, so $\theta_{A'}=1$.
Each local invertible has dimension $1$, and the module-dimension
formula gives $\dim_{\mathcal D}A'=N\dim_{\mathcal D}A\ne0$.
A connected \'etale algebra has nondegenerate multiplication pairing
against its normalized projection to the unit: the radical is a
right ideal, while its regular module is simple since its endomorphism
algebra is $\operatorname{Hom}(\mathbf1,A')=\mathbb C$ in a semisimple
module category; the pairing is nonzero on the unit. Thus $A'$ is a
rigid algebra in the sense of \cite[Definition 1.11]{ko}.
Theorems 1.17--1.18 and Remark 1.19 of that paper give a spherical
structure on $\mathcal D_{A'}$ with
\[
 \dim_{\mathcal D_{A'}}F_{A'}(X)=\dim_{\mathcal D}X.
\]
The induced spherical structure on its center has the dimension of
the underlying object. Hence the braided center equivalence preserves
the dimension of every simple. Two pivotal structures on a fusion
category differ by a tensor automorphism of the identity, scalar on
each simple; that scalar multiplies its nonzero pivotal dimension.
Equality of all simple dimensions therefore makes the two pivotal
structures equal. In a braided category the pivotal structure and
braiding determine the twist, proving the ribbon assertion.

The same dimension formulas in two stages give
$\dim F_H(R)=\dim R$. Its $N$ noninvertible summands form one orbit
under the odd-order invertibles, all of dimension $1$. Their common
dimension is consequently $\dim R/N=-t$.
\end{proof}

\paragraph{Scope of the identification.}
The theorem constructs the particular fusion category
$\mathcal H=\mathcal D_{A'}$ and identifies its full ribbon center.
Here ``untwisted'' means that its pointed associator is trivial.
It does not identify $\mathcal H$ with every category having these
fusion rules, or with a separately prescribed HI solution.
Such uniqueness cannot be inferred from modular data; distinct
order-fifteen HI categories with identical center modular data are
exhibited in \cite[Corollary 5.9]{gsy}.

\begin{corollary}\label{cor:our-center}
For the algebra $A=\mathbf1\oplus T$ constructed in $\mathcal D=\Par'$
above, the category $\mathcal D$ is ribbon equivalent to
$Z(\mathcal H_N)$, where $\mathcal H_N$ has the untwisted cyclic HI
fusion rules and $\dim\rho=-t$.
\end{corollary}
\begin{proof}
The branching and multiplicity budget prove that
$\mathcal D_A$ has exactly the invertibles $E=(\mathbb Z/N)^2$ and
$R$, with $R^2=\sum_E e+N^2R$ and $\dim R=-Nt$.
Its local ribbon category is the standard hyperbolic category,
and $\dim A=1+t^2\ne0$. Theorem~\ref{thm:ribbon-center} applies.
\end{proof}


\paragraph{Galois consequence.}
Choose a number-field form of the fusion category and its spherical
structure, and extend the automorphism
$\sqrt\mu\mapsto-\sqrt\mu$ to its algebraic closure. Applied to the
associativity and duality maps, it preserves the fusion rules and
the trivial pointed associator, and sends $-t$ to $\delta$.
All simple spherical dimensions of the conjugate are therefore
its Frobenius--Perron dimensions. Hence the conjugate is
pseudo-unitary. This is the same consequence already included in
\cite[Theorem A(1) and Section 1.3]{gsy}; pseudo-unitarity here means
equality of global and Frobenius--Perron dimensions. It does not
assert a positive-definite unitary structure or a Galois-conjugate
VOA construction.

\section{Modular data of the center}
\label{sec:gsy-modular}

The center equivalence in Corollary~\ref{cor:our-center} determines
a ribbon category. We now identify its complete modular data in the
notation of Gannon--Schopieray--Yadav \cite[Section 5 and Appendix F.1]{gsy}.
The comparison is uniform in odd $N$, including when $N^2+4$ is
composite or has repeated prime factors. In their notation the answer is
\[
 \mathcal{MD}_{-}\bigl(\Z/N;\Z/(N^2+4),\ b\longmapsto b^2/(N^2+4)\bigr).
\]
The subscript $-$ specifies the spherical dimension $\dim\rho=-t$.
The cyclic group and its quadratic form are additional information;
the spherical sign alone does not determine them in general.

\subsection{Quadratic forms and the exceptional labels}

Write $e(x)=\exp(2\pi ix)$ and $\mu=N^2+4$. In this section
quadratic forms are written additively in $\Q/\Z$, with their
$e$-values giving the twists. The local odd lattice category has metric group
\[
 E=(\Z/N)^2,\qquad q_E(x,y)=\frac{xy}{N}\pmod\Z,
 \qquad B_E((x,y),(x',y'))=\frac{xy'+x'y}{N}\pmod\Z.
\]
Here the coordinates are obtained from the difference-of-squares
coordinates in Part~IV by $x=a+b$, $y=a-b$. The letters $a,b$
below will instead label the exceptional cyclic group.

Recall $v=K_-^{-1}e_0$, with $v_i=(N-2i)/\mu$ and $v_N=-v_0$.
For each nonzero $a\in\Z/\mu$, choose the unique odd residue $s$
modulo $2\mu$ satisfying $s\equiv2a\pmod\mu$. Denote by
$X^{\mathrm{neu}}_{[a]}$ the exceptional simple with frame grid
\[
                  r\in sv+2\Z^N.
\]
The subscript $[a]$ means the orbit $\{a,-a\}$. Its dual-theta
label has parity $e_0$, which is exactly the neutral label forced by
\eqref{r:eq:SY}--\eqref{r:eq:SX} and $\dim J_\epsilon=(-1)^{\epsilon_0}$.
Changing $s$ by $2\mu$ preserves the grid, while replacing $s$ by
$-s$ gives the isomorphic opposite-grid module. This notation fixes
the neutral representative as well as its odd discriminant label.

The total conformal weight at $r=sv$ is
\[
             h(s)=\frac N4\left(\frac{s^2}{\mu}-1\right).
\]
Writing $s=2a+\mu k$, with $k$ odd, gives
\[
 h(s)-\frac{Na^2}{\mu}
       =Nak+\frac N4(\mu k^2-1)\in\Z,
\]
because $\mu\equiv1\pmod4$. Thus the exceptional twists are governed by
\begin{equation}\label{eq:gsy-natural-form}
 H_N^{\mathrm{met}}=\Z/\mu,\qquad
 q_N(a)=\frac{Na^2}{\mu}\pmod\Z,\qquad
 B_N(a,a')=\frac{2Naa'}{\mu}\pmod\Z.
\end{equation}
Nondegeneracy follows from $\gcd(2N,\mu)=1$.

\begin{lemma}[Uniform metric isometry]\label{lem:gsy-isometry}
For every odd $N\geq3$, put
\[
 u_N=\frac{N^2-N+2}{2}.
\]
Then multiplication by $u_N$ identifies
$(\Z/\mu,q_N)$ with $(\Z/\mu,q_0)$, where
$q_0(b)=b^2/\mu\pmod\Z$. In particular,
\begin{equation}\label{eq:gsy-gauss}
 \sum_{a\bmod\mu}e\bigl(q_N(a)\bigr)=\sqrt\mu.
\end{equation}
\end{lemma}
\begin{proof}
There is the integer identity
\[
 u_N^2-N=\mu\left(\frac{N-1}{2}\right)^2.
\]
Hence $u_N^2\equiv N\pmod\mu$. Since $N$ is a unit modulo $\mu$,
so is $u_N$, and $q_0(u_Na)=q_N(a)$ for every $a$.
The ordinary quadratic Gauss sum for the positive odd integer
$\mu\equiv1\pmod4$ is
$\sum_{b\bmod\mu}e(b^2/\mu)=\sqrt\mu$; the formula holds without
a square-free hypothesis. Applying the isometry proves
\eqref{eq:gsy-gauss}. This is the sign required in
\cite[Theorem 5.1]{gsy}, where $r_-=-\sqrt\mu$.
\end{proof}

\subsection{The complete modular-data theorem}

We use the normalized categorical Hopf-link matrix with
\begin{equation}\label{eq:gsy-normalization}
 S_{\one,\one}=\frac1{\dim\mathcal H_N}
              =\frac1{N(1+t^2)}>0.
\end{equation}
Write $T^{\mathrm{rib}}=\diag(\theta_Z)$ for the ribbon twist matrix.
This is the $T$ used in the center formulas of \cite{gsy}.
The object called $T$ earlier in the manuscript remains the simple
summand of $A=\one\oplus T$; the superscript distinguishes the matrix.

\begin{theorem}[Negative-branch GSY modular data]\label{thm:gsy-modular}
The ribbon categories $\mathcal D=\mathcal P_N'$,
$Z(\mathcal H_N)$, and $\Rep(E_N)$ have modular data
\begin{equation}\label{eq:gsy-identification}
 (S,T^{\mathrm{rib}})
 =\mathcal{MD}_{-}\bigl(\Z/N;\Z/\mu,q_0\bigr),
 \qquad q_0(b)=\frac{b^2}{\mu}\pmod\Z,
\end{equation}
under the following bijection of simple labels:
\begin{center}
\begin{tabular}{@{}lll@{}}
\toprule
Simple of $\mathcal D$ & GSY label & Ribbon twist\\
\midrule
$\one$ & $0_+$ & $1$\\
$T$ & $0_-$ & $1$\\
$Y_{[(x,y)]}$ & $\mathfrak a_{[x,y]}$ & $e(xy/N)$\\
$X^{\mathrm{neu}}_{[a]}$ & $\mathfrak d_{[u_Na]}$ & $e(Na^2/\mu)$\\
\bottomrule
\end{tabular}
\end{center}
The periodic labels range over
$((\Z/N)^2\setminus\{0\})/\{\pm1\}$ and the exceptional labels over
$(\Z/\mu\setminus\{0\})/\{\pm1\}$.
To display every entry, set
\[
 \alpha=\frac12\left(\frac1N+\frac1{\sqrt\mu}\right),
 \qquad
 \beta=\frac12\left(\frac1N-\frac1{\sqrt\mu}\right),
 \qquad r_G=\frac{N^2-1}{2},\quad r_H=\frac{\mu-1}{2}.
\]
In the ordered GSY basis $(0_+,0_-),(\mathfrak a),(\mathfrak d)$,
\begin{equation}\label{eq:gsy-full-S}
 S=\begin{pmatrix}
 \alpha&\beta&N^{-1}\mathbf1_{1\times r_G}&-\mu^{-1/2}\mathbf1_{1\times r_H}\\
 \beta&\alpha&N^{-1}\mathbf1_{1\times r_G}&\mu^{-1/2}\mathbf1_{1\times r_H}\\
 N^{-1}\mathbf1_{r_G\times1}&N^{-1}\mathbf1_{r_G\times1}&C_E&0\\
 -\mu^{-1/2}\mathbf1_{r_H\times1}&\mu^{-1/2}\mathbf1_{r_H\times1}&0&C_H
 \end{pmatrix},
\end{equation}
where
\begin{align}
 (C_E)_{[x,y],[x',y']}
   &=\frac2N\cos\left(\frac{2\pi(xy'+x'y)}N\right),
                                                    \label{eq:gsy-CE}\\
 (C_H)_{[b],[b']}
   &=\frac2{\sqrt\mu}\cos\left(\frac{4\pi bb'}\mu\right).
                                                    \label{eq:gsy-CH}
\end{align}
The diagonal entries of $T^{\mathrm{rib}}$ in this basis are
$1,1,e(xy/N),e(b^2/\mu)$. Thus the comparison fixes the tensor
unit, the full $S$-matrix, and the twists, rather than only their
spectra or the fusion rules.
\end{theorem}

\begin{proof}
It suffices to work in $\mathcal D$: the two ribbon equivalences
are Corollary~\ref{cor:our-center} and
Theorem~\ref{thm:parity-cancellation-voa}. We first keep the natural exceptional
coordinate $a$ of \eqref{eq:gsy-natural-form}.

\emph{Normalization and distinguished rows.}
The dimensions computed in Section~\ref{r:sec:budget} are
$1,t^2,1+t^2,-Nt$ for $\one,T,Y,X$, respectively.
Since $1+t^2=t\sqrt\mu$, the vacuum row in normalization
\eqref{eq:gsy-normalization} is
\[
 S_{\one,\one}=\alpha,\qquad S_{\one,T}=\beta,\qquad
 S_{\one,Y}=\frac1N,\qquad S_{\one,X}=-\frac1{\sqrt\mu}.
\]
The remaining distinguished entries also follow directly from the
branching matrix. Write $\mathcal E=\mathcal D_A^{\loc}$ and let
$b^0$ be restriction from its odd local simples to $\mathcal D$.
Its row at the local unit is $\one+T$, and its row at a nonzero
$e\in E$ is $Y_{[e]}$. In the categorical normalizations
$S^{\mathcal E}_{0,0}=1/N$ and
$S^{\mathcal D}_{\one,\one}=1/[N\dim A]$, one has
\[
              S^{\mathcal E}b^0=b^0S^{\mathcal D}.
\]
Indeed, factor the analytic branching identity
\eqref{r:eq:analytic} by the common parity matrix and use
\eqref{r:eq:index}. If $S_W^{\rm an}=c_W\widetilde S_W$ and
$S_L^{\rm an}=c_L\widetilde S_L$, that index identity is
$c_L=c_W\dim A$; the unnormalized neutral and local Hopf-link
matrices are $N\dim A\,S^{\mathcal D}$ and
$N S^{\mathcal E}$, respectively. The scalar factors therefore
cancel. The local unit row now gives
\[
 S_{\one,T}+S_{T,T}=\frac1N,\qquad
 S_{\one,Y}+S_{T,Y}=\frac2N,\qquad
 S_{\one,X}+S_{T,X}=0.
\]
Together with the vacuum row these imply
\[
 S_{T,T}=\alpha,\qquad S_{T,Y}=\frac1N,\qquad
 S_{T,X}=\frac1{\sqrt\mu}.
\]
These are precisely the tractable entries of
\cite[Section 5.1]{gsy}. The two distinguished labels are fixed by
their dimensions and by $A=\one\oplus T$; the periodic and
exceptional families have dimensions $1+t^2$ and $-Nt$.
The twists of $\one$ and $T$ are $1$ because $A$ is the conformal
extension algebra. The local lattice twists give $e(xy/N)$ for $Y$,
and the calculation preceding Lemma~\ref{lem:gsy-isometry} gives
$e(Na^2/\mu)$ for $X^{\mathrm{neu}}_{[a]}$.

\emph{Comparing the regular Fourier blocks.}
For a rational vector $w$, write
\[
 f_{\sigma,w}(\boldsymbol\tau)
 =\frac1{\eta_\Pi(\boldsymbol\tau)}
   \sum_{m\in\Z^N+w}
       e\left(\sum_i\tau_i(B_\sigma m)_i^2\right).
\]
Multiplicity one and the grid classification identify the refined
exceptional character with $f_{-,sv/2}$. Poisson summation gives
the finite transform
\begin{equation}\label{eq:gsy-grid-fourier}
 f_{\sigma,w}(-1/\boldsymbol\tau)
 =\frac1{\sqrt{\det(2K_\sigma)}}
  \sum_{z\in\Z^N/(2K_\sigma\Z^N)}
    e(-w^{\mathsf t}z)\,
    f_{\sigma,K_\sigma^{-1}z/2}(\boldsymbol\tau).
\end{equation}
Here $2K_\sigma w$ is integral for the grids being transformed.
The formula follows by Poisson summation for the quadratic matrix
$B_\sigma^{\mathsf t}\diag(\tau_i)B_\sigma$; its inverse sends
$z$ to the displayed dual grid. The determinant factor is
$2^{-N/2}/N$ in the periodic case and
$2^{-N/2}/\sqrt\mu$ in the antiperiodic case.
It holds on the convergent theta functions, hence on their formal
diagonal germs.

To pass to the normalization of $\mathcal D$, there is one common
scalar $\kappa$ such that
$S^{W,\mathrm{an}}_{Z,Z'}=\kappa S^{\mathcal D}_{Z,Z'}$ for neutral
simples $Z,Z'$. This follows from the analytic/Hopf-link comparison
in Section~\ref{sec:analytic-normalization} and the ribbon
factorization by $\mathcal P_N$. For neutral $Y$, its Fourier label
is $e_0$, so \eqref{r:eq:SY} and $S^{\mathcal D}_{\one,Y}=1/N$ give
\begin{equation}\label{eq:gsy-parity-normalization}
                         \kappa=-2^{-N/2}.
\end{equation}
In particular the analytic sign is determined, not replaced by a
positivity assumption.

For two neutral exceptional labels choose odd $s,s'$ with
$s\equiv2a$ and $s'\equiv2a'$ modulo $\mu$. The two opposite
target grids in \eqref{eq:gsy-grid-fourier} contribute
\[
 \frac{2^{1-N/2}}{\sqrt\mu}
       \cos\left(\frac{\pi Nss'}\mu\right)
 =-\frac{2^{1-N/2}}{\sqrt\mu}
       \cos\left(\frac{4\pi Naa'}\mu\right).
\]
For the equality, write $s=2a+\mu k$, $s'=2a'+\mu k'$ with
$k,k'$ odd: the additional phase in $Nss'/(2\mu)$ is
$N\mu kk'/2$ modulo integers, a half-integer. Division by
\eqref{eq:gsy-parity-normalization} therefore gives
\begin{equation}\label{eq:gsy-natural-XX}
 S^{\mathcal D}_{X^{\mathrm{neu}}_{[a]},X^{\mathrm{neu}}_{[a']}}
   =\frac2{\sqrt\mu}\cos\left(\frac{4\pi Naa'}\mu\right).
\end{equation}

For the periodic block, let $\xi_e$ be the odd lattice class
indexed by $e\in E$. The neutral character uses the theta coset
$\xi_e-\mathbf1/2$. Pairing an odd-order class with the
two-primary class $\mathbf1/2$ is integral, and
\[
 2(\mathbf1/2)^{\mathsf t}K_+(\mathbf1/2)=N/2.
\]
Thus the Fourier phase between two neutral periodic cosets is
$-e(-B_E(e,e'))$. Pairing the opposite grids and dividing by
$\kappa$ gives \eqref{eq:gsy-CE}. A periodic theta transform has
no antiperiodic regular support, and conversely, so $S_{Y,X}=0$.
These coefficient comparisons use the disjoint joint frame
supports and uniqueness of regular grid modules from
Sections~\ref{g:sec:generic-sector-classification}
and~\ref{g:sec:exceptional}, together with
Lemma~\ref{lem:ordinary-jets}. They do not assume linear
independence of scalar characters, nor do they distinguish labels
merely by their total conformal weight. They therefore also apply
at composite and nonsquare-free $\mu$.

Finally put $b=u_Na$. Lemma~\ref{lem:gsy-isometry} turns
\eqref{eq:gsy-natural-XX} into \eqref{eq:gsy-CH} and turns the
exceptional twist into $e(b^2/\mu)$. In Appendix~F.1 of \cite{gsy},
substituting $\sigma=-1$, $y=N/\sqrt\mu$, and
$B_H(b,b')=2bb'/\mu$ gives exactly
\eqref{eq:gsy-full-S}--\eqref{eq:gsy-CH}. Equation
\eqref{eq:gsy-gauss} supplies their required Gauss-sum normalization.
Every block and every twist agrees under the displayed bijection.
\end{proof}

\subsection{Ribbon twists and the VOA character normalization}

The theorem compares ribbon modular data. For an ordinary VOA trace,
the translation matrix includes the central-charge scalar:
\[
              T^{\mathrm{an}}=e(-c/24)T^{\mathrm{rib}}.
\]
Consequently, for the parity-cancelled realization $E_N$,
\[
              T^{E_N,\mathrm{an}}=e(2N/3)T^{\mathrm{rib}}.
\]
One must remove this scalar when comparing character transformations
with the center matrices in \cite{gsy}. It has cube $1$, consistently
with $c(E_N)=-16N\equiv0\pmod8$.
For $W_N$, the parity factor must also be retained:
\[
 T^{W_N,\mathrm{an}}
   =e(5N/24)\bigl(T^{\mathcal P_N,\mathrm{rib}}
                          \otimes T^{\mathrm{rib}}\bigr).
\]
The full character representation of $W_N$ is therefore not the
center representation by itself.

The vacuum row of \eqref{eq:gsy-full-S} recovers precisely
$1,t^2,1+t^2,-Nt$. In particular, the negative branch is fixed
with the vacuum at $0_+$; no interchange of $0_+$ and $0_-$ is
being used. All the labels are self-dual, so the categorical
normalization obeys $S^2=I$ and $(ST^{\mathrm{rib}})^3=I$.

\subsection{Matching the tabulated categories and the scope of the result}
\label{sec:small-comparison}

The compact GSY modular label for Theorem~\ref{thm:gsy-modular} is
$\mu^-_1$. The explicit isometry in Lemma~\ref{lem:gsy-isometry}
removes the need for separate square-root searches at small orders.
For completeness, the following are all rows of
\cite[Appendix F, Table 2]{gsy} with this modular label at each
listed $N$. A range of superscripts includes every row in that range.
The table is a comparison of exact modular labels, not a calculation
from the rounded numerical HI matrices printed there.
\begin{center}
\small
\begin{tabular}{@{}rrll@{}}
\toprule
$N$ & $\mu$ & Modular label & Matching GSY rows\\
\midrule
3 & 13 & $13^-_1$ & $\alpha(3)$\\
5 & 29 & $29^-_1$ & $\alpha(5)$\\
7 & 53 & $53^-_1$ & $\alpha(7)$\\
9 & 85 & $85^-_1$ & $\alpha^{(1)}(9)$\\
11 & 125 & $125^-_1$ & $\alpha_A^{(1)}(11)$\\
13 & 173 & $173^-_1$ & $\alpha(13)$\\
15 & 229 & $229^-_1$ & $\alpha^{(1)}(15),\alpha^{(2)}(15),\alpha^{(3)}(15)$\\
17 & 293 & $293^-_1$ & $\alpha(17)$\\
19 & 365 & $365^-_1$ & $\alpha^{(1)}(19)$\\
21 & 445 & $445^-_1$ & $\alpha^{(1)}(21),\alpha^{(3)}(21)$\\
23 & 533 & $533^-_1$ & $\alpha^{(1)}(23)$\\
25 & 629 & $629^-_1$ & $\alpha^{(1)}(25)$\\
27 & 733 & $733^-_1$ & $\alpha^{(1)}(27),\alpha^{(2)}(27),\alpha^{(3)}(27)$\\
29 & 845 & $845^-_1$ & $\alpha_A^{(1)}(29)$\\
31 & 965 & $965^-_1$ & $\alpha^{(1)}(31)$\\
33 & 1093 & $1093^-_1$ & $\alpha^{(1)}(33),\ldots,\alpha^{(5)}(33)$\\
35 & 1229 & $1229^-_1$ & $\alpha^{(1)}(35),\alpha^{(2)}(35),\alpha^{(3)}(35)$\\
37 & 1373 & $1373^-_1$ & $\alpha^{(1)}(37),\alpha^{(2)}(37),\alpha^{(3)}(37)$\\
39 & 1525 & $1525^-_1$ & $\alpha_A^{(2)}(39),\alpha_A^{(4)}(39)$\\
41 & 1685 & $1685^-_1$ & $\alpha^{(1)}(41)$\\
43 & 1853 & $1853^-_1$ & $\alpha^{(1)}(43)$\\
\bottomrule
\end{tabular}
\end{center}

Thus the realization has the modular data of the listed negative
spherical GSY centers. In particular it selects $85^-_1$, rather
than $85^-_3$, at $N=9$, and the cyclic $125^-_1$, rather than the
noncyclic $25^-_1\times5^-_1$, at $N=11$.
For arbitrary odd $N$, the uniform assertion is the exact datum
\eqref{eq:gsy-identification} in the GSY classification of center
modular data. Their general theorem permits other metric groups
and forms. It does not by itself assign $\mu^-_1$ to every category
in their uniform double-sine construction, and no such additional
assignment is needed or asserted here.

Nor does this theorem identify the associator of $\mathcal H_N$
with that of a separately prescribed GSY category. Already at
$N=15$, the three matching rows have the same modular data but
are pairwise not categorically Morita equivalent
\cite[Corollary 5.9]{gsy}. An equivalence with a named category
would require comparison of associativity data, or an equivalence
of centers retaining the canonical Lagrangian algebra. At $N=3$
the present theorem gives the negative-branch Haagerup modular data;
the classification results in
\cite{grossmansnyder,grossmanizumi} are not being used to promote
that equality to a named-category equivalence. The ribbon
equivalence with the center of the category constructed here
remains Corollary~\ref{cor:our-center}.


\section{An explicit VOA realization of the parity centralizer}
\label{sec:parity-cancellation}

We use ``vertex operator algebra'' in the lower-bounded, integer-graded
sense: negative conformal degrees are allowed, but there are only
finitely many of them and every homogeneous space is finite-dimensional.
The extension below is simple, self-contragredient, regular
(hence rational), and $C_2$-cofinite. Its regularity means that every
weak module is a direct sum of irreducible ordinary modules; its
$C_2$-cofiniteness is proved separately from this assertion.

\subsection{The parity coordinates}

Let $W=W_N$ and retain the parity subcategory
$\mathcal P\subset\operatorname{Rep}(W_N)$ from
Theorem~\ref{c:currents}.
Its parity modules are
\[
 J_\epsilon=
 \bigoplus_{\substack{r_i\geq1\cr r_i\equiv1+\epsilon_i\ ({\rm mod}\ 2)}}
       \bigotimes_{i=0}^{N-1}L_i(r_i,r_{i+1}),
 \qquad \epsilon\in P=(\mathbb Z/2)^N.
\]
The parity-current construction gives
$J_\epsilon\boxtimes_WJ_\eta\cong J_{\epsilon+\eta}$.
With the cycle matrix $K$ of diagonal $(N+2,2,\ldots,2)$ and
off-diagonal cycle entries $-1$, its self-braiding quadratic form,
twist, and dimension are
\begin{equation}\label{eq:parity-data}
 Q(\epsilon)=i^{\epsilon^{\mathsf t}K\epsilon},\qquad
 \theta_\epsilon=i^{\epsilon^{\mathsf t}K\epsilon+2N\epsilon_0},
 \qquad d_\epsilon=(-1)^{\epsilon_0}.
\end{equation}
For $N=3$ each off-diagonal entry is $-1$, so that
$x^{\mathsf t}Kx=\sum_i(x_i-x_{i+1})^2+Nx_0^2$ still holds.

Here is an explicit change of parity coordinates, needed to specify
the cancellation algebra. Set
\[
 (S\epsilon)_i=\epsilon_i+\epsilon_{i+1}+\epsilon_0
 \quad\text{in }\mathbb F_2.
\]
If $S\epsilon=0$, all successive differences equal $\epsilon_0$;
summing them and using oddness of $N$ gives $\epsilon_0=0$, and
then every coordinate is zero. Thus $S$ is invertible.
Let $T(\epsilon)$ be the number of cyclic transitions of the binary
word $\epsilon$; it is even. Then
\[
 \epsilon^{\mathsf t}K\epsilon=T(\epsilon)+N\epsilon_0,
 \qquad |S\epsilon|=
 \begin{cases}
 T(\epsilon),&\epsilon_0=0,\\
 N-T(\epsilon),&\epsilon_0=1.
 \end{cases}
\]
Substitution in \eqref{eq:parity-data} gives
\begin{equation}\label{eq:semion-coordinates}
 Q(\epsilon)=i^{|S\epsilon|},\qquad
 \theta_\epsilon=(-i)^{|S\epsilon|},\qquad
 d_\epsilon=(-1)^{|S\epsilon|}.
\end{equation}
The classification of pointed braided categories by their quadratic
forms, with the additional specified ribbon twist, therefore identifies
the parity ribbon category as
$\mathcal P\simeq(\mathcal S^-)^{\boxtimes N}$.
The nontrivial object of $\mathcal S^-$ has self-braiding $i$,
twist $-i$, and dimension $-1$.

\subsection{A minuscule conformal shift of the \texorpdfstring{$E_7$}{E7} lattice}

Choose the standard even unimodular $E_8$ lattice in $\mathbb R^8$,
put $\alpha=e_7+e_8$, and let
\[
 K_7=E_8\cap\alpha^\perp\cong E_7,\qquad
 u=e_6+\tfrac12e_7-\tfrac12e_8.
\]
Writing $\beta=e_6+e_7\in E_8$, we have $u=\beta-\alpha/2$.
Hence
\[
 u\in K_7^*\setminus K_7,\quad 2u\in K_7,
 \quad u^2=\tfrac32,
 \quad K_7^*/K_7\cong\mathbb Z/2.
\]
Define
\[
 U=(V_{K_7},\omega_{K_7}+u(-2)\mathbf1).
\]
The fields, vacuum and translation are those of the lattice vertex
algebra. Its Virasoro modes and momentum weights are
\begin{align*}
 L_u(n)&=L_{K_7}(n)-(n+1)u(n),&
 c(U)&=7-12u^2=-11,\\
 \operatorname{wt}_u(e^p)
 &=\frac{p^2}{2}-(u,p)
 =\frac{(p-u)^2}{2}-\frac34.
\end{align*}
The shifted VOA is integer graded since $u\in K_7^*$.
It is nonnegatively graded: for $p\in K_7$, the vector
$\beta-p\in E_8$ has inner product $1$ with $\alpha$, and therefore
has squared norm at least $2$. Its orthogonal projection $u-p$ has
squared norm at least $2-1/2=3/2$.
The positive definite square-completion formula also gives finite
weight spaces.

The two simple weak modules are the shifted versions of $V_{K_7}$
and $V_{K_7+u}$. Their vertex operators are unchanged.
The ordinary lattice VOA is regular, and each lattice module remains
ordinary for $L_u(0)$ by the same square-completion formula.
Thus $U$ is regular. It is $C_2$-cofinite because the state-field
map and hence its $C_2$-space are unchanged. Self-contragredience
follows from $2u\in K_7$
\cite[Theorems 3.1, 4.1, 4.4]{dongmason}.

Write $M=V_{K_7+u}$ with the shifted conformal structure.
For $p=u+a$, $a\in K_7$, its weights are
$a^2/2-3/4$ plus oscillator degrees. Thus it has a one-dimensional
lowest space of weight $-3/4$ and twist $i$.
The change of conformal vector preserves the braided category:
apply \cite[Proposition 2.17]{mcrae} to
$v=u(-1)\mathbf1$, for which $v_0\omega_{K_7}=0$;
both conformal gradings are nonnegative.
The lattice fusion and self-braiding are consequently
\[
 M\boxtimes_U M\cong U,
 \qquad c_{M,M}=e^{\pi i u^2}=-i.
\]
The order-two pointed ribbon category has
$d(M)=\theta_M/c_{M,M}=-1$.
Its associator on the triple of nontrivial labels is $-1$ in the
usual semion gauge. These data give the ribbon equivalence
\begin{equation}\label{eq:inverse-E7}
            \operatorname{Rep}(U)\simeq(\mathcal S^-)^{\mathrm{rev}}.
\end{equation}
In particular no theorem about the representation category of a
twisted HI VOA is needed to construct this auxiliary factor.

\subsection{The diagonal extension and all its weak modules}

\begin{theorem}\label{thm:parity-cancellation-voa}
Let $W_N$ be the VOA constructed above, with its established parity
ribbon factor, and put
$\mathcal D=\mathcal P'\subset\operatorname{Rep}(W_N)$.
Set $M^0=U$, $M^1=M$. Then
\[
 E_N=\bigoplus_{\epsilon\in P}
       J_\epsilon\otimes\bigotimes_{i=0}^{N-1}M^{(S\epsilon)_i}
\]
has a simple, self-contragredient, regular, $C_2$-cofinite VOA
extension structure with
\[
 c(E_N)=-16N,\qquad
 \operatorname{Rep}(E_N)\simeq\mathcal D
\]
as ribbon categories. In particular, $E_N$ is rational and has
exactly $N^2+3$ irreducible ordinary modules. Its least conformal
degree is $-(N-1)/2$, and its homogeneous spaces are finite-dimensional.
Thus it is a lower-bounded integer-graded VOA, and is not of CFT type.
\end{theorem}
\begin{proof}
Let $V=W_N\otimes U^{\otimes N}$.
Every tensor factor is nonnegatively graded and $C_2$-cofinite.
Their finite tensor product has these properties too. The Deligne
product theorem \cite[Theorem 5.1]{c:McRaeDeligne} gives
\[
 \operatorname{Rep}(V)
 \simeq\mathcal P\boxtimes\mathcal D\boxtimes\mathcal P^{\mathrm{rev}}.
\]
The displayed object $E_N$ is the diagonal regular algebra in the
first and third factors, tensored with the unit of $\mathcal D$.
Explicitly, all its self-braidings and twists cancel by
\eqref{eq:semion-coordinates}--\eqref{eq:inverse-E7}.
The diagonal pointed subcategory is Tannakian, and its regular
algebra is connected, commutative and separable.
Its multiplication is nonzero between every pair of its simple-current
components. It has dimension $|P|=2^N$, since each paired component
has dimension $1$.

The weights in this finite direct sum are integral and lower bounded,
with finite-dimensional homogeneous spaces, as verified below.
The extension correspondence
\cite[Theorem 7.5]{cmy} therefore supplies the VOA structure.
Its unit map embeds $V$ since $V$ is simple. Separability and
connectedness make its regular algebra module simple, which is
exactly simplicity of the extension VOA.

\emph{$C_2$-cofiniteness.}
Equip the same underlying vertex algebra as $V$ with the unshifted
lattice conformal vectors and write
\[
 V^{\mathrm{un}}=W_N\otimes V_{K_7}^{\otimes N}.
\]
This is a CFT-type, $C_2$-cofinite VOA. As a weak
$V^{\mathrm{un}}$-module, $E_N$ is the finite direct sum of the
irreducible modules
\[
 J_\epsilon\otimes
 \bigotimes_{i=0}^{N-1}V_{K_7+(S\epsilon)_i u},
 \qquad \epsilon\in P.
\]
Irreducibility follows from the tensor-product decomposition, and
the conformal shift has not changed the module vertex operators.
Each summand is $C_2$-cofinite as a $V^{\mathrm{un}}$-module by
\cite[Proposition 5.2]{abd}. Consequently the subspace
\[
 C_2^{V^{\mathrm{un}}}(E_N)
 =\operatorname{span}_{\mathbb C}
       \{v_{-2}e:v\in V^{\mathrm{un}},\ e\in E_N\}
\]
has finite codimension in $E_N$. The underlying vertex algebra
$V^{\mathrm{un}}$ embeds in $E_N$, so
\[
 C_2^{V^{\mathrm{un}}}(E_N)\subseteq C_2(E_N),
 \qquad
 \dim E_N/C_2(E_N)
 \leq\dim E_N/C_2^{V^{\mathrm{un}}}(E_N)<\infty.
\]
Here $V^{\mathrm{un}}$ is used only to apply the CFT-type module
theorem: its conformal vector need not agree with that of $E_N$.
The vertex operators, and hence the displayed $C_2$-spaces, are
unchanged. This proves $C_2$-cofiniteness with the stated conformal
structure on $E_N$.

\emph{The module category and regularity.}
The local-module category of the diagonal algebra in
$\mathcal P\boxtimes\mathcal P^{\mathrm{rev}}$ is
$\operatorname{Vec}$; this follows either from its Lagrangian
dimension or directly from the diagonal subgroup's orthogonal
complement. Hence the local algebra modules in
$\operatorname{Rep}(V)$ form the ribbon category $\mathcal D$.
The VOA/local-algebra correspondence is a braided tensor equivalence
\cite[Theorem 3.65]{ckm}; twists agree because both are
$e^{2\pi i L(0)}$ on the same underlying module.

This identifies the full ordinary category, and also proves
regularity rather than only semisimplicity within a chosen finite
category. Indeed the displayed Deligne product is semisimple. By
\cite[Proposition 2.11]{mcrae}, every cyclic weak $V$-module is
grading-restricted generalized, since $V$ is nonnegatively graded
and $C_2$-cofinite. It is therefore a finite direct sum of ordinary
simples. Every weak $V$-module is the sum of its cyclic submodules,
and hence a sum of simple submodules. Such a module is semisimple;
this proves regularity of $V$ in the weak-module sense.
Every weak $E_N$-module restricts to a direct
sum of ordinary simple $V$-modules, hence belongs to
$\operatorname{Ind}(\operatorname{Rep}(V))$.
The extension correspondence on direct limits
\cite[Theorems 7.5 and 7.7]{cmy} identifies it with a local
module for the same finite separable algebra there.
Here the passage to arbitrary weak modules can be checked on finite
pieces. Let $Q$ be such a local algebra module and let $X\subset Q$
be any finite $V$-subobject. The action map
\[
                       E_N\boxtimes_V X\longrightarrow Q
\]
has finite image, since its source lies in the finite category
$\operatorname{Rep}(V)$. Associativity makes that image an
$E_N$-submodule. Separability makes every finite algebra module
semisimple. These finite images cover $Q$, because $Q$ is a union
of finite $V$-subobjects and the algebra action is unital.
Consequently $Q$ is a sum of simple algebra submodules, hence a
direct sum of them. Each simple summand is finite as a $V$-object
and is ordinary for the common conformal vector: its total $L(0)$
is diagonalizable, lower bounded, and has finite-dimensional
eigenspaces. Locality passes to subobjects and direct summands.
Thus every weak $E_N$-module is a direct sum of irreducible
ordinary $E_N$-modules.
In particular every admissible module is completely reducible,
so $E_N$ is rational. The ribbon equivalence and the previously
established rank of $\mathcal D$ give its $N^2+3$ simple modules.

\emph{Self-contragredience.}
The conformal subalgebra $V=W_N\otimes U^{\otimes N}$ is
self-contragredient, since all its tensor factors are.
The finite $V$-module decomposition of $E_N$ contains $V$ as a
direct summand. Taking graded contragredients therefore shows that
the ordinary $E_N$-module $E_N^\vee$, restricted to $V$, contains
$V^\vee\cong V$. Induction--restriction adjunction
\cite[Lemma 2.61 and Theorem 3.65]{ckm} gives
\[
 \operatorname{Hom}_{E_N}(E_N,E_N^\vee)
 \cong\operatorname{Hom}_{V}(V,E_N^\vee)\ne0,
\]
because induction sends $V$ to $E_N$. The adjoint $E_N$-module is
irreducible by simplicity, and its graded contragredient is also
irreducible. A nonzero morphism between them is therefore an
isomorphism. Thus $E_N\cong E_N^\vee$ for its stated conformal
vector; this argument does not change that conformal vector or
assume CFT type.

\emph{Conformal degrees and central charge.}
Write $r=\mathbf1+v$ in a frame summand of $J_\epsilon$.
Its lowest total weight is
\[
 \frac14\sum_i(v_i-v_{i+1})^2+\frac N4v_0(v_0+2).
\]
The constraints $v_i\geq0$ and $v_i\equiv\epsilon_i\pmod2$
imply a squared difference of at least $1$ at every parity
transition and $v_0\geq\epsilon_0$; $v=\epsilon$ attains all
these bounds. Consequently
\[
 h_{\min}(J_\epsilon)
 =\frac{T(\epsilon)+3N\epsilon_0}{4}.
\]
Adding the auxiliary minimum $-3|S\epsilon|/4$ gives the integral
minimum of the corresponding extension component:
\[
 h_{\min}(E_{N,\epsilon})=
 \begin{cases}
 -T(\epsilon)/2,&\epsilon_0=0,\\
 T(\epsilon),&\epsilon_0=1.
 \end{cases}
\]
All higher weights differ from these by integers. The maximum
possible even transition number on an odd cycle is $N-1$, achieved
by a word with $\epsilon_0=0$.
This proves the stated least degree. Central charges add to
$-5N+N(-11)=-16N$.
\end{proof}

\begin{proof}[Completion of the proof of Theorem~\ref{thm:main}]
The cyclic construction, invariant form, and $C_2$ argument establish
the stated properties of $W_N$; Section~\ref{r:sec:projectivity}
proves its regularity. The parity factorization is
Section~\ref{sec:parity-factor}. The lattice branching and the vanishing
remainder in the dimension identity give
$|\operatorname{Irr}\mathcal D|=N^2+3$. The parity factorization
then gives $|\operatorname{Irr}(W_N)|=2^N(N^2+3)$.
Corollary~\ref{cor:our-center} constructs $\mathcal H_N$ and gives
the ribbon-center identification with $\dim\rho=-t$.
Theorem~\ref{thm:parity-cancellation-voa} supplies the explicit VOA $E_N$, its
simplicity, self-contragredience, $C_2$-cofiniteness, regularity
(and hence rationality), its central charge, its least degree, and
the required ribbon equivalence. Theorem~\ref{thm:gsy-modular}, with the metric
isometry of Lemma~\ref{lem:gsy-isometry}, proves the modular-data
identification in part (4). These are all the assertions of the theorem.
\end{proof}

\Needspace{12\baselineskip}
\appendix
\part*{Technical appendices}
\addcontentsline{toc}{part}{Technical appendices}
\section{Ordinary first images of proper cyclic intervals}
\label{app:ordinary-interval-images}
\label{neutral:app:proper-intervals}

Keep the hypothetical weight-one vacuum extension and put $H_*=2N+2$.
For a proper connected binary root support $J$ containing zero, let
$B_J$ be its primary and $r_J=(B_J)_{H_*-2}v$.

\begin{lemma}\label{neutral:lem:proper-first-images}
Every $r_J$ is a sum of ordinary simultaneous partial eigenvectors.
Its support consists only of coordinate spectra at the active frame
sites of $B_J$.
\end{lemma}

\begin{proof}
At a vacuum frame site, $L_i(-1)B_J=0$ and $h_i^{J}=0$.
Equations \eqref{neutral:eq:neutral-first-image}--
\eqref{neutral:eq:neutral-desc-first-image} therefore give
$H_ir_J=0$ exactly.  At its two endpoints,
Lemma~\ref{neutral:lem:first-image-cubic} makes the partial operator
semisimple on $r_J$.  Only internal $(3,3)$ factors can initially
have partial nilpotents.  For $J=\{0\}$ there are no such internal
factors, proving the base case.

First let $J=[0,k]$, with $0\leq k\leq N-2$, and write $B_k=B_J$.
For $0\leq a<k$, adjoin the ordinary interval $A_{a,k}=P_{[a+1,k]}$
to $B_a$.  Unique diagonal fusion and its level-one Ward coefficient
give
\begin{align*}
 (A_{a,k})_1B_a&=\kappa_{a,k}B_k,\\
 (A_{a,k})_0B_a&=\kappa_{a,k}
 \left(\sum_{i=a+1}^{k}L_i(-1)-(t+a)L_a(-1)\right)B_k,
                         \qquad\kappa_{a,k}\ne0.
\end{align*}
At a newly added site the right input is the vacuum, so its
intertwiner is $e^{zL_i(-1)}$ and has first coefficient one.  At the
shared site, the ordered $(1,3),(3,1)$ coefficient is
$-1/(\psi_a-1)=-(t+a)$.  This proves the displayed formulas, including
their signs.  The pairing of these coefficients with the output
partial weights is
\[
 \sum_{i=a+1}^{k}h_i^{(k)}-(t+a)h_a^{(k)}=1,
\]
by telescoping the internal weights
$2/((t+i)(t+i+1))$ and using $h_k^{(k)}=1+2/(t+k)$.

For any weight-two $A$ and weight-$H_*$ state $B$,
\eqref{neutral:eq:neutral-defect-product} and the zero-mode commutator
give $o(B)o(A)w=-o(A_0B)w/H_*$.  Applying this here gives the exact
transport relation
\begin{equation}\label{neutral:eq:proper-prefix-transport}
 \kappa_{a,k}\left(\sum_{i=a+1}^{k}H_i-(t+a)H_a-1\right)r_k
                       =-o(B_a)o(A_{a,k})w.
\end{equation}
Indeed, $o(L_i(-1)B_k)w=H_*(H_i-h_i^{(k)})r_k$.
Every ordinary interval preserves $U$, so if
$o(A_{a,k})w=\sum_i\alpha_iw_i$, its right-hand side is
\[
                    -\sum_i\alpha_i(h_i^{(a)}-H_i)r_a.
\]
By induction this is ordinary and supported at
$e_{N-1},e_0,\ldots,e_a$.

On a component $e_j$ with $0\leq j<k$, choose $a=j$.  The transport
operator has nonzero scalar value $-(t+j+1)$.  On $e_{N-1}$ every
cut has scalar value $-1$.  A polynomial in the commuting partial
operators with nonzero scalar part is invertible on the generalized
component, and its inverse takes an ordinary vector to an ordinary
vector.  All these components of $r_k$ are ordinary.
On the remaining component $e_k$, every right-hand side is zero.
Write $n_i=H_i-\delta_{ik}$.  Endpoint semisimplicity gives $n_kr_k=0$.
The cut $a=k-1$ gives $n_{k-1}r_k=0$, then $a=k-2$ gives
$n_{k-2}r_k=0$, and so on.  Vacuum sites and the fixed endpoint
were already treated.  This proves all prefixes.

For left growth, adjoin an ordinary interval $[a,b]$ to a heavy
interval starting at $b+1$ and wrapping through zero.  The analogous
level-one coefficients are one at $a-1,\ldots,b-1$ and $t+b+1$
at the shared site $b$.  The latter equals
$\psi_b/(\psi_b-1)$ for the ordered $(3,1),(1,3)$ input.
Their partial-weight pairing is
\[
 \left(1-\frac2{t+a}\right)
 +\sum_{i=a}^{b-1}\frac2{(t+i)(t+i+1)}
 +(t+b+1)\frac2{(t+b)(t+b+1)}=1.
\]
The transport operator is therefore
\[
                \widetilde D_{a,b}
                =\sum_{i=a-1}^{b-1}H_i+(t+b+1)H_b-1.
\]
For the suffix support $\{a,\ldots,N-1,0\}$, a component $e_j$
with $a\leq j\leq N-1$ uses the cut $b=j$, whose scalar is $t+j$.
The component $e_0$ has scalar $-1$.  Both are nonzero, and induction
on the shorter suffix proves ordinariness there.  On the new component
$e_{a-1}$, the right-hand sides vanish.  Endpoint semisimplicity
and the cuts $b=a,a+1,\ldots,N-1$ successively kill all its internal
nilpotents.  This proves every suffix together with root zero.

Finally take
$J=\{\ell,\ldots,N-1\}\cup\{0,\ldots,k\}$ with
$1\leq k<\ell-1$.  For each component $e_j$, $0\leq j\leq k$,
cut off the entire ordinary suffix $[\ell,N-1]$.  The left-growth
transport operator has scalar $-1$, and the remaining prefix image
is ordinary.  For each component $e_j$, $\ell-1\leq j\leq N-1$,
cut off the entire ordinary prefix $[1,k]$.  The right-growth
operator again has scalar $-1$, and its remaining suffix image is
ordinary.  These components exhaust the active sites, proving the
lemma without assuming a cyclic symmetry of the distinguished root.
\end{proof}



\section{Explicit cut polynomials}
\label{neutral:app:cut-polynomials}

This appendix defines the polynomials used in the Jacobian and
endpoint arguments by finite sums.  Thus all their stated values
and derivatives are identities of rational functions determined
inside the article.

\subsection{Local descendant evaluation}

For a local parameter $p$, put $h(p)=2(p-1)^2/p$.  For a PBW word
$\mathbf n=(n_1,\ldots,n_s)$ representing
$L(-n_1)\cdots L(-n_s)C$, define
\[
 Z_{\mathbf n}(x,y;p)
 =\prod_{k=1}^{s}(-1)^{n_k}
 \left(-x+n_ky+h(p)+\sum_{\ell>k}n_\ell\right).
\]
This is the repeated partial Zhu identity
\eqref{neutral:eq:neutral-partial-Zhu}, with the partial weight
updated after each inner descendant.  Let $V_d^+(p)$ and $V_d^-(p)$
be the local level-$d$ coefficients for ordered inputs
$(3,1),(1,3)$ and $(1,3),(3,1)$, respectively.  Through level three
they are the coefficients displayed in Section~\ref{r:sec:cyclic},
with $V_0^+=V_0^-=1$.  Write $Q_d^\pm(x,y;p)$ for the polynomial
obtained by substituting $Z_{\mathbf n}$ into $V_d^\pm$.

For three disjoint groups of sites, let $X,Y,Z$ denote the sums of
their left partial conformal operators, let $U,V,T$ denote independent
right variables, and let $h_X,h_Y,h_Z$ be the corresponding partial
weights of $C$.  For a polynomial $P$, write
$P_{abc}=P(U=a,V=b,T=c)$ and put
\begin{align*}
 \mathcal F(P)={}&(H-1)P_{000}
 +(h_X-X)(P_{100}-P_{000})\\
 &+(h_Y-Y)(P_{010}-P_{000})
 +(h_Z-Z)(P_{001}-P_{000}),\qquad H=2N.
\end{align*}
The right partial operators on $U=\operatorname{span}\{w_i\}$ are
orthogonal coordinate idempotents.  Since
$o(C)w=(H-1)r_C$ and $o(C)w_i=(h_i-H_i)r_C$, the action of the
descendant encoded by $P$ on $w$ is exactly $\mathcal F(P)r_C$.
There is no division by $r_C$ in this definition.

\subsection{The prefix cuts}

Fix $0\leq a\leq N-2$ and put $r=t+a$,
\[
 X=H_a,\qquad Y=H_{N-1},\qquad Z=\sum_{i<a}H_i,
\]
\[
 h_X=\frac2{r(r+1)},\qquad
 h_Y=\frac{2t^2}{1-t},\qquad h_Z=\frac2t-\frac2r.
\]
The ordered product $B_a\times A_a$, with $B_a=P_{[0,a]}$ and
$A_a=P_{[a+1,N-1]}$, has highest output $C$ with coefficient
$\kappa_a\ne0$.  Its degree-$d$ coefficient, divided by $\kappa_a$,
is represented by
\[
 P_d=\sum_{b+c+e=d}
 Q_b^+\!\left(X,U;\frac{r+1}{r}\right)
 Q_c^-\!\left(Y,V;\frac1{1-t}\right)
 \frac{(Z-T-h_Z)_{\underline e}}{e!},
                         \qquad0\leq d\leq3,
\]
where $(z)_{\underline e}=z(z-1)\cdots(z-e+1)$.
The last factor is the simultaneous translated unit intertwiner at
the prefix sites.  By \eqref{neutral:eq:neutral-defect-product},
\begin{equation}\label{neutral:eq:explicit-prefix-K}
 K_a=\sum_{\ell=0}^{3}
       \frac{\binom{H+1}{\ell}}{H+2-\ell}\,
                           \mathcal F(P_{3-\ell}),
 \qquad
 \kappa_aK_a(H)r_C=o(B_a)o(A_a)w.
\end{equation}
For $a=0$ the variable $Z$ is absent.  Direct differentiation of
this finite expression gives all four gradients in
Section~\ref{r:sec:cyclic}.  All denominators are nonzero at
$0<t<1/3$ and $r=t+a$: their possible zeros in a local parameter
are rational values excluded by its irrationality.

\subsection{Every ordinary interval and its heavy complement}

Let $A=P_{[a,b]}$, $1\leq a\leq b\leq N-1$, let $B=A^c$ be its
heavy complement, and put
\[
 r=t+a-1,\quad s=t+b,\qquad
 X=H_{a-1},\quad Y=H_b,\quad Z=\sum_{a\leq i<b}H_i.
\]
Here
\[
 h_X=\frac2{r(r+1)},\qquad h_Y=\frac2{s(s+1)},\qquad
 h_Z=\frac2{r+1}-\frac2s.
\]
Use the simultaneous disjoint-support normalization proved in
Appendix~\ref{r:app:first}.  For the product ordered as $A\times B$,
define
\[
 P_3^{a,b}=\sum_{d+e+f=3}
 Q_d^-\!\left(X,U;\frac{r+1}{r}\right)
 Q_e^+\!\left(Y,V;\frac{s+1}{s}\right)
                  \frac{(Z-T-h_Z)_{\underline f}}{f!}.
\]
Since $A$ has weight two, the defect identity and zero-mode
commutator give
\[
 o(B)o(A)w=-\frac1{H+2}o(A_0B)w.
\]
Consequently the full polynomial in the complementary-interval
equation is
\begin{equation}\label{neutral:eq:explicit-all-interval-K}
                 K_{a,b}=-\frac{\mathcal F(P_3^{a,b})}{H+2}.
\end{equation}
Evaluating and differentiating this formula at
$(X,Y,Z)=(0,0,0),(1,0,0),(0,1,0),(0,0,1)$ gives respectively
the off-support, left, right, and internal rows in
Appendix~\ref{r:app:first}.  The formula retains both left and
right Zhu variables at every intermediate step.

\subsection{Reconstructing the first derivatives}
\label{app:jet-map}
The four evaluation points in the preceding formula use the
coordinate order $(X,Y,Z)$. Their gradients are, respectively,
\[
 g_{\rm off},\quad g_{\rm left},\quad g_{\rm right},\quad g_{\rm inside}.
\]
For the prefix formula the same points give, after reordering
components to $(Z,X,Y)$,
\[
 G_{\rm suffix},\quad G_{\rm cut},\quad G_{\rm last},\quad G_{\rm prefix}.
\]
This is the gradient convention of Section~\ref{r:sec:cyclic}.

For example, write $P=P_3^{a,b}$ and evaluate the following
expression at $(X,Y,Z)=(0,0,0)$:
\begin{align*}
 \partial_XK_{a,b}=-\frac1{H+2}\big[&
 (H-1)\partial_XP_{000}-(P_{100}-P_{000})\\
 &+h_X\partial_X(P_{100}-P_{000})
 +h_Y\partial_X(P_{010}-P_{000})
 +h_Z\partial_X(P_{001}-P_{000})\big].
\end{align*}
Substituting the displayed $Z_{\mathbf n}$ and $V_d^\pm$ gives
$r(r^2+4r+1)/(3(r-1)(r+2))$, the first entry of $g_{\rm off}$.
The other entries use the same differentiation, with the selected
coordinate equal to one at an endpoint or interior component.
The prefix calculation uses
\eqref{neutral:eq:explicit-prefix-K} and $H=2/t-2t$.
The supplied program \texttt{uniform\_cut\_jets.py} implements
these PBW evaluations with first-order dual numbers over rational
function fields. It verifies the value and all three derivatives
at each of the eight points (32 exact identities); it does not
start from the asserted gradients. The subsequent endpoint
substitutions and matrix assembly are reproduced in
\texttt{first\_order\_all\_interval\_checks.py}.



\section{Maximum plateaux and the filtered Zhu span}
\label{app:maximum-plateaux}
\label{neutral:app:filtered-spanning}

The frame decomposition has highest vectors $v_m\in M_m$, with
\[
 h(m)=\sum_i(m_i-m_{i+1})^2+Nm_0(m_0+1),
                         \qquad m\in\mathbb Z_{\geq0}^N.
\]
Its diagonal multiplication has a nonzero top channel $M_m\times
M_n\to M_{m+n}$, and all other output labels $\ell$ obey
$\ell_i\leq m_i+n_i$ coordinatewise.

\begin{lemma}\label{neutral:lem:binary-Zhu-generation}
The Zhu algebra of $W$ is generated by the frame classes, all
connected proper binary interval primaries, and the full cyclic
primary $C$.  For a proper root subset $I$, the Zhu algebra of
$W_I=\bigoplus_{\operatorname{supp}(m)\subset I}M_m$ is generated by
the frame and the binary interval primaries supported in $I$.
If $A_0$ is the Zhu image of $W_{\{1,\ldots,N-1\}}$ and
$F_1W=\bigoplus_{m_0\leq1}M_m$, then
\[
 [F_1W]\subseteq A_0+
 \sum_{P\in\{C,\text{ proper heavy interval primaries}\}}A_0[P]A_0.
\]
\end{lemma}

\begin{proof}
We first prove a strong spanning assertion by induction on
$|m|=\sum_i m_i$.  If $\max_i m_i\geq2$, choose one connected
component of the maximum plateau, and let $u$ be its binary
indicator.  If the maximum plateau is the whole cycle, take $u$
to be the all-ones vector.  Direct expansion gives
\begin{align*}
 \Delta&=h(m)-h(u)-h(m-u)\\
 &=2\sum_i(u_i-u_{i+1})
       \bigl((m_i-u_i)-(m_{i+1}-u_{i+1})\bigr)
                                  +2Nu_0(m_0-u_0)\geq0.
\end{align*}
Each boundary term is nonnegative because a neighbor outside the
maximum plateau has height at most one less than its maximum.
The final term is nonnegative as well.
The product
\[
                       (v_u)_{-\Delta-1}v_{m-u}
\]
has a nonzero $v_m$ coefficient, and every other output label has
strictly smaller total label.  These other contributions are partial
Virasoro descendants of their highest vectors.  Induction expresses
$v_m$ in normal products of frame states and partial descendants of
binary highest vectors.  Every negative mode is a normal product
after a total derivative of its left factor.

The same spanning statement holds for partial descendants.  Each
$L_i(-1)=(\omega_i)_0$ is a derivation of every vertex product, so
it differentiates the normal-product factors.  For $n\geq2$,
$L_i(-n)a=(\omega_i)_{1-n}a$ adds a frame factor by a negative
product.  This proves closure under all partial Virasoro descendants.
It also proves the assertion within $W_I$, since the top-channel
construction and all lower output labels keep their support in $I$.

A disconnected binary support is the disjoint union of its separated
interval components.  Their unique diagonal output is their union,
its weight is the sum of their weights, and its highest coefficient
is nonzero.  All singular products between the separated components
vanish.  Their normal product consequently gives the union primary,
and their Zhu classes multiply to its class up to a nonzero scalar.
The full cyclic support is a separate connected case, with primary
$C$.  Repeated use of \eqref{neutral:eq:neutral-partial-Zhu} reduces
every partial descendant to its primary class and left/right frame
classes.  This proves the two Zhu-generation assertions.

For the filtered assertion, suppose $m_0\leq1$.  Whenever the maximum
is at least two, the chosen maximum plateau avoids root zero, so
$u_0=0$.  The first factor in the displayed negative product is
therefore in the open-chain algebra.  Moreover every lower output
label has $\ell_0=m_0$ exactly: at root coordinate zero the first
factor is the tensor unit in the aligned fusion rule.  The same
induction on $|m|$ thus expresses every state of $F_1W$ in normal
products with at most one binary factor whose support contains zero.
A terminal binary support has at most one connected component
containing zero, and that component is a proper heavy interval or
the whole cycle.  Partial derivatives and frame factors preserve
this count by the preceding derivation argument.

Now use a second, independent induction on total conformal weight
to pass to Zhu classes.  The identity
\[
 [a_{-1}b]=[a]*[b]-\sum_{j=1}^{\operatorname{wt}(a)}
          \binom{\operatorname{wt}(a)}j[a_{j-1}b]
\]
has strictly lower-weight singular terms.  In the spanning expression
just obtained, every such singular term remains in $F_1W$, because
one factor has root-zero label zero.  Its class is therefore covered
by the conformal-weight induction.  Higher negative modes reduce by
total derivatives.  Partial descendants reduce by
\eqref{neutral:eq:neutral-partial-Zhu}, adding only left and right
frame classes.  The right-hand side of the claimed inclusion is an
$A_0$-bimodule, so these reductions preserve it without commuting any
two noncommuting factors.  This proves the filtered assertion.
\end{proof}



\section{Uniform invertibility of the all-cut Jacobian}
\label{r:app:jac}

We prove that the gradients in Section~\ref{r:sec:cyclic} are
independent on the hyperplane \(\sum_i n_i=0\).  Fix the joint
eigenvalue tuple \(e_j\), let \(n=(n_i)\) lie in the kernel of the
corresponding Jacobian, and write
\[
 S_a=\sum_{i<a}n_i,\qquad z=n_{N-1},\qquad r=t+a,
\]
and denote the three appropriate gradient entries by \(A_a,B_a,D_a\).
The equations are
\[
 A_aS_a+B_an_a+D_az=0,\qquad
 S_{a+1}=(1-A_a/B_a)S_a-(D_a/B_a)z.
\]
The coefficient \(B_a\) is nonzero in each row.  This is immediate
for the suffix and cut rows; for the prefix and last rows it follows
from the signs established below.
If \(z=0\), induction from \(S_0=0\) gives \(n=0\).
Otherwise set \(u_a=S_a/z\); the required terminal value is
\(u_{N-1}=-1\).

\subsection{The last component}

Here
\[
 A=\frac{2t-3}{6t},\quad
 D=\frac{t-1}{3(t+1)},\quad
 B(r)=\frac{(r+1)((2t-3)r^2+(8t-9)r-4t+6)}
                 {6t(r-1)(r+2)}.
\]
Since \(0<t<1/3\), the first recurrence multiplier is
\[
 q_0=1-A/B(t)=
 \frac{t^2(2t+7)}{(t+1)(2t^2+7t-6)}\in(-1,0).
\]
For \(a\ge1\), \(r>1\) and
\[
 B(r)-A=
 \frac{r((2t-3)r^2+(8t-9)r+2t)}{6t(r-1)(r+2)}<0.
\]
Thus \(B(r)<A<0\), and \(q_a=1-A/B(r)\in(0,1)\).
The recurrence has fixed point \(-D/A\). Its terminal condition
would require
\[
       \prod_{a=0}^{N-2}q_a=1-A/D
                          =-\frac{3-t}{2t(1-t)}<-1.
\]
The left side lies in \((-1,0)\), a contradiction.
For the first multiplier, the inequality \(q_0>-1\) follows from
\(4t^3+16t^2+t-6<0\) on \(0<t<1/3\).

\subsection{The other components for ranks at least five}

Now \(0<t<1/5\). For \(a<j\), the suffix recurrence has multiplier
\((t+a)/(t+a+1)\) and fixed point \((1-t)/t\), hence
\[
                         u_j=\frac{j(1-t)}{t(t+j)}.
\]
At the cut row, with \(r=t+j\), substitution gives
\[
 u_{j+1}=
 \frac{(t-1)\left(6tr^2+(6t^3-t^2-5t+2)r
                         +2t^3-3t^2-9t+2\right)}
      {2t(r+1)(t+1)(2t-1)}.
\]
Both polynomial coefficients
\(6t^3-t^2-5t+2\) and \(2t^3-3t^2-9t+2\) are positive for
\(0<t<1/5\): they are decreasing on this interval and positive
at \(1/5\). Therefore \(u_{j+1}>0\).

Every subsequent prefix row has \(r>1\), multiplier
\[
    \frac{r(r^2+4r+1)}{(r+1)(r^2+4r-2)}\in(0,1),
\]
and common fixed point
\[
                 v_*=-\frac{(t-1)(2t^2+4t-1)}
                                  {t(t+1)(2t-1)}>0.
\]
Each step is a convex combination of its predecessor and \(v_*\).
Thus \(u_{N-1}>0\), contradicting the terminal value \(-1\).

\subsection{The two remaining rank-three cases}

Put \(t=(\sqrt{13}-3)/2\). At \(j=0\), exact reduction modulo
\(t^2+3t-1\) gives
\[
 u_1+1=\frac{53-175t}{2(16t-5)}<0,\qquad
 v_*+1=\frac{7t-2}{16t-5}<0.
\]
The next step is a convex combination of two values below \(-1\),
so cannot reach \(-1\). The signs follow from
\(2/7<t<53/175<5/16\).
At \(j=1\) the cut is the last row, and
\[
                 u_2=\frac{37-120t}{2(7t-2)}>0,
\]
using \(2/7<t<37/120\).
These inequalities can be checked exactly by evaluating
\(f(x)=x^2+3x-1\) at the rational endpoints:
\[
 f(2/7)=-3/49,\quad
 f(53/175)=9/30625,\quad
 f(37/120)=289/14400.
\]
All cases have zero homogeneous kernel, proving uniform
invertibility.

\section{Uniform elimination of the first-image obstruction}
\label{r:app:first}

We prove Theorem~\ref{r:thm:first} using the semisimplicity of the
proper-interval first images and the bound \(\mm^2r_{C,j}=0\).
First we compare complementary intervals to eliminate the two
possible endpoint components of the closing-root image.  We then
propagate the vanishing to every proper heavy interval.
All equations below are vector equations with scalar coefficients.

\subsection{Common normalization and complementary equations}

For disjoint binary root supports \(J,K\), aligned fusion has a
unique highest output on \(J\cup K\), with nonzero coefficient
\(\kappa_{J,K}\). Choose every local Virasoro intertwiner with
highest coefficient one. In a triple of pairwise disjoint binary
supports, each Virasoro site sees at most two nonvacuum inputs,
since its two labels depend on just the two adjoining roots.
The local associativity comparison therefore includes a tensor
unit. With its canonical left unit and translated right unit
\(e^{zL(-1)}\), its scalar is one. The translation contributes
descendants but no extra constant. Thus the \(\kappa\)'s obey
the ordinary cocycle identity on disjoint unions.

The only nontrivial crossing has local exponent
\(h_{33}-h_{31}-h_{13}=-2\), so its scalar is symmetric.
Iterated products of the singleton highest vectors are consequently
independent of order and parenthesization. They normalize all
\(\kappa_{J,K}\) to one while retaining prescribed singleton scales.
Choose ordinary-root scales so that their left endpoint coefficient
on \(w\) is one. The remaining free scale of root zero is chosen
so that \(C\) has leading charge term \(e^{-B}\); every full-support
disjoint product is linear in that root scale, so both normalizations
can be imposed simultaneously.

For an ordinary interval \(A=P_{[a,b]}\), \(1\le a\le b\le N-1\),
put
\[
 r=t+a-1,\quad s=t+b,\quad
 D(x)=x(x-1)(x+1)(x+2),\quad d(r,s)=\frac{s(s+1)}{D(r)}.
\]
The open-chain anticommutator induction gives
\begin{equation}\label{r:eq:ordinaryclosed}
 o(A)w=d(r,s)\left(r(r-1)w_{a-1}
                +2\sum_{i=a}^{b-1}w_i+(s+1)(s+2)w_b\right).
\end{equation}
For a singleton this is Lemma \ref{neutral:lem:ordinary-root-defect}. Adjoining the
next ordinary root multiplies old nonshared coefficients by
\((s+1)/(s-1)\); the shared and new endpoint coefficients become
\(2d(r,s)\) and \(d(r,s)(s+1)(s+2)\), proving the formula.

Let \(B=A^c\) be the heavy complementary interval. The exact
global equation is
\[
                    K_{a,b}(H)r_C=o(B)o(A)w.
\]
The polynomial vanishes at every \(e_j\). With
\(X=H_{a-1}\), \(Y=H_b\), \(Z=\sum_{a\le i<b}H_i\),
its gradients for an off-support, left-endpoint, or right-endpoint
component are respectively
\begin{align*}
 g_{\rm off}&=\left(
 \frac{r(r^2+4r+1)}{3(r-1)(r+2)},
 -\frac{(s+1)(s^2-2s-2)}{3(s-1)(s+2)},-\frac13\right),\\
 g_{\rm left}&=\left(
 -\frac{r(3r+1)}{6(r-1)},
 -\frac{(s+1)(3rs^2-3rs-6r+2s^2-4s-4)}{6(s-1)(s+2)},
 -\frac{3r+2}{6}\right),\\
 g_{\rm right}&=\left(
 -\frac{r(3r^2s+r^2+9rs+r-2)}{6(r-1)(r+2)},
 \frac{(s+1)(3s+2)}{6(s+2)},\frac{3s+1}{6}\right).
\end{align*}
For an internal component it is
\(g_{\rm inside}=(-r,s+1,1)/6\).
These follow from the level-three coefficients already displayed;
the finite polynomial prescription is given in
Appendix \ref{neutral:app:cut-polynomials}.

The proper first image is ordinary. Its right-side coefficient
on a component where it is supported is
\[
 d(r,s)\gamma_j(r,s)r_{B,j},\qquad
 \gamma_j=\begin{cases}
 2r+s^2+s-4,&j=a-1,\\
 r^2+r-2s-6,&j=b,\\
 r^2+r+s^2+s-4,&j\notin[a-1,b].
 \end{cases}
\]
This is the dot product of \eqref{r:eq:ordinaryclosed}'s coefficient
vector with \(h_i^B-\delta_{ij}\).
Since \(\mm^2r_{C,j}=0\), these equations use only the gradients,
with \(n_i=(H_i-\delta_{ij})r_{C,j}\) and \(\sum_i n_i=0\).

\subsection{The endpoint components for ranks at least five}

Write \(\sigma=t+N-1=t^{-1}-1\). On component \(e_0\), set
\(u=r_{R_0,0}\). Repeated proper transport gives
\[
 r_{B,0}=\begin{cases}
 D(\sigma)u/D(s),&a=1,\\[2pt]
 \dfrac{2r(r+1)D(\sigma)}{t(t+1)^2(t+2)D(s)}u,&a\ge2.
 \end{cases}
\]
These formulas follow by removing ordinary tips on each side of
zero, using the nonzero scalar parts of the previously established
transport operators. Here is the scalar recursion explicitly.
Write $r_{B,j}=T_j(a,b)u$, initialized by
$T_j(1,N-1)=1$. Removing one tip gives, whenever the indicated
intervals are nonempty,
\begin{align*}
 \frac{T_0(a+1,b)}{T_0(a,b)}
 &=\begin{cases}2/[t(t+1)],&a=1,\\(r+2)/r,&a\ge2,\end{cases}\\
 \frac{T_0(a,b-1)}{T_0(a,b)}&=\frac{s+2}{s-2}.
\end{align*}
These are scalar equalities, not divisions by the possibly zero
vector $u$. They follow by evaluating the proper-transport
operator in \eqref{neutral:eq:proper-prefix-transport} at $e_0$,
using the ordinary-interval coefficient vector
\eqref{r:eq:ordinaryclosed}; the left-tip operator has shared-site
coefficient $s+1$ and coefficient one at each new site.
For example, at an endpoint component $j\in\{0,N-1\}$ the
right-growth scalar is
\[
 -\frac{1+2/r-\delta_{j,a-1}}{-1-r\delta_{j,a-1}}.
\]
The old heavy interval has partial weight $1+2/r$ at $a-1$ and
zero at $a$; the ordinary root has coefficient one at $a-1$.
For left growth the scalar is
\[
 -\frac{(s+1)(s+2)}{(s-1)(s-2)}\,
 \frac{1-2/(s+1)-\delta_{j,b}}{(s+1)\delta_{j,b}-1}.
\]
Here the old heavy weight at $b$ is $1-2/(s+1)$, and the first
fraction is the right coefficient of the adjoining ordinary root.
These two scalar evaluations also give the recurrences for the
last component below. Products telescope by
$D(s)/D(s-1)=(s+2)/(s-2)$ and $r(r+1)$, giving the stated
$T_0(a,b)$. Set
\[
 z=\frac{2D(\sigma)}{t(t+1)^2(t+2)}u,\qquad
 A_*(x)=(x-1)(x+2),\quad
 F(x)=x+3+\frac6{A_*(x)},\quad
 G(x)=x-2-\frac6{A_*(x)}.
\]
For every pair of off-component frame endpoints
\(1\le i<k\le N-1\), the global equation becomes
\begin{equation}\label{r:eq:triangle}
 F(t+i)n_i-G(t+k)n_k-\sum_{i<\ell<k}n_\ell
 =3z\left(\frac1{A_*(t+i)}+\frac1{A_*(t+k)}\right).
\end{equation}
Add the equations for \((i,k)\) and \((k,\ell)\), then subtract
the one for \((i,\ell)\). Since
\[
 F(x)-G(x)+1=\frac{6x(x+1)}{A_*(x)},
\]
we obtain \(n_k=z/[(t+k)(t+k+1)]\) for \(2\le k\le N-2\).
For \(N\ge5\), the set \(\{2,\ldots,N-2\}\) is nonempty.
Using a known internal site in \eqref{r:eq:triangle} extends the
formula to \(k=1,N-1\), since \(F(t+1)\ne0\) and
\(G(\sigma)>0\).
The sum relation then gives
\(n_0=z(t-1/(t+1))\). Substitution in any remaining interval
with \(a=1\) leaves
\[
                    -\frac{(3t+1)(t^2+1)}{6(t-1)}z=0.
\]
The scalar substitution behind this residual is
\begin{align*}
 &g_{\rm left}(t,s)\cdot
 \left(t-\frac1{t+1},\frac1{s(s+1)},\frac1{t+1}-\frac1s\right)\\
 &\qquad-\frac{(t+1)(2t+s^2+s-4)}{2(t-1)(s-1)(s+2)}
 =-\frac{(3t+1)(t^2+1)}{6(t-1)}.
\end{align*}
The coefficient is nonzero, so \(u=0\) and every \(n_i=0\)
on this component.

On component \(e_{N-1}\), put \(u=r_{R_0,N-1}\). The transport
formulas are
\[
 r_{B,N-1}=\begin{cases}
 \dfrac{r(r+1)}{t(t+1)}u,&b=N-1,\\[4pt]
 \dfrac{2r(r+1)(\sigma+1)(\sigma+2)}{t(t+1)D(s)}u,&b\le N-2.
 \end{cases}
\]
For completeness, the corresponding scalar recurrences are
\[
 \frac{T_{N-1}(a+1,b)}{T_{N-1}(a,b)}=\frac{r+2}{r},\qquad
 \frac{T_{N-1}(a,b-1)}{T_{N-1}(a,b)}=
 \begin{cases}
 \dfrac{2(\sigma+2)}{(\sigma-2)(\sigma-1)\sigma},&b=N-1,\\[3pt]
 \dfrac{s+2}{s-2},&b\le N-2.
 \end{cases}
\]
They use the same operator at $e_{N-1}$ and the same initial value.
Set \(z=2(\sigma+1)(\sigma+2)u/[t(t+1)]\).
The off-component equations are again \eqref{r:eq:triangle}, now
for \(0\le i<k\le N-2\). Triangle differences give the same
formula for \(n_k\) at \(1\le k\le N-3\). It extends to the
two other off-component sites using \(F(t)<0\) and
\(G(t+N-2)>0\). The sum fixes \(n_{N-1}\). A remaining interval
ending at \(N-1\) gives
\[
             -\frac{(3\sigma+2)(\sigma^2+2\sigma+2)}
                         {6(\sigma+2)}z=0.
\]
In this case the corresponding scalar identity, with $s=\sigma$, is
\begin{align*}
 &g_{\rm right}(r,s)\cdot
 \left(\frac1{r(r+1)},\frac1s-s-1,\frac1{r+1}-\frac1s\right)\\
 &\qquad-\frac{s(r^2+r-2s-6)}{2(r-1)(r+2)(s+2)}
 =-\frac{(3s+2)(s^2+2s+2)}{6(s+2)}.
\end{align*}
Again \(u=0\) and every \(n_i=0\). The two endpoints are the
only allowed components of \(r_{R_0}\), so \(r_{R_0}=0\).

\subsection{Exact rank-three endpoint calculation}

For \(N=3\), use the four unknown vectors
\((n_0,n_1,n_2,u)\), with \(u=r_{R_0,j}\).
The row order is the sum relation followed by the interval
equations \([1,1],[1,2],[2,2]\). Substituting the same gradients
and transport formulas, and reducing modulo \(t^2+3t-1\), gives
the matrices below. Explicitly, for $(a,b)=(1,1),(1,2),(2,2)$,
choose $g=(g_X,g_Y,g_Z)$ to be the left, right, or off row according
as $j=a-1$, $j=b$, or neither. With $\mathbf e_i$ the coordinate
row vectors, the corresponding matrix row is
\[
 \left(g_X\mathbf e_{a-1}+g_Y\mathbf e_b+
       g_Z\sum_{a\le i<b}\mathbf e_i,
       -d(r,s)\gamma_j(r,s)T_j(a,b)\right).
\]
The first row is $(1,1,1,0)$. This gives every entry without a
separate rank-three ansatz:
{\small
\[
M_0=\begin{pmatrix}
1&1&1&0\\
-(5t-4)/18&(19t+11)/6&0&-(22t+73)/3\\
-(5t-4)/18&-(3t+2)/6&1/6&(42t+139)/3\\
0&(t+10)/3&(t+2)/9&-4(53t+175)/9
\end{pmatrix},
\]
\[
M_2=\begin{pmatrix}
1&1&1&0\\
-(t+1)/9&-(t-7)/3&0&-4(17t+56)/9\\
1/6&(3t+7)/6&(5t+19)/18&-(102t+337)/3\\
0&-(19t+46)/6&(5t+19)/18&(122t+403)/3
\end{pmatrix}.
\]
}
Their exact determinants are
\[
       \det M_0=-\frac{376t+1253}{27},\qquad
       \det M_2=\frac{4016t+13265}{27}.
\]
Both are nonzero for \(t>0\). Thus \(r_{R_0}=0\) also for
\(N=3\). The displayed matrices follow by substituting the three
intervals into \eqref{neutral:eq:explicit-all-interval-K} and the
transport relations, then reducing modulo \(t^2+3t-1\).

\subsection{Propagation and the final Jacobian step}

Suppose a proper heavy interval \(B\) satisfies \(r_B=0\), so
\(o(B)U=0\). Adjoin an ordinary interval \(A\) to form another
proper heavy interval \(B'\). Its unique order-two OPE has
\[
 A_1B=\kappa B',\qquad
 A_0B=\kappa\sum_i c_iL_i(-1)B',\qquad
 \sum_i c_i h_i^{B'}=1.
\]
For right growth, \(c_i\) is \(-(t+a)\) at the shared site,
one at each new site, and zero elsewhere. For left growth, its
shared-site value is \(t+b+1\), again with new-site values one.
Put \(Q=\sum_i c_iH_i\). On every top, the partial Zhu formula
and the usual mode commutator give
\[
 [o(A),o(B)]=o(A_0B)+o(A_1B)=\kappa[Q,o(B')].
\]
The scalar terms cancel because their partial-weight pairing is
one. The left side kills \(U\), as \(o(A)U\subset U\).

The first image \(r_{B'}\) is ordinary. Project the identity,
applied to \(w_i\), to a target component \(e_j\):
\[
           (c_j-c_i)(h_i^{B'}-\delta_{ij})r_{B',j}=0.
\]
On the active support, shared and new sites have different
coefficients. For any target \(j\), choose an active \(i\) with
\(c_i\ne c_j\); then \(i\ne j\). Every active partial weight
is nonzero: internally it is \(2/[(t+i)(t+i+1)]\), and at an
endpoint it is an irrational \((3,1)\) or \((1,3)\) weight.
Hence \(r_{B',j}=0\).

Starting at \(R_0\) and adjoining roots on either side reaches
every proper heavy cyclic interval, proving \(r_B=0\) for all of
them. Every right side in \eqref{r:eq:K} is now zero. Appendix
\ref{r:app:jac}'s invertible Jacobian therefore gives
\(\mm r_{C,j}=0\). This proves Theorem~\ref{r:thm:first} for every odd \(N\ge3\).


\section{The weight-zero top has vacuum open-chain action}
\label{r:app:h0-open-chain}

We use the frame decomposition of \(W\), its nonzero diagonal
multiplication, the elementary cubic null relations, and
Appendix~\ref{neutral:app:filtered-spanning}.
Let $M$ be a finite-dimensional representation of $A(W)$ whose
commuting partial conformal operators $H_i$ have joint generalized
eigenvalue zero and satisfy
\[
 \sum_{i=0}^{N-1}H_i=0.
\]
These hypotheses hold on the relevant neutral aligned simple tops of lowest
weight zero.  They also hold on a vacuum extension by a distinct such simple:
the total conformal nilpotent factors through a homomorphism from the vacuum
to the distinct simple and hence vanishes.  We do not assert the result for
arbitrary logarithmic tops with a nonzero total conformal nilpotent.

\subsection{Eliminating the partial nilpotents}

For the ordinary elementary root $R_j$, $1\leq j<N$, put
\[
 a=1+\frac2{t+j},\qquad b=2-a,
 \qquad t^2+Nt-1=0,\quad t>0.
\]
Its two active partial weights are $a,b$ at sites $j,j-1$.
The local cubic relation is
\[
 (x-y)\bigl((x-y)^2-2\psi(x+y)+2\psi-1\bigr)=0.
\]
The second factor is invertible on left and right multiplication by a
nilpotent $H_i$, since $2\psi-1\ne0$.  Thus $[H_i,R_j]=0$ for every $i$.
Let $\mathcal R\subseteq\End(M)$ be the commutative algebra
generated by the $H_i$, and let $\mathfrak m=(H_0,\ldots,H_{N-1})$
be its nilpotent maximal ideal.  The operators $R_j$ commute with
$\mathcal R$ by the preceding null relation.  The partial Zhu formula
\begin{equation}
 [L_i(-n)v]=(-1)^n\bigl(-H_i*[v]+n[v]*H_i+h_i[v]\bigr),\qquad n\geq1,
 \label{r:eq:h0-partial-zhu}
\end{equation}
gives $[D_aR_j]=-aR_j$.  In particular $[R_j,D_aR_j]=0$.

Only the vacuum and root self-channel occur in $R_j*R_j$ and in
\[
 [R_j,D_aR_j]=[(R_j)_0D_aR_j]+[(R_j)_1D_aR_j].
\]
All states used have weight at most four, whereas the next aligned channel
has lowest weight eight.  Write $\nu_j,\sigma_j\ne0$ for the vacuum and
self-channel coefficients.  Set
\begin{align*}
 u(h)&=\frac{h(h+1)}{3(h-4)(3h-1)},
 &d(a)&=(a-4)(a+2)(2a-3)(2a-1),\\
 \kappa(a)&=-\frac{a(a-2)(5a^2-10a+2)}{4d(a)},
 &\eta(a)&=-\frac{7a^2(a-2)^2(a-1)}{8d(a)}.
\end{align*}
The commutator has the form $P(H)+Q(H)R_j=0$, where $P(0)=0$
and $Q(0)=2\sigma_j\eta(a)\ne0$.  Since $Q(H)$ is invertible
in $\mathcal R$, this expresses $R_j$ as
$-Q(H)^{-1}P(H)\in\mathfrak m$.
Modulo $\mathfrak m^2$, the product and half the commutator give
\begin{align*}
 0&\equiv\nu_j\bigl(u(a)H_j+u(b)H_{j-1}\bigr)
             +\sigma_j\kappa(a)R_j,\\
 0&\equiv\nu_j\bigl(-b\,u(a)H_j+a\,u(b)H_{j-1}\bigr)
             +\sigma_j\eta(a)R_j.
\end{align*}
Eliminating $R_j$ uses the exact identity
\[
 \bigl((\eta+b\kappa)u(a),(\eta-a\kappa)u(b)\bigr)
 =\frac{a^2(a-2)^2}{24d(a)}(a+1,a-3).
\]
All denominators, as well as $\sigma_j$, $\eta(a)$ and the common
factor $a^2(a-2)^2/(24d(a))$, are nonzero: each possible zero would
force $a$ to be rational.  No division by $\kappa(a)$ is needed.
We obtain
\begin{equation}
 (t+j+1)H_j-(t+j-1)H_{j-1}\in\mathfrak m^2.
 \label{r:eq:h0-recurrence}
\end{equation}
Iteration gives
\[
 H_j\equiv\frac{t(t+1)}{(t+j)(t+j+1)}H_0\pmod{\mathfrak m^2}.
\]
All coefficients are positive.  Their sum, together with $\sum_iH_i=0$,
gives $H_0\in\mathfrak m^2$, then $\mathfrak m=\mathfrak m^2$.
Nilpotence forces $\mathfrak m=0$.  Therefore
\begin{equation}
 H_i=0\quad(0\leq i<N),\qquad R_j=0\quad(1\leq j<N).
 \label{r:eq:h0-frame-zero}
\end{equation}

\paragraph{Derivation of the coefficients.}
For a local self-intertwiner with highest weight $h$ in all three positions,
the first evaluated coefficients at left and right Zhu weights zero are
\[
 f_0=1,\qquad f_1=-\frac h2,\qquad
 f_2=\frac{h(h+1)(h^2+3h-2)}{4(h+2)(2h-1)}.
\]
The central charge is $c=(-3h^2+7h-2)/(h+1)$.  The level-two Ward solution is
\[
 \frac{h(h+1)}{(h+2)(2h-1)}L(-2)v
 +\frac{(h-2)(h+1)}{4(h+2)(2h-1)}L(-1)^2v.
\]
For a right-input derivative,
$f^D_n=(h-n)f_n-(h+n-1)f_{n-1}$, with $f^D_0=h$.
Thus $[z^2](1+z)^2F_aF_b=\kappa$ and
$[z^2](1+z)F^D_aF_b=2\eta$.
For the vacuum channel the linear coefficients are $g_2=2h/c$,
$g_3=-2h/c$, and $g_4=3\alpha+2\beta$, where
\[
 \alpha=\frac{3h(c-2h+4)}{c(5c+22)},\qquad
 \beta=\frac{2h(5h+1)}{c(5c+22)}.
\]
Thus $g_4-g_2=u(h)$.  Differentiation gives the vacuum linear pair
$2(-b\,u(a),a\,u(b))$.  These rational identities use
levels at most four and follow by substituting the displayed Ward
coefficients and collecting the indicated powers of $z$.

\subsection{All ordinary intervals vanish}

Let $P=P_{[j,k]}$, $1\leq j\leq k<N$, be an ordinary interval primary of
weight two.  Its endpoint weights are
\[
 a=1+\frac2{t+k},\qquad b=1-\frac2{t+j}.
\]
By \eqref{r:eq:h0-frame-zero} and \eqref{r:eq:h0-partial-zhu},
\begin{equation}
 \frac1a[P,D_aP]-\frac1b[P,D_bP]=0.
 \label{r:eq:h0-endpoint-difference}
\end{equation}
We induct on the interval length.  Every output has labels in
$\{0,1,2\}$, supported in $[j,k]$, and lowest weight at most four.
Strict-support outputs and their descendants vanish by induction and the
maximum-plateau strong spanning of the smaller-support open-chain
subalgebras.  A full-support output is either the binary interval itself,
or has one interior height-two plateau and unchanged endpoint labels one.
Indeed, the weight is the sum of squared adjacent label differences:
endpoint height two costs at least six, and each separated interior
height-two plateau adds two over binary weight two.

The extra weight-four channels occur only as highest vectors in
$P_0D_iP$.  Their coefficients in the two commutators are respectively
$a$ and $b$ times the same canonical channel coefficient, so they cancel
in \eqref{r:eq:h0-endpoint-difference}.  For the self-channel, if $F_2$ is
the degree-two coefficient in the product of all local self series, then
the coefficient in $[P,D_aP]$, divided by $a$, is
\[
 F_2-2g(a)-\tfrac12,\qquad
 g(h)=\frac{f_2(h)}h-\frac h8
     =\frac{5h^2+4h-4}{8(h+2)(2h-1)}.
\]
All interior data cancel between endpoints.  Thus
$2\sigma_P(g(b)-g(a))P=0$, where $\sigma_P\ne0$ is the self-channel
coefficient.  The exact factorization is
\[
 g(a)-g(b)=
 \frac{(a-b)(7ab-2a-2b+4)}
 {8(a+2)(2a-1)(b+2)(2b-1)}.
\]
The denominators and $a-b$ are nonzero.  The remaining numerator,
after multiplying by $(t+j)(t+k)$ and using $t^2+Nt-1=0$, is
\[
 7(j+k-N)t+7jk-10(k-j)-21.
\]
If the expression vanished, irrationality of $t$ would force both
its coefficient of $t$ and its constant term to vanish. For $j=1$
the constant term is $-3k-11<0$. For $j\geq2$, using $k\geq j$,
the constant term is $k(7j-10)+10j-21>0$ (already $4k-1>0$
when $j=2$). It never vanishes. Hence every ordinary interval
acts by zero.  The full ordinary open-chain Zhu image therefore acts by
its vacuum character on the entire top $M$.


\section{Completion of the weight-zero Zhu calculation}
\label{r:app:h0-completion}

Retain the hypotheses and conclusion of Appendix~\ref{r:app:h0-open-chain}.
We now show that the proper heavy intervals also act trivially.
Only the full-cycle class $C$ remains; its quadratic relation will
determine the quotient governing lowest weight zero.

\subsection{The known ground character and maximal heavy intervals}

Lemma~\ref{neutral:lem:second-ground-character} constructs a simple
ground character $T_0$ on which every proper interval vanishes and $C$
acts by $2\lambda\ne0$.  The argument uses only the quotient $A/W$,
charge support and the invariant pairing; semisimplicity of the lattice
restriction is not required.

Fix $1\leq j<N$ and put $q=t+j$, $Q=R_j$,
$B_j=v_{\mathbf c-e_j}$, $K_j=v_{\mathbf c+e_j}$, where
$\mathbf c=(1,\ldots,1)$ labels $C$.  The exact aligned channels of
$Q\times C$ are $C,B_j,K_j$, with weights $2N,2N+2,2N+2$.
Since $Q=0$ and $H_i=0$, the identities $Q*C=0$ and $[Q,D_jC]=0$
give
\begin{align*}
 0&=\mu_-B_j+\mu_+K_j+\alpha C,\\
 0&=\mu_-d_-B_j+\mu_+d_+K_j+\beta C.
\end{align*}
Here $\mu_\pm\ne0$ are canonical highest-channel coefficients, and all
partial descendants of $C$ have reduced to scalar multiples of $C$.
Right differentiation gives
$d_\pm=h_j(Q)+h_j(C)-h_j(B_j\text{ or }K_j)$.  The two output weights are
\[
 h_j(B_j)=\frac{q-1}{q+1},\qquad
 h_j(K_j)=\frac{(q+2)(q+3)}{q(q+1)},
 \qquad d_--d_+=\frac6q\ne0.
\]
Eliminating $K_j$ expresses $B_j$ as a fixed scalar multiple of $C$,
universally on these tops.  Evaluate this identity on $T_0$:
$B_j=0$ and $C=2\lambda\ne0$, so the scalar is zero.  Thus every maximal
proper heavy interval vanishes.  This argument does not assume that
$K_j$ vanishes on $T_0$.

\subsection{Descending through all proper heavy intervals}

Let $B$ be any nonmaximal proper heavy interval.  Extend its right endpoint
by an ordinary root $j$ to a proper heavy interval $B'$.  Suppose by
descending induction that $B'=0$.  The three outputs of $R_j\times B'$
replace its endpoint label by $0,1,2$.  The first two are $B,B'$, of
weight $H=2N+2$; the third has weight $H+4$.  It cannot occur in
$R_j*B'$, whose states have weight at most $H+2$.  After partial Zhu
reduction,
\[
 0=R_j*B'=\kappa_-\mu_-B+\kappa_0\mu_0B'
            =\kappa_-\mu_-B,
 \qquad
 \kappa_-=-\frac{q+2}{(q+3)(3q+2)}\ne0,
 \quad q=t+j>0.
\]
Therefore $B=0$.

The scalar requires only level two.  At site $j$ the two inputs have
weight $1+2/q$ and the output is the vacuum, whose evaluated series is
$1$.  At site $j-1$ the first, second, and output weights are
\[
 p=\frac{q-2}{q},\qquad x=\frac2{q(q-1)},\qquad
 w=\frac{q+1}{q-1}.
\]
The evaluated local series is
\[
 1-z+\frac{3q^2+10q+4}{(q+3)(3q+2)}z^2+O(z^3).
\]
Indeed, writing the second descendant as
$\alpha L(-2)v+\beta L(-1)^2v$, the Ward equations are
\begin{align*}
 3\alpha+2(2w+1)\beta
    &=\frac{(w+p-x+1)(w+p-x)}{2w},\\
 (4w+c/2)\alpha+6w\beta&=w+2p-x,
 \qquad c=13-6\left(\frac q{q-1}+\frac{q-1}q\right).
\end{align*}
Its Zhu evaluation is $w\alpha+w(w+1)\beta$.
Multiplication of the series by $(1+z)^2$ and extraction of its degree-two
coefficient gives $\kappa_-$.  All other first inputs are the tensor unit.
This proves that every proper heavy interval vanishes.

\subsection{The universal cyclic quadratic and its consequences}

The full Zhu algebra is generated by the partial frame classes, all proper
interval highest classes, and $C$.  The preceding results leave only $C$.
The exact lowest-weight formula
\[
 h(m)=\sum_i(m_i-m_{i+1})^2+Nm_0(m_0+1)
\]
shows that a frame summand with $m_0\geq2$ starts at weight at least $6N$.
Every state in $C*C$ has weight at most $4N<6N$, so $C*C$ lies in the
root-zero filtration piece $F_1W$.  The filtered maximum-plateau Zhu
spanning lemma gives
\[
 [F_1W]\subseteq A_0+
 \sum_{P\text{ proper heavy or }P=C}A_0[P]A_0,
\]
where $A_0$ is the ordinary open-chain image.  Its vacuum action and the
heavy annihilation show that $C^2=\gamma\mathbf1+\delta C$ with fixed
scalars.  The vacuum character gives $\gamma=0$, and $T_0$ gives
$\delta=2\lambda$.  Consequently
\[
 C(C-2\lambda)=0,
 \qquad
 A(W_N)/\langle H_0,\ldots,H_{N-1}\rangle
 \cong\frac{\mathbb C[X]}{X(X-2\lambda)}
 \cong\mathbb C\oplus\mathbb C.
\]
The vacuum and $T_0$ supply both distinct characters.  Thus the
surjection from $\mathbb C[X]/(X(X-2\lambda))$ to the quotient is
an isomorphism.

Every relevant finite top with generalized partial spectrum zero and
exact total conformal operator zero first reduces to this quotient by
Appendix~\ref{r:app:h0-open-chain}.  It is therefore a direct sum of the
known vacuum and $T$ characters.  The Zhu correspondence excludes any
additional neutral aligned simple of lowest weight zero.  In a vacuum
extension by a distinct such simple, the top extension splits over this
semisimple quotient; the previously proved injection of vacuum extensions
into Zhu-top extensions then splits the full module extension.  Vacuum
self-extensions split by
Proposition~\ref{neutral:prop:two-lowest-weights}.

Thus every vacuum extension by a simple target of lowest weight
zero splits.


\section{The heavy-root vacuum coefficient excludes a weight-one defect}
\label{r:app:heavy-vacuum-obstruction}

This calculation uses the closing-root first-image annihilation
from Theorem~\ref{r:thm:first}.
Suppose a nonsplit vacuum extension has a simple target of lowest weight
one.  Lemma~\ref{neutral:lem:ordinary-root-defect} gives nonzero vectors
$w_j$ with
\[
 U=\operatorname{span}\{w_0,\ldots,w_{N-1}\},\qquad
 H_iw_j=\delta_{ij}w_j.
\]
Theorem~\ref{r:thm:first} gives $o(R_0)U=0$ for the closing root
in every odd rank, including $N=3$.  Its total and active local
weights are
\[
 H=2N+2,\qquad a=1+\frac2t,\qquad b=1-2t,
 \qquad a+b=H,\quad t^2+Nt-1=0.
\]
Every partial Virasoro descendant of $R_0$ also annihilates $U$: iterate
\[
 [L_i(-n)v]=(-1)^n\bigl(-H_i*[v]+n[v]*H_i+h_i[v]\bigr)
\]
and use $H_iU\subset U$.  The aligned channels of $R_0\times R_0$ are
the vacuum, $R_0$, and $2e_0$.  The last has lowest weight $6N+8$, above
the maximum weight $2H=4N+4$ in $R_0*R_0$.  Thus its self-channel
annihilates $U$, and it remains to evaluate the vacuum channel on $w_0$.

\subsection{The local evaluated vacuum block}

Let $h$ be an irrational local $(3,1)$ or $(1,3)$ weight, with
\[
 c=\frac{-3h^2+7h-2}{h+1}.
\]
Its highest vector has the null relation
\begin{equation}
 \bigl(L(-1)^3-2(h+1)L(-2)L(-1)+h(h+1)L(-3)\bigr)v_h=0.
 \label{r:eq:heavy-null}
\end{equation}
Normalize the vacuum intertwiner as
\[
 I(v_h,z)v_h=z^{-2h}\sum_{n\geq0}v_nz^n,\qquad v_0=\mathbf1.
\]
Evaluate these vacuum descendants in the highest-weight-one Virasoro
module, with normalized highest vector and dual $\xi,\xi'$, and put
\[
 V_h(z)=\sum_{n\geq0}\langle \xi',o(v_n)\xi\rangle z^n,
 \qquad F(x)=(x-1)^{-2h}V_h(x-1).
\]
This is the formal four-point block with weights $1,h,h,1$ at
$\infty,x,1,0$, in the vacuum channel between the two $h$ insertions.
It is defined by the formal vacuum intertwiner and the ordinary Virasoro
module action, and has normalization $V_h(0)=1$.

Set
\begin{align*}
 A(x)&=\frac{x-1}{x}-\frac{x}{x-1},\\
 B(x)&=\frac{1+2h}{x}-\frac{1+2h}{x-1}
                  +\frac h{(x-1)^2}+\frac1{x^2},\\
 C(x)&=-\frac{x-1}{x^2}+\frac{x}{(x-1)^2},\\
 D(x)&=-\frac{2h}{x^2}+\frac{2h}{(x-1)^2}
                  -\frac{2h}{(x-1)^3}-\frac2{x^3}.
\end{align*}
The null relation gives
\begin{equation}
 F'''-2(h+1)AF''+
 \bigl(-2(h+1)B+h(h+1)C\bigr)F'+h(h+1)DF=0.
 \label{r:eq:heavy-bpz}
\end{equation}
For a direct Ward verification, allow the other insertions to vary at
$y,s$.  At $y=1,s=0$, translation and scaling give
\[
 \partial_yF=-xF'-2hF,\qquad
 \partial_sF=(x-1)F'+2hF.
\]
The $L(-2)$ insertion is $AF'+B_0F$, with
\[
 B_0=\frac1{x^2}+\frac h{(x-1)^2}+\frac{2h}{x}-\frac{2h}{x-1}.
\]
Differentiating the global identities for the first descendant gives
$AF''+BF'$ for $L(-2)L(-1)$.  Since
$L(-3)=[L(-1),L(-2)]$, its insertion is $CF'+DF$, with
$C=A'+B_0-B$ and $D=B_0'$.  This proves
\eqref{r:eq:heavy-bpz}, including the outer-weight-one terms.

\subsection{Uniform solution of the vacuum channel}

Put $u=(x-1)^2/x$, $\theta=u\,d/du$, and substitute
$F=(x-1)^{-2h}G(u)$.  After removing $(x-1)^{-2h}$, the left side of
\eqref{r:eq:heavy-bpz} equals exactly
\begin{equation}
 \frac{4(x+1)}{x^2(x-1)^3}
 \left[
 \theta\left(\theta-\frac h2\right)
       \left(\theta-\frac{1+3h}{2}\right)
 +\frac u4(\theta-h-1)(\theta-h)(\theta+2)
 \right]G(u).
 \label{r:eq:heavy-transformed}
\end{equation}
One can verify the identity directly by expanding the bracket as an
operator in $d/du$.  Its coefficients on $G''',G'',G',G$, respectively,
are
\[
 u^3(1+u/4),\qquad
 \frac{u^2}{2}(5-4h+(2-h)u),\qquad
 \frac u4(3h^2-7h+2+h(h-5)u),\qquad
 \frac u2h(h+1).
\]
The chain rule with $u'=(x^2-1)/x^2$ gives exactly the four coefficients
of the left-hand side after its common factor has been removed.
Construct the ordinary power-series solution $G(0)=1$.  Its recurrence
has no zero denominator because $h$ is irrational, and gives
\begin{equation}
 G(u)=\sum_{n\geq0}d_nu^n,\qquad
 d_n=\left(-\frac14\right)^n
 \frac{(-h-1)_n(-h)_n(2)_n}
 {(1-h/2)_n((1-3h)/2)_n\,n!},
 \label{r:eq:heavy-dn}
\end{equation}
where $(x)_n=x(x+1)\cdots(x+n-1)$.  The exponents of \eqref{r:eq:heavy-bpz} at $x=1$ are
$-2h,-h,h+1$, with nonintegral pairwise differences because $h$ is
irrational.  The normalized Frobenius solution with exponent $-2h$
is therefore unique.  Thus the series just constructed is the vacuum
block, and its dependence on $u$ follows from the equation and
normalization.
Consequently
\[
 V_h(z)={}_3F_2\!\left(
 \begin{matrix}-h-1,-h,2\\1-h/2,(1-3h)/2\end{matrix};
 -\frac{z^2}{4(1+z)}\right).
\]
For comparison with the low-level Ward equations, $d_1=2h/c$ and
\[
 d_2=\frac{h^2(h+1)}{(h-4)(h-2)(3h-1)}
\]
agree with the level-two and level-four vacuum Ward equations.


\subsection{The coefficient cannot vanish}

On $w_0$ the two active partial Zhu weights are $1,0$.  At weight zero,
every positive vacuum descendant evaluates to zero; the second local
series is exactly $1$.  Therefore, after dividing out the nonzero
root vacuum coefficient $\nu_0$, the vacuum contribution to
$R_0*R_0$ is
\[
 P_H(h)=[z^{2H}](1+z)^HV_h(z),\qquad h=a=1+2/t.
\]
Using \eqref{r:eq:heavy-dn},
\[
 (1+z)^HV_h(z)=\sum_{n\geq0}d_nz^{2n}(1+z)^{H-n}.
\]
For $n<H$ the term is a polynomial of degree $H+n<2H$; for $n>H$
its lowest degree is $2n>2H$.  Hence only $n=H$ contributes:
\begin{equation}
 \boxed{\quad
 P_H(h)=\left(-\frac14\right)^H
 \frac{(-h-1)_H(-h)_H(2)_H}
 {(1-h/2)_H((1-3h)/2)_H\,H!}\ne0.
 \quad}
 \label{r:eq:heavy-nonzero}
\end{equation}
Every possible zero or pole of these finite Pochhammer products would
make $h$ rational, whereas $1+2/t$ is irrational.  The scalar $\nu_0$
is nonzero by the nondegenerate invariant form on the multiplicity-one
self-contragredient summand \(M_{e_0}\).

Finally apply $R_0*R_0$ to $w_0\ne0$.  Its left side is zero because
$R_0U=0$.  Its self-channel descendants annihilate $U$, and the higher
channel is excluded by weight.  Its vacuum channel equals
$\nu_0P_H(a)w_0\ne0$, a contradiction.  Thus no nonsplit vacuum
extension with a lowest-weight-one simple target exists.  Together with the weight-zero analysis and
Proposition~\ref{neutral:prop:two-lowest-weights}, this proves that
every extension of the vacuum by a simple module splits.

\section{The uniform symbolic frame certificate}
\label{app:certificate}

The screening-frame construction has one computer-assisted algebraic
step. The charge bound in Section~\ref{c:app:frame} reduces the
Virasoro relations to proper intervals, charge $B$, and finitely many
tails $B+I$. The proper-interval and tail equations are proved there.
The program reproduced below checks the remaining charge-$B$
identities with $N$, $t$, and the ordered site positions left
indeterminate.

A coefficient involves at most two frame indices and three oscillator
indices. The program enumerates their equality and order patterns;
its interval-count formula includes the end sites. After the
denominators prescribed by the frame formulas are cleared, every
residual numerator expands to zero. The frame proof separately
establishes that those denominators are nonzero at the parameters
used in the theorem. Together with the mathematical reduction to
these patterns, this makes the calculation uniform in $N$.

The program was extracted from this manuscript and rerun on
23 October 2026. Its output was
\texttt{UNIFORM\_SYMBOLIC\_PASS}, with the following counts:
\begin{center}
\begin{tabular}{@{}lrrrrr@{}}
\toprule
Distinct site indices & 1 & 2 & 3 & 4 & 5\\
Coefficient identities & 7 & 40 & 69 & 46 & 10\\
\bottomrule
\end{tabular}
\end{center}
All 172 residual numerators vanish identically. The calculation uses
exact symbolic arithmetic, with no numerical substitution for the rank
or site positions. It checks these identities, not the surrounding
representation-theoretic arguments.

To reproduce the calculation, save the listing as
\nolinkurl{certify_uniform_epsilon.py} and run it with Python~3 and
SymPy. It prints a JSON report and writes the same report beside the
script. The program requires no external mathematical data and makes
no network requests. The accompanying source-and-checks archive
contains this program and the supplementary programs, together
with their execution reports. In particular,
\nolinkurl{uniform_cut_jets.py} checks the eight first jets
from Appendix~\ref{neutral:app:cut-polynomials},
\nolinkurl{first_order_all_interval_checks.py} assembles the endpoint
matrices from those derivatives and the transport formulas, and
\nolinkurl{direct_charge_B_n3.py} provides the independent Fock-mode
check described in Section~\ref{c:app:frame}. The remaining programs
check the BPZ, Zhu, interval, heavy-root, and dimension identities.
The archive's README states each check's scope; successful algebraic
checks do not certify the surrounding representation theory.

The TeX source is self-contained: it uses no external
\verb|\input| file or bibliography database. The symbolic certificate
is not a proof-assistant verification of the manuscript.

\subsection{Complete program}

\begin{lstlisting}
"""Uniform symbolic verification of the charge-B Virasoro-frame equations.

N, t, and all distinct site positions are indeterminates. This enumerates
order/equality types of at most five site indices, NOT numerical ranks.
Interval sums use the exact counting formula explained in the proof note.
"""
from collections import Counter
from functools import lru_cache
from itertools import combinations_with_replacement, product
from pathlib import Path
import json
import time
import sympy as s
# Exact polynomial helpers copied from the original helper module.
# Keeping them here removes the unrelated numerical exploration dependency.
def poly_add(*args):
    out = {}
    for p in args:
        for k,v in p.items():
            out[k] = out.get(k,0)+v
    return {k:v for k,v in out.items() if v != 0}

def scale(p,c):
    return {k:c*v for k,v in p.items() if c*v != 0}

def mul(p,q):
    out={}
    for a,x in p.items():
        for b,y in q.items():
            k=tuple(u+v for u,v in zip(a,b))
            out[k]=out.get(k,0)+x*y
    return {k:v for k,v in out.items() if v != 0}


N,t=s.symbols('N t')


def ba(*args):
    ans={}
    for p in args:
        for key,c in p.items():ans[key]=ans.get(key,0)+c
    return {key:c for key,c in ans.items() if c != 0}

def bm(p,q):
    ans={}
    for (M,H),a in p.items():
        for (M2,H2),b in q.items():
            key=(M|M2,H|H2);ans[key]=ans.get(key,0)+a*b
    return {key:c for key,c in ans.items() if c != 0}

def bs(p,c):return {key:v*c for key,v in p.items() if v*c != 0}

BC={(frozenset(),frozenset()):s.S.One}


def setup(m):
    positions=s.symbols(f'j0:{m}')
    r=[t+j for j in positions];p=[1/(x+1) for x in r]
    psi=[(x+1)/x for x in r];R=(t+N)/t
    def bv(kind,i,offset=0):
        point=(i,offset)
        return {(frozenset([point]) if kind=='M' else frozenset(),
                 frozenset([point]) if kind=='H' else frozenset()):s.S.One}
    mom=[];u=[];v=[]
    for i in range(m):
        mi,mnext=bv('M',i),bv('M',i,1)
        hi,hnext=bv('H',i),bv('H',i,1)
        mom.append(ba(mi,bs(mnext,-r[i]/(r[i]+1))))
        u.append(ba(bs(mi,1/r[i]),bs(mnext,-1/(r[i]+1))))
        gi=ba(BC,bs(mi,-1),bs(hi,R-1))
        gn=ba(BC,bs(mnext,-1),bs(hnext,R-1))
        v.append(ba(bs(gi,1/r[i]),bs(gn,-1/(r[i]+1))))
    comp=[ba(bs(BC,p[i]),bs(mom[i],-1)) for i in range(m)]
    def val(point):return positions[point[0]]+point[1]
    def interval_sum(pol):
        terms=[]
        for (M,H),c in pol.items():
            upper=min(H) if H else (m,0)
            upperval=val(upper) if H else N
            if M:
                lo,hi=min(M),max(M)
                count=0 if upper<=hi else val(lo)*(upperval-val(hi))
            else:count=upperval*(upperval-1)/2
            if count!=0:terms.append(c*count)
        return s.Add(*terms)
    @lru_cache(None)
    def source(i,j,indices):
        left=bm(u[i],v[j]);right=bm(v[i],u[j])
        for k in indices:
            left=bm(left,mom[k]);right=bm(right,comp[k])
        return interval_sum(ba(left,right))

    variables=3*m
    one={(0,)*variables:s.S.One}
    def atom(i,k):
        key=[0]*variables;key[3*i+k-1]=1;return {tuple(key):s.S.One}
    atoms={(i,k):atom(i,k) for i in range(m) for k in (1,2,3)}
    def D(pol,i):
        ans=mul(scale(atoms[i,1],p[i]),pol)
        for key,c in pol.items():
            for k in (1,2):
                power=key[3*i+k-1]
                if power:
                    kk=list(key);kk[3*i+k-1]-=1;kk[3*i+k]+=1
                    ans=poly_add(ans,{tuple(kk):k*power*c})
        return ans
    def Dt(pol):return poly_add(*(D(pol,i) for i in range(m)))
    def action(pol,i,mode):
        if mode==-1:return D(pol,i)
        if mode==0:return {key:c*sum(k*key[3*i+k-1] for k in (1,2,3)) for key,c in pol.items()}
        ans={}
        for key,c in pol.items():
            if mode==1 and key[3*i+1]:
                kk=list(key);kk[3*i+1]-=1;kk[3*i]+=1
                ans=poly_add(ans,{tuple(kk):2*c*key[3*i+1]})
            if mode==2:
                if key[3*i]==2:
                    kk=list(key);kk[3*i]-=2
                    ans=poly_add(ans,{tuple(kk):2*psi[i]*c})
                if key[3*i+1]:
                    kk=list(key);kk[3*i+1]-=1
                    ans=poly_add(ans,{tuple(kk):-2*psi[i]*p[i]*c})
        return ans
    corrections=[]
    for i in range(m):
        pol=scale(atoms[i,2],(1+R)*p[i]/2)
        for j in range(m):
            diag=(R*p[i]-r[i]*p[i]**2)/2 if j==i else -(R if j<i else 1)*r[j]*p[j]/(2*r[i]*(r[i]+1))
            pol=poly_add(pol,scale(mul(atoms[j,1],atoms[j,1]),diag))
            if j!=i:
                cross=(R if j<i else 1)*p[i]*p[j]
                pol=poly_add(pol,scale(mul(atoms[i,1],atoms[j,1]),cross))
        corrections.append(pol)
    def tag(pol,k):return {(k,)+key:c for key,c in pol.items()}
    @lru_cache(None)
    def linear(i,j):
        ans={}
        for k in range(4):
            ans=poly_add(ans,tag(scale(action(corrections[j],i,k-1),s.S.One/s.factorial(k)),k))
            base=action(corrections[i],j,k-1)
            for zpower in range(k,-1,-1):
                if zpower!=k:base=Dt(base)
                coefficient=-(-1)**k*s.binomial(k,zpower)/s.factorial(k)
                ans=poly_add(ans,tag(scale(base,coefficient),zpower))
        if i==j:
            ans=poly_add(ans,tag(scale(Dt(corrections[i]),-1),0),tag(scale(corrections[i],-2),1))
        return ans
    return source,linear,variables


def cases(m):
    labels=range(m)
    # (lambda power, oscillator list, exponential-series coefficient)
    shapes=[(3,(),s.Rational(1,6))]
    for k in labels:
        shapes.extend([(2,((k,1),),s.Rational(1,2)),
                       (1,((k,2),),s.Rational(1,2)),
                       (0,((k,3),),s.Rational(1,3))])
    for k,l in combinations_with_replacement(labels,2):
        shapes.append((1,((k,1),(l,1)),s.Rational(1,2) if k==l else s.S.One))
    for k,l in product(labels,repeat=2):
        shapes.append((0,((k,2),(l,1)),s.Rational(1,2)))
    for inds in combinations_with_replacement(labels,3):
        denominator=s.prod(s.factorial(c) for c in Counter(inds).values())
        shapes.append((0,tuple((k,1) for k in inds),1/denominator))
    for i in labels:
        for j in range(i,m):
            for zpower,osc,factor in shapes:
                if {i,j}|{k for k,_ in osc} != set(labels):continue
                yield i,j,zpower,osc,factor


def verify(max_sites=5):
    results=[];start=time.time()
    for m in range(1,max_sites+1):
        source,linear,variables=setup(m)
        checked=0
        for i,j,zpower,osc,factor in cases(m):
            key=[0]*variables
            for k,degree in osc:key[3*k+degree-1]+=1
            lhs=linear(i,j).get((zpower,)+tuple(key),0)
            rhs=factor*source(i,j,tuple(sorted(k for k,_ in osc)))
            residual=s.expand(s.together(lhs+rhs).as_numer_denom()[0])
            if residual!=0:
                raise AssertionError({'sites':m,'i':i,'j':j,'lambda_power':zpower,
                                      'oscillators':osc,'residual':str(s.factor(residual))})
            checked+=1
        results.append({'distinct_site_indices':m,'coefficient_order_types':checked})
        print(json.dumps({'stage':'PASS','distinct_sites':m,'identities':checked,'seconds':round(time.time()-start,2)}),flush=True)
    return {'status':'UNIFORM_SYMBOLIC_PASS','parameters':'Independent N,t,j0<...<jk; no numerical rank substitution',
            'interval_moment_formula':'min(M)*max(0,min(H or N)-max(M)); for empty M, U*(U-1)/2',
            'all_order_types':results,'total_identities':sum(x['coefficient_order_types'] for x in results)}


if __name__=='__main__':
    result=verify()
    Path(__file__).with_suffix('.json').write_text(json.dumps(result,indent=2)+'\n')
    print(json.dumps(result,indent=2))

\end{lstlisting}

\Needspace{6\baselineskip}
\phantomsection
\addcontentsline{toc}{section}{References}
\begin{thebibliography}{99}
% Compact reference type.
\small
\raggedright

\bibitem{izumi}
M. Izumi,
\emph{The structure of sectors associated with Longo--Rehren inclusions.
II: Examples}, Rev. Math. Phys. 13 (2001), no. 5, 603--674.
\url{https://doi.org/10.1142/S0129055X01000818}.

\bibitem{evansgannon2011}
D. E. Evans and T. Gannon,
\emph{The exoticness and realisability of twisted Haagerup--Izumi
modular data}, Commun. Math. Phys. 307 (2011), 463--512.
\url{https://arxiv.org/abs/1006.1326}.

\bibitem{evansgannon2017}
D. E. Evans and T. Gannon,
\emph{Non-unitary fusion categories and their doubles via endomorphisms},
Adv. Math. 310 (2017), 1--43.
\url{https://arxiv.org/abs/1506.03546}.

\bibitem{huangcyclic}
T.-C. Huang,
\emph{Cyclic Haagerup--Izumi fusion categories at every odd order},
2026. \url{https://arxiv.org/abs/2609.15986}.

\bibitem{gsy}
T. Gannon, A. Schopieray and H. Yadav,
\emph{On Haagerup--Izumi fusion categories}, 2026.
\url{https://arxiv.org/abs/2609.25185v1}.
Theorem 5.1 and Appendix F.1 specify the center modular data;
Appendix F, Table 2 gives the negative-branch comparison labels.
Section 6 gives the forward HI-to-near-group construction and
states the additional central data needed for an inverse.
Corollary 5.9 gives the order-fifteen nonuniqueness example.


\bibitem{c:Koshida}
S. Koshida, \emph{Module categories of the generic Virasoro VOA and quantum groups}, SIGMA 21 (2025), 039; \url{https://sigma-journal.com/2025/039/sigma25-039.pdf}.


\bibitem{virasoro}
T. Creutzig, C. Jiang, F. Orosz Hunziker,
D. Ridout and J. Yang,
\emph{Tensor categories arising from the Virasoro algebra},
Adv. Math. 380 (2021), 107601.
\url{https://arxiv.org/abs/2002.03180}.
Numbering follows arXiv version 2.


\bibitem{gluing}
T. Creutzig, S. Kanade and R. McRae, \emph{Gluing vertex algebras}, Adv. Math. 396 (2022), 108174; \href{https://arxiv.org/abs/1906.00119}{arXiv:1906.00119}. Numbering follows arXiv version 2.


\bibitem{c:DLMRegular}
C. Dong, H. Li and G. Mason, \emph{Regularity of rational vertex operator algebras}, Adv. Math. 132 (1997), 148--166; \href{https://arxiv.org/abs/q-alg/9508018}{arXiv:q-alg/9508018}, Theorem 3.16.


\bibitem{c:LiForm}
H. Li, \emph{Symmetric invariant bilinear forms on vertex operator algebras}, J. Pure Appl. Algebra 96 (1994), 279--297.


\bibitem{c:TK}
A. Tsuchiya and Y. Kanie, \emph{Fock space representations of the Virasoro algebra---intertwining operators}, Publ. Res. Inst. Math. Sci. 22 (1986), 259--327; \href{https://doi.org/10.2977/prims/1195178069}{doi:10.2977/prims/1195178069}. Proposition 2.7 and its proof in Section 2.5; \href{https://ems.press/content/serial-article-files/42320}{open publisher PDF}.


\bibitem{c:McRaeDeligne}
R. McRae, \emph{Deligne tensor products of categories of modules for vertex operator algebras}, \href{https://arxiv.org/abs/2304.14023}{arXiv:2304.14023}.


\bibitem{cmy}
T. Creutzig, R. McRae and J. Yang,
\emph{Direct limit completions of vertex tensor categories},
Commun. Contemp. Math. 24 (2022), no. 2, 2150033.
\url{https://arxiv.org/abs/2006.09711}.
We use Theorems 7.5 and 7.7 of arXiv version 2.


\bibitem{c:MPS}
S. Morrison, E. Peters and N. Snyder, \emph{Categories generated by a trivalent vertex}, Selecta Math. (N.S.) 23 (2017), no. 2, 817--868; \href{https://arxiv.org/abs/1501.06869}{arXiv:1501.06869}.


\bibitem{huang}
Y.-Z. Huang, \emph{Cofiniteness conditions, projective covers and the logarithmic tensor product theory}, J. Pure Appl. Algebra 213 (2009), 458--475; \href{https://arxiv.org/abs/0712.4109}{arXiv:0712.4109}.


\bibitem{c:Moriwaki}
Y. Moriwaki, \emph{Quantum coordinate ring in WZW model and affine vertex algebra extensions}, Selecta Math. 28 (2022), 68; \href{https://arxiv.org/abs/2111.11357}{arXiv:2111.11357}.


\bibitem{c:Miyamoto}
M. Miyamoto, \emph{$C_1$-cofiniteness and fusion products for vertex operator algebras}, in \emph{Conformal Field Theories and Tensor Categories}, Math. Lect. Peking Univ., Springer, Heidelberg, 2014, 271--279; \href{https://arxiv.org/abs/1305.3008}{arXiv:1305.3008}.


\bibitem{huangyang}
Y.-Z. Huang and J. Yang,
\emph{Logarithmic intertwining operators and associative algebras},
J. Pure Appl. Algebra 216 (2012), no. 6, 1467--1492.
We use the corrected arXiv version 3, Lemma 5.5 and
Propositions 5.6--5.8.
\url{https://arxiv.org/abs/1104.4679}.
See also the authors' erratum:
\url{https://sites.math.rutgers.edu/~yzhuang/rci/math/papers/HY-erratum.pdf}.


\bibitem{zhu}
Y. Zhu,
\emph{Modular invariance of characters of vertex operator algebras},
J. Amer. Math. Soc. 9 (1996), no. 1, 237--302.
\url{https://doi.org/10.1090/S0894-0347-96-00182-8}.
Theorem 2.2.2 gives the correspondence between simple modules
and simple modules for the Zhu algebra.


\bibitem{alsw}
R. Allen, S. Lentner, C. Schweigert, and S. Wood,
\emph{Duality structures for representation categories of vertex
operator algebras and the Feigin--Fuchs boson},
Selecta Math. (N.S.) 31 (2025), article 36.
Numbering follows arXiv version 2.
\url{https://arxiv.org/abs/2107.05718}.


\bibitem{ep}
P. Etingof and D. Penneys,
\emph{Rigidity of non-negligible objects of moderate growth in
braided categories}, Forum Math. Pi 14 (2026), e7.
Theorem numbering follows arXiv:2412.17681v3.
\url{https://arxiv.org/abs/2412.17681}.


\bibitem{mcrae}
R. McRae,
\emph{On rationality for $C_2$-cofinite vertex operator algebras},
Camb. J. Math. 14 (2026), no. 1, 1--115.
\url{https://arxiv.org/abs/2108.01898}.
Numbering follows arXiv version 3.


\bibitem{abd}
T. Abe, G. Buhl, and C. Dong,
\emph{Rationality, regularity, and \(C_2\)-cofiniteness},
Trans. Amer. Math. Soc. 356 (2004), no. 8, 3391--3402.
\url{https://arxiv.org/abs/math/0204021}.
Proposition 5.2 gives $C_2$-cofiniteness of irreducible weak modules;
its CFT-type hypothesis is applied to $V^{\mathrm{un}}$, not to $E_N$.


\bibitem{dgno}
V. Drinfeld, S. Gelaki, D. Nikshych and V. Ostrik,
\emph{On braided fusion categories I},
Selecta Math. 16 (2010), 1--119.
\url{https://arxiv.org/abs/0906.0620}.
We use Theorem 4.18, Proposition 4.19 and Remark 4.21.


\bibitem{mimachi}
K. Mimachi, \emph{The irreducibility of the monodromy representation
associated with the Dotsenko--Fateev equation},
Commun. Math. Phys. 405 (2024), article 71,
\href{https://doi.org/10.1007/s00220-024-04942-7}{doi:10.1007/s00220-024-04942-7},
especially equation (1.1) and Theorem 11.


\bibitem{km}
M. Krauel and M. Miyamoto,
\emph{A modular invariance property of multivariable trace functions
for regular vertex operator algebras}, J. Algebra 444 (2015), 124--142.
\url{https://arxiv.org/abs/1501.00124}.
Lemmas 4--5 and equations (2.2), (2.5) fix the normalization of the
formal Ward recursion used here.


\bibitem{ko}
A. Kirillov, Jr., and V. Ostrik,
\emph{On a $q$-analog of the McKay correspondence and the ADE
classification of $\widehat{\mathfrak{sl}}_2$ conformal field theories},
Adv. Math. 171 (2002), 183--227.
\url{https://arxiv.org/abs/math/0101219}.
Numbering follows arXiv version 3: Definition 1.11,
Theorems 1.17--1.18, Remark 1.19 and Lemma 1.20.


\bibitem{ckm}
T. Creutzig, S. Kanade and R. McRae,
\emph{Tensor categories for vertex operator superalgebra extensions},
Mem. Amer. Math. Soc. 295 (2024), no. 1472.
\url{https://arxiv.org/abs/1705.05017}.
We use Lemma 2.61 and Theorem 3.65 of arXiv version 2.


\bibitem{cmsy}
T. Creutzig, R. McRae, K. Shimizu, and H. Yadav,
\emph{Commutative algebras in Grothendieck--Verdier categories,
rigidity, and vertex operator algebras},
Commun. Contemp. Math. 28 (2026), no. 10, 2550091.
Numbering follows arXiv version 2 (18 October 2025).
\url{https://arxiv.org/abs/2409.14618}.


\bibitem{dmno}
A. Davydov, M. M\"uger, D. Nikshych and V. Ostrik,
\emph{The Witt group of non-degenerate braided fusion categories},
J. Reine Angew. Math. 677 (2013), 135--177.
\url{https://arxiv.org/abs/1009.2117}.
We use Section 3.3, Section 3.6, Proposition 3.16, Corollary 3.30,
Remark 3.31 and Corollary 4.1.


\bibitem{egno}
P. Etingof, S. Gelaki, D. Nikshych and V. Ostrik,
\emph{Tensor Categories}, Mathematical Surveys and Monographs 205,
American Mathematical Society, 2015.
\url{https://math.mit.edu/~etingof/egnobookfinal.pdf}.
We use Proposition 9.2.2.


\bibitem{dongmason}
C. Dong and G. Mason,
\emph{Shifted vertex operator algebras},
Math. Proc. Cambridge Philos. Soc. 141 (2006), 67--80.
\url{https://arxiv.org/abs/math/0411526}.
We use Theorems 3.1, 4.1 and 4.4.



\bibitem{evanssfold}
D. E. Evans, \emph{A vertex operator algebra for the S-fold realisation
of a non-unitary Haagerup--Izumi modular fusion category}, 2026.
\url{https://arxiv.org/abs/2610.26341}.
Section 6.2 proposes the cyclic Virasoro sum used here.

\bibitem{cgduality}
T. Creutzig and D. Gaiotto, \emph{Vertex algebras for S-duality},
Commun. Math. Phys. 379 (2020), 785--845.
\url{https://arxiv.org/abs/1708.00875}.

\bibitem{flohr1993}
M. Flohr, \emph{$W$-algebras, new rational models and completeness
of the $c=1$ classification}, Commun. Math. Phys. 157 (1993), 179--212.
\url{https://arxiv.org/abs/hep-th/9207019}.

\bibitem{flohr1994}
M. Flohr, \emph{More curiosities at effective $c=1$},
Mod. Phys. Lett. A 9 (1994), 1071--1082.
\url{https://arxiv.org/abs/hep-th/9312097}.

\bibitem{gangkim}
D. Gang and D. Kim, \emph{Generalized non-unitary Haagerup--Izumi
modular data from 3D S-fold SCFTs}, JHEP 03 (2023), 185.
\url{https://arxiv.org/abs/2211.13561}.

\bibitem{gangkimlee}
D. Gang, D. Kim and S. Lee, \emph{A non-unitary bulk-boundary
correspondence: Non-unitary Haagerup RCFTs from S-fold SCFTs},
SciPost Phys. 17 (2024), 064.
\url{https://arxiv.org/abs/2310.14877}.

\bibitem{jeongsohn}
K. Jeong and H. Sohn, \emph{Non-unitary Haagerup-like TQFTs and
RCFTs from generalized S-fold SCFTs}, 2026.
\url{https://arxiv.org/abs/2608.11946}.

\bibitem{evansgannon2014}
D. E. Evans and T. Gannon, \emph{Near-group fusion categories and
their doubles}, Adv. Math. 255 (2014), 586--640.
\url{https://arxiv.org/abs/1208.1500}.

\bibitem{riesenc8}
A. Riesen, \emph{The Haagerup VOA at central charge eight},
announcement, 2026. \url{https://arxiv.org/abs/2610.13933}.

\bibitem{evanswords}
D. E. Evans, \emph{Haagerup--Izumi modular data from hyperbolic
words: S-fold TQFTs, central charge, torsion and the Bethe vacua},
2026. \url{https://arxiv.org/abs/2610.18512}.

\bibitem{grossmanizumi}
P. Grossman and M. Izumi, \emph{Drinfeld centers of fusion categories
arising from generalized Haagerup subfactors},
Quantum Topol. 13 (2022), 593--668.
\url{https://arxiv.org/abs/1501.07679}.

\bibitem{grossmansnyder}
P. Grossman and N. Snyder, \emph{Quantum subgroups of the Haagerup
fusion categories}, Commun. Math. Phys. 311 (2012), 617--643.
\url{https://arxiv.org/abs/1102.2631}.

\end{thebibliography}
\end{document}
